\documentclass[11pt,a4paper]{article}
\usepackage[T1]{fontenc}
\usepackage{lmodern}
\pdfextension mapfile{+lm.map}
\pdfextension mapfile{+cm.map}
\pdfextension mapfile{+symbols.map}
\pdfextension mapfile{+euler.map}
\usepackage[margin=27mm,headheight=15pt]{geometry}
\usepackage{amsmath,amssymb,amsthm,mathtools,bm}
\usepackage{tikz-cd,longtable,booktabs,array,enumitem,microtype}
\usepackage[authoryear,round]{natbib}
\usepackage{xcolor,hyperref}
\hypersetup{colorlinks=true,linkcolor=blue!40!black,citecolor=blue!40!black,urlcolor=blue!40!black}
\usepackage{fancyhdr}
\pagestyle{fancy}\fancyhf{}\fancyhead[L]{\small RESEARCH\_SA05}\fancyhead[R]{\small Proof audit}\fancyfoot[C]{\thepage}
\newcommand{\Q}{\mathbb Q}\newcommand{\Z}{\mathbb Z}\newcommand{\R}{\mathbb R}\newcommand{\C}{\mathbb C}\newcommand{\F}{\mathbb F}
\newcommand{\Gm}{{\mathbb G_m}}\newcommand{\Sha}{\mathop{\mathrm{Sha}}}
\DeclareMathOperator{\Gal}{Gal}\DeclareMathOperator{\Hom}{Hom}\DeclareMathOperator{\End}{End}\DeclareMathOperator{\Spec}{Spec}\DeclareMathOperator{\Rep}{Rep}
\newcommand{\et}{\mathrm{\acute et}}\newcommand{\cris}{\mathrm{cris}}
\DeclareMathOperator{\Aut}{Aut}\DeclareMathOperator{\Br}{Br}\DeclareMathOperator{\Pic}{Pic}\DeclareMathOperator{\coker}{coker}
\DeclareMathOperator{\im}{im}\DeclareMathOperator{\rk}{rk}\DeclareMathOperator{\CH}{CH}\DeclareMathOperator{\Sym}{Sym}
\newtheorem{theorem}{Theorem}[section]\newtheorem{proposition}[theorem]{Proposition}\newtheorem{lemma}[theorem]{Lemma}\newtheorem{corollary}[theorem]{Corollary}
\theoremstyle{definition}\newtheorem{definition}[theorem]{Definition}\newtheorem{example}[theorem]{Example}\newtheorem{question}[theorem]{Question}
\theoremstyle{remark}\newtheorem{remark}[theorem]{Remark}
\newcommand{\status}[1]{\par\noindent\textbf{Status: #1.}\par}
\newcommand{\ms}[1]{\textnormal{[manuscript #1]}}
\providecommand{\mathscr}[1]{\mathcal{#1}}
\newcommand{\needcheck}[1]{\par\noindent\textbf{Unresolved proof gate.} #1\par}
\setlength{\emergencystretch}{2em}
\allowdisplaybreaks
\title{Grothendieck's finite-field motivic gerbe and the scope of his absolute cohomology\\\large RESEARCH\_SA05: a proof audit of cotes 11, 12, and 15}
\author{AI-produced by GPT-6 Astra (OpenAI, ChatGPT Pro mode)}
\date{19 November 2026}
\begin{document}
\maketitle
\begin{abstract}
\textbf{INCOMPLETE.} The programme stopped under the required two-review rule; the accompanying ledger identifies every unresolved item. We distinguish the finite-field motivic group, its connected component, and the gerbe obtained by rigidifying the Artin--Tate fibre functor. We prove and audit statements against the supplied AI transcriptions and primary literature. No original manuscript scans were supplied, so disputed readings are never certified palaeographically. No priority is claimed. Every assertion retains its stated hypotheses and its actual proof status.
\end{abstract}
\tableofcontents
\clearpage
\section{A. The finite-field gerbe: bands, local classes, and descent}
\label{sec:A}

\subsection{Conventions and the exact hypothesis}
We write $\overline{\Q}\subset\C$ and fix an embedding into
$\overline{\Q}_p$. The Tate object $\Q(1)$ has weight $-2$.
The Frobenius endomorphism $\pi_X$ of a variety over $\F_q$ is induced
by raising its coordinates to their $q$th powers; its pullback on cohomology
has the usual Weil eigenvalues, and $\pi_{\Q(-1)}=q$.
On the geometric-point fibre functor of the Artin category, this is
\emph{geometric} Galois Frobenius. For local Brauer groups we use the
arithmetic normalization
\[
 \operatorname{inv}_K(L/K,\sigma_{\rm ar},a)
       =\frac{\operatorname{ord}_K(a)}{[L:K]}
 \quad(L/K\text{ unramified}).
\]
Here the cyclic algebra has $uxu^{-1}=\sigma_{\rm ar}(x)$ and
$u^{[L:K]}=a$. In particular the endomorphism algebra of a simple
isocrystal of slope $s$ has invariant $-s$. Reversing the reciprocity
normalization reverses all nonarchimedean invariants below; it does not
change the classes represented by the stated cyclic algebras.
For a gerbe we use arrows $\phi_\sigma:{}^\sigma\omega\to\omega$
and cocycle $c_{\sigma,\tau}=\phi_\sigma{}^\sigma\phi_\tau
\phi_{\sigma\tau}^{-1}$. With this convention the class evaluated
at a character is the \emph{opposite} of the corresponding simple
object's endomorphism-algebra class. Indeed the semilinear operators
on its geometric fibre make it a left module over the crossed product
of $c$, and its descended endomorphisms are the centralizer of that
crossed product. The two central-simple-algebra classes sum to zero.

Our conditional category $\mathcal M_q$ is the category of rational
numerical pure motives over $\F_q$, with all Tate twists and the
commutativity constraint corrected by the K\"unneth projectors so that
its realizations have ordinary, rather than super, tensor symmetry.
We assume the full Tate conjecture for every smooth projective variety
and every product over every finite extension of $\F_q$: the order of
the pole of the zeta function in codimension $r$ is the dimension of
numerical cycle classes, with the corresponding cycle-class
surjectivity and semisimplicity at the eigenvalue $1$. This is the
hypothesis used in \citet[\S1, especially Propositions 1.15 and
1.17]{AMilne1994}; it gives full faithfulness after extension of
coefficients for both prime-to-$p$ and crystalline realizations.
Where a polarization is invoked we additionally specify the Hodge
standard positivity for the abelian motives involved. No assertion
about arbitrary Chow motives follows from these hypotheses.

The category is semisimple and Tannakian, and its geometric fibre
functors have a commutative band $G$. Thus $G$ is a group over $\Q$
even when no $\Q$-valued fibre functor exists. We use $H^i(k,T)$ for
continuous protorus cohomology, with the fpqc interpretation of torsors
and banded gerbes; the finite-level derived-limit exact sequence is
retained throughout. It is not our definition that $H^i$ commutes with
inverse limits.

\subsection{A.1. Correcting the Weil-number group and the rigidification}

\begin{theorem}[The band and its canonical decompositions]
\label{A:band}
\status{Proved under the stated Tate hypothesis; the group calculation is unconditional}
Put
\[
 \Gamma=W(p)=\{\lambda\in\overline\Q^{\times}:\
  p^N\lambda\text{ is integral for some }N\geq0,
  \quad |\tau\lambda|=p^{w(\lambda)/2}\ (\tau),\quad
  w(\lambda)\in\Z\}.
\]
Then $G=D(\Gamma)$, where $D$ is the contravariant equivalence from
Galois modules to groups of multiplicative type. Let $\mu_\infty$
be the group of roots of unity, $\Gamma^0=\ker(w)$, and
$M=\Gamma^0/\mu_\infty$. There are canonical identifications
\begin{align*}
 X^*(G^0)&=\Gamma/\mu_\infty=M\oplus\Z[\sqrt p],\\
 G/G^0&=\widehat\Z,\qquad G^0=H^0\times\Gm,
 \qquad H^0=D(M),\\
 X^*(H)&=\Gamma'=\Gamma/\langle\mu_\infty,p\rangle
            =M\oplus\Z/2\Z,\\
 H&=H^0\times\mu_2.
\end{align*}
The maps $\epsilon:G\to\Gm$ and $i:\Gm\to G$ have character
maps $1\mapsto p$ and $\lambda\mapsto w(\lambda)$, respectively,
and $\epsilon i(t)=t^2$. The restriction $i|_{\mu_2}$ splits $H$.
The group $M$ is a $\Q$-vector space, and complex conjugation acts
as $-1$. Consequently $H^0$ is a protorus, an inverse limit of tori
compact at the real place. It is not an algebraic torus of finite type.
\end{theorem}
\begin{proof}
The Tate full-faithfulness assertion identifies endomorphisms of every
tensor construction on a motive, after extension to $\Q_\ell$, with
the centralizer of its semisimple Frobenius. The stabilizer of these
tensor invariants is the Zariski closure of Frobenius. It is
commutative and diagonalizable. Descent therefore makes the band a
group of multiplicative type. Its characters are the Frobenius
eigenvalues, with their actual $\Gal(\overline\Q/\Q)$ action.

All eigenvalues satisfy the displayed conditions. Conversely take
$\lambda\in\Gamma$ and choose $N$ so that $a=p^N\lambda$ is integral
and its weight $m=w(\lambda)+2N$ is positive. Any $m$th root $b$ of
$a$ is integral: it is integral over the ring of algebraic integers,
which is integrally closed. Every complex conjugate of $b$ has
absolute value $\sqrt p$. Thus $b$ is a Weil $p$-integer of weight
one. Honda--Tate supplies an abelian variety over $\F_p$ with $b$
as a Frobenius eigenvalue. The eigenvalue $b^m p^{-N}=\lambda$
then occurs in a tensor power with a Tate twist. The case of weight
zero integral $a$ can instead be treated by Artin motives, since
Kronecker's theorem makes $a$ a root of unity. This proves the whole
character group, not merely a subgroup of it. This also derives the
motive version of Honda--Tate used in
\citet[Proposition 2.6 and Remark 2.7]{AMilne1994}.

The torsion in $\Gamma$ is exactly $\mu_\infty$. The dual
$D(\mu_n)$ is the constant cyclic group $\Z/n\Z$: its geometric
characters have precisely the cyclotomic Galois action. Passing to
the limit gives the constant profinite quotient $\widehat\Z$.
For $\lambda$ of weight $w$, the formula
\[
 [\lambda]\longmapsto
 \bigl([\lambda p^{-w/2}],w\bigr)
\]
is an isomorphism $\Gamma/\mu_\infty\to M\oplus\Z$.
The sign of the square root and every Galois change of that root
are roots of unity, so the formula is canonical on these quotients
and Galois equivariant. Under it $p$ maps to $(0,2)$.
All the asserted exact sequences, the torsion $\Z/2\Z$ in
$\Gamma'$, and the splitting by $i|_{\mu_2}$ follow.

If $u\in\Gamma^0$, every algebraic $n$th root of $u$ is again
in $\Gamma^0$. The quotient by roots of unity is torsion free:
$u^n$ torsion implies $u$ torsion. It is therefore uniquely divisible.
For every embedding, $\overline{u}=u^{-1}$, giving the action $-1$.
A finite Galois-stable lattice in $M$ is the character lattice of a
compact real torus, so these lattices give the asserted inverse limit.
Finally let $a$ be a root of $X^2-X+p$. The imaginary quadratic field
$\Q(a)$ has $p$ split, including $p=2$ (discriminant $-7$).
The element $u=a/\bar a$ has normalized $p$-adic valuations $1,-1$
and is not torsion. Its rational powers give a copy of $\Q$ in $M$;
thus $M$ is not finitely generated, and $H^0$ is not of finite type.
\end{proof}

\begin{remark}[The transcription is not a group definition]
\label{A:monoid}
\status{Refuted literally; corrected by Theorem~\ref{A:band}}
The transcription of 15:5 reads ``entiers alg\'ebriques'' and
``$\nu\in\mathbf Z^+$''. Those elements form a monoid, not the
character group of a rigid category: $p$ is in it and $p^{-1}$ is not.
We have explicitly replaced that text by its group of fractions.
Both positive and negative integral weights occur. Weight zero does
not imply root of unity once denominators at $p$ are allowed.
Also the arrow transcribed as $i|_{\mu_2}:\mu_2\to H^0$ must land
in $H$; its composite with $H\to\mu_2$ is the identity, so it
cannot land in $H^0$.
\end{remark}

\begin{proposition}[Artin versus Artin--Tate rigidification]
\label{A:rigidification}
\status{Refuted as literally formulated in the prompt; corrected statement proved}
Let $F^0$ be the intrinsic fibre functor on
$\operatorname{Niv}^0(\mathcal M)^0$, the Artin motives. The gerbe
of fibre functors extending this $F^0$, with an identification on
that subcategory, has band $G^0$, not $H$. To obtain band $H$ one
must also fix a trivialization of the Tate line. Equivalently one
extends $F^0$ to the Artin--Tate subcategory by $F^0(\Q(1))=\Q$.
Its group is
\[
 G/H=D(\mu_\infty\oplus\Z p)
          =\widehat\Z\times\Gm.
\]
\end{proposition}
\begin{proof}
For any quotient $G\to A$, the automorphisms of a fibre functor
acting identically on the subcategory $\Rep(A)$ form its kernel.
The Artin characters are precisely $\mu_\infty$, giving kernel
$G^0$. Adding the Tate character adds $p\in\Gamma$ and gives $H$.
The two character subgroups intersect trivially, which proves the
last product decomposition. The manuscript itself writes
$F^0(\Q(1))=\Q$ on 15:3 even though $\Q(1)$ is not Artin; that
line explains the missing enlargement, rather than making it true
that Artin rigidification alone has band $H$.
\end{proof}

\begin{proposition}[Geometric base change and the universal Frobenius]
\label{A:P}
\status{Proved under the stated Tate hypothesis; identification of character modules unconditional}
There is a canonical isomorphism
\[
 \Gamma/\mu_\infty\xrightarrow{\sim}
 W(p^\infty):=\varinjlim_{n\mid n'}
   \bigl(W(p^n),\lambda\mapsto\lambda^{n'/n}\bigr).
\]
Thus $G^0=P$, the Weil-number protorus of Langlands--Rapoport and
Milne, and $H^0=P_0$, its weight-zero factor. The base-change
functor $\mathcal M_p\to\mathcal M_{\overline{\F}_p}$ induces
the inclusion $P\hookrightarrow G$ on bands. The inclusion
$\Gamma\hookrightarrow\overline\Q^\times$ is Galois equivariant
and hence defines a canonical $\mathfrak f\in G(\Q)$.
\end{proposition}
\begin{proof}
Send $\lambda\in W(p)$ to its germ at $n=1$. Its germ vanishes
exactly when some positive power is $1$, giving kernel
$\mu_\infty$. Given a representative $a\in W(p^n)$, choose an
$n$th root $b$. All its conjugates have absolute value $p^{w(a)/2}$,
and a sufficiently large power of $p$ makes it integral. Thus
$b\in W(p)$ and its germ is that of $a$. Choices differ by a root
of unity. This proves the assertion and all equivariance claims.
Over $\F_{p^n}$ Frobenius becomes the $n$th power, so the same
map is exactly the character map of geometric base change.
Evaluation $\lambda\mapsto\lambda$ is a Galois-equivariant
homomorphism into $\overline\Q^\times$, which is the definition
of a $\Q$-point of $D(\Gamma)$.
\end{proof}

\begin{proposition}[Centres and crystalline endomorphisms]
\label{A:endomorphisms}
\status{Proved under the stated Tate hypothesis}
For a simple $N\in\mathcal M_q$ the algebra $A_N=\End(N)$ is a
central division algebra over $E_N=\Q[\mathfrak f_N]$, and
\[
 A_N\otimes\Q_p\simeq
 \End_{F\text{-}\mathrm{Isoc}(\F_q)}(T_pN).
\]
The right side means endomorphisms commuting with crystalline
Frobenius, including its semilinear structure. It does not mean all
endomorphisms of the underlying vector space. If $d_N^2=[A_N:E_N]$,
then $\operatorname{rank}(N)=d_N[E_N:\Q]$.
\end{proposition}
\begin{proof}
Simplicity in a semisimple category gives a division algebra.
Full faithfulness at any $\ell\ne p$ identifies its scalar extension
with the centralizer of semisimple Frobenius. The double-centralizer
theorem says that the centre is $\Q_\ell[\mathfrak f_N]$;
faithful scalar extension gives $E_N$ over $\Q$. The crystalline
identity is the $p$-adic full-faithfulness statement of
\citet[Proposition 1.17]{AMilne1994}, applied to $N$ with itself.
A maximal commutative subalgebra of the centralizer has dimension
$\operatorname{rank}(N)$ over the realization field. A maximal
commutative subalgebra of $A_N$ has dimension $d_N[E_N:\Q]$,
which proves the rank formula (ibid., Proposition 2.8).
For abelian varieties the full-faithfulness input is Tate's
endomorphism theorem, so this portion is unconditional. The proof
of Theorem~\ref{A:band} explains precisely the additional tensor
argument needed to pass from abelian varieties to arbitrary pure
motives under the full Tate hypothesis.
\end{proof}

\subsection{A.2. Crystalline realization and the local classes}

\begin{theorem}[The warning about $\Z$ and $\widehat\Z$]
\label{A:crystalline}
\status{Proved; Lead A2 is false in its proposed nonneutral interpretation over $\F_p$}
Over $\F_p$, an isocrystal is a finite-dimensional $\Q_p$-vector
space with an invertible \emph{linear} operator. Forgetting the
operator is an exact faithful tensor functor to $\Q_p$-vector
spaces. Thus $T_p$, after this forgetting, really is a fibre functor
over $\Q_p$. Its Frobenius action generally does not extend
continuously from $\Z$ to $\widehat\Z$. Over
$\overline{\F}_p$ the realization instead takes values in the
$\Q_p$-linear category $\mathcal R_p$ of isocrystals on
$B=W(\overline{\F}_p)[1/p]$. Its fibre-functor gerbe has band
$D(\Q)$ and class $\beta_p$ characterized by
\[
 \operatorname{inv}_p\bigl(s_*(\beta_p)\bigr)=s\pmod\Z,
 \qquad s\in\Q=X^*(D(\Q)).
\]
It is not neutral over $\Q_p$.
\end{theorem}
\begin{proof}
The Witt Frobenius on $W(\F_p)=\Z_p$ is the identity, which proves
the first assertion. The Tate object $\Q(-1)$ has $F=p$.
A continuous representation of the compact group $\widehat\Z$
on a finite-dimensional $\Q_p$-space has compact image. The
image of the subgroup $\Z$ acting by $p^n$ is unbounded. This
already proves the warning on 15:3, including at $p=2$.
More generally an invertible operator extends continuously to
$\widehat\Z$ precisely when the closure of the cyclic subgroup
it generates is compact; equivalently it preserves a lattice, or
all its eigenvalues are $p$-adic units. Necessity follows by taking
the bounded orbit of a lattice and its inverse; sufficiency follows
from the profiniteness of $\mathrm{GL}_r(\Z_p)$.

Over $\overline{\F}_p$, the Dieudonn\'e--Manin simple isocrystal
of slope $a/b$, $(a,b)=1$, has basis $e_0,\ldots,e_{b-1}$ and
$F(e_i)=e_{i+1}$, $F(e_{b-1})=p^a e_0$, semilinearly for the
arithmetic Witt Frobenius. Its endomorphism algebra is the cyclic
division algebra generated by the unramified degree-$b$ extension
and a cyclic shift $u$ satisfying
\[
 u x u^{-1}=\sigma_{\rm ar}^{-1}(x),\qquad u^b=p^a.
\]
For example the unramified scalar $x$ acts on $e_i$ by
$\sigma_{\rm ar}^i(x)$, which verifies the relation directly.
Its invariant is $-a/b$ with our normalization. The slope simples
and their tensor products give the geometric characters $\Q$;
the opposite-algebra computation in our conventions gives the
displayed gerbe invariant $a/b$. This is
the concrete form of \citet[Theorem 2.9, Proposition 2.14, and
Example 3.28]{AMilne1994}. Slope $1/2$ has nonsplit quaternion
endomorphisms, so no $\Q_p$-fibre functor can exist. This concerns
the geometric-base-field category, and does not contradict the
first paragraph for $\F_p$.
\end{proof}

\begin{proposition}[The Artin torsor and the connecting map]
\label{A:eta}
\status{Proved under the stated Tate hypothesis; the Artin calculation is unconditional}
For the Artin--Tate rigidification of Proposition~\ref{A:rigidification},
let $\eta_p=\operatorname{Isom}^{\otimes}(T_p|_{\mathrm{Artin}},F^0_p)$,
with the right action on the source; extend it by the trivial Tate
factor. Under
\[
 H^1(\Q_p,G/H)=H^1(\Q_p,\widehat\Z)
       =\Hom_{\rm cts}(G_{\Q_p},\widehat\Z),
\]
it is the unramified quotient. Arithmetic local Frobenius maps to
$+1$ when the canonical motivic Frobenius is the generator $+1$.
The opposite isomorphism torsor has value $-1$ and gives
$\xi_p=-\partial(\eta_p^{\mathrm{op}})$ instead.
The class of the rigidified gerbe is
\[
 \xi_p=\partial(\eta_p),\qquad
 \partial:H^1(\Q_p,G/H)\longrightarrow H^2(\Q_p,H).
\]
\end{proposition}
\begin{proof}
For $\Spec\F_{p^n}$ the crystalline algebra is
$W(\F_{p^n})[1/p]$, the unramified degree-$n$ extension of
$\Q_p$; the standard Artin algebra is $\Q_p^n$. The torsor
of tensor identifications is therefore the torsor of its embeddings.
These torsors are compatible as $n$ varies. More precisely, over
$\overline\Q_p$ write $U(a)=(\tau(a))_\tau$ for the splitting
of the unramified algebra. If $F$ is its Witt Frobenius, then
$\sigma_{\rm ar}U=UF$. Thus the right-torsor cocycle
$U^{-1}\sigma_{\rm ar}U$ is $F$, the crystalline realization of
$\mathfrak f$, and has value $+1$. This computes the quotient
with its orientation, not just up to sign. The $\Gm$ factor contributes no
$H^1$ by Hilbert 90.

The full $G$-gerbe has the crystalline fibre functor as a
$\Q_p$-object. Imposing an identification with $F^0$ is the
lifting gerbe of this $G/H$-torsor for
$1\to H\to G\to G/H\to1$. The class of a lifting gerbe is
its connecting class: choose local lifts of the torsor's transition
functions; an arrow $\phi_\sigma:{}^\sigma(T_p,U)\to(T_p,U)$
satisfies $U\phi_\sigma={}^\sigma U$. Its failure to satisfy
the cocycle identity is the connecting $H$-valued $2$-cocycle.
This gives the formula with no replacement
of crystalline cohomology by a continuous Galois representation.
\end{proof}

\begin{proposition}[The real obstruction, with the Serre argument repaired]
\label{A:real}
\status{Proved for abelian motives, and hence under the stated hypothesis for all motives}
The full category over $\F_p$ has no $\R$-valued fibre functor.
At infinity the rigidified gerbe has class
\[
 \xi_\infty=(0,1/2)\in
 H^2(\R,H^0)\times H^2(\R,\mu_2)=0\times\Z/2\Z.
\]
The nonzero class is represented by the real Weil extension,
whose lift $j$ of complex conjugation satisfies $j^2=-1$.
\end{proposition}
\begin{proof}
First, Honda--Tate applied to the Weil number $\sqrt p$ gives a
simple abelian surface over $\F_p$. Its endomorphism division
algebra has centre $\Q(\sqrt p)$, local invariant $1/2$ at
each real embedding, and zero at finite places. At $p$ the formula
is $(1/2)[\Q_p(\sqrt p):\Q_p]=1\equiv0$, so this works also
at $p=2$. The algebra has degree two over its centre and the
motive has rank four. Over $\R$ its algebra is
$\mathbb H\times\mathbb H$. A faithful module for that algebra
has real dimension at least eight, whereas a fibre functor must
give dimension four. This is a contradiction.

This repairs the preliminary reading's use of a supersingular
elliptic curve over $\F_p$. Its \emph{geometric} quaternion
endomorphisms need not be defined over $\F_p$; one cannot simply
put that quaternion algebra into its endomorphism ring there.
The original elliptic proof becomes valid after extension to
$\F_{p^2}$, where Honda--Tate supplies a curve with scalar
Frobenius $p$ and all geometric endomorphisms. An $\R$-fibre
functor over $\F_p$ would extend to this category: its restriction
to all Artin motives is standard since
$\Hom(\Gal(\C/\R),\widehat\Z)=0$, and the resulting split
Artin algebra $\R^2$ supplies either extension-of-constants
idempotent. Its fibre on this elliptic curve would be two-dimensional
and admit the definite quaternion algebra, again impossible.

The odd-line argument on 15:11 is also valid after specifying
polarizability. A neutral real category with this band would have
the real character $\sqrt p$, a rank-one object of odd weight.
Its polarization must be alternating and nondegenerate, impossible
on a one-dimensional real space. Equivalently the Tate character
on the connected band is twice a real character: the disconnected
kernel $H=H^0\times\mu_2$ records exactly this forbidden odd
line. Disconnectedness \emph{by itself} is not an obstruction for
an arbitrary polarizable Artin category; the connected band, the
Tate character, and odd-weight polarization are essential to this
first argument.

The computation $H^2(\R,H^0)=0$ is proved in
Theorem~\ref{A:hasse} below. Since any real fibre functor has the
standard Artin restriction and any one-dimensional real Tate line
can be trivialized, nonneutrality of the full category implies
nonneutrality of the rigidified gerbe. Finally
$H^2(\R,\mu_2)\simeq\Br(\R)[2]=\Z/2\Z$, giving the stated
class. The cocycle $j^2=-1$ is precisely its image under
$\mu_2\hookrightarrow\Gm$; on weight $w$ the realization has
semilinear square $(-1)^w$.
\end{proof}

\subsection{A.3. The protorus Hasse calculation}

\begin{lemma}[Compact real tori satisfy the degree-two Hasse principle]
\label{A:compacthasse}
\status{Proved unconditionally}
If $T/\Q$ is a torus compact over $\R$, with character lattice
$L$, then $\Sha^2(\Q,T)=0$. In fact
\[
 H^2(\Q,T)\xrightarrow{\sim}\bigoplus_v H^2(\Q_v,T).
\]
\end{lemma}
\begin{proof}
Complex conjugation $c$ acts by $-1$ on $L$. Choose a finite
Galois splitting field $E$. Continuity and torsion-freeness give
$H^1(\Q,L)=H^1(\Gal(E/\Q),L)$, since a continuous homomorphism
from the profinite group $G_E$ to the discrete lattice has finite,
and hence zero, image. In the finite image acting on $L$, $c=-1$
is central and $L^{\langle c\rangle}=0$. Inflation--restriction
therefore injects $H^1(\Q,L)$ into $H^1(\R,L)$.
In particular $\Sha^1(\Q,L)=0$.

Apply \citet[Theorem I.4.20(a)]{AMilneADT}, with the set of all
places and $M=L$, so that the Cartier dual is $T$. Its perfect
pairing of $\Sha^1(\Q,L)$ with $\Sha^2(\Q,T)$ proves
injectivity. The last portion of the exact sequence in
Theorem I.4.20(b) has cokernel dual to $L^{G_\Q}$.
This group is zero because already $L^c=0$. Hence the localization
map is surjective as well. Every application is to a finite-rank
lattice; no inverse-limit assertion is hidden in this step.
\end{proof}

\begin{theorem}[The global class really is locally determined]
\label{A:hasse}
\status{Proved unconditionally for the band in Theorem~\ref{A:band}}
For $H^0=D(M)$ one has
\[
 \Sha^2(\Q,H^0)=0,\qquad H^2(\R,H^0)=0.
\]
For a finite place $v$, there is a canonical identification
\[
 H^2(\Q_v,H^0)=\Hom\bigl(M^{G_{\Q_v}},\Q/\Z\bigr).
\]
More precisely, for any cofinal countable increasing sequence of
Galois-stable lattices $L_n\subset M$, and $T_n=D(L_n)$,
\[
 H^2(\Q,H^0)=\varprojlim_n H^2(\Q,T_n)
            =\varprojlim_n\bigoplus_v H^2(\Q_v,T_n).
\]
These equalities are conclusions of the proof, not definitions of
cohomology. The assertion in A.3 that $\Sha^2$ of an arbitrary
protorus is simply a limit of finite-level $\Sha^2$ is not a
valid general rule.
\end{theorem}
\begin{proof}
The field $\overline\Q$ is countable, so such a countable exhaustion
exists. The inclusions of lattices dualize to surjective torus
transition maps. Fix $n$ and a splitting field $E_n$ for $T_n$.
Restriction to $E_n$ kills $H^1(\Q,T_n)$ by Hilbert 90; restriction
followed by corestriction is multiplication by $d=[E_n:\Q]$.
Thus $d$ kills that group. Since $M$ is uniquely divisible, it
contains $(1/d)L_n$. For large $m$ the inclusion $L_n\subset L_m$
factors through $(1/d)L_n$, and the resulting torus map factors
through multiplication by $d$ on a copy of $T_n$. Consequently
$H^1(\Q,T_m)\to H^1(\Q,T_n)$ is zero. The inverse system of
these $H^1$ groups is pro-zero, so its inverse limit and its first
derived inverse limit vanish.

The exact sequence of \citet[Proposition 2.8]{AMilne2003} is
\[
 0\longrightarrow\varprojlim{}^1 H^1(\Q,T_n)
 \longrightarrow H^2(\Q,H^0)
 \longrightarrow\varprojlim H^2(\Q,T_n)\longrightarrow0.
\]
The left term has now actually been killed. A class locally zero
maps to a locally zero class at every finite level and is zero
there by Lemma~\ref{A:compacthasse}; hence it is zero globally.
The displayed full inverse-limit description follows from the same
lemma.

The identical argument works over $\Q_v$, choosing a finite
splitting extension and its degree. Local Tate duality
\citep[Theorem I.2.4]{AMilneADT} identifies
$H^2(\Q_v,T_n)$ with $\Hom(L_n^{G_{\Q_v}},\Q/\Z)$.
Invariants commute with this increasing union because each invariant
element already belongs to a finite-rank stable lattice. Applying
$\Hom(-,\Q/\Z)$ proves the finite-place formula. At infinity
$T_n$ is a product of norm-one circles, whose degree-two group is
zero: in cyclic cohomology the invariant circle quotient by the
norm is trivial, or use the norm sequence
$1\to U(1)\to\operatorname{Res}_{\C/\R}\Gm\to\Gm\to1$
and $H^1(\R,\Gm)=0$. The real $H^1$ inverse system is killed
by the same division argument, giving $H^2(\R,H^0)=0$.
\end{proof}

\begin{theorem}[Explicit global and local classes]
\label{A:class}
\status{Proved unconditionally for the constructed gerbe; motivic identification under the stated hypotheses}
Write $v_p(p)=1$, and for $m\in M$ choose a Weil-unit representative
$u$. The number $v_p(m):=v_p(u)$ is independent of that choice.
There is exactly one class
\[
 \xi_0=(a_p,b_{p,\infty})\in H^2(\Q,H^0)\times H^2(\Q,\mu_2)
\]
with local components
\begin{center}
\begin{tabular}{lll}
\toprule
Place & $H^0$ component & $\mu_2$ component\\
\midrule
$\ell\ne p,\infty$ & $0$ & $0$\\
$p$ & $m\mapsto v_p(m)\pmod\Z$, $m\in M^{G_{\Q_p}}$ & $1/2$\\
$\infty$ & $0$ & $1/2$\\
\bottomrule
\end{tabular}
\end{center}
Here $b_{p,\infty}$ is the quaternion Brauer class ramified
exactly at $p$ and infinity, viewed in $H^2(\Q,\mu_2)$ by Kummer.
The compact component $a_p$ has infinite order. The finite local
components determine both factors of $\xi_0$.
\end{theorem}
\begin{proof}
At each finite level $T_n$, restrict the displayed $p$-adic
functional to $L_n^{G_{\Q_p}}$ and put zero at every other place.
Lemma~\ref{A:compacthasse} gives a unique global class with these
local values. Uniqueness makes the finite-level classes compatible.
The derived-limit calculation in Theorem~\ref{A:hasse} lifts them
to exactly one $a_p$. This supplies existence without a motivic
conjecture. The Brauer reciprocity exact sequence supplies exactly
one $2$-torsion class with invariants $1/2$ at $p,\infty$ and zero
elsewhere; Kummer identifies it with $H^2(\Q,\mu_2)$ since
$\Pic(\Q)=0$. If all finite invariants are prescribed, their
sum fixes the infinite invariant, proving the last assertion.

For comparison with the motivic gerbe, geometric base change gives
the isocrystal slope morphism
\[
 D(\Q)\longrightarrow P_{\Q_p},\qquad
 X^*(P)\longrightarrow\Q,\quad [\lambda]\longmapsto v_p(\lambda).
\]
On the splitting $X^*(P)=M\oplus\Z[\sqrt p]$, its character map
is $(m,n)\mapsto v_p(m)+n/2$. Theorem~\ref{A:crystalline}
therefore gives the compact local functional $v_p$ and the
weight-line invariant $1/2$. At $\ell\ne p$, geometric
\'etale cohomology is the fibre functor extending the standard
Artin--Tate functor, giving neutrality. The real computation is
Proposition~\ref{A:real}. The resulting global class is the unique
one just constructed.

One may check the connecting-map normalization directly. For
$m\in M^{G_{\Q_p}}$, some $u^n=a$ is fixed by local Galois:
$\sigma u/u$ ranges in a finite group of roots of unity. Lifting
the unramified $\widehat\Z$ cocycle through $D(\Gamma^0)$
produces the cyclic class with invariant
$\eta_p(\sigma_{\rm ar})\operatorname{ord}_p(a)/n=v_p(u)$.
This is the same class as $\partial\eta_p$ in
Proposition~\ref{A:eta}, not a second unrelated local prescription.

For infinite order, use $u=a/\bar a$ from the last paragraph of
Theorem~\ref{A:band} and choose the $p$-adic embedding with
$v_p(u)=1$. The class $[u]/n$ lies in $M^{G_{\Q_p}}$ and has
valuation $1/n$. The functional thus takes the values $1/n$
for all $n\geq1$; no nonzero integer annihilates it.
\end{proof}

\begin{proposition}[An explicit arithmetic description of the character module]
\label{A:valuations}
\status{Proved unconditionally}
For finite Galois CM fields $K\subset\overline\Q$, let $X_K$
be their sets of primes above $p$. Then
\[
 M\simeq\varinjlim_K
 \{f:X_K\to\Q:\ f(cw)=-f(w)\},
\]
where transition maps pull back functions along maps of primes,
and $[u]$ maps to $w\mapsto\operatorname{ord}_w(u)/
\operatorname{ord}_w(p)$. In particular the local groups and the
functional in Theorem~\ref{A:class} are specified by explicit
rational permutation modules.
\end{proposition}
\begin{proof}
A weight-zero Weil number generates a CM field unless it is real,
in which case it is $\pm1$. On its Galois closure, complex
conjugation is central because $c(u)=u^{-1}$ holds for every
conjugate. Its valuations are supported at $p$ and are
anti-invariant under $c$. If all are zero, it is an algebraic unit
all of whose conjugates have absolute value one, hence a root of
unity. This proves injectivity.

Conversely scale a rational anti-invariant function by an integer
so that the corresponding exponents, including ramification
indices, are integral. They define a fractional ideal $I$ supported
at $p$ with $cI=I^{-1}$. A positive power $I^h=(b)$ is principal.
Then $b/\bar b$ has exactly twice these scaled valuations and
all its complex absolute values are one. Taking the requisite
root, and then its class modulo roots of unity, recovers the given
function. This operation is compatible with extension of fields
and proves surjectivity.
\end{proof}

\subsection{A.4. The comparison with the pseudomotivic gerbe}

\begin{theorem}[The correct comparison, and the strict-band counterexample]
\label{A:comparison}
\status{Strict equivalence with band $H$ or $G$ refuted; corrected equivalence proved under the stated hypotheses}
Let $\mathscr H$ be the Artin--Tate rigidified gerbe of class
$\xi_0$. Its pushout along $H\hookrightarrow P=G^0$ is the
pseudomotivic $P$-gerbe of
\citet[Theorem 3.29 and Theorem 3.31]{AMilne1994}, with the
inverse-limit and adelic corrections of
\citet[Theorem 6.3 and Definition 6.6]{AMilne2003}.
Equivalently it is the gerbe obtained by rigidifying only the
Artin fibre functor $F^0$, and under the motivic hypotheses it
is the gerbe of motives over $\overline{\F}_p$.

There is no band-preserving equivalence of the $H$-gerbe itself
or of the full $G$-gerbe with that $P$-gerbe. Their bands have
component groups $\mu_2$, $\widehat\Z$, and $1$, respectively.
The morphisms that do exist are pushouts
\[
 \mathscr H\longrightarrow\mathscr P
       \longrightarrow\mathscr G,
 \qquad H\hookrightarrow P\hookrightarrow G.
\]
\end{theorem}
\begin{proof}
The character isomorphism of Proposition~\ref{A:P} identifies the
connected band with $P$, not merely with an abstract isogenous
group. Under $H=H^0\times\mu_2\hookrightarrow
P=H^0\times\Gm$, the local compact class is unchanged and the
$\mu_2$ class becomes the real/quaternion weight class.
At $p$ this is exactly the pushout of $\mathcal R_p^\vee$ by
its slope morphism; at infinity it is exactly the pushout of the
real realization gerbe by the weight cocharacter. At the other
places it has the standard \(\ell\)-adic object. These are the
local morphisms in Milne's definition, with the above signs.
The uniqueness theorem for this \emph{adelically supplied} system
is \citet[Theorem 6.3]{AMilne2003}; its proof uses
Propositions 3.5, 3.6, 3.10, and 4.3, rather than incorrectly
discarding all derived limits. Hence this is an equivalence of
systems with local data, not merely a match of names.

The base-change tensor functor has band inclusion $P\to G$ by
Proposition~\ref{A:P}. Artin descent identifies its image gerbe
with the rigidification of the Artin functor: on any finite
extension the Artin algebra becomes split and its idempotents
recover the extension-of-constants category; compatible
idempotents give the algebraic closure. Adding a trivialization
of the Tate line takes the further fibre of $P\to\Gm$, which
is $\mathscr H$. This proves both identifications of gerbes.

Any equivalence of banded gerbes induces an isomorphism of their
bands. An isomorphism preserves the identity component and its
component quotient. The three quotients just computed are
nonisomorphic. This is a complete counterexample to the literal
claim in Lead A1 and A.4, for every prime $p$, including $2$.
It also proves why the corrected pushout statement is the
appropriate replacement.
\end{proof}

\begin{proposition}[Dictionary between gerbes and groupoids]
\label{A:dictionary}
\status{Proved}
Let $\mathscr P$ be a $P$-gerbe over $\Q$ and choose
$x\in\mathscr P(\overline\Q)$. Its groupoid over
$\Spec(\overline\Q)\times_{\Spec\Q}\Spec(\overline\Q)$
is the sheaf of isomorphisms between the two pullbacks of $x$.
Its diagonal kernel is $P_{\overline\Q}$. For each place $v$,
a realization functor induces a morphism of such groupoids;
its map on kernels is $1\to P_{\Q_\ell}$ for $\ell\ne p$,
$D(\Q)\to P_{\Q_p}$ at $p$, and $\Gm\to P_\R$ at infinity.
Conversely fpqc descent for this groupoid recovers $\mathscr P$.
Changing $x$ changes the groupoid by an isomorphism whose
conjugacy ambiguity is exactly the change of the chosen
identification of objects.
\end{proposition}
\begin{proof}
Composition of isomorphisms supplies multiplication and the
identity of $x$ supplies the unit; the cocycle condition is
associativity of composition. Restriction to the diagonal gives
$\Aut(x)=P_{\overline\Q}$. An object on an arbitrary base is
locally an object isomorphic to a pullback of $x$, together with
its descent isomorphisms. Those isomorphisms are exactly a torsor
for the groupoid satisfying its cocycle identity. This constructs
inverse functors and proves the converse, rather than identifying
a groupoid with its kernel alone. A tensor realization gives a
contravariant morphism on fibre-functor gerbes. The kernel maps
are read from weight and slope, as calculated above, so all
morphisms in the dictionary type-check. This is the dictionary
used in \citet[\S3, Examples 3.26--3.28]{AMilne1994} and
\citet[\S6]{AMilne2003}.
\end{proof}

\subsection{A.5. Unconditional existence and the remaining enrichment issues}

\begin{theorem}[Unconditional $H$-gerbe and its automorphisms]
\label{A:unconditional}
\status{Proved unconditionally as a statement about banded gerbes}
The class in Theorem~\ref{A:class} defines an $H$-gerbe over
$\Q$ with exactly the prescribed local classes. It is unique up
to $H$-equivalence. No nonzero class can be added without changing
some local class. The group of isomorphism classes of its
band-preserving self-equivalences is
\[
 H^1(\Q,H)=\Q^\times/\Q^{\times2}\times
  \frac{\prod_n T_n(\Q)}
  {\{(a_n f_n(a_{n+1})^{-1})_n:(a_n)_n\in\prod_nT_n(\Q)\}},
\]
for the explicit lattice exhaustion in Theorem~\ref{A:hasse};
$f_n:T_{n+1}\to T_n$ is its transition. The group of automorphisms
of any one such equivalence is $H(\Q)=\{\pm1\}$.
\end{theorem}
\begin{proof}
Banded commutative gerbes are classified by $H^2$; the descent
cocycle construction of Proposition~\ref{A:eta} realizes any such
class. Existence and uniqueness were proved in
Theorem~\ref{A:class}. Twisting a band-preserving equivalence by
an $H$-torsor gives every other one, with $H(\Q)$ as its
$2$-automorphism group. This is the elementary torsor form of
the gerbe classification in \citet[\S5]{AMilne2003}.

The same pro-zero argument used for $H^1(\Q,T_n)$ gives
$H^1(\Q,H^0)=\varprojlim{}^1 T_n(\Q)$. The displayed quotient
is the explicit countable-system definition of that derived
limit, and Kummer gives the other factor. Finally any rational
point of $H^0$ evaluates on $m/n$ in a fixed number field
(the stabilizer of $m/n$ is the stabilizer of $m$). Its value on
$m$ is therefore divisible by every positive integer inside that
number field's multiplicative group. The intersection of all
power subgroups there is $1$, by valuations, Dirichlet units,
and the finite root-of-unity group. Thus $H^0(\Q)=1$.
\end{proof}

\begin{proposition}[The bare category and the category with local data]
\label{A:enrichment}
\status{Bare classification proved; classification of arbitrary framed local data by the bare class refuted}
Fix the band diagram $H\hookrightarrow G$, the grading $i$,
and the augmentation $\varepsilon$, with the Artin--Tate
rigidification specified in Proposition~\ref{A:rigidification}.
Then $\xi_0$ determines the underlying Tannakian category with
these structures up to compatible equivalence, and conversely.
It does not in general determine a family of chosen local
neutralizations extending the fixed rigidification.

Write $S=\{\ell:\ell\ne p,\ \ell\text{ finite}\}$. The precise
replacement for that stronger classification is relative
cohomology, defined by the shifted cone of
$R\Gamma(\Q,H)\to\prod_{\ell\in S}R\Gamma(\Q_\ell,H)$.
For a fixed locally trivializable class its fibre is a torsor under
\[
 \coker\left[H^1(\Q,H)\longrightarrow
              \prod_{\ell\in S}H^1(\Q_\ell,H)\right].
\]
The same statement with a restricted adelic product applies when
an integral compatibility condition is imposed.
\end{proposition}
\begin{proof}
The gerbe of fibre functors and the category of its finite-rank
representations recover each other by tensor descent. The band
diagram recovers the specified subcategory and its rigidification;
$i$ recovers the grading and $\varepsilon$ the Tate augmentation.
Equivalence classes of $H$-gerbes are exactly $H^2(\Q,H)$, so
these constructions prove both directions of the bare assertion.
They do not choose objects in any local gerbe.

Once an object of a neutral local gerbe has been chosen, every
other object is obtained by twisting by an $H$-torsor, uniquely
up to its isomorphism class in $H^1(\Q_\ell,H)$. A global
band-preserving equivalence changes this family by the
localizations of a global $H$-torsor. This proves the displayed
cokernel description directly. Equivalently a relative
$2$-cocycle is a global $2$-cocycle together with a local
$1$-cochain whose coboundary is its localization; quotienting
by global and local coboundaries gives exactly that description.

Here is a fully explicit obstruction already at $p=2$. Twist a
fixed neutralization at $3$ by $3\in\Q_3^\times/\Q_3^{\times2}$,
in the $\mu_2$ factor, and leave every other odd-prime
neutralization unchanged. This changes no local $H^2$ class.
If a global equivalence identified the two families, its
$\mu_2$-torsor would be represented by
$a\in\{3,-3,6,-6\}$ modulo rational squares: its odd
valuations are prescribed, and only the valuation at $2$ is
unrestricted. The classes $3$ and $-3$ are nonsquares at $5$;
$6$ is a nonsquare at $7$; and $-6$ is a nonsquare at $13$.
None has the prescribed trivial localizations. The compact
factor cannot repair this, since $H=H^0\times\mu_2$ and the
projection to the second factor is functorial. These residue
checks are certified in the accompanying program. The change
has finite support, so the counterexample also survives the
usual almost-everywhere integral requirement.

This refutes the stronger reading in A.3(c), rather than the
carefully separated assertion actually transcribed on 15:7.
That page first says ``la classe de $\mathcal H$ comme
$H$-gerbe'' for $\mathcal M,\mathrm{grad},\mathrm{aug},F_0$,
then introduces $H^2(S\bmod T,H)$ after including the $F_\ell$.
If the chosen identifications with $F^0$ are forgotten, one
must further quotient by their local automorphism actions; this
is a different moduli problem, and the framed counterexample
does not assert otherwise. For the $P$-gerbe with the coherent
adelic object and the two local realization morphisms, the
specific uniqueness theorem is \citet[Theorem 6.3]{AMilne2003},
with $\Q^\times$ as the residual $2$-automorphism group.
That theorem includes more than the calculation $\Sha^2=0$.
\end{proof}

\begin{proposition}[Forgetting the local comparison framings]
\label{A:unframed-adelic}
\status{Proved for coherent adelic systems with the prescribed Artin restrictions}
Fix the category and global Artin--Tate rigidification of
Proposition~\ref{A:enrichment}. Suppose that the prime-to-$p$
realizations form one object over the finite adeles away from $p$,
and that this adelic object's Artin restriction is isomorphic to
$F^0\otimes\mathbb A_f^p$, not just separately at each place. At $p$
fix the Artin restriction of the ordinary crystalline $\Q_p$
fibre functor, namely the unramified Artin torsor computed in
Proposition~\ref{A:eta}. If the local comparison isomorphisms with
these restrictions are \emph{not} part of the data, all such
coherent systems on the fixed gerbe are equivalent. One may also
include the prescribed real realization morphism.

Thus the underlying class suffices for the equivalence class of
this unframed coherent system. This assertion and the framed
counterexample in Proposition~\ref{A:enrichment} concern different
specified data. The equivalence itself is not asserted to be
unique as an equivalence preserving the global $H$-rigidification.
\end{proposition}
\begin{proof}
Put $A=G/P=\widehat\Z$. Over any one of the finite coefficient
fields, the isomorphism torsor between two full $G$-fibre functors
has class in $H^1(k,G)$. Their Artin restrictions are isomorphic,
so its image in $H^1(k,A)$ is zero. Exactness for
$1\to P\to G\to A\to1$ therefore lifts that class to
$H^1(k,P)$. A chosen comparison of the Artin restrictions effects
this lift; the result after forgetting the comparison does not
depend on that choice. The argument also works over the
prime-to-$p$ adelic coefficient ring: the stipulated coherent
Artin comparison supplies a reduction of the adelic isomorphism
torsor to $P$.

For fixed-band real realization morphisms into the $G$-gerbe,
their difference is a $G$-torsor by the same descent argument for
the sheaf of isomorphisms of morphisms. It lifts to a $P$-torsor
because $H^1(\R,\widehat\Z)=0$: a continuous homomorphism
from the group of order two to the torsion-free group
$\widehat\Z$ is trivial. If no real morphism is specified, choose
the zero local difference there. The resulting local differences
give a class in $H^1(\mathbb A,P)$, with exactly the adelic
meaning used by \citet[\S6.2]{AMilne2003}.

Propositions 3.6 and 3.10 of that source prove
\[
 H^1(\Q,P)\xrightarrow{\sim}H^1(\mathbb A,P).
\]
This is the specific global-to-adelic theorem, not a consequence
of the bare degree-two Hasse principle. It provides a global
$P$-torsor whose localizations give the differences. Moreover
\[
 H^1(\Q,H)\longrightarrow H^1(\Q,P)
\]
is surjective: in $H=H^0\times\mu_2\to P=H^0\times\Gm$
the first factor is the identity and $H^1(\Q,\Gm)=0$.
Lift the global torsor along this map and twist the global
$H$-equivalence by it. Its underlying $G$-equivalence then
identifies every unframed local realization and still preserves
the global rigidification. This proves existence in the required
direction, including the nontrivial Artin torsor at $p$.

There remains at least the ambiguity from
$H^1(\Q,\mu_2)=\Q^\times/\Q^{\times2}$: its image in
$H^1(\Q,P)$ is zero, while it can change the global
$H$-equivalence. Hence uniqueness of its isomorphism class is not
part of this assertion. For the $P$-system itself, the sharper
uniqueness and residual $2$-automorphisms are exactly those of
Milne's Theorem 6.3, already stated above.
\end{proof}

\begin{proposition}[What the Lefschetz category actually supplies]
\label{A:lefschetz}
\status{Proved unconditionally; equality of its band with $P$ is refuted}
Let $\operatorname{LMot}(\overline{\F}_p)$ be the category generated
by abelian varieties, with correspondences generated by divisor
classes and with the ordinary Tannakian commutativity constraint.
It is a semisimple Tannakian category. Its band is
\[
 L=\mathop{\prod}_{\Pi}^{\Gm} L_\Pi,\qquad
 X^*(L_\Pi)=\Z^{\Pi}/
 \{f:c f=f,\ \textstyle\sum_{\pi\in\Pi}f(\pi)=0\},
\]
where $\Pi$ runs through the geometric simple isogeny classes,
identified with Galois orbits of integral weight-one Weil germs.
The fibre product identifies every multiplier
$[\pi]+[c\pi]$. In particular $L$ is commutative, but it is
not the Weil-number protorus $P$ in general. Compatible powers of Frobenius generate
a subgroup $P\hookrightarrow L$. Prime-to-$p$ realization and
crystalline realization, the weight and Tate morphisms, and the
real obstruction are available unconditionally. The claimed
identification of the entire character group with Weil numbers,
full Frobenius-centralizer faithfulness for arbitrary tensor
objects, and the resulting formula for $H^0$ cannot simply be
transferred to this smaller-correspondence category.
\end{proposition}
\begin{proof}
The Tannakian assertion is \citet[Theorem 1.7]{AMilne1999};
Proposition 1.8 identifies the band on the category generated by
an abelian variety $A$ with its Lefschetz group, and Corollary 1.9
passes to the full category. Over the algebraic closure of a finite
field all abelian varieties have many endomorphisms. The displayed
character module is exactly Lemma 4.2, and its fibre product is
Theorem 4.3 there. These are unconditional applications of
Honda--Tate and the divisor theorem, not the full Tate conjecture.
The natural map on characters sends a formal generator indexed by
$\pi$ to its Weil germ; its kernel consists of the additional
relations among actual Frobenius numbers. This defines $P\to L$.
Every Weil germ is generated by weight-one integral germs and
Tate twists, by the root argument in Theorem~\ref{A:band}; hence
the character map is surjective and the group map is injective.
The following example proves that its kernel on characters need
not vanish.
\end{proof}

\begin{example}[An explicit extra Frobenius relation]
\label{A:lefschetz-counterexample}
\status{Refuted: $L=P$ for all finite-field Lefschetz motives}
Take $p=5$, $\alpha=1+2i$, and
$\beta=(3+\sqrt{-11})/2$. They are Frobenius numbers of the
elliptic curves
\[
 E_\alpha:y^2=x^3+x,\qquad
 E_\beta:y^2=x^3+4x+2
 \quad\text{over }\F_5,
\]
which have four and three rational points, respectively. Choose
$\lambda^2=\alpha\beta$. This is an integral weight-one
$5$-Weil number, with irreducible polynomial
\[
 x^8-6x^6+15x^4-150x^2+625.
\]
Honda--Tate supplies a simple abelian fourfold $B/\F_5$ with this
Frobenius number. In
\[
 h^1(B)^{\otimes2}\otimes h^1(E_\alpha)^\vee
                      \otimes h^1(E_\beta)^\vee
\]
there is a geometric eigenline on which Frobenius is $1$, but
the Lefschetz group acts nontrivially. Thus its character is zero
on $P$ and nonzero on $L$.
\end{example}
\begin{proof}
Both quadratic numbers have norm $5$; taking a square root of
their product preserves the required absolute value and
integrality. Elimination gives the displayed polynomial. Exact
factorization over $\Q$ and both finite-field point counts are
recorded by \texttt{check\_lefschetz\_counterexample.py}.
There is also a conceptual degree check: in
$K=\Q(i,\sqrt{-11})$ the valuations of $\alpha\beta$ at the
four split primes over $5$ are $(2,1,1,0)$, so it is not a square.
Let $L_1=K(\lambda)$, of degree eight over $\Q$. At the two
$K$-primes where $v_5(\alpha\beta)=1$, the extension $L_1/K$
is ramified quadratic; the local degree times slope is
$2(1/2)=1$. At the other primes the slopes are $0$ or $1$.
All Honda--Tate invariants over $\F_5$ therefore vanish, giving
division index one and dimension $[L_1:\Q]/2=4$.
The four conjugates of $\alpha\beta$ have distinct valuation vectors; consequently
no quotient of two of them is a root of unity. All powers
therefore still generate $K$. Over $\F_{25}$ its slopes are
$1,1/2,1/2,0$. Honda--Tate gives division index $2$ and dimension
$2[K:\Q]/2=4$. This is still dimension four, so no multiplicity appears on
base change to $\F_{25}$. Over every further even extension the
same slopes and centre persist, so $B$ remains simple geometrically.
Its geometric centre is $K$, distinct from the two elliptic
centres, and its Lefschetz group is the norm-similitude torus
of $K/K^+$. These conclusions also verify every hypothesis of the
fibre-product description in Proposition~\ref{A:lefschetz}.

Over an algebraic closure of the coefficient field, choose the
eigencharacters labelled by $\lambda,\alpha,\beta$. The tensor
character $2[\lambda]-[\alpha]-[\beta]$ evaluates on Frobenius
as $\lambda^2/(\alpha\beta)=1$. For the geometric band we
use Frobenius squared over $\F_{25}$, where it centralizes all
geometric endomorphisms; its value on the character is still $1$.
In the $B$-factor of $L$,
vary one embedding by $t$ and its conjugate by $t^{-1}$,
leave the other conjugate pair fixed, and leave the two elliptic
factors fixed. The common multiplier is $1$, so this is a
cocharacter of $L$. It evaluates the tensor character as $t^2$.
Thus the character is nonzero in $X^*(L)$, as asserted.
\end{proof}

\begin{proposition}[The unconditional gerbes used in Kisin's paper]
\label{A:kisin}
\status{Known construction; the signed dictionary proved and independently checked}
Kisin's quasi-motivic gerbe has a band $Q$ defined by a universal
condition on real and $p$-adic cocharacters. Its pushout by
$Q\twoheadrightarrow P$ is the pseudomotivic gerbe, with the
same $P$ and the same local realization morphisms as above.
This construction is unconditional. It supplies a Tannakian
category of representations of that gerbe; identifying that
category with all numerical motives is a separate claim with the
hypotheses of A.1. The quasi-motivic band $Q$ itself is not a
replacement name for $G$, $H$, or $P$.
\end{proposition}
\begin{proof}
For a finite Galois field $L/\Q$, Kisin defines
\[
 Q^L=(\operatorname{Res}_{L(\infty)/\Q}\Gm
       \times\operatorname{Res}_{L(p)/\Q}\Gm)/\Gm.
\]
Lemma 3.1.4 gives its universal cocharacter condition; the
transition maps of 3.1.5 include their local-degree factors.
At the other finite places the source gerbe is trivial, at
infinity it is the real fundamental extension, and at $p$ it is
the inverse limit of the local fundamental extensions, of band
$D(\Q)$; these are 3.1.6. Theorem 3.1.9 gives existence and
uniqueness, with the specified almost-everywhere integral
condition of 3.1.8 and compatibility of the Galois band action.

For CM fields, 3.5.1--3.5.5 define the quotient $P$ by Weil
characters, with the finite-level integrality condition on the
local norms. That condition gives finite-rank lattices; it
exhausts the Weil-germ group after the CM field is enlarged.
Indeed a fixed rational valuation denominator is cleared by
enlarging the local degree, and its weight remains unchanged.
Lemma 3.5.3 sends both universal cocharacters to the weight and
slope cocharacters of $P$. The definition in 3.5.5 then takes
the pushout of the gerbe. The signs require an explicit
local-source identification. Kisin's equation (3.5.2), with
$|p|_p=p^{-1}$, pairs his $p$-adic cocharacter with a Weil
character by the negative normalized valuation. In 3.3.3 he
chooses a lift $d_\sigma$ of arithmetic Witt Frobenius whose
$n$th power at the unramified degree-$n$ level is $p^{-1}$.
Thus that source fundamental class has invariant $-1/n$ in
our normalization. Our isocrystal gerbe has invariant $+1/n$.
The inversion $[-1]:D(\Q)\to D(\Q)$ identifies the two source
gerbes: it takes $+1/n$ to $-1/n$. Precomposing Kisin's slope
morphism with this inversion changes $-s$ into $+s$.
Equivalently, pushing his negative fundamental class through
his negative valuation cocharacter gives our positive local
class. The target band $P$ is unchanged. At infinity inversion
does not change the order-two Brauer class. Consequently this
is the groupoid/gerbe operation in
Proposition~\ref{A:dictionary}, with the specified source-band
identification, not a silent identification of opposite signs.
Kisin notes explicitly in 3.1.5, footnote 13, that his auxiliary
map has the inverse convention to Langlands--Rapoport's; the
band and local-class dictionary must include this choice.
All these source statements were opened in
\citet[\S\S3.1 and 3.5]{AKisin2017}.

Kisin's application to points of abelian-type Shimura varieties
is Theorem 0.3, expanded as Theorem 4.6.7, with $p>2$ in this
paper and a conjugation ambiguity in the adelic action. It
does not assert the full Tate conjecture for arbitrary varieties
and it does not justify dropping that hypothesis in A.1.
\end{proof}

\begin{remark}[The literature date changes the abelian scope]
\label{A:2026-abelian}
\status{Known from the opened October 2026 theorem; the implication is stated explicitly}
The dated primary preprint of \citet[Theorem I]{AvanHoftenSempliner2026}
proves the rational Hodge conjecture for CM abelian varieties.
Its \S1.1.2--1.1.3 applies Milne's Hodge-to-Tate implication
to all abelian varieties over finite fields. With these cited
results the full-correspondence category of abelian motives,
and all its tensor powers, has the Weil band $P$
unconditionally. This does not change the Lefschetz category:
the latter was defined by restricting its morphisms to those
generated by divisors, and Example~\ref{A:lefschetz-counterexample}
still proves that its band can be larger. Nor does this theorem
settle the Tate conjecture for arbitrary smooth projective
varieties. The preprint status and exact version are retained in
the bibliography rather than represented as an earlier theorem.
\end{remark}

\IfFileExists{sections/part_b.tex}{\section{Absolute cohomology: the finite-field identification and its arithmetic limits}\label{sec:part-b}
\providecommand{\BW}{\mathrm W}
\providecommand{\BM}{\mathrm M}
\providecommand{\Bcts}{\mathrm{cont}}
\providecommand{\Bf}{\mathrm f}
\providecommand{\Bc}{\mathrm c}
\providecommand{\Bcorank}{\operatorname{corank}}
\providecommand{\BPic}{\operatorname{Pic}}
\providecommand{\BCl}{\operatorname{Cl}}
\providecommand{\Br}{\operatorname{Br}}
\providecommand{\BCH}{\operatorname{CH}}
\providecommand{\BRHom}{\operatorname{RHom}}
\providecommand{\BRGamma}{\operatorname{R\Gamma}}

\subsection{Conventions and the audit of the inherited claims}
The symbols in this section distinguish five different theories.  We write
\(H^i_{\BM}(X,\Q(n))\) for rational motivic cohomology in the sense of
Voevodsky (equivalently, the appropriate Adams eigenspace of rational
\(K\)-theory for the regular schemes considered below),
\(H^i_{\et}(X,\Z(n))\) for etale motivic hypercohomology,
\(H^i_{\Bcts}(X,\Q_\ell(n))\) for continuous \(\ell\)-adic cohomology,
\(H^i_{\BW}\) for the specified Weil--etale theory, and
\(H^1_{\Bf}\) for the Bloch--Kato Selmer subgroup.  A subscript \(\Bf\)
in degrees other than one requires the choice of a Selmer complex; it is
not defined just by taking subgroups of ordinary cohomology.  Tate twists
are normalized so that \(\Q(1)\) has weight \(-2\).

For a finite field \(k=\F_q\), arithmetic Frobenius is
\(\phi_q:x\mapsto x^q\), and \(W_k=\langle\phi_q\rangle\cong\Z\).
The generator \(e_q\in H^1(W_k,\Z)\) sends \(\phi_q\) to one.
When using geometric Frobenius on geometric cohomology we use its inverse
consistently; changing that convention changes the sign of \(e_q\), not
the multiplication-by-degree formula below.  In number-ring formulae,
\(\bar X\) means the Artin--Verdier compactification.  Omitting this bar
is a substantive change of theory, even when \(X\) is proper over
\(\Spec\Z\).

The attached transcriptions, rather than scans, are the witnesses used
here.  Thus an asserted correction of a printed formula below is a
mathematical correction of the transcription; it does not constitute a
new paleographic reading.  The reading's R5 is correct about the rational
class in degree three only after the theory is specified.  Its integral
tables omit two nonzero groups of order two.  R6 correctly detects the
failure of unrestricted arithmetic \(\ell\)-adic comparison, but its
positive finite-field formulation needs the additional rational-versus-
numerical-equivalence condition explained below.

\subsection{B.1: the actual finite-field theorem and the missing hypothesis}

\begin{proposition}[Restriction on the arithmetic class]\label{prop:B-e}
\status{Proved}
For \(k_m=\F_{q^m}\), restriction sends
\[
 e_q\longmapsto m e_{q^m}.
\]
For the point, Weil--etale motivic cohomology is
\[
H^i_{\BW}(\Spec\F_q,\Z(n))=
\begin{cases}
\Z,&n=0,\ i=0,1,\\
\Z/(q^{|n|}-1),&n\ne0,\ i=1,\\
0,&\text{otherwise},
\end{cases}\tag{B.1}
\]
where for negative twists the etale convention is
\(\Z(n)=\Q/\Z'(n)[-1]\), with the characteristic-primary part omitted.
After tensoring with \(\Q\), only the two groups with \(n=0\) survive.
\end{proposition}
\begin{proof}
The subgroup \(W_{k_m}\subset W_k\) is generated by \(\phi_q^m\).
The restricted homomorphism has value \(m\) on this generator.  For
\(n=0\), use the length-one resolution of the trivial module over
\(\Z[t,t^{-1}]\), giving \([\Z\xrightarrow{0}\Z]\).
For \(n>0\), the finite-field motivic calculation gives the cyclic group
\(K_{2n-1}(\F_q)=\Z/(q^n-1)\) in degree one and no other group.
For \(n<0\), Frobenius-minus-one is surjective on
\(\Q/\Z'(n)\), and its kernel has order \(q^{|n|}-1\).
The coefficient complex is shifted by minus one.  Weil and etale
cohomology agree for torsion complexes, so the asserted negative-twist
calculation follows.  These comparisons are the torsion and rational
parts of Geisser's Theorem~7.1~\cite{BGeisser2004}.
\end{proof}

The expression \(\Z/(1-q^n)\) in 15:173 is therefore an ordinary cyclic
group only for nonnegative \(n\).  For \(n<0\), its meaningful version is
\(\Z[1/p]/(1-q^n)\), canonically isomorphic to
\(\Z/(q^{|n|}-1)\).  Similarly, the table of continuous cohomology of
\(\widehat\Z\) with coefficients \(\Z_\ell(n)\) presupposes
\(\ell\ne p\); the characteristic-primary coefficient theory is not
the same cyclotomic etale sheaf.

\begin{theorem}[The precise comparison gate]\label{thm:B-Geisser-gate}
\status{Proved}
Let \(X/\F_q\) be smooth and projective of dimension \(d\).  In the
notation of Geisser's Conjectures~8.1--8.3, let \(K(X,n)\) denote the
integral Weil--etale to continuous-cohomology comparison, \(L(X,n)\)
finite generation, and \(T(X,n)\) the conjunction of the Tate cycle-class
conjecture in codimension \(n\), semisimplicity at eigenvalue one, and
equality of rational and numerical equivalence in that codimension.
Then
\[
 K(X,n)+K(X,d-n)\ \Longrightarrow\ L(X,n)
 \ \Longrightarrow\ K(X,n)\ \Longrightarrow\ T(X,n).
 \tag{B.2}
\]
Assuming \(T(Y,r)\) for every smooth projective variety over finite fields
and every \(r\), all the comparisons \(K(X,n)\) and finite-generation
statements \(L(X,n)\) follow.  Moreover, independently of these
conjectures,
\[
 H^i_{\BW}(X,\Q(n))\simeq
 H^i_{\BM}(X,\Q(n))\oplus H^{i-1}_{\BM}(X,\Q(n)).
 \tag{B.3}
\]
Cup product with \(e_q\) acts on these summands by
\( (a,b)\mapsto(0,a)\).
\end{theorem}
\begin{proof}
Equation (B.2) is precisely Geisser's Theorem~8.4, with the three
clauses of his Conjecture~8.3 retained; its last clause must not be
dropped.  Equation (B.3) and the cup-product matrix are Theorem~7.1(c),
deduced there from the computation of the derived pushforward between
the Weil and etale topoi.  Rational etale motivic cohomology agrees with
rational motivic cohomology here.  Thus these are statements about the
actual specified complexes, rather than dimension counts for numerical
cycles~\cite[Theorems~7.1 and~8.4]{BGeisser2004}.
\end{proof}

This identifies an exact defect in the proposed B.1 premise.  The Tate
conjecture and semisimplicity alone supply the expected
\emph{continuous}\(\ell\)-adic groups; they do not, by the cited theorem,
supply the missing equality of rational and numerical equivalence.
Conversely the requested integral comparison already implies that
equality.  No counterexample to that equality is asserted here.  The
unresolved implication must therefore remain part of B.1's status.

For clarity, let \(V^r=H^r_{\et}(\bar X,\Q_\ell(n))\), with
\(\ell\ne p\).  Hochschild--Serre gives
\[
0\longrightarrow V^{i-1}_{\phi_q}
\longrightarrow H^i_{\Bcts}(X,\Q_\ell(n))
\longrightarrow (V^i)^{\phi_q}\longrightarrow0.\tag{B.4}
\]
Purity puts a possible eigenvalue one only in \(r=2n\).
Consequently the groups can be nonzero only in degrees \(2n,2n+1\).
Semisimplicity at one makes the natural invariants-to-coinvariants map
an isomorphism.  Under the Tate conjecture its source is the image of
codimension-\(n\) cycles.  If rational and numerical equivalence agree,
(B.2)--(B.3) give exactly the manuscript's cycle spaces in these two
degrees and the isomorphism by cup product with \(e_q\).  The odd
identification requires the chosen normalization of the arithmetic
generator; this is the valid content of ``pas canoniquement.''

\begin{proposition}[A contradiction already at the point]\label{prop:B-continuity}
\status{Refuted; corrected statement proved}
There is no cohomology theory which simultaneously has the finite-field
Weil--etale value \(H^1(\F_q,\Q)=\Q e_q\), the restriction in
Proposition~\ref{prop:B-e}, continuity under the filtered union of finite
extensions, and vanishing of higher cohomology over \(\overline{\F}_q\).
\end{proposition}
\begin{proof}
All transition maps in the directed system of rational degree-one
groups are multiplication by a positive integer, hence isomorphisms.
Its colimit is a one-dimensional rational space, not zero.  Integrally
the same colimit of copies of \(\Z\), with maps multiplication by the
extension degrees, is \(\Q\).  The contradiction uses no conjecture.
Ordinary etale cohomology of the geometric point instead has no higher
\(\Q_\ell\)-cohomology: it is a different theory and does not repair this
incompatible list of requirements.
\end{proof}

The unqualified demand that actual Weil--etale cohomology satisfy
\emph{every} assertion of 12:60--73 is thus refuted.  The two assertions
must be separated in any replacement.  Nor can the degree-one group be
identified with \(\operatorname{Ext}^1\) in the semisimple category of
pure numerical motives: that extension group vanishes for the unit,
whereas \(H^1_{\BW}(\F_q,\Q)=\Q\).  A theory of extensions must specify
a larger, nonsemisimple arithmetic category or a different derived
construction.

\paragraph{What is unconditionally identified, and what uniqueness can mean.}
For curves all twists satisfy the comparison: this is Geisser's
Theorem~9.5 (dimension at most one belongs to its stated range).
Theorem~9.3 gives the codimension-one comparison whenever the divisor
Tate conjecture holds, in particular for abelian varieties.  For a smooth
projective geometrically connected curve \(C/\F_q\), the two relevant
integral rows can also be checked without appealing to the general
conjectural mechanism:
\[
\begin{array}{c|rrrr}
 &i=0&i=1&i=2&i=3\\ \hline
H^i_{\BW}(C,\Z)&\Z&\Z&\BPic^0(C)^D&(\F_q^\times)^D\\
H^i_{\BW}(C,\Z(1))&0&\F_q^\times&\BPic(C)&\Z .
\end{array}\tag{B.8}
\]
Here \(\Z(1)=\Gm[-1]\), and all unlisted degrees vanish.  This is
Lichtenbaum's Theorem~6.1 after that shift.  The rational rows are
therefore \(\Q\) in degrees zero and one for twist zero and in degrees
two and three for twist one~\cite[Theorem~6.1]{BLichtenbaum2005}.

\begin{theorem}[All twists for motives of abelian type, using the current theorem]
\label{thm:B-abelian-type}
\status{Proved from the cited current theorem}
Let \(X/\F_q\) be smooth projective, and suppose that its rational
Chow motive belongs to the thick rigid category generated by abelian
varieties and Artin motives.  Then, for every \(i,n\) and every prime
\(\ell\),
\[
 H^i_{\BW}(X,\Z(n))\otimes\Z_\ell
 \xrightarrow{\sim}H^i_{\Bcts}(X,\Z_\ell(n)),
\]
and all the integral Weil--etale groups are finitely generated.
For \(\ell=p\), the right side is the logarithmic Hodge--Witt
comparison, rather than a cyclotomic etale sheaf in characteristic
\(p\).  In particular,
\[
 H^i_{\BW}(X,\Q(n))=
 \begin{cases}
 CH^n(X)_\Q,&i=2n,\ 2n+1,\\
 0,&\text{otherwise},
 \end{cases}
\]
where cup product by the normalized \(e_q\) identifies the two
copies.  Thus the previously known divisor and curve cases extend to
all twists on all abelian varieties, and on all varieties of abelian
type.  This statement does not assert that every Chow motive is of
abelian type.
\end{theorem}
\begin{proof}
The external input dated 2026 is van Hoften--Sempliner's
Theorem~I, the Hodge conjecture for CM abelian varieties.  Milne's
Theorem~7.1 then proves the Tate conjecture for every abelian variety
over every finite field~\cite[Theorem~I and \S1.1.3]{AvanHoftenSempliner2026}
\cite[Theorem~7.1]{AMilne1999}.  We use that theorem as an opened
literature result, not as a proof established in this research.

Tate classes pass to direct summands by applying the algebraic
projector to a representing cycle.  They pass to Tate twists and
finite sums, while a tensor product of abelian motives is again the
motive of an abelian variety, up to the finite-extension and twist
operations already allowed.  For a variety over a finite extension,
Galois averaging of an algebraic cycle, followed by pushforward,
descends the required invariant cycle with rational coefficients.
Consequently the Tate statement passes to the whole stated category.
Its Frobenius realizations are semisimple: this holds on \(H^1\)
of an abelian variety, passes to its exterior algebra and all tensor
constructions, and is preserved by projectors and finite extension.
This is also Kahn's Lemma~1.9.

There is an essential extra step beyond Tate surjectivity.  Kahn's
Theorem~1.4 places these Chow motives in the Kimura finite-dimensional
category, and Proposition~1.5 makes the kernel of the numerical
endomorphism map nilpotent.  His Theorem~1.10 therefore identifies
rational and numerical equivalence here.  Corollary~2.2 gives the
vanishing of positive rational \(K\)-groups.  With these hypotheses
verified, his Theorem~3.6 proves both the \(\ell\ne p\) and
the logarithmic \(p\)-adic comparisons; Corollary~3.8(c) supplies
finite generation.  The rational formula follows either from his
Lemma~3.7 and Corollary~3.8(d), or by substituting rational
motivic vanishing into Geisser's splitting (B.3).
All these are exact statements for the same category, so no general
Beilinson assumption has entered~\cite[Theorem~1.10, Corollary~2.2,
Theorem~3.6, Corollary~3.8]{BKahn2003}.
\end{proof}

\begin{proposition}[The remaining universal assertion is exactly a phantom-motive gate]
\label{prop:B-phantom-gate}
\status{Proved equivalence; the equivalent universal assertion remains unresolved}
Assume Tate surjectivity and semisimplicity at eigenvalue one for
every smooth projective variety over \(\F_q\), in every
codimension.  Then the following universal statements are equivalent:
\begin{enumerate}
\item The integral Weil--etale comparison holds for all varieties,
degrees and twists, and all these integral groups are finitely generated.
\item Rational equivalence equals numerical equivalence on all
rational Chow groups.
\item There is no nonzero rational Chow motive \(P\) whose image
\(\bar P\) in numerical motives is zero.
\item Every rational Chow motive belongs to the category generated
by abelian and Artin motives.
\end{enumerate}
In particular, the missing assertion in B.1 can be expressed without
an unspecified regulator or unspecified comparison map: it is the
vanishing of numerical Chow phantoms.  The Tate and semisimplicity
hypotheses alone have not been shown to imply that vanishing here.
\end{proposition}
\begin{proof}
The precise equivalences, including integral comparison and finite
generation, are Kahn's Theorems~5.1 and~5.2, especially clauses
(iv bis), (iv ter), (v), (vi ter), and (vi quater).  Here is the
categorical obstruction in explicit form.  Strong Tate, which is
the stated combination of cycle surjectivity and semisimplicity,
implies that every numerical motive is of abelian type.  Thus, for
a Chow motive \(M\), there is an abelian-type Chow motive \(N\)
and inverse numerical isomorphisms
\(\bar f:\bar M\to\bar N\),
\(\bar g:\bar N\to\bar M\).
Fullness of the numerical quotient lifts them to \(f,g\).
The element \(u=1_N-fg\) is numerically zero.  Nilpotence on
the abelian-type motive \(N\) gives \(u^r=0\) for some \(r\).
Replace \(f\) by
\[
 f'=(1_N+u+\cdots+u^{r-1})f.
\]
Then \(f'g=1_N\), so \(gf'\) is an idempotent on \(M\).
The rational Chow category splits idempotents, hence
\[
 M\simeq N\oplus P,\qquad \bar P=0.
\]
This constructs the exact leftover motive.  If numerical phantoms
vanish, every \(M\) is abelian-type and
Theorem~\ref{thm:B-abelian-type} applies.  Conversely universal
rational--numerical equality makes the numerical quotient faithful:
if \(\bar P=0\), then \(1_P=0\), so \(P=0\).
The equivalence of that equality with universal comparison is also
Geisser's Theorem~8.4 after the Tate and semisimplicity hypotheses
have been fixed.  Kahn's argument identifies precisely which
additional step the two stated hypotheses do not supply
\cite[Theorems~5.1--5.2 and the final proof on pp.~999--1000]{BKahn2003}.
\end{proof}

\begin{proposition}[A smaller index set for the remaining Chow comparison]
\label{prop:B-diagonal-Chow-gate}
\status{Proved reduction; the restricted cycle-injectivity assertion remains unresolved}
Keep the universal Tate and eigenvalue-one semisimplicity hypotheses
of Proposition~\ref{prop:B-phantom-gate}.  Fix \(\ell\ne p\).
The universal comparison in that proposition is equivalent to
injectivity of the cycle map only for the pairs
\[
 \dim X=2m-1,\qquad
 \operatorname{CH}^{m}(X)_\Q\otimes\Q_\ell
       \longrightarrow H^{2m}(\bar X,\Q_\ell(m)),
 \qquad m\ge2.                                      \tag{B.9a}
\]
Here all smooth projective varieties in the displayed dimensions,
over all finite extensions under consideration, are quantified.
There is also a dimension-preserving version: to obtain
Beilinson's injectivity in every codimension for varieties of
dimension at most \(D\), it suffices to prove it for
\[
 3\le d\le D,\qquad 2\le i\le\lceil d/2\rceil.
 \tag{B.9b}
\]
The first still undecided ordinary-Chow assertion in either range
is codimension two on threefolds.  Under the fixed Tate hypotheses,
that assertion is equivalent to
\[
 H^4(F_0\mathbb P^4_{\overline{\F}_q},\Q_\ell(3))^
       {\Gal(\overline{\F}_q/\F_q)}=0,              \tag{B.9c}
\]
where \(H^a(F_0Y,\Q_\ell(b))\) means the colimit of
continuous \(\ell\)-adic cohomology over dense open subsets
of \(Y\).  It is not defined by tensoring the Milnor
\(K\)-group of the function field with \(\Q_\ell\).
\end{proposition}
\begin{proof}
Under universal Tate, numerical and \(\ell\)-homological
equivalence coincide, so the Beilinson clause can be stated as
cycle-class injectivity.  Codimensions zero and one are already
injective.  Balkan--Schreieder's Corollary~1.3 proves (B.9b)
under exactly the fixed Tate and semisimplicity hypotheses.
The mechanism matters: their Theorem~7.1(2) translates the
low-codimension assertions into vanishing of
\[
 H^{2j}(F_0\mathbb P^{d+1}_{\overline{\F}_q},
        \Q_\ell(j+1))^{\Gal(\overline{\F}_q/\F_q)}
\]
for \(j\le\lceil d/2\rceil\).  For larger \(j\),
\(2j>d+1\); a cofinal system of affine opens has no
cohomology above its dimension \(d+1\).  Thus those remaining
groups vanish automatically.  The same theorem then gives
Beilinson injectivity in the upper half of the codimensions.
This is the full implication used from the source, and does not
assume a hard Lefschetz theorem for Chow groups
\cite[Corollary~1.3, Theorem~7.1 and its proof in \S7.4]
{BBalkanSchreieder2026}.

To reduce (B.9b) to (B.9a), let \(X\) have dimension \(d\)
and let \(2\le i\le\lceil d/2\rceil\).  Set
\[
 t=d-2i+1\ge0,\qquad Y=X\times\mathbb P^t,
 \qquad m=i+t=d-i+1.
\]
Then \(\dim Y=2m-1\).  If \(h\) is the hyperplane
class of \(\mathbb P^t\), the map
\[
 z\longmapsto \operatorname{pr}_X^*z\cdot
                   \operatorname{pr}_{\mathbb P^t}^*h^t
\]
is a split injection from \(\operatorname{CH}^i(X)_\Q\)
to \(\operatorname{CH}^m(Y)_\Q\): its composite with
\((\operatorname{pr}_X)_*\) is the identity, since
\(\deg h^t=1\).  A homologically trivial class maps to
a homologically trivial class, so (B.9a) makes it zero and
pushforward makes the original class zero.  We have proved
(B.9b) for each \(D\), and hence universal Beilinson.
Proposition~\ref{prop:B-phantom-gate} finishes the equivalence.
The converse is immediate.  Finally, (B.9c) is precisely
Theorem~7.1(2) with \(d=3,i=2\), as stated explicitly
in Remark~7.2 of the source.
\end{proof}

This reduces the index set of Chow assertions; it does not claim
that proving the first threefold case proves all higher cases.
The projective-space construction can increase dimension: a
middle-codimension correspondence on a surface squared has
\((d,i)=(4,2)\), and is sent to \((5,3)\), not to a
threefold.  In particular, the already known ordinary Chow-group
equalities on an individual surface do not remove its possible
Chow-phantom endomorphism.  Geisser's Frobenius argument for
Parshin's conjecture uses the motivic endomorphism category;
replacing its universal rational=numerical hypothesis by equality
only on that surface is not justified
\cite[Propositions~2.3, 3.1--3.2 and Theorem~3.3]{BGeisser1998}.

\begin{proposition}[Rational varieties as a restricted test class]
\label{prop:B-rational-blowup-gate}
\status{Proved test-class reduction; the same Chow obstruction remains}
Under the hypotheses of Proposition~\ref{prop:B-diagonal-Chow-gate},
the remaining universal comparison is equivalent to cycle-class
injectivity only for
\[
 \operatorname{CH}^{2m}(Y)_\Q\otimes\Q_\ell
       \longrightarrow H^{4m}(\bar Y,\Q_\ell(2m)),
 \qquad \dim Y=4m-1,\quad m\ge2,
\]
where \(Y\) ranges over smooth projective rational varieties over
all the finite constant-field extensions under consideration.
It even suffices to test the varieties obtained by blowing up
\(\mathbb P^{4m-1}\) along a smooth projective centre of
dimension \(2m-1\).
In particular, the codimension-two threefold problem embeds as a
split summand of the codimension-four problem on rational
sevenfolds.  This does not reduce it to rational threefolds.
\end{proposition}
\begin{proof}
Let \(X\) be smooth projective of dimension \(d=2m-1\).
After a finite extension of the ground field, there is a closed
immersion \(X\hookrightarrow\mathbb P^{2d+1}\).
Here is the required projection argument, including the field
qualification.  Start with a very ample embedding in
\(\mathbb P^N\); if \(N\le2d+1\), use a linear inclusion.
Otherwise the closure of the union of secant lines has dimension
at most \(2d+1\), since the distinct-pair parameter space has
dimension \(2d\).  The union of embedded tangent lines has
dimension at most \(2d\), by the projectivized tangent bundle
of dimension \(2d-1\).  Thus a general linear centre of
dimension \(N-2d-2\) avoids both loci.  Such centres form a
nonempty open subset of a Grassmannian, so one exists over a
finite extension of the finite ground field.  Projection from
it separates geometric points and tangent vectors.  It is
therefore universally injective and unramified; properness then
makes it a closed immersion
\cite[Lemma~41.7.2]{BStacksBlowupTests}.

Write \(Y=\operatorname{Bl}_X\mathbb P^{4m-1}\), let
\(j:E\hookrightarrow Y\) be the exceptional divisor and
\(\pi:E\to X\) its projection.  The centre is regular of
codimension \(2m\).  Consequently \(Y\) is smooth
projective and rational, and
\[
 E=\operatorname{Proj}_X\operatorname{Sym}
                  (\mathcal I/\mathcal I^2),\qquad
 \xi=c_1(\mathcal O_E(1)),\qquad
 \mathcal O_Y(E)|_E=\mathcal O_E(-1).
\]
Smoothness may also be checked in the usual local charts of the
blow-up of a smooth coordinate subspace.  The exceptional-bundle
description and the sign of its normal bundle are
\cite[Lemmas~31.33.4 and 42.59.11]{BStacksBlowupTests}.
Define explicit maps on Chow groups by
\[
 \begin{aligned}
 I_m:\operatorname{CH}^m(X)&\longrightarrow
                   \operatorname{CH}^{2m}(Y),
 &I_m(z)&=j_*(\xi^{m-1}\pi^*z),\\
 R_m:\operatorname{CH}^{2m}(Y)&\longrightarrow
                   \operatorname{CH}^{m}(X),
 &R_m(w)&=-\pi_*(\xi^{m-1}j^*w).
 \end{aligned}
\]
The Cartier-divisor self-intersection identity gives
\(j^*j_*u=-\xi u\).  The projective-bundle pushforward
identity gives \(\pi_*(\xi^{2m-1}\pi^*z)=z\), since
\(\pi\) has relative dimension \(2m-1\).  Hence
\[
 R_m I_m(z)=\pi_*(\xi^{2m-1}\pi^*z)=z.
\]
These identities hold integrally before tensoring
\cite[Lemmas~42.30.3 and 42.36.1]{BStacksBlowupTests}.
All maps commute with cycle classes.  A homologically trivial
class on \(X\) therefore maps to one on \(Y\); injectivity
for the latter and the displayed retraction kill the original
class.

Finally, for a finite constant extension \(k'/k\), restriction
and corestriction on Chow groups compose to multiplication by
\([k':k]\).  The latter is invertible after tensoring with
\(\Q_\ell\), so the result over \(k'\) descends to
\(k\)~\cite[Lemma~42.15.2]{BStacksBlowupTests}.
Proposition~\ref{prop:B-diagonal-Chow-gate} now proves the
asserted universal implication.  Its converse is immediate.
\end{proof}

This is a restriction of the geometric test class, with explicit
split correspondences; it does not prove cycle injectivity on
that class.  The arbitrary original variety survives as the
blow-up centre.  Thus the substantive phantom obstruction has
been preserved, and this encoding is not counted as a new
resolution of an open gate.

\begin{proposition}[A surface-support criterion for the first Chow gate]
\label{prop:B-surface-supported-CH2}
\status{Proved; the unrestricted threefold gate remains unresolved}
Let \(X/k\) be smooth, projective and geometrically integral, where
\(k\) is finite, and put \(n=\dim X\geq2\).
Suppose that, after a finite extension of \(k\), there is a closed
subscheme \(S\subset X\), of dimension at most two, for which
\[
 [\eta_X]\in\operatorname{im}\bigl(
 CH_0(S_{k(X)})_{\Q}\longrightarrow
 CH_0(X_{k(X)})_{\Q}\bigr).
 \tag{B.S1}\label{eq:B-generic-surface-support}
\]
Then rational and numerical equivalence agree on codimension-two
cycles on \(X\).  In particular, for every \(\ell\ne\operatorname{char}k\),
\[
 CH^2(X)_{\Q}\otimes_{\Q}\Q_\ell\longrightarrow
 H^4(X_{\overline k},\Q_\ell(2))
 \quad\hbox{is injective}.
\]
Condition~\eqref{eq:B-generic-surface-support} is equivalent to
universal rational support on this fixed \(S\): the pushforward
\(CH_0(S_L)_{\Q}\to CH_0(X_L)_{\Q}\) is surjective for every
field extension \(L/k\).  Thus it is a condition over function
fields, not merely a condition on zero-cycles over the finite base
field.  For threefolds satisfying it, rational and numerical
equivalence agree in every ordinary Chow degree.
\end{proposition}
\begin{proof}
All assertions descend through a finite constant extension by
restriction and corestriction, so make the extension in the hypothesis.
Clearing denominators in~\eqref{eq:B-generic-surface-support},
spreading its finitely many cycles and rational functions over a
nonempty open of the first factor, and using Chow localization gives
\[
 N\Delta_X=\Theta_1+\Theta_2\quad\hbox{in }CH^n(X\times X),
 \qquad
 \operatorname{supp}\Theta_1\subset X\times S,
 \quad\operatorname{supp}\Theta_2\subset D\times X,
\]
where \(N>0\) and \(D\subset X\) is a divisor containing the
closed complement of that open.  This is the generic-point proof
of the Bloch--Srinivas decomposition, so the support implication
does not require an additional conjecture.  The corresponding
curve-support factorization is explicitly proved in
\citet[Lemma~5.1 and its proof, p.~841]{BFuVial2023}.
Conversely, after any extension \(L/k\), move a zero-cycle on the
smooth variety \(X_L\) away from \(D_L\).  The action of
\(\Theta_2\) is zero on that cycle, while the action of
\(\Theta_1\) lands in the image of \(CH_0(S_L)\).
Division by \(N\) proves universal rational support.  The other
implication follows by taking \(L=k(X)\) and the generic point.

Transpose the decomposition and write
\[
 N\Delta_X=\Gamma_1+\Gamma_2,
 \qquad
 \operatorname{supp}\Gamma_1\subset S\times X,
 \quad\operatorname{supp}\Gamma_2\subset X\times D.
 \tag{B.S2}\label{eq:B-transposed-surface-support}
\]
Replace the irreducible components of \(S\) and \(D\) by smooth
projective alterations, separately when necessary
\citep[Theorem~4.1 and Remark~4.2]{c_dejong1996}.  Denote their
maps to \(X\) by \(a:S'\to X\) and \(b:D'\to X\), with
\(\dim S'=s\leq2\) and \(\dim D'=n-1\) on each component.
Proper surjective pushforward on Chow groups is surjective with
rational coefficients: for an integral cycle downstairs, choose a
closed point of its generic fibre, take its closure upstairs, and
divide by the resulting finite generic degree.  Consequently the
two supported correspondences lift, component by component, to
\[
 \widetilde\Gamma_1\in CH^{s}(S'\times X)_{\Q},
 \qquad
 \widetilde\Gamma_2\in CH^{n-1}(X\times D')_{\Q}.
\]
Their actions on codimension-two Chow groups have the precise
factorizations
\[
\begin{aligned}
 (\Gamma_1)_*&=(\widetilde\Gamma_1)_*a^*:
 CH^2(X)_{\Q}\longrightarrow CH^2(S')_{\Q}
                 \longrightarrow CH^2(X)_{\Q},\\
 (\Gamma_2)_*&=b_*(\widetilde\Gamma_2)_*:
 CH^2(X)_{\Q}\longrightarrow CH^1(D')_{\Q}
                 \longrightarrow CH^2(X)_{\Q}.
\end{aligned}
\]
The pullback \(a^*\) exists as the refined pullback for a
morphism between smooth varieties; the dimension shifts above
also follow directly by composing the graph correspondences.
All these maps preserve numerical equivalence, by the projection
formula, and preserve homological equivalence.

Here are the elementary finite-field vanishing inputs.  On any
smooth projective geometrically integral variety \(Y\) over a
finite field, every degree-zero zero-cycle is torsion: choose a
smooth projective geometrically integral curve through its finite
support, using \citet[Corollary~3.4, with Theorem~3.3 and its
proof]{DPoonen2004Bertini}; its class on that curve belongs to the
finite group of rational points of the Jacobian.  Pushing forward
annihilates a positive multiple of the original cycle.  A finite
constant extension handles geometric components if needed.
Thus the numerically trivial part of \(CH^2(S')_{\Q}\) is zero:
it is zero by dimension when \(s<2\), and is the degree-zero
part of \(CH_0(S')_{\Q}\) when \(s=2\).  Numerically trivial
divisors on \(D'\) are algebraically trivial up to torsion; the
remaining Picard-zero points over a finite field are torsion too.
Hence \(CH^1(D')_{\Q,\mathrm{num}}=0\).  The divisor assertion
is also the exact input in the proof of Fu--Vial's Lemma~5.1;
here no hypothesis on \(H^3(X)\) is needed.

For numerically trivial \(z\in CH^2(X)_{\Q}\), both actions vanish.
Equation~\eqref{eq:B-transposed-surface-support}
gives \(Nz=0\), hence \(z=0\).  Homologically trivial cycles
are numerically trivial.  To obtain the stated injectivity even
after tensoring with \(\Q_\ell\), recall that numerical Chow
groups are finite dimensional: a full-rank matrix of rational
intersection numbers embeds its rank in the finite-dimensional
cohomological pairing.  Choose a rational basis for \(CH^2(X)_{\Q}\)
and complementary cycles giving an invertible rational intersection
matrix.  That same matrix remains invertible over \(\Q_\ell\),
so the realized basis remains linearly independent over \(\Q_\ell\).
For a threefold the remaining ordinary degrees are degree zero,
divisors, and zero-cycles, already covered by the preceding
elementary statements.
\end{proof}

\begin{corollary}[Geometrically uniruled threefolds]
\label{cor:B-uniruled-threefold}
\status{Proved ordinary Chow equality; no full motivic comparison inferred}
Every smooth projective geometrically uniruled threefold over a
finite field satisfies rational equals numerical equivalence in
every ordinary Chow degree.  In particular the first unresolved
codimension-two threefold test in
Proposition~\ref{prop:B-diagonal-Chow-gate} may be restricted to
threefolds which are not geometrically uniruled.
\end{corollary}
\begin{proof}
After a finite constant extension, geometric uniruledness gives a
dominant rational map \(\phi:T\times\mathbf P^1\dashrightarrow X\),
where \(T\) is an integral surface.  The two sides have the same
dimension, so the map is generically finite.  After another finite
constant extension choose \(a\in\mathbf P^1(k)\) for which
\(\phi\) is defined at the generic point of \(T\times\{a\}\).
Let \(S\) be the closure of its image; then \(\dim S\leq2\).

Over \(K=k(X)\), choose a closed point in the generic fibre of
the domain of \(\phi\).  Its residue field \(K'/K\) is finite,
and the corresponding point of \(T\) is generic: a proper closed
subset of \(T\), multiplied by \(\mathbf P^1\), has dimension
at most two and cannot dominate \(X\).  The resulting rational
map \(\mathbf P^1_{K'}\dashrightarrow X_{K'}\) extends to a
morphism because the source is a regular curve and the target is
proper.  On this curve the selected \(K'\)-point and \(a\) are
rationally equivalent.  Their images show that
\([\eta_X]_{K'}\) is represented by a zero-cycle on \(S_{K'}\).
Corestriction proves~\eqref{eq:B-generic-surface-support} over
\(K\), after division by \([K':K]\).  The proposition applies,
and finite-constant descent finishes the proof.
\end{proof}

The support hypothesis must not be weakened to
\(CH_0(X_{\overline{\F}_p})_{\Q}=\Q\): that statement already
holds for every smooth projective geometrically integral
finite-field variety by the curve argument above.  Nor does the
proposition establish Parshin's higher-Chow vanishing, comparison
for all twists of an individual such variety, or the higher odd
dimensional tests in Proposition~\ref{prop:B-diagonal-Chow-gate}.
The criterion is not subsumed by our abelian-type theorem: for an
arbitrary smooth projective surface \(S\), \(S\times\mathbf P^m\)
has universal support on \(S\times\{a\}\), while its surface
Chow summand is not in general known to be of abelian type.

\begin{proposition}[Two failures of interchange in a proposed symbol proof]
\label{prop:B-limit-diagnostics}
\status{Proved counterexamples to two limit arguments; the B1 gate is unchanged.}
The following two implications, which might otherwise be used to attack
Proposition~\ref{prop:B-diagonal-Chow-gate}, are false.
\begin{enumerate}
\item Put \(k=\overline{\F}_2\), \(L=k(t)\), and \(\ell=3\).
The natural map
\[
 \varinjlim_{U\subset\mathbf P^1_k}
 H^1_{\mathrm{cont}}(U,\Z_3(1))
 \longrightarrow H^1_{\mathrm{cont}}(\Spec L,\Z_3(1)),
\]
where \(U\) runs over dense opens, is injective and not surjective.
It remains nonsurjective after tensoring with \(\Q_3\), and an
explicit omitted class is invariant under \(\Gal(k/\F_2)\).
\item There is a torsion-free \(\Z\)-module \(M\) with an automorphism
\(F\), such that every element of \(M\otimes\Q\) has finite
\(F\)-orbit, but the rationalized \(3\)-adic completion has a nonzero
\(4\)-eigenvector.  Thus local finiteness before completion does not
exclude a non-root-of-unity eigenvalue after completion.
\end{enumerate}
\end{proposition}
\begin{proof}
For (1), the opens \(U=\mathbf A^1_k\setminus D\), with \(D\) finite,
are cofinal.  Their coordinate rings are localizations of the PID
\(k[t]\), so \(\operatorname{Pic}(U)=0\).  Kummer theory gives
\[
 H^1_{\mathrm{cont}}(U,\Z_3(1))
   =\bigoplus_{a\in D}\Z_3,
\]
with coordinates the valuations at \(t-a\).  Here \(k^\times\) is
\(3\)-divisible, and the transition maps on
\(H^0(U,\mu_{3^r})\) are surjective.  Consequently the
\(\varprojlim{}^1\) term in continuous cohomology vanishes.  These are
continuous coefficient systems, not discrete constant
\(\Z_3\)-sheaves
\cite[(3.1), Theorem~3.2 and (3.27), pp.~216--217, 222]{BJannsen1988}.
On the field, \(L^\times/k^\times\) is freely generated by the linear
polynomials \(t-a\).  Therefore the corresponding calculation gives
\[
\begin{split}
 H^1_{\mathrm{cont}}(\Spec L,\Z_3(1))
 &=\varprojlim_r L^\times/(L^\times)^{3^r}\\
 &=\left\{(v_a)\in\prod_{a\in k}\Z_3:
       \text{for every }r,\ v_a\in3^r\Z_3
       \text{ for all but finitely many }a\right\}.
\end{split}
\]
The open colimit is the finite-support subgroup
\(\bigoplus_{a\in k}\Z_3\), proving injectivity.

Enumerate the monic irreducible polynomials \(P_n\in\F_2[t]\),
\(n\geq1\), by degree and then by binary coefficient encoding.  Define
\[
 c_r=\prod_{1\leq n<r}P_n(t)^{3^n}
       \quad\text{in }L^\times/(L^\times)^{3^r}.
\]
Since \(c_{r+1}/c_r=P_r(t)^{3^r}\), these classes are compatible.
Their limit has valuation \(3^n\ne0\) at each geometric root of
\(P_n\), so it has infinite support.  Each \(c_r\) belongs to
\(\F_2(t)^\times\); the resulting Kummer class is Galois invariant.
No nonzero scalar in \(\Q_3\) makes its support finite.  This proves
all assertions in (1), including the rational one.

For (2), take the group of locally constant integer-valued functions
\[
 M=C_{\mathrm{lc}}(\Z_3,\Z),\qquad (Fh)(x)=h(x+1),
\]
where \(\Z\) is discrete.  Each function factors through some
\(\Z/3^s\), and a finite sum of rational multiples of such functions
does too.  Thus \(F^{3^s}h=h\) for some \(s\), and
\(Fh=4h\) in \(M\otimes\Q\) implies
\((4^{3^s}-1)h=0\), hence \(h=0\).

Choosing integer lifts on the finitely many clopen fibres gives
\[
 M/3^rM=C_{\mathrm{lc}}(\Z_3,\Z/3^r),\qquad
 \widehat M_3=C(\Z_3,\Z_3).
\]
The second identity follows by taking compatible reductions and is
\(F\)-equivariant.  The elementary identity
\[
 v_3(4^{3^s}-1)=s+1
\]
follows by induction: if \(u=1+3^a b\), \(a\geq1\) and \(3\nmid b\),
then
\[
 u^3-1=3^{a+1}b\bigl(1+3^a b+3^{2a-1}b^2\bigr),
\]
and the last factor is a \(3\)-adic unit.  Consequently the functions
\[
 h_r(x)=4^{\,x\bmod 3^{r-1}}\pmod{3^r}
\]
are compatible and satisfy \(Fh_r=4h_r\), including at the period's
wraparound.  Their limit \(h\in\widehat M_3\) satisfies
\(Fh=4h\) and \(h(0)=1\).  It is not \(3\)-primary torsion and
remains a nonzero eigenvector after tensoring with \(\Q_3\).
\end{proof}

The second example uses inflation maps between finite permutation
lattices.  It does not invalidate the following distinct calculation.
If \(P=\bigoplus_O\Z[O]\) for fixed finite \(F\)-orbits \(O\),
its \(3\)-adic completion embeds in \(\prod_O\Z_3[O]\).
For a \(q\)-eigenvector, projection to an orbit of size \(d\) gives
\((q^d-1)v_O=0\).  If \(q\in\Q_3^\times\) is not a root of unity,
all projections vanish, even after rationalization.  Thus a residue
argument must examine the actual transition maps and the unramified
kernel; the fixed-orbit permutation calculation itself remains valid.

Neither example identifies its completed group with the degree-four
birational cohomology in B1.  They show precisely why finite-coefficient
norm-residue comparison, or finite Frobenius orbits of individual
symbols, cannot supply the missing interchange without uniform
ramification control.  The exact finite programs are
\begin{quote}
\small
\path{part_b/code/kummer_limit_counterexample.py}\\
\path{part_b/code/frobenius_completion_counterexample.py}
\end{quote}
with their correspondingly named \texttt{.txt} outputs in
\texttt{part\_b/output}.  They check the stated finite congruences;
the infinite assertions are proved above.

\begin{proposition}[The Chow obstruction on an explicit fixed open]
\label{prop:B-fixed-open-ramification-obstruction}
\status{Proved exact comparison formula; no colimit vanishing or survival assertion}
Let \(k=\F_q\), let \(\ell\ne\operatorname{char}k\), and put
\(Y=\mathbb P^4_k\).  Let \(D_1,D_2\subset Y\) be smooth
hypersurfaces of positive degree meeting transversely in a smooth
surface \(S\).  Write \(U=Y\setminus(D_1\cup D_2)\) and
\(\Gamma=\Gal(\bar k/k)\).  Assume the divisor Tate conjecture
for \(S\).  Define
\[
\begin{aligned}
 V&=\operatorname{CH}^1(S)_\Q\otimes\Q_\ell,\\
 C&=\bigl(\operatorname{CH}^2(D_1)_\Q
             \oplus\operatorname{CH}^2(D_2)_\Q\bigr)
                 \otimes\Q_\ell,\\
 G:V&\longrightarrow C,
       &G(z)&=((i_1)_*z,-(i_2)_*z),
\end{aligned}
\]
where \(i_a:S\hookrightarrow D_a\) are the inclusions.  Let
\[
 \operatorname{cl}_C:C\longrightarrow
 H^4(\bar D_1,\Q_\ell(2))\oplus H^4(\bar D_2,\Q_\ell(2))
\]
be the ordinary cycle map.  There is an isomorphism
\[
\begin{split}
 \operatorname{coker}\bigl[
   H_M^4(U,\Q(3))\otimes\Q_\ell
       &\longrightarrow H^4(\bar U,\Q_\ell(3))^\Gamma
                         \bigr]\\
       &\simeq\operatorname{im}G\cap\ker\operatorname{cl}_C.
\end{split}
\]
Here \(H_M\) is ordinary rational motivic cohomology and the
right-hand cohomology is continuous \(\ell\)-adic cohomology.
The identification is canonical after fixing the displayed
Mayer--Vietoris orientation and the usual cycle maps.
\end{proposition}
\begin{proof}
Put \(D=D_1\cup D_2\).  The Mayer--Vietoris triangle for
supports and purity give
\[
\begin{split}
 R\Gamma(\bar S,\Q_\ell(1))[-4]
   &\longrightarrow
      R\Gamma(\bar D_1,\Q_\ell(2))[-2]
        \oplus R\Gamma(\bar D_2,\Q_\ell(2))[-2]\\
   &\longrightarrow R\Gamma_{\bar D}(\bar Y,\Q_\ell(3)).
\end{split}
\]
For the support triangle, apply ordinary open Mayer--Vietoris
to the cover
\((Y\setminus D_1)\cup(Y\setminus D_2)=Y\setminus S\),
compare with restriction from \(Y\), and take the three
mapping cones in functorial cochain models.  The same construction
works motivically; no descent across the singular closed immersion
\(D\hookrightarrow Y\) is assumed.  The purity shifts are those
for the regular immersions of codimensions two and one,
respectively~\cite[Theorem~3.1.1, p.~323]{c_riou2014}.

The relevant support sequence is
\[
\begin{aligned}
 H^1(\bar S,\Q_\ell(1))
 &\longrightarrow\bigoplus_{a=1}^2H^3(\bar D_a,\Q_\ell(2))
 \longrightarrow H^5_{\bar D}(\bar Y,\Q_\ell(3))\\
 &\longrightarrow H^2(\bar S,\Q_\ell(1))
 \longrightarrow\bigoplus_{a=1}^2H^4(\bar D_a,\Q_\ell(2)).
\end{aligned}
\]
The first two terms are pure of weight \(-1\).  Hence the
kernel of the map from the support group to the last Gysin
kernel has \(F-1\) invertible, where \(F\) is geometric
Frobenius.  Both its \(\Gamma\)-invariants and its first
\(\Gamma\)-cohomology vanish.  Taking invariants therefore
identifies the support invariants with the Gysin kernel on
invariant subspaces.  This uses the proper smooth weight
theorem, without a separate semisimplicity assumption
\cite[Corollary~3.3.9, p.~207]{c_deligneweil1980}.

Next, localization gives
\[
 H^4(\bar Y,\Q_\ell(3))\longrightarrow
 H^4(\bar U,\Q_\ell(3))\longrightarrow
 H^5_{\bar D}(\bar Y,\Q_\ell(3))\longrightarrow
 H^5(\bar Y,\Q_\ell(3))=0.
\]
The image of the first term is pure of weight \(-2\), so
the same invariant and coinvariant argument identifies
\(H^4(\bar U,\Q_\ell(3))^\Gamma\) with the support
invariants.  Divisor Tate on \(S\), together with rational
Kummer injectivity for its Picard group, identifies
\[
 V\simeq H^2(\bar S,\Q_\ell(1))^\Gamma.
\]
Thus the continuous target in the statement is exactly
\(\ker(\operatorname{cl}_C G)\subset V\).

For ordinary motivic cohomology, the analogous support sequence
has the segment
\[
 H^5_D(Y,\Q(3))\longrightarrow
 H_M^2(S,\Q(1))\longrightarrow
 H_M^4(D_1,\Q(2))\oplus H_M^4(D_2,\Q(2)).
\]
The last three groups are the corresponding rational Chow
groups.  Its image in \(\operatorname{CH}^1(S)_\Q\)
is consequently the kernel of the rational signed Gysin map.
Motivic localization also gives
\[
 H_M^4(U,\Q(3))\longrightarrow H^5_D(Y,\Q(3))
                   \longrightarrow H_M^5(\mathbb P^4_k,\Q(3)).
\]
The last group is zero: the projective-bundle formula expresses it as
\[
 \bigoplus_{r=0}^4 H_M^{5-2r}(k,\Q(3-r))=0,
\]
by the ordinary finite-field motivic calculation.  These are
the usual motivic purity, localization and projective-bundle
maps~\cite[(14.5.1), Theorem~15.12, Corollary~15.13,
Theorem~15.15]{BMVW2006}.
The ordinary cycle map commutes with these maps.  Since
\(\Q_\ell\) is flat over \(\Q\), its image in \(V\)
is therefore exactly \(\ker G\).  Finally, applying \(G\)
induces the elementary isomorphism
\[
 \ker(\operatorname{cl}_C G)/\ker G
       \simeq\operatorname{im}G\cap\ker\operatorname{cl}_C,
\]
which proves the assertion.
\end{proof}

The formula locates a concrete obstruction: divisor Tate on the
surface supplies an algebraic representative, but the outgoing
Gysin relation is initially only a cohomological relation on the
threefold boundary components.  It need not yet be a rational
Chow relation.  We have neither proved that this intersection
vanishes nor constructed a nonzero element of it.
Moreover, failure of motivic surjectivity on this fixed \(U\)
does not show survival under every further shrinking: a class
could become zero on a smaller open without having had a motivic
lift on the original one.  The proposition thus neither solves
nor refutes the projective-space colimit gate.  It assumes only
the displayed smooth transverse pair, without asserting that
every open admits such a compactification.

\subsubsection*{Ordinary Milnor elements and infinitely divisible residues}

The following intermediate result uses the refined weight bound in
\citet[Proposition~3.7, p.~16]{BCTSchreiederDivisibility2026} to extend the
conclusion of \citet[Lemma~4.1, pp.~11--12]{BSchreiederMilnorTate2026} beyond
its stated degree and unramified-vanishing hypotheses.  No novelty claim
is made, and this is not a resolution of the invariant-colimit gate.

\begin{proposition}[Kernel of the ordinary generic Milnor regulator]
\label{prop:B-Milnor-divisible-kernel}
\status{Proved intermediate statement; the universal B.1 gate is unchanged}
Let $k=\overline{\mathbf F}_p$, let $K/k$ be a finitely generated field
extension, and let $\ell\ne p$.  Define ordinary generic continuous
cohomology by
\[
 H^n_\eta(K,\mathbf Z_\ell(n))
   =\varinjlim_{\varnothing\ne U\subset X}
      H^n_{\mathrm{cont}}(U,\mathbf Z_\ell(n)),
\]
where $X$ is a model of $K$ and the limit can be taken over its smooth
dense opens.  For the ordinary cycle map
\[
 \rho_n:K_n^M(K)\longrightarrow
                 H^n_\eta(K,\mathbf Z_\ell(n)),
\]
one has, for every $n\ge0$,
\[
 \ker\rho_n=\bigcap_{r\ge1}\ell^rK_n^M(K).
 \tag{B.M1}\label{eq:B-Milnor-kernel}
\]
For $n\ge1$, suppose moreover that $X/k$ is smooth projective and that
$\alpha\in K_n^M(k(X))$ has infinitely $\ell$-divisible residue at every
prime divisor of $X$.  Then $\alpha$ is infinitely $\ell$-divisible.
The second assertion deliberately excludes $n=0$.
\end{proposition}

\begin{proof}
The cycle map is defined on finite-type opens before passing to the
colimit.  It can be constructed from products of Kummer classes of
symbols, or from the compatible motivic cycle map in
\citet[\S6.1, equations~(6.2)--(6.4)]{BCTSchreiederDivisibility2026}.
It commutes with residues, finite-extension restriction and transfer,
and reduction modulo $\ell^r$.  Reduction on the field is the
norm-residue isomorphism.  Consequently $\rho_n(\alpha)=0$ implies
$\alpha\in\ell^rK_n^M(K)$ for every $r$.

We prove the reverse implication by induction on $n$, simultaneously
for all finitely generated extensions of $k$.  The case $n=0$ is the
injection $\mathbf Z\hookrightarrow\mathbf Z_\ell$.
Assume the assertion in degree $n-1$.  First suppose that $K$ has a
smooth projective model $X/k$, and let $\alpha$ be infinitely
$\ell$-divisible.  Every residue $\partial_D\alpha$ is infinitely
$\ell$-divisible in $K_{n-1}^M(k(D))$.  The induction hypothesis for
the residue field $k(D)$ therefore gives
\[
 \partial_D\rho_n(\alpha)
       =\rho_{n-1}(\partial_D\alpha)=0
\]
in ordinary generic continuous cohomology for every prime divisor $D$.

Localization and codimension-one purity place the rationalized class
in the ordinary refined unramified group
\[
 H^n_{0,\mathrm{nr}}(X,\mathbf Q_\ell(n))
 =\operatorname{im}\bigl(
 H^n(F_1X,\mathbf Q_\ell(n))\longrightarrow
 H^n(F_0X,\mathbf Q_\ell(n))\bigr).
\]
Here is the finite-support detail.  Start with one open on which the
class is defined.  It misses finitely many prime divisors.  The zero
generic residue on each of them vanishes already on a dense open of
that divisor.  Delete the finitely many remaining closed subsets of
codimension at least two in $X$, including singular and intersection
loci.  The localization sequence then lifts the class across all
prime divisors.  This argument involves an ordinary filtered
colimit; it makes no assertion about commuting that colimit with an
inverse limit.  The precise refined-unramified definition is
\citet[equations~(2.7)--(2.8), p.~11]{BCTSchreiederDivisibility2026}.

By the cited Proposition~3.7, the weights of this rational group lie
in
\[
 [-n,\max(-n,-2)],
\]
which is strictly negative for $n\ge1$.  Choose a finite field over
which $X$ and $\alpha$ are defined.  The ordinary Milnor element
$\alpha$ has a finite Frobenius orbit, as does its regulator.  The
finite-dimensional span of that orbit has only roots of unity as
Frobenius eigenvalues and thus cannot be a nonzero submodule of this
strictly negative-weight group.  Hence $\rho_n(\alpha)$ is zero
rationally.  The ordinary integral generic continuous group is
torsion-free by \citet[Theorem~2.2 and Corollary~2.3,
pp.~6--7]{BSchreiederMilnorTate2026}; the change of twist used for
torsion-freeness is allowed because $k$ contains all $\ell$-power
roots of unity.  Thus $\rho_n(\alpha)=0$ integrally.

For a general $K$, choose a projective model over a finite field and
a smooth projective alteration
\citep[Theorem~4.1 and Remark~4.2]{c_dejong1996}.
After choosing a geometric component and, if necessary, enlarging
the finite field of definition, its function field $L$ is a finite
extension of $K$.  The pullback of $\alpha$ is infinitely
$\ell$-divisible and its regulator vanishes by the just-proved
smooth-projective case.  Restriction followed by corestriction is
multiplication by $[L:K]$ on ordinary generic continuous cohomology:
this is finite-\'etale transfer on suitable dense opens for the
separable part, and the purely inseparable part is handled by
universal-homeomorphism invariance and its degree.  Therefore
\[
 [L:K]\rho_n(\alpha)=0.
\]
Torsion-freeness gives $\rho_n(\alpha)=0$.  This proves the induction
statement for arbitrary residue fields as well as for $K$ itself.
No degree-prime-to-$\ell$ condition is needed in this argument.

Finally, take an ordinary $\alpha$ on a smooth projective $X$ whose
residues are infinitely $\ell$-divisible, without initially assuming
the same of $\alpha$.  The proved degree-$(n-1)$ assertion again
kills every residue of its regulator.  The same localization,
negative-weight and finite-orbit argument gives
$\rho_n(\alpha)=0$.  Formula~\eqref{eq:B-Milnor-kernel} then gives
infinite $\ell$-divisibility.  This last argument requires $n\ge1$:
in degree zero all residues vanish whereas $K_0^M=\mathbf Z$ is not
divisible.
\end{proof}

\begin{corollary}[The separated Milnor quotient is free]
\label{cor:B-Milnor-free-quotient}
\status{Proved intermediate structural consequence; no new universal gate}
In the notation of Proposition~\ref{prop:B-Milnor-divisible-kernel}, set
\[
 D_\ell(K_n)=\bigcap_{r\ge1}\ell^rK_n^M(K),\qquad
 M_n(K)=K_n^M(K)/D_\ell(K_n).
\]
Then $D_\ell(K_n)$ is $\ell$-divisible and $M_n(K)$ is a free abelian
group of at most countable rank.  Moreover, $D_\ell(K_n)$ is independent
of the prime $\ell\ne p$.  Thus one canonical quotient $M_n(K)$
induces the canonical isomorphisms of completions
\[
 \widehat{K_n^M(K)}_\ell\simeq\widehat{M_n(K)}_\ell
 \qquad(\ell\ne p),
\]
equivariant for the Frobenius action from any finite-field model.
No saturated residue embedding or equivariant permutation basis is
asserted, and a common integral module does not by itself make its
completed Frobenius spectra independent of the prime.
\end{corollary}

\begin{proof}
Since $D_\ell(K_n)=\ker\rho_n$ and the regulator target is torsion-free,
$M_n(K)$ is torsion-free.  If $a\in D_\ell(K_n)$, choose $b$ with
$\ell b=a$.  Then $\ell\rho_n(b)=0$, so $\rho_n(b)=0$ and
$b\in D_\ell(K_n)$.  Thus $D_\ell(K_n)$ is $\ell$-divisible.

We prove freeness by induction on $n$ for all finitely generated $K/k$.
The case $n=0$ is $M_0(K)=\mathbf Z$.  For arbitrary $K$, take a finite
extension $L/K$ with a smooth projective model, supplied by an
alteration as in the proposition.  Restriction induces an injection
$M_n(K)\hookrightarrow M_n(L)$: if the restricted class has zero
regulator, transfer kills $[L:K]\rho_n(a)$, and torsion-freeness gives
$\rho_n(a)=0$.  It therefore suffices to prove freeness for the
function field of a smooth projective $X$.

For $n\ge1$, the residue assertion of the proposition gives an injection
\[
 M_n(k(X))\lhook\joinrel\longrightarrow
       \bigoplus_{D\in X^{(1)}}M_{n-1}(k(D)).
\]
Each ordinary Milnor element has only finitely many nonzero residues,
so this is a direct sum.  By the induction hypothesis its summands
are free abelian; hence so is the direct sum, and every subgroup of
a free abelian group is free.  All the fields and indexing sets here
are countable because $\overline{\mathbf F}_p$ is countable.  This
completes the induction.

Finally, divisibility of $D_\ell(K_n)$ gives, for every $r$, the
canonical equivariant isomorphism
\[
 K_n^M(K)/\ell^rK_n^M(K)\simeq M_n(K)/\ell^rM_n(K).
\]
Taking inverse limits proves the completion assertion.  None of the
argument makes a residue embedding $\ell$-saturated or chooses a
Frobenius-stable basis.

For independence of the kernel, let $\ell,\ell'\ne p$.  The image of
$D_\ell(K_n)$ in the free abelian group
$K_n^M(K)/D_{\ell'}(K_n)$ is $\ell$-divisible.  Such an image is zero:
an element of a free abelian group has finitely many integer
coordinates and cannot be divisible by every power of $\ell$ unless
it vanishes.  Thus $D_\ell(K_n)\subset D_{\ell'}(K_n)$, and symmetry
gives equality.  The common subgroup is divisible by every prime
different from $p$, and its common quotient gives all the displayed
completion isomorphisms.
\end{proof}

\paragraph{The remaining obstruction.}
The result concerns ordinary Milnor elements, whose Frobenius
orbits are finite.  It does not identify all continuous generic
classes, or all completed Milnor classes, with such elements.
In the first B.1 gate, a residue of
$H^4(F_0\mathbf P^4_k,\mathbf Q_\ell(3))$ lies in degree three
and twist two on a boundary threefold.  Proposition~3.7 gives the
unramified weight interval $[-1,0]$ there, so weight zero is not
excluded.  Replacing that twist by the diagonal twist three would
change the problem.  Moreover, Proposition~\ref{prop:B-limit-diagnostics}
shows that completion of a locally finite Frobenius module can
create non-root-of-unity eigenvectors; its abstract function module
is not asserted to be a Milnor module.

The additional weight input in the proof is essential to this
argument.  Inverse-limit divisibility statements alone do not
correct residues.  For example, the complex
\[
 C^0=\mathbf Z\ \hookrightarrow C^1=\mathbf Z[1/\ell],\qquad C^2=0
\]
has $H^0(C)=0$, divisible $H^1(C)=\mathbf Z[1/\ell]/\mathbf Z$,
zero primary torsion in its terms, and $H^1(C/\ell^r)=0$ for every
$r$.  Nevertheless, $1\in C^0$ is not infinitely $\ell$-divisible
while its image in $C^1$ is.  This is a formal obstruction to that
inference, not a geometric counterexample.


There is a useful limited uniqueness statement.  Fix a semisimple
\(\Q\)-linear category of pure motives with simple unit \(\mathbf1\)
and \(\End(\mathbf1)=\Q\), and assume its realization detects precisely
the unit-isotypic part as the Frobenius-one eigenspace.  Any
\(\Q\)-linear cohomological functor on its bounded derived category
with the prescribed continuous-cohomology values on pure objects is,
noncanonically, the functor
\[
 M\longmapsto\Hom(\mathbf1,M)\oplus
 \bigl(\Hom(\mathbf1,M)\otimes L\bigr)[-1],\qquad \dim_\Q L=1.
 \tag{B.9}
\]
Indeed every object is a direct sum of simple objects; comparison kills
all nonunit simples in every degree and prescribes the two groups of
the unit.  Every complex in a semisimple category splits into its
cohomology, so there are no further extension choices for this additive
derived functor.  A chosen \(e\in L\setminus0\) normalizes the
isomorphism.  Without \(e\), multiplication by \(a\in\Q^\times\) on
the odd factor is a nontrivial natural automorphism preserving \(H^0\).
This proves uniqueness only in this specified additive sense, for this
category.  It is not a uniqueness theorem for topoi, integral complexes,
products, six operations, or rational cycle classes.  Identifying (B.9)
with actual motivic Weil--etale cohomology requires the gate in
Theorem~\ref{thm:B-Geisser-gate}; the manuscript's unrestricted list is
already inconsistent by Proposition~\ref{prop:B-continuity}.

\subsection{B.2: Tate twists on number rings and the exact regulator defects}
Let \(K\) have \(r_1\) real and \(r_2\) complex places, let \(S\) be
a finite set of finite places, and put
\[
 X=\Spec\mathcal O_{K,S},\quad T=S\cup S_\ell,\quad
 U=X[1/\ell]=\Spec\mathcal O_{K,T},\quad
 r_n=\begin{cases}r_2&n\text{ even},\\r_1+r_2&n\text{ odd}.
 \end{cases}
\]
All the cohomology in the next table is \emph{full continuous}
cohomology of \(U\).  In degree one it is equivalently the cohomology
of the Galois group of the maximal extension unramified outside
\(T\) and infinity.

There is a scope correction to R6's proposed comparison example.
Its ranks one and zero for \(\Q(0)\) compare two constant-coefficient
theories on the fixed scheme \(\Spec\Z[1/p]\).  If, as B.2
explicitly requests, the scheme is instead replaced by
\(X[1/\ell]\) for each coefficient prime, the degree-one ranks in
twist zero are both one.  A counterexample with exactly this latter
convention is \(X=\Spec\Z[1/p]\) and twist one: (B.11) gives
rank one for \(\ell=p\), but rank two for \(\ell\ne p\).
The extra generator is the newly inverted prime \(\ell\).
Thus full cohomology still fails the requested common rational
comparison, but its verification must use the same base convention.

\begin{theorem}[All full rational etale groups, with the defect left visible]
\label{thm:B-full-tate}
\status{Proved}
For every prime \(\ell\),
\[
\begin{array}{c|ccc}
n&H^0_{\Bcts}(U,\Q_\ell(n))&
\dim_{\Q_\ell}H^1_{\Bcts}(U,\Q_\ell(n))&
\dim_{\Q_\ell}H^2_{\Bcts}(U,\Q_\ell(n))\\ \hline
0&\Q_\ell&r_2+1+\delta_{K,\ell}(0)&\delta_{K,\ell}(0)\\
1&0&r_1+r_2-1+|T|&|T|-1\\
n\ge2&0&r_n&0\\
n<0&0&r_n+\delta_{K,\ell}(n)&\delta_{K,\ell}(n).
\end{array}\tag{B.10}
\]
All other degrees vanish.  For \(n=1\), more precisely,
\[
 H^1_{\Bcts}(U,\Q_\ell(1))=
 \mathcal O_{K,T}^\times\otimes\Q_\ell,\qquad
 H^2_{\Bcts}(U,\Q_\ell(1))=
 \ker\bigl(\Q_\ell^T\xrightarrow{\sum}\Q_\ell\bigr).
 \tag{B.11}
\]
For \(n\ge2\), the Chern-class map identifies degree one with
\(K_{2n-1}(\mathcal O_K)^{(n)}\otimes\Q_\ell\).
The defect \(\delta_{K,\ell}(0)\) is exactly the Leopoldt defect.
For \(n<0\), writing \(m=1-n\), it is exactly the dimension of
the kernel of the \(\ell\)-adic regulator
\[
 K_{2m-1}(\mathcal O_K)^{(m)}\otimes\Q_\ell
 \longrightarrow\bigoplus_{v\mid\ell}K_v .\tag{B.12}
\]
In particular, replacing every defect in (B.10) by zero is equivalent,
at the indicated \((K,\ell,n)\), to Leopoldt for \(n=0\), or to
Schneider's regulator-injectivity conjecture at \(m=1-n\) for \(n<0\).
\end{theorem}
\begin{proof}
For twist one, Kummer theory identifies degree one with units because
the ideal class group is finite.  The Brauer reciprocity sequence
identifies the rational Tate module of \(\Br(U)\) with the kernel of
the sum of the finite local invariants; real places contribute only
groups killed by two.  This proves (B.11).

For positive twists at least two, the rational Chern-class comparison
and the finiteness of the degree-two integral cohomology give the
assertions; the rank is Borel's \(r_n\).  For odd \(\ell\) a precise
modern reference is Nickel's Theorem~3.2; at two use Rognes--Weibel's
Theorems~0.4 and~0.6, after rationalization.  Their real-place errors
are finite two-groups, hence do not alter this rational assertion.
Localization away from \(\ell\) does not change either group when
\(n\ne1\), since the additional residue-field terms have twist
\(n-1\ne0\) and are rationally zero~\cite[Theorem~3.2]{BNickel2022}
\cite[Theorems~0.4 and~0.6]{BRognesWeibel2000}.

For \(n\le0\), Poitou--Tate duality identifies
\[
 H^2_{\Bcts}(U,\Q_\ell(n))^\vee=
 \ker\left(H^1_{\Bcts}(U,\Q_\ell(1-n))
 \longrightarrow\bigoplus_{v\in T}H^1(K_v,\Q_\ell(1-n))\right).
 \tag{B.13}
\]
No local degree-two term intervenes: its dual is
\(H^0(K_v,\Q_\ell(1-n))=0\).  If \(n=0\), the valuations first
restrict this kernel to \(\mathcal O_K^\times\otimes\Q_\ell\), and
the remaining map is the vector of local \(\ell\)-adic logarithms.
This is exactly the definition of the Leopoldt defect.  If \(n<0\),
the positive-twist comparison and the inverse Bloch--Kato exponential
identify the map with (B.12).  These descriptions also show independence
of the extra places in \(S\).  The global Euler characteristic is
\(-r_n\), and \(H^0\) is one-dimensional only in twist zero; hence
the degree-one dimensions follow.  Rational cohomological dimension is
two, including at \(\ell=2\), since the additional real-place
cohomology is two-torsion.  Finally, Nickel's Conjecture~3.5,
Remarks~3.7--3.10 and Lemma~3.11 identify precisely this strict
localization kernel with Schneider's conjecture; no nonvanishing of
an unspecified regulator has been assumed~\cite[\S3.6]{BNickel2022}.
\end{proof}

\paragraph{The local conditions must depend on the arithmetic open.}
For \(L/\Q_\ell\) finite, \(d_L=[L:\Q_\ell]\), the dimensions of
the Bloch--Kato local subgroups are
\[
\begin{array}{c|rrrr}
 n&\dim H^1_e&\dim H^1_f&\dim H^1_g&\dim H^1\\ \hline
n<0&0&0&0&d_L\\
0&0&1&1&d_L+1\\
1&d_L&d_L&d_L+1&d_L+1\\
n>1&d_L&d_L&d_L&d_L .
\end{array}\tag{B.14}
\]
The twist-zero finite line consists of unramified characters; the
twist-one finite subgroup is \(\mathcal O_L^\times\otimes\Q_\ell\),
while its geometric subgroup is all of
\(L^\times\widehat\otimes\Q_\ell\).  For a finite place
\(v\nmid\ell\), the full \(H^1(K_v,\Q_\ell(n))\) has dimension
one exactly for \(n=0,1\), and zero otherwise; its unramified subgroup
has dimension one exactly for \(n=0\).  The first table is
Bloch--Kato's Example~3.9, with their Proposition~3.8 providing the
orthogonality statements.  The second follows directly from inertia
cohomology and the eigenvalue of Frobenius~\cite[Proposition~3.8,
Example~3.9]{BBlochKato1990}.

Define \(H^1_{f,X}(K,V)\) using \(H^1_f(K_v,V)\) at finite
\(v\notin S\) and \(H^1_g(K_v,V)\) at \(v\in S\); this is
exactly the convention of Bloch--Kato's Definition~5.1.  Away from
\(\ell\), the latter condition imposes no restriction for the Tate
representations at issue.  Imposing \(f\) also at a removed place
above \(\ell\) would incorrectly discard that place's unit valuation.
Imposing no condition on twist zero there would retain precisely the
ramified characters responsible for R6.

\begin{proposition}[An unconditional rational structure in degree one]
\label{prop:B-Hf-tate}
\status{Proved}
For every integer \(n\), all primes \(\ell\), and the above
geometric local conditions, the natural realization maps give
\[
 H^1_{f,X}(K,\Q_\ell(n))\simeq V_X(n)\otimes\Q_\ell,
\quad
 V_X(n)=
 \begin{cases}
 \mathcal O_{K,S}^\times\otimes\Q,&n=1,\\
 K_{2n-1}(\mathcal O_K)^{(n)}\otimes\Q,&n\ge2,\\
 0,&n\le0.
 \end{cases}\tag{B.15}
\]
In particular the degree-one assertion for negative twists does not
require Schneider's conjecture.  That conjecture concerns the full
groups in (B.10), not the group in (B.15).
\end{proposition}
\begin{proof}
For \(n=1\), the Kummer map in (B.11) and (B.14) keep exactly
the \(S\)-units.  For \(n\ge2\), all local finite conditions are
already satisfied, so Theorem~\ref{thm:B-full-tate} applies.  For
\(n=0\), classes satisfying these conditions are continuous
\(\Q_\ell\)-valued characters unramified at every finite place.
Global class field theory factors them through a finite ideal class
group (and finite real-place factors); hence they vanish.  For \(n<0\),
both local \(f\) and \(g\) are zero.  The strict global kernel is
dual, by (B.13) with the twists interchanged, to the degree-two group
in positive twist \(1-n\ge2\), which is zero by the positive-twist
theorem.  This proves the assertion without setting
\(\delta_{K,\ell}(n)\) to zero.
\end{proof}

\paragraph{Degrees other than one and the scope of a corrected requirement 4).}
On the proper number ring there is a precise higher-degree correction.
Take the Bloch--Kato Selmer complex.  At a place over \(\ell\),
the local finite complex for the crystalline Tate representation is
\[
 C_f(K_v,V)=
 [D_{\rm cris}(V)\xrightarrow{\ ((1-\varphi),\mathrm{can})\ }
 D_{\rm cris}(V)\oplus t_V],\qquad
 t_V=D_{\rm dR}(V)/\operatorname{Fil}^0,
\]
in degrees zero and one.  Away from \(\ell\), use
\([V^{I_v}\xrightarrow{1-\operatorname{Frob}_v}V^{I_v}]\).
These complexes have cohomology \(H^0(K_v,V)\) and
\(H^1_f(K_v,V)\), and no other cohomology.  The global complex is
the mapping fibre of the global cochain complex plus these local
complexes mapping to the direct sum of the full local complexes.
Fontaine's \S\S4.4--4.5 give these constructions and their duality;
the real-place Tate errors are two-torsion and disappear in this
rational calculation~\cite[pp.~214--216]{BFontaine1992}.

Here is the duality argument in the degrees actually used.  Local
Tate duality and Bloch--Kato's orthogonality of finite subspaces
give an isomorphism of the quotient cone with
\[
 \operatorname{Cone}\bigl(C_f(K_v,V)\longrightarrow
   R\Gamma(K_v,V)\bigr)
 \simeq R\operatorname{Hom}_{\Q_\ell}
       (C_f(K_v,V^\vee(1)),\Q_\ell)[-2].
\]
Indeed its degree-one quotient is dual to the complementary finite
subspace, and its degree-two local invariant is dual to
\(H^0(K_v,V^\vee(1))\); all other groups vanish.
To keep the issue of choices explicit, set
\[
\begin{split}
 C(V)&=\operatorname{coker}\left(H^1(U,V)\longrightarrow
       \bigoplus_{v\in T}H^1(K_v,V)/H^1_f(K_v,V)\right),\\
 K(V)&=\ker\left(H^2(U,V)\longrightarrow
       \bigoplus_{v\in T}H^2(K_v,V)\right).
\end{split}
\]
The Selmer cone and the Poitou--Tate sequence, respectively, give
\[
\begin{gathered}
0\longrightarrow C(V)\longrightarrow H_f^2(K,V)
 \longrightarrow K(V)\longrightarrow0,\\
0\longrightarrow C(V)\longrightarrow
 H_f^1(K,V^\vee(1))^\vee
 \longrightarrow K(V)\longrightarrow0.
\end{gathered}
\]
These sequences split as vector spaces.  Thus the degree-two
isomorphism below exists, but is not asserted to be canonical
from the finite subgroups alone.  The extreme-degree
identifications follow directly from global invariants and the
local invariant sum.  We obtain
\[
 H_f^0=H^0(K,V),\quad H_f^1=H_f^1(K,V),\quad
 H_f^2\simeq H_f^1(K,V^\vee(1))^\vee,\quad
 H_f^3\simeq H^0(K,V^\vee(1))^\vee.\tag{B.16}
\]
For Tate motives, (B.15)--(B.16) supply a common rational space in
every degree, including the dual unit and dual Borel-regulator spaces;
the higher comparison isomorphisms use the specified splittings.
Their graded groups agree with Flach--Morin's rational compactified
complex (B.6), now for \(\mathcal O_K\).  This also proves why a comparison with
\emph{ordinary} motivic cohomology in every degree would be false:
\(H^3_{\BM}(\Spec\Z,\Q(1))=0\), whereas (B.16) gives \(\Q_\ell\).

For an arithmetic open, merely specifying the subgroups in degree one
does not specify the higher Selmer complex.  In particular, at a
removed \(\ell\)-adic place in twist one one must retain the local
degree-two invariant as well as \(H^1_g=H^1\) to reproduce the
localization triangle.  The rational absolute complex with localization
has, in addition to the degree-one \(S\)-unit group,
\[
 H^2(X,\Q(1))=\ker(\Q^S\xrightarrow{\sum}\Q),\qquad
 H^3(X,\Q(1))=\operatorname{coker}(\Q^S\xrightarrow{\sum}\Q).
 \tag{B.17}
\]
For \(n\ne1\) it has the same groups as the compactified proper
number ring.  This follows by applying the finite-field computation
to the Gysin boundary in twist \(n-1\).  It is a construction of
a specified rational absolute theory; it is not an assertion that
every expression denoted \(H^i_f\) in the literature has these values.

\begin{corollary}[The common rational spaces in every Tate degree]
\label{cor:B-Hf-all-degrees}
\status{Proved for the specified absolute and Selmer complexes}
For \(X=\Spec\mathcal O_{K,S}\), the following table lists all
possibly nonzero common \(\Q\)-spaces \(H^i_{\rm abs}(X,\Q(n))\).
The dual in the table is the \(\Q\)-linear dual.
\[
\begin{array}{c|c|l}
 n&i&H^i_{\rm abs}(X,\Q(n))\\ \hline
 0&0&\Q\\
 0&2&(\mathcal O_K^\times\otimes\Q)^\vee\\
 1&1&\mathcal O_{K,S}^\times\otimes\Q\\
 1&2&\ker(\Q^S\xrightarrow{\sum}\Q)\\
 1&3&\operatorname{coker}(\Q^S\xrightarrow{\sum}\Q)\\
 n\ge2&1&K_{2n-1}(\mathcal O_K)^{(n)}\otimes\Q\\
 n<0&2&(K_{1-2n}(\mathcal O_K)^{(1-n)}\otimes\Q)^\vee .
\end{array}
\tag{B.17a}
\]
Every unlisted group is zero.  For every prime \(\ell\), including
two, these spaces tensor to the cohomology of the specified Selmer
complex with the geometric open-place conditions.  Neither Leopoldt
nor Schneider is assumed in this assertion.
\end{corollary}
\begin{proof}
For the proper number ring, substitute (B.15) into (B.16).  For the
open, the local change from the finite to the geometric condition
is zero in every twist except one.  In twist one it adds the local
valuation and invariant in degrees one and two.  The map
\(\mathcal O_{K,S}^\times\otimes\Q\to\Q^S\) is surjective:
its integral cokernel embeds in the finite ideal class group.
The map between the remaining invariant lines is the sum map by
Brauer reciprocity.  This proves exactly (B.17), including
\(S=\varnothing\), when its cokernel is \(\Q\).  Tensoring the
construction with \(\Q_\ell\) gives the local complexes just
described.  The negative-twist assertion uses the dual of the proved
positive-twist group, not injectivity of the regulator (B.12).
\end{proof}

The precise degree-one replacement for the manuscript's requirement
4) is the Bloch--Kato realization conjecture: for a geometric motive
\(M\) over \(K\), define the motivic extension space over an
arithmetic open by the image from a regular model (and use
homologically trivial cycles in the middle case).  Its etale regulator
must induce
\[
 H^1_{\rm mot,f,X}(K,M)\otimes\Q_\ell
 \xrightarrow{\sim}H^1_{f,X}(K,M_\ell)\quad\text{for every }\ell.
 \tag{B.18}
\]
For motives of the form \(h^m(Y)(r)\), this is exactly the
model-image and realization assertion in Bloch--Kato's
Conjecture~5.3(i)--(iii), with the local convention of Definition~5.1.
In Fontaine's direct precursor to his joint work with Perrin-Riou,
\S6.5 expresses the proper-base comparison as \emph{f-admissibility}:
the spaces \(H_f^0(K,M)\) and \(H_f^1(K,M)\) are finite-dimensional,
and their realization maps \(c^0_\ell,c^1_\ell\) are isomorphisms
after tensoring with \(\Q_\ell\), for every \(\ell\).
His \S6.6 conjectures this for the motivic category, and \S6.7
spells it out for the geometrically defined motivic spaces
\cite[pp.~221--223]{BFontaine1992}.  This is a directly opened
primary-source formulation; the later full Fontaine--Perrin-Riou
article was not obtained in this run.
A statement in all degrees additionally needs a chosen arithmetic
motivic category and the dual completion (B.16)--(B.17), as well as
compatibility with localization and duality.  Over finite fields the
corresponding absolute complex is the Weil complex of B.1; its cycle
comparison retains all three clauses of Theorem~\ref{thm:B-Geisser-gate}.
Bloch--Kato's degree-one conjecture alone does not imply those extra
finite-field clauses.  The precise conjectural package below states
them separately and proves the resulting two specializations;
it does not infer them from an unspecified ``finite condition.''

\paragraph{Scholbach's construction.}
Scholbach's Theorem~6.8 identifies an intermediate-extension motivic
group with \(H^1_{f,\backslash\mathrm{crys}}\), under the stated
axioms for mixed motives and their weights.  His Theorems~6.11 and~6.13
identify it with the model-image definition.  These numbers were
checked against the opened arXiv version.  Remark~6.2 explicitly
excludes the crystalline condition at places over \(\ell\) from the
paper's treatment.  The derived \(\Hom\) definition and assumed six
operations support analogues of requirements 1)--3) in its stated
motivic category; the existence and properties of the motivic
\(t\)-structure are among the assumptions.  They do not prove (B.18)
with the missing local conditions restored.  Moreover ordinary
motivic cohomology in that construction has
\(H^1_{\BM}(\F_q,\Q)=0\) and
\(H^3_{\BM}(\Spec\Z,\Q(1))=0\), so it cannot itself be the
full absolute theory with the odd Weil--etale classes.  This is an
explicit distinction, rather than an unidentified comparison
map~\cite[Remark~6.2 and Theorems~6.8, 6.11, 6.13]{BScholbach2015}.

The other supplied sources do not remove this distinction.  In the
published version of Scholbach's \emph{Special L-values}, Axiom~1.6
states the mixed-motive assumptions, and Theorem~4.5 identifies the
finite-field perfectness assertion with rational=numerical equivalence
plus Beilinson--Parshin vanishing~\cite{BScholbachSpecial2017}.
Adrien Morin's supplied constructible-sheaf paper has a different
coefficient scope: Theorem~2.8 treats twists \(n\le0\), under
its stated perfectness hypothesis (or using Weil--eh cohomology and
resolution), and Theorem~8.3 specifies when a motivic application is
of that constructible kind.  These results do not supply the omitted
positive-twist or arbitrary pure-motive realization
\cite{BAdrienMorin2024}.

\begin{definition}[A precise conjectural replacement for requirement 4)]
\label{def:B-corrected-comparison-package}
\status{Conjectural formulation; its stated specializations are proved below}
The following is a comparison \emph{package}, rather than a claim that
Bloch--Kato's degree-one conjecture already constructs an absolute
category.  Its data include enhanced, stable, \(\Q\)-linear
categories \(\mathcal A(X)\), with exact unital functors from
constructible geometric motives, for finite-type schemes
\(X/\Z\).  These are the proposed arithmetic completions;
their rational cohomology complexes are
\[
 R\Gamma_{\rm abs}(X,M)
   =R\Hom_{\mathcal A(X)}(\mathbf1_X,M),
 \qquad H^i(X,M)=H^iR\Gamma_{\rm abs}(X,M),
\]
and compatible \(\ell\)-adic arithmetic realization complexes
\(\mathsf S_\ell(X,M)\).  The complexes have finite-dimensional
cohomology, and the regulator is required to be a quasi-isomorphism
\[
 c_{\ell,X,M}:R\Gamma_{\rm abs}(X,M)\otimes_\Q\Q_\ell
       \xrightarrow{\ \sim\ }\mathsf S_\ell(X,M).
 \tag{B.18a}
\]
Here the word ``arithmetic'' is fixed by the following requirements;
it does not mean unrestricted cohomology of \(X[1/\ell]\).
\begin{enumerate}
\item On a smooth projective variety over \(k=\F_q\), the rational
complex is the rational motivic Weil--etale complex.  For
\(\ell\ne p\) its realization is
\[
 \mathsf S_\ell(X,M)=
 \operatorname{Cone}\bigl(1-\phi_q:
 R\Gamma(\bar X,M_\ell)\longrightarrow
 R\Gamma(\bar X,M_\ell)\bigr)[-1].
\]
Arithmetic Frobenius and its action are chosen consistently as above.
At \(p\), Tate coefficients use the logarithmic Hodge--Witt
realization.  The conjectural cycle hypothesis is \emph{all three}
clauses of Geisser's \(T(Y,r)\), for every smooth projective
\(Y\) over finite fields: Tate surjectivity, semisimplicity at one,
and rational=numerical equivalence.  The normalized arithmetic class
is \(e_q\), with restriction \(e_q\mapsto m e_{q^m}\).
\item On a proper number ring, prescribe a pure generic motive
\(M/K\) together with a geometric extension \(M_X\), or
explicit intermediate-extension data, representing the finite
condition.  The generic motive is not assumed to be lisse on
the whole proper ring.  The degree-zero and degree-one rational spaces and maps
are Fontaine's \(A^a(M)=H_f^a(K,M)\) and
\(c^a_\ell\), for \(a=0,1\).  For both \(M\) and
\(M^*=M^\vee(1)\), these are required to be finite-dimensional
and to satisfy f-admissibility.  In the geometrically defined case,
\(A^1\) is the model-image group of Bloch--Kato, with the
homologically trivial cycle group in its middle-degree case.
The higher rational spaces are the dual completion
\[
 A^2(M)=A^1(M^*)^\vee,\qquad
 A^3(M)=A^0(M^*)^\vee,
\]
and the realization is the Selmer complex (B.16).  Its local
degree-one condition is unramified away from \(\ell\), and
crystalline \(f\) at retained places above \(\ell\).
An arithmetic open uses the geometric \(g\) condition at omitted
places above \(\ell\), and the entire local change complex,
including its degree-two term.  At real places the compact-support
complex includes the stated Tate modification; the rational
archimedean completion is part of the data.
\item For an open--closed decomposition \(j:U\hookrightarrow X\),
\(i:Z\hookrightarrow X\), both arithmetic theories have the
localization triangles for \(i^!\) and \(j^*\), and the
compact-support triangles for \(j_!\) and \(i^*\).
The maps (B.18a) commute with these triangles, their boundary maps,
proper pushforward, pullback where defined, and finite-extension
restriction and trace.  For the structural map \(f:X\to\Spec\Z\),
the definition is compatible with
\[
 R\Gamma_{\rm abs}(X,M)=
 R\Gamma_{\rm abs}(\Spec\Z,Rf_*M),
\]
and similarly for \(\mathsf S_\ell\) and for compact support.
The existence of these compatible extensions to arbitrary
constructible geometric coefficients and arbitrary finite-type
bases is an \emph{additional conjectural assertion}.  It is not
deduced merely by naming local subgroups of \(H^1\).
\end{enumerate}
The first two clauses specify the local and global theories that
the extension must preserve.  The third states the missing gluing
and functoriality assertion expressly.  Thus this is an existential
conjecture about complexes, maps and compatibilities, rather than
an assertion that equal dimensions suffice.  The arithmetic
degree-zero and degree-one clauses are Bloch--Kato--Fontaine;
the finite-field and gluing clauses are the additional requirements
needed for the present absolute-cohomology question
\cite[Conjecture~5.3]{BBlochKato1990}
\cite[\S\S6.5--6.7]{BFontaine1992}
\cite[Conjecture~8.3 and Theorem~8.4]{BGeisser2004}.
\end{definition}

\begin{proposition}[The two specializations of the corrected package]
\label{prop:B-comparison-specializations}
\status{Proved conditional implications; the Tate number-ring specialization is unconditional}
The finite-field clauses of
Definition~\ref{def:B-corrected-comparison-package} imply every
\emph{compatible} comparison and cycle-space assertion proved in
B.1, including the odd classes and the extension-degree factor.
Its number-ring clauses specialize, for every \(M=\Q(n)\),
to the unconditional table in
Corollary~\ref{cor:B-Hf-all-degrees}.  No implication from the
arithmetic degree-one clause alone to the finite-field Beilinson
clause is used.
\end{proposition}
\begin{proof}
For finite fields, the universal three-part cycle hypothesis is
exactly the hypothesis of Geisser's Theorem~8.4.  It gives the
integral comparisons and finite generation; Theorem~7.1 gives
the rational motivic splitting and the cup-product matrix for
\(e_q\).  Purity and semisimplicity then give the two cycle
spaces in degrees \(2n,2n+1\), and restriction of the group
homomorphism \(\phi_q\mapsto1\) gives the factor \(m\).
These are precisely the proofs of
Theorem~\ref{thm:B-Geisser-gate} and
Proposition~\ref{prop:B-e}; they apply to projectors and Tate
twists by functoriality.  The contradictory algebraically closed
vanishing assertion is excluded, as
Proposition~\ref{prop:B-continuity} requires.

For number rings, the maps \(c^0_\ell,c^1_\ell\) for Tate
motives are the invariant, Kummer and Borel--Soul\'e maps already
proved in B.2.  Their dual-twist counterparts are among the same
Tate calculations.  Applying the two Poitou--Tate exact sequences
and choosing their vector-space splittings gives (B.16), hence
every degree of (B.17a).  Localization adds only the two
finite-residue-field lines when \(n=1\), giving (B.17).
This proves the required arithmetic specialization for all
\(n\), all \(\ell\), and all finite \(S\).  For a general
prescribed dual pair the same implication is exactly
Proposition~\ref{prop:B-general-rational-gate}; it is conditional
on its \(c^0,c^1\) hypotheses.  None of these deductions
establishes the conjectural gluing clause or the full rational
duality requested in B.5.
\end{proof}

\begin{proposition}[Scholbach's ordinary Hom theory is not the full absolute theory]
\label{prop:B-Scholbach-finite-support}
\status{Proved counterexample to the full identification; the source's conditional comparison is retained}
The ordinary motivic cohomology functor used in Scholbach's
\emph{f-cohomology and motives over number rings} does not satisfy
the full finite-field part of the corrected comparison package.
This failure already occurs on an explicitly specified object
of its stated category over \(\Spec\Z\).
\end{proposition}
\begin{proof}
Fix primes \(p\ne\ell\), let
\(i:\Spec\F_p\hookrightarrow\Spec\Z\), and take
\(N=i_*\mathbf1_{\F_p}\).
Scholbach's Definition~1.3 includes this geometric motive, and
Definition~1.7 defines its cohomology by ordinary derived
\(\Hom\).  Adjunction in Axiom~1.1 and the \(K\)-theory
identification in Axiom~1.8 give
\[
 \begin{split}
 H^1_{\rm Scholbach}(\Spec\Z,N)
 &=\Hom_{DM(\F_p)}(\mathbf1,\mathbf1[1])\\
 &=H^1_{\BM}(\F_p,\Q(0))
   =K_{-1}(\F_p)^{(0)}_\Q=0.
 \end{split}
\]
The last equality is the negative \(K\)-group vanishing of a
field; it is also the degree-one motivic point calculation.
On the other hand, proper pushforward along \(i\) and the
Weil point computation give
\[
 H^1_{\rm abs}(\Spec\Z,N)=\Q e_p,
 \qquad
 H^1_{\Bcts}(\F_p,\Q_\ell)
   =\Hom_{\Bcts}(\widehat\Z,\Q_\ell)=\Q_\ell.
\]
For the last equality, a continuous homomorphism is determined
by the image of \(1\), and every element of \(\Q_\ell\)
defines such a homomorphism through
\(\widehat\Z\twoheadrightarrow\Z_\ell\).
Thus the realization map has zero source and a nonzero target.
Discrete etale \(\Q_\ell\)-cohomology is not being used.
This counterexample concerns the proposed full identification,
not Scholbach's conditional Theorems~6.8, 6.11 and~6.13 for
generic f-cohomology.  Their axioms and their explicitly omitted
crystalline condition remain exactly as stated in the preceding
source audit~\cite[Definitions~1.3, 1.7; Axioms~1.1, 1.8;
Remark~6.2]{BScholbach2015}.
\end{proof}


\subsection{B.3: abelian schemes, local coranks, and the certified counterexample}
Let \(A\) be an abelian scheme on
\(X=\Spec\mathcal O_K[1/N]\), with \(\ell\mid N\), and write
\(A^t\) for its dual.  Let
\(\mathcal S_\ell(A^t)=\varprojlim_a\operatorname{Sel}_{\ell^a}(A^t/K)\)
be the pro-\(\ell\) Selmer group, and let \(D\) denote Pontryagin
duality to \(\Q_\ell/\Z_\ell\).  At real places the local degree-zero
group used in this duality is the component group, as in local Tate
duality.

\begin{theorem}[The exact support sequence and corank]\label{thm:B-Sha-corank}
\status{Proved}
There is an exact sequence
\[
\begin{split}
0\longrightarrow\Sha(A/K)[\ell^\infty]
\longrightarrow H^1_{\et}(X,A)[\ell^\infty]
\longrightarrow\bigoplus_{v\mid N\ \mathrm{or}\ v\mid\infty}
 H^1(K_v,A)[\ell^\infty]
\longrightarrow\mathcal S_\ell(A^t)^D .
\end{split}\tag{B.19}
\]
There is no assertion of surjectivity onto the last term.
Put \(d=[K:\Q]\dim A\),
\(s=\Bcorank_{\Z_\ell}\Sha(A/K)[\ell^\infty]\), and let
\[
 t=\operatorname{rank}_{\Z_\ell}\operatorname{im}
 \left(\mathcal S_\ell(A^t)\longrightarrow
 \prod_{v\mid\ell}A^t(K_v)^{\wedge}_\ell\right).
\]
Then
\[
 \Bcorank_{\Z_\ell}H^1_{\et}(X,A)[\ell^\infty]=d+s-t.
 \tag{B.20}
\]
If \(\Sha(A/K)[\ell^\infty]\) is finite, then \(t\) is the
dimension of the \(\ell\)-adic closure of the Mordell--Weil group
\(A^t(K)\) inside its local points.  It must not in general be
replaced by the abstract rank \(r=\operatorname{rank}A(K)\).
\end{theorem}
\begin{proof}
The generic-fibre map embeds \(H^1(X,A)\) into \(H^1(K,A)\).
Its image consists of classes locally trivial at the finite places
outside \(N\): over those henselian integer rings proper smooth
extension and Lang's theorem give \(H^1(\mathcal O_v,A)=0\).
The converse follows by extending a locally trivial torsor across the
missing codimension-one points and descent on the Dedekind base.
The unconditional Cassels--Tate sequence of Gonz\'alez-Avil\'es and
Tan, Main Theorem, is
\[
0\longrightarrow\Sha(A)[\ell^\infty]\longrightarrow H^1(K,A)[\ell^\infty]
\longrightarrow\bigoplus_v H^1(K_v,A)[\ell^\infty]
\longrightarrow\mathcal S_\ell(A^t)^D\longrightarrow0.
\]
Taking the inverse image of the direct sum supported on \(N\) and
infinity proves (B.19).  This source is needed because replacing
the last term by \(A^t(K)^{\wedge,D}_\ell\) without a finiteness
assumption is unjustified~\cite[Main Theorem and equation~(6)]{BGonzalezTan2007}.

Local Tate duality identifies the dual of each local \(H^1\) with
\(A^t(K_v)^{\wedge}_\ell\).  Its rank is
\([K_v:\Q_\ell]\dim A\) for \(v\mid\ell\), and zero at all
other finite or real places.  The dual of the last map of (B.19)
is precisely the local-point map defining \(t\), so its image has
corank \(t\).  Additivity of coranks in the resulting short exact
sequences gives (B.20).  Finally the Selmer Kummer sequences give
\[
0\longrightarrow A^t(K)^{\wedge}_\ell
\longrightarrow\mathcal S_\ell(A^t)
\longrightarrow T_\ell\Sha(A^t/K)\longrightarrow0.
\]
The dual abelian varieties are isogenous via a polarization, so their
Sha groups have the same corank.  If that corank is zero, these
cofinitely generated groups are finite and their Tate modules vanish.
This proves the last assertion.
\end{proof}

The distinction between \(t\) and \(r\) is already effective over \(\Q\).
The second curve in the attached computation is
\[
 E':y^2+y=x^3+x^2-2x,
 \qquad \Delta(E')=389.
\]
The certified output is \(\operatorname{ellrank}(E')=[2,2,0,
[[-2,0],[4,8]]]\), with trivial torsion.  Its two-Selmer group
therefore has dimension two and \(\Sha(E')[2]=0\).  On
\(\Spec\Z[1/(2\cdot389)]\), we have \(d=1,s=0,r=2\), while
\(t=1\).  To see the last equality without a numerical logarithm,
a non-torsion rational point cannot become torsion in \(E'(\Q_2)\),
and its multiple in the formal-group neighborhood has nonzero
\(2\)-adic logarithm; the closure it generates has rank one.
Consequently (B.20) gives corank zero.  The putative formula
\(d-r+s\) would give minus one and is false.  The same
\texttt{x011.gp} file records all these data.

\begin{proposition}[Identical global ranks, different arithmetic coranks]
\label{prop:B-paired-corank}
\status{Refuted rank-only determination; explicit corrected coranks proved}
Put \(S=\Spec\Z[1/(2\cdot11\cdot37\cdot389)]\), and let
\[
\begin{aligned}
 E_{11}&:y^2+y=x^3-x^2-10x-20,\\
 E_{37}&:y^2+y=x^3-x,\\
 E_{389}&:y^2+y=x^3+x^2-2x.
\end{aligned}
\]
The abelian schemes
\(A=E_{389}\times E_{11}\) and \(B=E_{37}\times E_{37}\)
over the same base \(S\) have
\[
\begin{gathered}
 \dim A=\dim B=2,\qquad
 \operatorname{rank}A(\Q)=\operatorname{rank}B(\Q)=2,\\
 \Sha(A/\Q)[2^\infty]=\Sha(B/\Q)[2^\infty]=0.
\end{gathered}
\]
Nevertheless their arithmetic cohomology groups have distinct
coranks:
\[
 \operatorname{corank}_{\Z_2}H^1(S,A)[2^\infty]=1,
 \qquad
 \operatorname{corank}_{\Z_2}H^1(S,B)[2^\infty]=0.
\]
Thus even fixing the number field, coefficient prime, base,
dimension, Mordell--Weil rank and entire primary Sha group does
not determine the requested arithmetic corank.  The local closure
contribution \(t\) in Theorem~\ref{thm:B-Sha-corank} is essential.
\end{proposition}
\begin{proof}
The complete algebraic calculation is
\texttt{part\_b/code/paired\_corank.gp}, with full output in
\texttt{part\_b/output/paired\_corank.txt}.  It uses PARI~2.19.1
and no analytic-rank or BSD calculation.  Global reduction gives
conductors \(11,37,389\) and discriminants \(-11^5,37,389\),
respectively, so all three are elliptic schemes over \(S\).
The exact outputs are
\[
\begin{array}{c|c|c|c}
 E&\operatorname{ellrank}(E)&E(\Q)_{\rm tors}&
             \dim_{\F_2}\operatorname{Sel}_2(E/\Q)\\ \hline
 E_{11}&[0,0,0,[]]&\Z/5&0\\
 E_{37}&[1,1,0,[[-1,0]]]&0&1\\
 E_{389}&[2,2,0,[[-2,0],[4,8]]]&0&2.
\end{array}
\]
For precision, the third \texttt{ellrank} entry is
\(s_0=\dim_{\F_2}\Sha[2]/2\Sha[4]\); it is not in
general \(\dim\Sha[2]\).  The official algorithm description
states that the two-Selmer dimension \(C\), rational two-torsion
dimension \(T\), and \(s_0\) are unconditional, with
\[
 r_2=C-T-s_0,\qquad
 C=T+\operatorname{rank}E(\Q)+\dim_{\F_2}\Sha[2].
\]
Here the lower and upper rank bounds coincide, \(T=s_0=0\),
so these identities give \(\Sha[2]=0\) in every case
\citep[entries \texttt{ellrank} and \texttt{ell2cover}]{BPari2026}.

There is a direct additional certificate of that last assertion.
The program outputs a basis of all everywhere locally soluble
two-covers, with its exact rational maps to the elliptic curve.
The quartics are empty for \(E_{11}\), are
\[
       v^2=u^4-4u^3+6u^2+1
\]
for \(E_{37}\), and are
\[
\begin{aligned}
 v^2&=u^4-4u^3-2u^2+8u+1,\\
 v^2&=u^4-4u^3+10u^2+1
\end{aligned}
\]
for \(E_{389}\).  Every listed basis cover has the rational
point \((u,v)=(0,1)\).  Under the recorded covering maps these
points give \((-1,0)\), \((4,8)\), and \((-2,0)\), respectively.
The program verifies each map by an exact polynomial identity.
Thus every two-Selmer basis class has trivial image in Sha.
The two-descent exact sequence proves \(\Sha[2]=0\) directly,
and a primary torsion group with no elements of order two is zero.
Products preserve both the Mordell--Weil decomposition and Sha,
which proves all the stated global equalities.

It remains to compute the local closure ranks, without a numerical
logarithm.  Each \(E(\Q_2)\) contains an open formal subgroup
isomorphic to \(4\Z_2\), by the formal logarithm on depth
\(r=2>v_2(2)/(2-1)=1\)
\citep[Theorem~IV.6.4(b), p.~132, and Proposition~VII.2.2,
p.~191]{BSilverman2009Local}.
It has finite index because \(E(\Q_2)\) is compact.
A non-torsion rational point remains non-torsion over \(\Q_2\):
an equation \([m]P=0\) after field extension is the same equation
over \(\Q\).  A nonzero multiple in the formal subgroup therefore
has a nonzero logarithm.  The closure of its integer multiples is
its logarithm times \(\Z_2\), of rank one.  Since the ambient
local group itself has rank one, the closures of \(E_{37}(\Q)\)
and \(E_{389}(\Q)\) each have rank one; the finite group
\(E_{11}(\Q)\) has closure rank zero.

The rational and local point groups of a product are products,
and closure commutes with a finite product.  Consequently
\(t(A)=1+0=1\), whereas \(t(B)=1+1=2\).  The product
principal polarizations identify these varieties with their duals.
Because their primary Sha groups vanish, the pro-Selmer groups
in Theorem~\ref{thm:B-Sha-corank} are precisely the completions
of these rational-point groups.  At the other inverted finite
primes the local \(2\)-primary cohomology has corank zero.
Both varieties have \(d=2\) and \(s=0\), so formula (B.20)
gives \(2-1=1\) and \(2-2=0\), as claimed.  The latter group
is finite by the finite-cotype assertion already proved there.
\end{proof}


\begin{example}[The counterexample in R6, without a BSD assumption]\label{ex:B-X011}
\status{Refuted; corrected statement proved}
Take
\[
 E=X_0(11):\quad y^2+y=x^3-x^2-10x-20.
\]
Its discriminant is \(-11^5\), its conductor is \(11\), and it has
good reduction over \(\Z[1/11]\).  The reproducible algebraic
two-descent in \texttt{part\_b/code/x011.gp} gives
\[
 \operatorname{ellrank}(E)=[0,0,0,[]],\qquad
 E(\Q)_{\rm tors}\cong\Z/5,
\]
with torsion generator \((5,5)\).  Thus its rank is certified to be
zero; this uses no analytic-rank computation or BSD.  PARI's official
description of \texttt{ellrank} specifies that the lower and upper
bounds and the two-descent data are unconditional.  The model of
\(X_0(11)\) is also derived explicitly in Elkies's October~13, 2021
lecture notes~\cite{BElkies2021,BPari2026}.

For any prime \(\ell\), let \(X=\Spec\Z[1/(11\ell)]\).
The Kummer sequence gives
\[
0\longrightarrow E(\Q)\otimes\Q_\ell
\longrightarrow H^1_{\Bcts}(X,V_\ell E)
\longrightarrow V_\ell H^1_{\et}(X,E)\longrightarrow0.
\]
The first term is zero.  The global Euler characteristic of
\(V_\ell E\), whose complex-conjugation eigenspaces both have
dimension one, is \(-1\).  Its \(H^0\) is zero, so
\(\dim H^1=1+\dim H^2\ge1\).  It follows that
\(H^1_{\et}(X,E)[\ell^\infty]\) is infinite for every \(\ell\).
For \(\ell=2\) the zero two-Selmer calculation additionally gives
\(\Sha(E)[2]=0\), hence \(\Sha(E)[2^\infty]=0\); then (B.20)
gives exact corank one.  All hypotheses of the boxed assertion on
12:75 hold, including invertibility of \(\ell\) on \(X\), and its
conclusion fails.
\end{example}

\paragraph{The corrected boxed assertion and its precise remaining conjecture.}
Set
\[
 H^1_{\rm loc\,0}(X,A)[\ell^\infty]=
 \ker\left(H^1(X,A)[\ell^\infty]\longrightarrow
 \bigoplus_{v\mid\ell}H^1(K_v,A)[\ell^\infty]\right).
\]
The local groups at the remaining finitely many places in (B.19) are
finite, so this kernel is finite if and only if
\(\Sha(A/K)[\ell^\infty]\) is finite.  The corrected statement is
therefore a Sha-finiteness conjecture, rather than a theorem obtained
by putting \(\ell\) in the denominator.  At the rational Tate-module
level, local Kummer theory identifies the finite condition with
the image of \(A(K_v)\widehat\otimes\Q_\ell\) (Bloch--Kato,
Example~3.11).  The global Selmer Kummer sequence becomes
\[
0\longrightarrow A(K)\otimes\Q_\ell
\longrightarrow H_f^1(K,V_\ell A)
\longrightarrow V_\ell\Sha(A/K)\longrightarrow0.\tag{B.21}
\]
Thus the rational-points assertion, with the restored finite local
conditions, is equivalent exactly to the finite boxed kernel.
This is the content of Bloch--Kato's Proposition~5.4, here derived
from Kummer theory.  The motive written \(T_{\rm mot}(A)\) on
12:74 has realization \(V_\ell A\) and weight \(-1\), so the
transcription's weight \(+1\) must also be corrected.

\begin{proposition}[An arithmetic relative-curve counterexample to 12:79]
\label{prop:B-Brauer-arithmetic}
\status{Refuted; corrected statement proved}
Let \(S=\Spec\Z[1/22]\) and let \(f:X\to S\) be the smooth
proper genus-one curve given by the above model of \(X_0(11)\), with
its zero section.  Then
\[
 \left(\Br(X)/f^*\Br(S)\right)[2^\infty]
 \simeq H^1_{\et}(S,E)[2^\infty]
\]
has \(\Z_2\)-corank one.  Thus surjectivity modulo finite groups
of \(\Br(S)\to\Br(X)\) is false for arithmetic bases, even for
a smooth proper relative curve with a section.
\end{proposition}
\begin{proof}
The zero section splits the relative degree map, giving
\(R^1f_*\Gm=\BPic_{X/S}\simeq E\times\Z\) as etale sheaves.
Proper base change identifies \(R^2f_*\mu_{2^a}\) with the
constant sheaf \(\Z/2^a\); the degree-one section class maps to its
generator.  The relative Kummer sequence therefore shows that
\((R^2f_*\Gm)[2^a]=0\) for every \(a\).  The Leray
spectral sequence, and the splitting on pullback supplied by the
section, now give on the torsion subgroup the usual exact edge
sequence
\[
 0\longrightarrow\Br(S)\longrightarrow\Br(X)
 \longrightarrow H^1(S,\BPic_{X/S})
 \longrightarrow H^3(S,\Gm),
\]
as far as its \(2\)-primary component is concerned.  Here
\(H^1(S,\Z)=0\), since a continuous homomorphism from the profinite
etale fundamental group to the discrete torsion-free group \(\Z\)
is zero.  Moreover \(H^3(S,\Gm)=0\): Milne's
Proposition~II.2.1 gives its description as the cokernel of the
sum of local invariants at the omitted places, and the summand at
the omitted finite place \(2\) already surjects onto \(\Q/\Z\).
This is stated explicitly in Remark~II.2.2(a)~\cite{AMilneADT}.
All groups in \(H^1(S,E)\) are torsion, and \(\Br(X)\) is
torsion because \(X\) is regular.  Taking primary components of
the resulting exact sequence is therefore exact.  The preceding
certified two-descent and (B.20) give the asserted corank one.
\end{proof}

\paragraph{The finite-field Hochschild--Serre correction.}
Let now \(S/\F_q\) be smooth and geometrically connected,
\(A/S\) an abelian scheme, \(\ell\ne p\),
\(\Gamma=\Gal(\overline{\F}_q/\F_q)\), and
\(C=A(\bar S)\).  The actual five-term sequence is
\[
0\longrightarrow H^1(\Gamma,C)\longrightarrow H^1(S,A)
\longrightarrow H^1(\bar S,\bar A)^\Gamma
\xrightarrow{d_2}H^2(\Gamma,C).\tag{B.29}
\]
The last term deleted on 12:76--77 cannot be discarded solely
because \(\Gamma\) has torsion cohomological dimension one: \(C\)
is not a torsion module.  For a lattice \(M\) with continuous
\(\Gamma\)-action, the exact calculation is
\[
0\longrightarrow M^\Gamma\otimes\Q/\Z
\longrightarrow(M\otimes\Q/\Z)^\Gamma
\longrightarrow H^1(\Gamma,M)\longrightarrow0,
\qquad
H^2(\Gamma,M)=(M\otimes\Q/\Z)_\Gamma.\tag{B.30}
\]
Indeed \(M\otimes\Q\) has zero higher continuous cohomology:
every cochain factors through a finite quotient and averaging is
exact over \(\Q\).  The rational-to-torsion exact sequence then
proves (B.30).  In particular \(H^1(\Gamma,M)\) is finite,
whereas \(M=\Z\) gives \(H^2(\Gamma,M)=\Q/\Z\).
The equality with \((M\otimes\Q/\Z)^\Gamma\) asserted for
\(H^1(\Gamma,M)\) on 12:76 misses the first term of (B.30).

Modulo the constant trace, Lang--Neron gives the finite-rank group
\(M\) on the page.  Lang's theorem and torsion cohomological
dimension one dispose of the trace's higher cohomology, so
\(H^1(\Gamma,C)\) is finite; they do not dispose of \(d_2\).
Write
\[
 W=H^1_{\Bcts}(\bar S,V_\ell A),\qquad
 M_\ell=C^{\wedge}_\ell\otimes_{\Z_\ell}\Q_\ell
\]
where \(C^{\wedge}_\ell=\varprojlim_a C/\ell^a C\), so the
divisible trace is removed.  Geometric Kummer
gives
\[
 0\longrightarrow M_\ell\longrightarrow W
 \longrightarrow V_\ell H^1(\bar S,\bar A)
 \longrightarrow0.\tag{B.31}
\]
Its invariant sequence includes the boundary
\((V_\ell H^1(\bar S,\bar A))^\Gamma\to(M_\ell)_\Gamma\).
Thus taking invariants in the bottom row of 12:77 is not
automatically right exact either.  These are two manifestations of
the same obstruction.

\begin{lemma}[Splitting off divisor classes on a smooth open variety]
\label{lem:B-divisor-splitting}
\status{Proved}
Let \(Y/\F_q\) be smooth, separated and of finite type, and let
\(N_Y\) be the image of \(\BPic(\bar Y)\otimes\Q_\ell\) in
\(H^2(\bar Y,\Q_\ell(1))\), for \(\ell\ne p\).  Then
\(N_Y\) is a \(\Gamma\)-equivariant direct summand.  No
semisimplicity or Tate conjecture for the complementary summand is
assumed.
\end{lemma}
\begin{proof}
First suppose \(Y\) is an open subset of a smooth projective
\(Y'\).  On \(Y'\), the intersection form, formed with an ample
line bundle, is nondegenerate on the rational Neron--Severi group.
Its orthogonal complement therefore splits the divisor-class
subspace in \(H^2(\bar Y',\Q_\ell(1))\).  This is the precise
linear argument in the proof of Yuan's Theorem~1.2; it uses the
Hodge index theorem, not the divisor Tate conjecture~\cite{BYuan2020}.
The kernel of restriction to \(Y\) is spanned by the classes of
the codimension-one components of \(Y'\setminus Y\), by the
Gysin sequence.  It lies inside that divisor summand.  Quotienting
the splitting therefore splits \(N_Y\) inside the image of
restriction; every line bundle on the regular open \(Y\) extends
as a divisor to the regular \(Y'\).  The remaining quotient of
\(H^2(\bar Y,\Q_\ell(1))\) has weights at least one, whereas
the image of restriction has weight zero.  The decomposition into
the relatively prime Frobenius factors of these two weights splits
the image as well.  The Gysin sequence and weight assertion are
given in Balkan--Schreieder's proof of Proposition~4.10 and
Lemma~3.11~\cite{BBalkanSchreieder2026}.

In general an alteration supplies a smooth \(V\), open in a
smooth projective variety, and a proper generically finite
\(\tau:V\to Y\), of degree \(a>0\).  The trace on cohomology
satisfies \(\tau_*\tau^*=a\), and it takes divisor classes to
divisor classes.  These assertions are the projection formula and
cycle compatibility in Balkan--Schreieder's Lemma~2.2.  If \(r_V\)
is the equivariant retraction constructed above, then
\[
 r_Y=a^{-1}\tau_*r_V\tau^*
 :H^2(\bar Y,\Q_\ell(1))\longrightarrow N_Y
\]
restricts to the identity on \(N_Y\).  This proves the result.
The alteration and trace argument works for separated smooth
schemes after taking a proper compactification and a projective
alteration; it does not require \(Y\) itself to be proper.
\end{proof}

\begin{proposition}[The finite-field chain over a regular finite-type base]
\label{prop:B-finite-chain-proper}
\status{Proved}
For the smooth separated base \(S\) above, the natural map
\[
H^1(S,A)[\ell^\infty]\longrightarrow
H^1(\bar S,\bar A)[\ell^\infty]^\Gamma
\]
has finite kernel and cokernel.  Consequently the following four
conditions are equivalent:
\begin{enumerate}
\item \(H^1(S,A)[\ell^\infty]\) is finite;
\item the Kummer map \(A(S)\otimes\Q_\ell\to
H^1_{\Bcts}(S,V_\ell A)\) is an isomorphism;
\item the same map to \(W^\Gamma\) is an isomorphism;
\item \(H^1(\bar S,\bar A)[\ell^\infty]^\Gamma\) is finite.
\end{enumerate}
These are equivalences, without an assumption that their equivalent
finiteness assertions hold.  If \(S\) is projective and \(A\) is
constant, these assertions of finiteness hold unconditionally.
\end{proposition}
\begin{proof}
Put \(B=A^t\), and choose an integer \(a>1\).  Multiplication
\([a]\) on \(B/S\) acts by \(a^j\) on \(R^jf_*\Q_\ell\).
The distinct eigenvalues force all Leray differentials to vanish
and split the degree-two filtration.  The \(a\)-eigenspace of
\(H^2(\bar B,\Q_\ell(1))\) is consequently
\(H^1(\bar S,R^1f_*\Q_\ell(1))=W\).
On relative Picard groups, \([a]^*\) acts by \(a\) on
\(\BPic^0_{B/S}=A\), and by \(a^2\) on relative
Neron--Severi classes; on line bundles pulled back from the base
it acts by one.  The zero section makes
\(\BPic(B)\to\BPic_{B/S}(S)\) surjective.  Thus the
\(a\)-eigenspace of the divisor-class image \(N_B\) is exactly
the Kummer subspace \(M_\ell\) of (B.31).  This identification
also follows at the cocycle level: the line bundle attached to a
section of \(A\) is the pullback of the Poincare bundle, and
choosing its local \(\ell^r\)-division sections gives precisely
the section's Kummer cocycle in \(A[\ell^r]=R^1f_*\mu_{\ell^r}\).

Let \(e_1\) denote the polynomial projector onto the
\(a\)-eigenspace.  Lemma~\ref{lem:B-divisor-splitting} gives a
\(\Gamma\)-equivariant retraction \(r\) onto \(N_B\).
The map \(e_1re_1:W\to M_\ell\) is the identity on
\(M_\ell\), even though \(r\) need not commute with \([a]^*\).
This splits (B.31) equivariantly and makes its invariant sequence
right exact.  All projectors are rational, so the argument also
applies to \(\ell=2\).

Lang's theorem and Lang--Neron identify
\(A(S)\otimes\Q_\ell=(M_\ell)^\Gamma\).
The kernel in continuous Hochschild--Serre is
\(H^0(\bar S,V_\ell A)_\Gamma=0\).  Indeed evaluation at a
closed point of \(S\) embeds this global-section space in the
Tate module of that abelian fibre; a power of Frobenius on it has
weight \(-1\), excluding eigenvalue one.  The diagram of global
Kummer and the now exact invariant sequence of (B.31) therefore
identifies
\(V_\ell H^1(S,A)\) with
\((V_\ell H^1(\bar S,\bar A))^\Gamma\).
The finite sheaf \(A[\ell]\) and Kummer prove cofiniteness of
both torsion groups.  Their natural comparison consequently has
finite kernel and cokernel, and all four conditions are equivalent.
This proves, in particular, that the missing transgression in
(B.29) has finite image in the stated regular-base setting.

For constant \(A=A_0\times S\), sections modulo the finite
constant group identify, up to a finite ambiguity when choosing an
Albanese base point, with \(\Hom(\operatorname{Alb}_S,A_0)\).
Moreover
\(W=\Hom(V_\ell\operatorname{Alb}_S,V_\ell A_0)\).
Tate's homomorphism theorem for finite-field abelian varieties
identifies its invariants with those homomorphisms tensored with
\(\Q_\ell\).  Condition (3) follows~\cite[Main Theorem, \S1]{BTate1966}.
\end{proof}

Keller's Theorem~2.7.19 and Remark~2.7.20(d) identify the same
proper-base finiteness assertion with equality of algebraic and
analytic rank.  The height pairing is nondegenerate in the required
case; it cannot be silently omitted from the theorem's printed
condition before that remark is applied~\cite{BKeller2019}.
Thus 12:80 is a proved equivalence with BSD in this scope, not a
proof of BSD for an arbitrary nonconstant abelian scheme.

\begin{proposition}[Relative curves without a section]
\label{prop:B-relative-curve}
\status{Proved}
Let \(S/\F_q\) be smooth, separated and connected of finite type,
and let \(f:Y\to S\) be a smooth proper curve with geometrically
connected fibres and relative Jacobian \(J\).  For \(\ell\ne p\),
\[
 V_\ell\bigl(\Br(Y)/\operatorname{im}\Br(S)\bigr)
 \simeq V_\ell H^1(S,J).
 \tag{B.32}
\]
The kernel of \(\Br(S)[\ell^\infty]\to\Br(Y)[\ell^\infty]\)
is finite.  Consequently the following are equivalent:
\begin{enumerate}
\item \(\Br(S)\to\Br(Y)\) is surjective modulo a finite
\(\ell\)-primary group;
\item \(H^1(S,J)[\ell^\infty]\) is finite;
\item \(H^1(\bar S,\bar J)[\ell^\infty]^\Gamma\) is finite.
\end{enumerate}
If \(f\) has a section, the natural edge map gives the stronger
integral isomorphism
\[
 (\Br(Y)/\Br(S))[\ell^\infty]
 \simeq H^1(S,J)[\ell^\infty].
\]
\end{proposition}
\begin{proof}
Choose a closed point of degree \(d>0\) on the generic fibre.
Its closure in the regular scheme \(Y\) is a Cartier divisor;
its line bundle \(L\) has degree \(d\) on every geometric
fibre, because degree in a proper smooth family of curves is locally
constant.  The closure is not asserted to be finite over \(S\).
Cup product with \(c_1(L)\) induces the isomorphism
\(R^0f_*\Q_\ell\to R^2f_*\Q_\ell(1)\), namely
multiplication by \(d\).  Deligne's Proposition~2.4, equation
(2.5), therefore applies directly; for a curve this is the only
nontrivial Lefschetz condition.  In particular its Leray filtration
in degree two is split over \(\Q_\ell\), with successive terms
\[
 H^2(S,\Q_\ell(1)),\qquad
 H^1(S,V_\ell J),\qquad \Q_\ell.
\]
The first term is the image of \(f^*\), and the last term is
split by \(c_1(L)/d\)
\cite[Proposition~2.4, (2.5), pp.~111--112]{c_deligne1968}.

Put \(P=\BPic_{Y/S}=R^1f_*\Gm\).  Etale locally the
curve has a section, so there is an exact sequence of sheaves
\(0\to J\to P\xrightarrow{\deg}\Z\to0\).
The image of \(L\) has degree \(d\), hence the boundary
\(\Z\to H^1(S,J)\) has finite cyclic image of order dividing
\(d\).  Moreover \(H^1(S,\Z)=0\).  Thus
\(H^1(S,J)\to H^1(S,P)\) has finite kernel and is surjective.
The obstruction to lifting a section of \(P\) to a line bundle
on \(Y\) is in
\(\ker(\Br(S)\to\Br(Y))\).  This kernel is killed by
\(d\): restrict to the function field of the regular integral
\(S\), then to the chosen degree-\(d\) point, and apply
Brauer corestriction.  Injectivity of \(\Br(S)\to\Br(k(S))\)
returns the assertion to \(S\).  Its \(\ell\)-primary part
is finite, because Kummer and finiteness of
\(H^2(S,\mu_\ell)\) make \(\Br(S)[\ell^\infty]\)
cofinitely generated.

It remains to identify the middle algebraic subspace, rather than
infer anything from the rank of a Leray term alone.  The relative
Chern class of a degree-zero relative line bundle is its Kummer
class in \(H^1(S,V_\ell J)\): this follows by choosing etale
local lifts and \(\ell^r\)-division sections and comparing the
resulting transition cocycles.  The bounded Brauer obstruction just
proved permits every section of \(J\) to be lifted after
multiplication by \(d\).  The divisor-class image in the
middle term is therefore exactly \(J(S)\otimes\Q_\ell\).
In the first term it is the divisor-class image from \(S\),
and in the last term it is all of \(\Q_\ell\), generated by
\(L\).  Equivalently, one can apply this argument to the
\(\ell\)-adic completions of the Picard groups and then tensor
with \(\Q_\ell\); finite kernels and the bounded obstruction
contribute no term.  Quotienting the three Leray terms by these
Chern-class images, the two Kummer sequences give
\[
 \frac{V_\ell\Br(Y)}{V_\ell\Br(S)}
 \simeq
 \frac{H^1(S,V_\ell J)}{J(S)\otimes\Q_\ell}
 \simeq V_\ell H^1(S,J).
\]
This proves (B.32).  All torsion groups in this comparison are
cofinitely generated on their \(\ell\)-primary parts, again by
Kummer and finite etale cohomology.  Thus vanishing of these rational
Tate modules is equivalent to finiteness, proving (1) equivalent
to (2); Proposition~\ref{prop:B-finite-chain-proper} gives (3).

When there is a section, \(P=J\times\Z\).  On strict local
bases the curve has a section, and proper base change and Kummer
show \((R^2f_*\Gm)[\ell^\infty]=0\).  The resulting
Leray edge sequence on primary torsion has target
\(H^1(S,P)\).  Its next differential lands in the kernel of
\(H^3(S,\Gm)\to H^3(Y,\Gm)\), which is zero because
restriction to the section splits pullback.  This proves the final
integral assertion without assuming \(H^3(S,\Gm)=0\).
\end{proof}

\begin{remark}[A source error avoided in the preceding proof]\label{rem:B-Keller-closure-error}
\status{Refuted; corrected statement proved}
Keller's Lemma~2.17 claims that the closure of a generic-fibre point
of a proper relative curve is finite over an integral base.  This
is false in the printed generality.  Over any finite field take
\(S=\mathbb A^2_{x,y}\), \(Y=\mathbb P^1_{[u:v]}\times S\),
and the generic point \([u:v]=[x:y]\).  Its closure is
\(uy-vx=0\), the blowup of the origin, whose fibre over
\((0,0)\) is the whole \(\mathbb P^1\).  On its two standard
charts the coordinate rings are \(k[x,t]\), \(y=xt\), and
\(k[y,s]\), \(x=ys\); it is integral and has no additional
irreducible component.  This explicitly contradicts the lemma's
quasi-finiteness step.  The degree of the associated line bundle,
which is all the proof above uses, remains one on every fibre.
Thus the section-free part of the preceding proposition is not
inherited from that lemma or from its use in Corollary~2.18 and
Theorem~4.27~\cite[Lemma~2.17]{BKeller2016}.
\end{remark}

The proposition proves the ``via des jacobiennes'' equivalence on
12:79 over its stated regular base, without adding a section
hypothesis.  It does not prove the equivalent finiteness assertions
for every nonconstant Jacobian.  When \(S\) is a smooth projective
curve, \(\Br(S)=0\), and the assertion at a fixed
\(\ell\ne p\) is precisely finiteness of
\(\Br(Y)[\ell^\infty]\), equivalently its divisor Tate
conjecture by Yuan's Theorems~1.2 and 1.3(4).
No finiteness of the whole torsion group is inferred merely from
finiteness of individual primary parts.  The separate argument in
Theorem~\ref{thm:B-Brauer-prime-uniform} proves finiteness of the
whole prime-to-\(p\) quotient under the same surface-Tate
condition.  The unqualified assertion
over nonproper bases is refuted below by
Theorem~\ref{thm:B-Brauer-characteristic}.

For a nonproper smooth curve base, Qin's Propositions~2.4--2.5 and
Lemma~2.6 provide an independent proof of the comparison, including
possible bad reduction at the omitted points~\cite{BQin2024}.
The argument above controls higher-dimensional nonproper regular
bases too, by splitting only the algebraic subspace, a weaker and
proved property than semisimplicity of all of \(W\).
\paragraph{The chain over every finite-type base.}
The smooth-base splitting above is enough to remove regularity from
the Kummer comparison itself.  One needs injectivity on sections,
not descent of a family of cohomological retractions along a
hypercover.

\begin{lemma}[Sections and their completion on a reduced base]
\label{lem:B-sections-reduced}
\status{Proved}
Let \(S/\F_q\) be reduced and of finite type, without a
separatedness assumption, and let \(A/S\) be an abelian scheme.
Put \(C_S=A(\bar S)\).  Then \(C_S/(C_S)_{\rm tors}\)
is finitely generated, its Galois action has finite image, and
the natural map gives an isomorphism
\[
 C_S\otimes_{\Z}\Q_\ell
 \xrightarrow{\sim}
 (C_S)^{\wedge}_\ell\otimes_{\Z_\ell}\Q_\ell
 \qquad (\ell\ne p).
\]
There exists a finite disjoint union \(V\) of smooth separated
finite-type \(\F_q\)-schemes with a jointly surjective map
\(v:V\to S\).  Pullback embeds \(C_S\) in \(C_V\)
and embeds their quotients by torsion.
\end{lemma}
\begin{proof}
Take a finite affine cover of \(S\).  Alter each reduced
irreducible component of each member of the cover.  De Jong's
Theorem~4.1 and Remark~4.2 give a smooth separated alteration,
since \(\F_q\) is perfect; their disjoint union is the
required \(V\)~\cite{c_dejong1996}.  The construction and
surjectivity persist after extension to \(\overline{\F}_q\).
If a section becomes zero on \(\bar V\), its equalizer with
the zero section is a closed subscheme of \(\bar S\) whose
underlying set is all of \(\bar S\).  Its ideal is therefore
contained in the nilradical and is zero.  Thus pullback is
injective.  Applying the same argument to an integer multiple of
a section shows injectivity modulo torsion as well.

For each connected component of \(\bar V\), restriction to
its generic point embeds its sections into the points of an
abelian variety over a finitely generated regular extension of
\(\overline{\F}_q\).  Lang--Neron, in the exact form of
Conrad's Theorems~2.1 and~7.1, says that the quotient by the
constant trace is finitely generated~\cite{BConrad2006}.  All
\(\overline{\F}_q\)-points of an abelian variety over
\(\overline{\F}_q\) are torsion: the variety and each point
descend to a finite field.  The quotient of its function-field
points by all torsion is consequently finitely generated.  A
subgroup of a finitely generated abelian group is finitely
generated.  Applying this twice, first to \(C_V\) and then to
\(C_S\), proves the assertion about sections modulo torsion.
Representatives of finitely many generators are defined over one
finite extension of \(\F_q\), which proves the finite-image
assertion.

Evaluation at one geometric point of each connected component of
\(\bar S\) injects \(C_S[\ell^\infty]\) into a finite
sum of groups \(A_{\bar s}[\ell^\infty]\).  Indeed a section
of a finite etale group scheme which vanishes at one point of a
connected scheme vanishes everywhere.  Thus the Pontryagin dual
of \(C_S[\ell^\infty]\) is a quotient of a finite free
\(\Z_\ell\)-module; it is finitely generated.  Equivalently,
this torsion group is a sum of finitely many copies of
\(\Q_\ell/\Z_\ell\) and a finite group.  Its
\(\ell\)-adic completion has zero rationalization.  The
prime-to-\(\ell\) torsion is \(\ell\)-divisible and has
zero completion.  Finally the extension of the torsion subgroup
by the finite free group \(C_S/(C_S)_{\rm tors}\) splits as
an extension of abelian groups.  Computing completions in that
splitting proves the displayed natural isomorphism; its statement
does not depend on the splitting chosen.
\end{proof}

\begin{theorem}[The finite-field chain for arbitrary finite-type schemes]
\label{thm:B-finite-chain-all-bases}
\status{Proved}
Let \(S/\F_q\) be any scheme of finite type, let \(A/S\)
be an abelian scheme, and let \(\ell\ne p\), including
\(\ell=2\) when \(p\ne2\).  No regularity, reducedness,
properness or separatedness of \(S\) is required.  Set
\[
 W_S=H^1_{\Bcts}(\bar S,V_\ell A),\qquad
 N_S=A(\bar S)\otimes\Q_\ell.
\]
The geometric Kummer sequence
\[
 0\longrightarrow N_S\longrightarrow W_S
 \longrightarrow V_\ell H^1_{\et}(\bar S,\bar A)
 \longrightarrow0
\]
has a \(\Gamma\)-equivariant splitting.  The natural map
\[
 H^1_{\et}(S,A)[\ell^\infty]\longrightarrow
 H^1_{\et}(\bar S,\bar A)[\ell^\infty]^\Gamma
\]
has finite kernel and cokernel.  The four conditions in
Proposition~\ref{prop:B-finite-chain-proper} are therefore
equivalent for every such \(S\).
\end{theorem}
\begin{proof}
First suppose that \(S\) is reduced and take \(V\) as in
Lemma~\ref{lem:B-sections-reduced}.  Kummer maps commute with
pullback.  Write \(i:N_S\hookrightarrow N_V\) for the
injection of that lemma and \(v^*:W_S\to W_V\) for
cohomological pullback.  The proof of
Proposition~\ref{prop:B-finite-chain-proper}, component by
component, supplies a \(\Gamma\)-equivariant retraction
\(r_V:W_V\to N_V\).  The action of \(\Gamma\) on
\(N_V\) factors through a finite group \(G\).  Choose any
linear retraction \(b:N_V\to N_S\) of \(i\) and average:
\[
 b_G=\frac1{|G|}\sum_{g\in G}g\,b\,g^{-1}.
\]
The scalar \(|G|^{-1}\) is available in \(\Q_\ell\),
even when \(\ell\mid |G|\).  Hence
\[
 r_S=b_G\,r_V\,v^*:W_S\longrightarrow N_S
 \quad\hbox{satisfies}\quad r_S|_{N_S}=1.
\]
This proves the splitting.  It does not require \(v^*\) to be
injective on the whole of \(W_S\), nor compatible retractions
on intersections of the covering alterations.

For a nonreduced \(S\), filter its nilpotent ideal by powers,
so that successive closed immersions have square-zero ideals
\(I\).  On the small etale site, smoothness of \(A\) gives
a surjection from the sheaf of its sections to the sheaf on the
thickening's reduction.  Its kernel for a square-zero step is
\(\operatorname{Lie}(A)\otimes I\), an additive sheaf killed
by \(p\).  The cohomology exact sequence and induction show
that the kernels and cokernels on \(H^0\) and \(H^1\) are
of bounded \(p\)-power exponent.  Such maps are isomorphisms
on \(\ell\)-primary torsion and after tensoring sections
with \(\Q_\ell\).  The finite etale sheaves
\(A[\ell^a]\) and their cohomology are unchanged by a
nilpotent thickening.  Thus the Kummer sequence and its splitting
reduce to those already proved for \(S_{\rm red}\).

It remains to verify the passage from geometric splitting to the
stated finite comparison.  Norms from finite constant extensions
give
\[
 A(S)\otimes\Q_\ell=N_S^\Gamma.
\]
Indeed a finite set of section representatives is defined over
one finite extension, and averaging over that extension descends
the invariant rational combination; injectivity is immediate.
Furthermore
\(H^0(\bar S,V_\ell A)_\Gamma=0\).  To see this,
evaluate at a finite set of geometric points, one on each
connected component, lying over closed points of \(S\).
This is injective on sections of a lisse sheaf.  After a finite
constant extension it is equivariant into a finite sum of Tate
modules of abelian varieties over finite fields.  Their
geometric Frobenius eigenvalues have weight \(-1\), so no
positive power has eigenvalue one.  Continuous
Hochschild--Serre therefore identifies
\[
 H^1_{\Bcts}(S,V_\ell A)=W_S^\Gamma.
\]
Taking invariants in the split geometric Kummer sequence and
comparing with global Kummer gives the natural isomorphism
\[
 V_\ell H^1_{\et}(S,A)
 \xrightarrow{\sim}
 \bigl(V_\ell H^1_{\et}(\bar S,\bar A)\bigr)^\Gamma.
\]
The group \(H^1(S,A)[\ell]\) is a quotient of the finite
group \(H^1(S,A[\ell])\), and the same holds over
\(\overline{\F}_q\).  Finite-coefficient etale finiteness
also holds for a nonseparated finite-type scheme by a finite
affine cover and its Cech spectral sequence.  Consequently both
\(\ell\)-primary groups are cofinitely generated.  Taking
invariants preserves this property, and the Tate module of the
invariants is the invariant Tate module, since invariants commute
with inverse limits.  A map between cofinitely generated
\(\ell\)-primary torsion groups with an isomorphism on
rational Tate modules has finite kernel and cokernel.  This proves
the comparison and all four equivalences.
\end{proof}

The transcription of 12:79 itself explicitly retains the condition
that the base be regular.  The relative-curve proposition above
therefore proves the requested Jacobian--Brauer equivalence in
that stated scope.  The present theorem removes the former
singular-base gap in the Kummer chain; it does not prove its
equivalent finiteness conditions for every nonconstant abelian
scheme.


\begin{example}[Why the positive-degree argument retains its regularity hypothesis]
\label{ex:B-singular-infinite-period}
\status{Proved counterexample to an extension beyond the stated manuscript scope}
There is a reduced projective surface \(S/\F_2\) and a smooth
proper genus-one scheme \(f:X\to S\) for which every global line
bundle has relative degree zero.  Thus the positive-degree class
used in Proposition~\ref{prop:B-relative-curve} cannot be asserted
on an arbitrary singular base.  This example does not assert that
the Brauer quotient is infinite.

Take
\[
 E:\ y^2z+yz^2+x^3+x^2z=0\ \subset T=\mathbf P^2_{\F_2},
 \qquad O=[0:1:0],
 \qquad S=T\amalg_E T.
\]
The derivatives of the cubic are \(x^2,z^2,y^2+x^2\), with no
common projective zero, so \(E\) is an elliptic curve.  The
pushout is the explicit finite \(T\)-scheme
\[
 S=\operatorname{Spec}_T(\mathcal O_T\oplus\mathcal O_T(-3)),
 \qquad
 R[t]/(t(t-f_E))\simeq R\times_{R/(f_E)}R
\]
on each trivializing affine, where multiplication on the second
summand uses the cubic section \(f_E\).  It is reduced because
the displayed algebra embeds in the product of two domains, and
projective because it is finite over \(T\).

Let \(A=E\times S\).  The closed-cover sequence on the small
etale site is
\[
 0\longrightarrow A\longrightarrow i_{1*}A_T\oplus i_{2*}A_T
 \xrightarrow{\mathrm{res}_1-\mathrm{res}_2}i_{E*}A_E
 \longrightarrow0.
\]
For surjectivity it suffices to pass to a strict henselian local
ring.  A section into the smooth curve \(E\) on its closed
quotient lifts across the resulting henselian pair; see the Stacks
Project, Lemma~15.13.3~\cite{BStacksSmoothLift}.  The kernel is
the fibre-product description of \(S\).  Every map
\(\mathbf P^2\to E\) is constant, since its restriction to
each projective line is constant, whereas the rigidity lemma gives
\[
 A(T)=E(\F_2),\qquad
 A(E)=E(\F_2)\oplus\operatorname{End}_{\F_2}(E).
\]
The cohomology boundary therefore injects
\(\operatorname{End}_{\F_2}(E)\) into \(H^1_{\et}(S,A)\).
In particular \(\alpha=\delta(1_E)\) has infinite order:
\([m]:E\to E\) has degree \(m^2\) and is nonconstant for every
nonzero integer \(m\).

To represent this class by a scheme, glue two copies
\(X_i=E\times T\) along \(E\times E\), using the identity
on one side and \((a,t)\mapsto(a+t,t)\) on the other.
Both attaching maps are closed immersions, so the pushout is a
scheme by the Stacks Project, Lemma~37.67.8; Lemma~37.67.9(4)
gives smoothness over \(S\)~\cite{BStacksClosedPushout,BStacksPushoutMorphisms}.  The same local lifting just used
trivializes the translation cocycle on an etale cover.  Thus this
scheme is the \(A\)-torsor of class \(\pm\alpha\), locally
isomorphic to \(E\times S\), and is proper with geometrically
connected genus-one fibres.

Under the principal polarization \(E\simeq E^\vee\), translation
by \(a\) acts on \(\operatorname{Pic}^d(E)\) by translation
by \(\pm da\): represent a divisor by its sum in \(E\).
Consequently \(\operatorname{Pic}^d_{X/S}\) is a torsor of
class \(\pm d\alpha\).  A global line bundle of degree \(d\)
would give a section of that torsor, forcing \(d\alpha=0\),
hence \(d=0\).  This proves the assertion, and also shows that
\(X\to S\) is not projective.  The cubic smoothness, its five
\(\F_2\)-points, and the displayed fibre-product algebra were
checked exactly by
\texttt{part\_b/code/singular\_genus\_one\_checks.py}; its output
is preserved with the computations.  The infinite-order and torsor
claims are proved above, not inferred from the finite computation.
\end{example}

\begin{example}[A fully specified nonzero degree-two group on 12:78]\label{ex:B-H2-constant}
\status{Proved}
Let \(C/\F_2\) be the good reduction of the displayed
\(X_0(11)\) model, and let \(A=C\times C\to C\) be its
constant Jacobian scheme.  Take \(\ell=3\).  Direct enumeration
gives \(\#C(\F_2)=5\), so Frobenius on \(V_3C\) has
polynomial \(t^2+2t+2\).  Its two roots are distinct, and the
centralizer of a companion matrix has dimension two.  Kummer,
the constant-case result just proved, and Hochschild--Serre give
\[
 V_3H^2_{\et}(C,A)
 \simeq H^2_{\Bcts}(C,V_3A)
 \simeq H^1_{\Bcts}(\bar C,V_3A)_\Gamma
 \simeq\End_{\F_2}(C)\otimes\Q_3,
\]
of dimension two.  The geometric degree-two term has weight one,
so contributes no invariants; in the degree-one term,
\(H^1(\bar C,\Q_3)\otimes V_3C=(V_3C)^\vee\otimes V_3C\).
Semisimplicity identifies its coinvariants with its invariants,
and Tate's Main Theorem identifies the latter with endomorphisms.
The point count and exact centralizer calculation are reproduced in
\texttt{part\_b/code/finite\_field\_checks.py}.

This verifies the intended nonvanishing in R6.  If a dual Tate module
is replaced in the transcription by the Tate module of the dual
abelian variety, a twist must be retained:
\((V_\ell J)^\vee=V_\ell(J^t)(-1)\).
An integral isomorphism also cannot be read off from a displayed
\(\Q_\ell\)-vector-space calculation without tracking finite
Kummer kernels.
\end{example}

\begin{theorem}[The prime-to-characteristic relative-curve gate is a surface gate]
\label{thm:B-relative-surface-gate}
\status{Proved equivalence; the equivalent surface Tate assertion remains unresolved}
Fix primes \(p\ne\ell\).  The following statements are equivalent.
\begin{enumerate}[label=(\roman*)]
\item The Tate conjecture for divisors holds for every smooth
projective surface over \(\F_p\), including connected surfaces
whose field of constants is a finite extension of \(\F_p\).
\item For every finite field \(k/\F_p\), every smooth separated
finite-type \(k\)-scheme \(S\), and every smooth projective
morphism \(f:X\to S\) with geometrically connected curve fibres,
the cokernel of
\[
 f^*:\Br(S)[\ell^\infty]\longrightarrow\Br(X)[\ell^\infty]
\]
is finite.  Here \(\Br=H^2_{\mathrm{et}}(-,\Gm)_{\mathrm{tors}}\)
is the cohomological Brauer group.
\item The assertion in (ii) holds just for smooth curve families
with a section over nonempty opens of \(\mathbb P^1_k\), for
all finite \(k/\F_p\).
\item The single birational invariant vanishes:
\[
 H^3(F_0\mathbb P^3_{\overline{\F}_p},\Q_\ell(2))^
 {\operatorname{Gal}(\overline{\F}_p/\F_p)}=0.
\]
Here \(F_0\) denotes the colimit over dense open subsets, as
in Proposition~\ref{prop:B-diagonal-Chow-gate}.
\end{enumerate}
Thus the universal regular separated-base relative-curve question
in its intended \(\ell\ne p\) form is already equivalent to a
fixed-dimensional surface conjecture.  This theorem does not
identify the characteristic-primary assertion with that conjecture,
and does not extend its surface equivalence to arbitrary singular
or nonseparated bases.
\end{theorem}

\begin{proof}
We first prove (i)\(\Rightarrow\)(ii), retaining the arithmetic
cohomology term that can be nonzero on a nonproper base.
By \citet[Theorems~6 and~8, pp.~88--90]{BKahn2023Divisors},
hypothesis (i) implies
\[
 \operatorname{Pic}(Z)\otimes\Q_\ell
 \twoheadrightarrow H^2(Z_{\bar k},\Q_\ell(1))^{\Gamma_k}
 \tag{B.24a}
\]
for every smooth separated finite-type \(Z/k\).
The surface reduction uses
\citet[Proposition~4.1, Lemma~4.2 and Theorem~4.3,
author-version pp.~25--27]{BMorrow2019Variational}.
For clarity, its threefold step chooses a smooth ample surface
after finite extension whose Frobenius-one generalized eigenspace
has no larger dimension; the crystalline variational divisor
theorem then identifies the algebraic eigenspaces.  We use Kahn's
published iteration in higher dimension, rather than silently
assuming a characteristic-zero Picard Lefschetz argument in
positive characteristic.  Kahn's passage to smooth opens uses a
divisor-correspondence retraction and alteration trace, so it
does not add an unproved general semisimplicity hypothesis.

Write \(\Gamma=\operatorname{Gal}(\bar k/k)\).
For a smooth finite-type \(Z/k\), continuous Hochschild--Serre
in degree two gives the natural exact sequence
\[
 0\longrightarrow
 H^1\bigl(\Gamma,H^1(Z_{\bar k},\Q_\ell(1))\bigr)
 \longrightarrow H^2(Z,\Q_\ell(1))
 \longrightarrow H^2(Z_{\bar k},\Q_\ell(1))^\Gamma
 \longrightarrow0.
 \tag{B.24b}
\]
Finite-coefficient cohomology is finite in these degrees, so
passing from finite coefficients to the continuous complexes
has the needed Mittag--Leffler property.  The rational Kummer
sequence maps \(H^2(Z,\Q_\ell(1))\) onto
\(V_\ell\Br(Z)\), and rational divisor classes map to zero.
Using (B.24a), we may subtract an arithmetic divisor class from
any class in the middle of (B.24b) to make its geometric
restriction zero.  Hence
\[
 H^1\bigl(\Gamma,H^1(Z_{\bar k},\Q_\ell(1))\bigr)
 \twoheadrightarrow V_\ell\Br(Z).
 \tag{B.24c}
\]
We have not claimed that the source vanishes: on \(Z=\Gm\)
it is a copy of \(\Q_\ell\).

The geometric Leray sequence of \(f\) has
\[
 0\longrightarrow H^1(S_{\bar k},\Q_\ell)
 \longrightarrow H^1(X_{\bar k},\Q_\ell)
 \longrightarrow C\longrightarrow0,
 \qquad
 C\subseteq H^0(S_{\bar k},R^1f_*\Q_\ell).
 \tag{B.24d}
\]
The right-hand global-section representation is pure of weight
one.  Indeed, on each geometric connected component evaluation
at a geometric point is injective for a lisse sheaf.  Choose
closed points on all components after a finite constant-field
extension.  Evaluation embeds the section space into a finite
sum of \(H^1\)'s of smooth projective fibres, equivariantly for
that finite-index Galois group.  The Weil bounds on those fibres
give weight one, and taking powers of Frobenius shows that the
same weight holds before extension.  Thus \(C(1)\) has weight
\(-1\).  In particular \(1-\phi\) is invertible on it, and
both \(H^0(\Gamma,C(1))\) and \(H^1(\Gamma,C(1))\) vanish.
The cohomology sequence of (B.24d) therefore gives
\[
 H^1\bigl(\Gamma,H^1(S_{\bar k},\Q_\ell(1))\bigr)
 \xrightarrow{\sim}
 H^1\bigl(\Gamma,H^1(X_{\bar k},\Q_\ell(1))\bigr).
\]
Naturality of Kummer and Hochschild--Serre, together with
(B.24c) for \(X\), now proves
\[
 V_\ell\Br(S)\twoheadrightarrow V_\ell\Br(X).
\]
Both \(\ell\)-primary Brauer groups have finite cotype: their
\(\ell\)-torsion is a quotient of the finite group
\(H^2(-,\mu_\ell)\).  A group of finite cotype is a finite
group extended by a finite sum of \(\Q_\ell/\Z_\ell\)'s.
Surjectivity of the rational Tate-module map makes its divisible
part surjective, so the original cokernel is finite.  This proves
(ii), and (ii)\(\Rightarrow\)(iii) is immediate.

Assume (iii), and let \(Y/k\) be a smooth projective
geometrically connected surface.  Over \(\bar k\), choose a
pencil in a sufficiently ample system with smooth generic member
and a nonempty transverse finite base locus.  The required
Bertini conditions are open conditions on finitely many
coefficients; a choice over \(\bar k\) is therefore defined
over a finite extension \(k'/k\).  Enlarge it so that at least
one base point is rational, and blow up the base locus.  We
obtain a smooth projective surface \(\widetilde Y/k'\), a
proper morphism \(h:\widetilde Y\to\mathbb P^1_{k'}\), and
a section \(s\) supplied by that exceptional divisor.
Over a nonempty open \(U\subset\mathbb P^1_{k'}\) the
restriction \(X=\widetilde Y\times_{\mathbb P^1}U\to U\)
is a smooth projective geometrically connected curve family.
No assertion about the singularity type of its omitted fibres
is needed.

The map
\[
 \Br(\widetilde Y)[\ell^\infty]
 \longrightarrow
 \Br(X)[\ell^\infty]/h^*\Br(U)[\ell^\infty]
 \tag{B.24e}
\]
is injective.  If a class \(\alpha\) restricts to \(h^*\beta\),
pull back by the section over \(U\).  Since the section extends
over \(\mathbb P^1\) and
\(\Br(\mathbb P^1_{k'})=\Br(k')=0\), we obtain \(\beta=0\).
Then \(\alpha\) vanishes on a dense open and is zero, because
the Brauer group of a regular integral scheme injects into that
of its function field
\cite[Corollary~1.8]{BGrothendieckBrauerII1966}.
By (iii), the right side of (B.24e) is finite, hence so is its
left side.  Pullback from \(Y_{k'}\) to its blowup is injective
on Brauer groups: both restriction maps embed in the same
function-field Brauer group and commute with pullback.
Thus \(\Br(Y_{k'})[\ell^\infty]\) is finite.

Finiteness descends from \(k'\) to \(k\), with no assumption
that \([k':k]\) is prime to \(\ell\).  Restriction followed
by corestriction is multiplication by \([k':k]\); the kernel of restriction
on \(\Br(Y)[\ell^\infty]\) is contained in
\(\Br(Y)[\ell^v]\), where \(v=v_\ell([k':k])\), a finite
group by Kummer.  Its image is also finite, proving the claim.
Finally, for proper \(Y\), \(\operatorname{Pic}^0(Y)(k)\) is
finite and its N\'eron--Severi group is finitely generated.  Hence
the rational Kummer kernel is exactly
\(\operatorname{Pic}(Y)\otimes\Q_\ell\).
Moreover \(H^1(Y_{\bar k},\Q_\ell(1))\)
has weight \(-1\).  Sequence (B.24b) and rational Kummer then
identify vanishing of \(V_\ell\Br(Y)\) with the Tate
surjectivity for divisors on \(Y\).  This gives (i) for every
constant field.  A connected smooth surface over \(\F_p\)
with a larger field of constants is handled by that field and
the corresponding induced Galois representation.  Disconnected
surfaces are handled componentwise.

Finally, (i)\(\Leftrightarrow\)(iv) is the exact surface case
of \citet[Theorem~7.1(1) and Remark~7.2]{BBalkanSchreieder2026}.
Their Proposition~4.6 removes a separate eigenvalue-one
semisimplicity hypothesis: in degree two it follows from the
divisor Tate assertion.  The constant-field quantifier agrees
with (i), because the geometric cohomology of a variety over
\(\F_{p^r}\), viewed over \(\F_p\), is an induced
representation and its \(\Gamma_{\F_p}\)-invariants are the
original \(\Gamma_{\F_{p^r}}\)-invariants.  No vanishing of
the displayed group has been proved here.
\end{proof}

This equivalence reduces the regular separated relative-curve
gate to surfaces; it does not prove their Tate conjecture.
The prime-to-\(p\) equivalence chain on arbitrary finite-type
bases in Theorem~\ref{thm:B-finite-chain-all-bases} remains a
separate statement with its original full scope.

\paragraph{The same surface gate for all abelian schemes and all bases.}
The universal finiteness assertion in the manuscript is broader
than the relative-curve assertion just considered.  Its base may
be singular, nonreduced or nonseparated.  The following descent
argument shows that this added generality does not enlarge the
remaining divisor-Tate gate.

\begin{lemma}[Rational sections descend through proper covers]
\label{lem:B-proper-section-descent}
\status{Proved}
Let \(S\) be a noetherian \(\F_p\)-scheme, let
\(q:T\to S\) be proper and surjective, and let \(A/S\)
be an abelian scheme.  The sequence
\[
0\longrightarrow A(S)\otimes\Q
\longrightarrow A(T)\otimes\Q
\xrightarrow{\operatorname{pr}_1^*-\operatorname{pr}_2^*}
A(T\times_S T)\otimes\Q
\]
is exact.  The same equalizer statement holds after scalar
extension from \(\Q\) to \(\Q_\ell\).
\end{lemma}
\begin{proof}
First suppose \(S\) reduced, and take an actual section
\(b:T\to A\) whose two pullbacks agree as morphisms.
Let \(Z\) be its scheme-theoretic image in \(A\), with
maps \(j:Z\hookrightarrow A\) and \(g:Z\to S\).
The morphism \(g\) is proper and \(T\to Z\) is
surjective, by Stacks Lemma~29.42.10~\cite{BStacksProperImage}.
On every geometric fibre, any pair of points of \(T\) has
the same image under \(b\), by the equalizer hypothesis.
Since proper surjectivity persists under base change, each
geometric fibre of \(g\) has exactly one point.  Thus \(g\)
is quasi-finite and radicial as well as proper and surjective.
It is a finite universal homeomorphism; the exact criteria are
Stacks Lemmas~29.10.2 and~37.44.1~\cite{BStacksRadicial,BStacksProperFinite}.

There are an integer \(r\ge0\) and a morphism
\(h:S\to Z\) such that
\[
 h g=F_Z^r,\qquad g h=F_S^r.
\]
This is Kelly's Proposition~3 and Corollary~20, proved through
Lemma~18 and Corollary~19~\cite{BKelly2018}.
For reduced \(S\), it also follows directly from
Stacks Lemma~29.47.9: every local element of \(\mathcal O_Z\)
has a \(p\)-power in \(\mathcal O_S\), and finite
presentation and quasi-compactness give one exponent for a finite
set of algebra generators~\cite{BStacksFrobeniusPower}.
The local power maps glue.

Write
\(A^{(p^r)}=A\times_{S,F_S^r}S\).
The pair \((jh,1_S)\) defines a section \(c\) of
\(A^{(p^r)}/S\), and its pullback to \(T\) is
\(F_{A/S}^r b\).  Applying the iterated Verschiebung gives
\[
 a=V_{A/S}^r c\in A(S),\qquad q^*a=[p^r]b,
\]
because \(V^rF^r=[p^r]\); see Conrad's Example~4.4
and the Frobenius--Verschiebung discussion~\cite{BConrad2006}.
Thus an actual equalizer section descends after multiplication
by a power of \(p\).

For a rational equalizer element, multiply first by an integer
to obtain an actual section.  Its actual Cech difference is
torsion, since its rational difference vanishes.  Multiplication
by the order of that difference reduces to the preceding case.
Injectivity over reduced \(S\) follows because the equality
locus of two sections into the separated scheme \(A\) is
closed: surjective pullback makes its defining ideal nilpotent,
and reducedness makes that ideal zero.

For nonreduced \(S\), the finite filtration of its nilpotent
ideal has successive additive section kernels
\(\operatorname{Lie}(A)\otimes I\).  Its cohomology exact
sequences show that reduction changes section groups only by
bounded \(p\)-primary kernels and cokernels, as in
Theorem~\ref{thm:B-finite-chain-all-bases}.
Reduce \(T\) and its fibre product as well.  Rationalizing
therefore reduces the asserted sequence to the reduced case.
Finally \(\Q_\ell\) is flat over \(\Q\), so scalar
extension preserves the equalizer.  This last step avoids any
attempt to clear arbitrary \(\Q_\ell\)-coefficients by an
integer.
\end{proof}

\begin{theorem}[The universal all-base finiteness gate]
\label{thm:B-all-bases-surface-gate}
\status{Proved equivalence; the surface Tate assertion is unresolved}
Fix distinct primes \(p\) and \(\ell\).  The following
universal assertions are equivalent.
\begin{enumerate}
\item Divisor Tate holds for every smooth projective surface
over every finite extension of \(\F_p\).
\item For every finite-type \(\F_p\)-scheme \(S\) and
every abelian scheme \(A/S\), the group
\(H^1_{\et}(S,A)[\ell^\infty]\) is finite.  No
regularity, reducedness, properness or separatedness hypothesis
is imposed on \(S\).
\item The assertion in (2) holds for the Jacobians of smooth
projective geometrically connected curve families with a section
over nonempty opens of \(\mathbf P^1\), over all finite
extensions of \(\F_p\).
\end{enumerate}
These assertions are also equivalent to the birational
degree-three vanishing in
Theorem~\ref{thm:B-relative-surface-gate}.
\end{theorem}
\begin{proof}
The implication (2)\(\Rightarrow\)(3) is immediate.
Proposition~\ref{prop:B-relative-curve} and
Theorem~\ref{thm:B-relative-surface-gate} give
(3)\(\Rightarrow\)(1).  We prove (1)\(\Rightarrow\)(2).

First let \(T/k\) be smooth and separated of finite type,
where \(k/\F_p\) is finite, and let \(A/T\) be abelian.
Set \(B=A^t\), with structure map \(\pi:B\to T\).
The Kahn--Morrow implication used in
Theorem~\ref{thm:B-relative-surface-gate} gives divisor Tate
on the smooth separated total space \(B\).
For any integer \(m>1\), multiplication \([m]\) has
eigenvalue \(m^q\) on \(R^q\pi_*\Q_\ell\).
The distinct eigenvalues split the degree-two Leray filtration;
its \(m\)-eigenspace is
\[
 e_1 H^2(\bar B,\Q_\ell(1))
 =H^1(\bar T,R^1\pi_*\Q_\ell(1))
 =H^1(\bar T,V_\ell A).
\]
The divisor identification proved in
Proposition~\ref{prop:B-finite-chain-proper} identifies the
\(m\)-eigenspace of the algebraic divisor image with
\(N_T=A(\bar T)\otimes\Q_\ell\): the Poincare bundle
realizes sections, the zero section kills their Brauer obstruction,
and base Picard and relative Neron--Severi classes have eigenvalues
one and \(m^2\), respectively.  Applying \(e_1\) to
divisor Tate therefore proves
\[
 \kappa_T:N_T^\Gamma
 \xrightarrow{\sim}H^1(\bar T,V_\ell A)^\Gamma.
\]
Injectivity here is geometric Kummer, with the completion
identification in Lemma~\ref{lem:B-sections-reduced}.

Now let \(S/k\) be reduced, separated and of finite type.
Take a proper surjection \(q:T\to S\) with \(T\)
smooth separated of finite type: take the finitely many reduced
irreducible components, compactify, apply a projective alteration,
and restrict above \(S\).  De Jong's Theorem~4.1 and
Remark~4.2 supply the required smooth alterations~\cite{c_dejong1996}.
If a finite extension of constants is used, retain it in the
proper cover; assertion (1) includes all such finite extensions.
Put \(T_a=T^{\times_S(a+1)}\) and
\(\mathcal V=V_\ell A\).  Proper Cech descent gives
\[
 E_1^{a,b}=H^b(\bar T_a,\mathcal V)
 \ \Longrightarrow\ H^{a+b}(\bar S,\mathcal V).
\]
This follows from Stacks Lemmas~85.36.3--85.36.6 at finite
\(\ell\)-power coefficients~\cite{BStacksProperDescent}.
Finite-coefficient etale cohomology is finite in this situation,
so the inverse systems satisfy Mittag--Leffler.  Taking continuous
cohomology and tensoring with \(\Q_\ell\) preserves the
displayed first-quadrant spectral sequence in total degree one.
Its five-term sequence identifies the pullback kernel on
\(H^1\) with \(E_2^{1,0}\).

Each \(H^0(\bar T_a,\mathcal V)\) is pure of weight
\(-1\), including when \(T_a\) is singular: after a
finite constant extension, evaluation at one closed point of
each geometric connected component embeds it equivariantly into
a finite sum of Tate modules of abelian varieties over finite
fields.  Consequently \(E_2^{1,0}\), a subquotient of the
degree-zero row, has no \(\Gamma\)-invariants.

Take \(c\in H^1(\bar S,\mathcal V)^\Gamma\).
The smooth-base case gives a unique \(b\in N_T^\Gamma\)
such that \(q^*c=\kappa_T(b)\).
The two pullbacks of \(c\) to \(\bar T_1\) agree.
Kummer injectivity on the possibly singular \(T_1\), supplied
by Theorem~\ref{thm:B-finite-chain-all-bases}, therefore
gives \(\operatorname{pr}_1^*b=\operatorname{pr}_2^*b\).
Lemma~\ref{lem:B-proper-section-descent}, applied over
\(\bar k\) and then extended from \(\Q\) to
\(\Q_\ell\), descends \(b\) uniquely to
\(a\in N_S\).  Uniqueness makes \(a\) invariant.
The difference \(c-\kappa_S(a)\) is then invariant in
the weight-minus-one pullback kernel, hence zero.  Thus
invariant Kummer surjectivity holds over every reduced separated
base.

Nilpotent thickening does not change finite etale torsion sheaves
or their cohomology, and changes section groups only by bounded
\(p\)-primary groups.  This removes reducedness.  For a
nonseparated finite-type \(S\), use a finite affine open
cover.  Every finite intersection is quasi-compact and separated.
The ordinary open-cover Cech spectral sequence again has its
degree-one pullback kernel in the pure weight-minus-one
degree-zero row.  The pulled-back section descends because
sections are a Zariski sheaf and tensoring with \(\Q\)
is exact for the finite products in its equalizer, followed by
scalar extension to \(\Q_\ell\).
The same subtraction argument proves invariant Kummer
surjectivity on \(S\).

Finally Theorem~\ref{thm:B-finite-chain-all-bases}, including
its finite-cotype assertion, identifies this invariant Kummer
surjectivity with the actual finiteness of
\(H^1(S,A)[\ell^\infty]\).  This proves (2).
\end{proof}

This is a reduction of the full finite-type abelian-scheme
finiteness assertion, not a proof of its equivalent surface Tate
condition.  Extending the relative-curve Brauer quotient
statement to singular bases requires separate control of
non-torsion Picard-degree and base \(H^2(S,\Gm)\) terms.
The integral torsor-lattice and unit-cup arguments below provide
that control and prove the same surface-Tate equivalence in
the full arbitrary-base scope.  The characteristic-primary
counterexample below is separate.

\begin{proposition}[Finite generation of the localized torsor group]
\label{prop:B-integral-torsor-lattice}
\status{Proved conditional strengthening; surface Tate remains unresolved}
Assume the equivalent surface-Tate conditions of
Theorem~\ref{thm:B-all-bases-surface-gate}.  For every
finite-type \(\F_p\)-scheme \(S\), every abelian scheme
\(A/S\), and every \(\ell\ne p\),
\[
 H^1_{\et}(S,A)\otimes\Z_{(\ell)}
       \quad\hbox{is finitely generated over }\Z_{(\ell)}.
                                                               \tag{B.24i}
\]
Thus the non-torsion part on a singular base is controlled as
well as its \(\ell\)-primary torsion.  If \(f:X\to S\)
is a smooth proper family of geometrically connected curves,
then
\[
 H^1_{\et}(S,\operatorname{Pic}_{X/S})[\ell^\infty]
       \quad\hbox{is finite},
\]
even when its degree-one Picard torsor has infinite order.
\end{proposition}
\begin{proof}
First let \(S\) be reduced and separated.  Choose a proper
surjection \(q:T\to S\) with \(T\) smooth as in
Theorem~\ref{thm:B-all-bases-surface-gate}, and put
\[
 K=\ker\bigl(H^1_{\et}(S,A)\longrightarrow
                        H^1_{\et}(T,A)\bigr).
\]
For \([P]\in K\), choose a section of its torsor over
\(T\).  The difference of the two pulled-back sections is
an element of \(A(T\times_ST)\) satisfying the ordinary
\v Cech one-cocycle identity.  Changing the section changes
this cocycle by a coboundary, and contracted product adds
cocycles.  We obtain a homomorphism
\[
 K\longrightarrow H^1\bigl(A(T_\bullet)\bigr),     \tag{B.24j}
\]
where the right side denotes cohomology of the actual
cosimplicial groups of sections, not a presumed proper-cover
descent theorem for \(\et\) torsors.

The kernel of (B.24j) is \(p\)-primary.  If the cocycle is
a coboundary, translate the chosen section to make its two
pullbacks agree exactly.  Its schematic image \(Z\) in
\(P\) then has one geometric point in each fibre over
\(S\), by the same argument as
Lemma~\ref{lem:B-proper-section-descent}.
Here \(P\) can be an algebraic space: it is proper over
\(S\) by descent from the abelian scheme, so its closed
image \(Z\) is proper and quasi-finite over \(S\).
Zariski's main theorem makes \(Z\to S\) finite, and a
finite algebraic space over a scheme is a scheme
\cite[Lemmas~76.34.3 and 67.45.3]{BStacksFiniteSpaces}.
Thus \(Z\to S\) is a finite universal homeomorphism of
schemes.  Kelly's Frobenius factorization supplies a morphism
\(S\to Z\) over a fixed power of absolute Frobenius.
Composing with \(Z\to P\) gives a section of the twist
\(P^{(p^r)}\).  Relative Frobenius identifies that twist
with the pushout of \(P\) along
\(F^r_{A/S}:A\to A^{(p^r)}\), as is checked on an
\(\et\) cover trivializing the torsor.  Pushing its section
along \(V^r\) trivializes the pushout by
\(V^rF^r=[p^r]\).  Hence \(p^r[P]=0\), proving the
claim~\cite[Proposition~3 and Corollary~20]{BKelly2018}.
Exact localization of (B.24j) gives an injection
\[
 K[1/p]\lhook\joinrel\longrightarrow
             H^1\bigl(A(T_\bullet)[1/p]\bigr).     \tag{B.24k}
\]

For every finite-type \(\F_p\)-scheme \(U\), the group
\(A(U)[1/p]\) is finitely generated over \(\Z[1/p]\).
For reduced \(U\), restriction to its finitely many generic
points is injective by separatedness.  Lang--N\'eron makes
each group of rational points of the resulting generic
abelian variety finitely generated: the trace over its finite
constant field has only finitely many rational points
\cite[Theorems~2.1 and 7.1]{BConrad2006}.
For nonreduced \(U\), the kernel of restriction to
\(U_{\mathrm{red}}\) has a finite filtration by additive
Lie-sheaf sections and is killed by a power of \(p\).
After inverting \(p\), the group is a submodule of the
finitely generated reduced section module.  No surjectivity
of global sections across a nilpotent thickening is needed
for this finite-generation argument.  The right side of
(B.24k) is a subquotient of \(A(T_1)[1/p]\), hence is
finitely generated, as is its submodule \(K[1/p]\).

On a regular noetherian integral scheme \(W\), restriction
\(H^1(W,A)\to H^1(k(W),A)\) is injective.  The exact
source for arbitrary dimension is Stacks Project
Remark~110.74.2; its Picard argument extends a generically
defined line bundle on the regular dual abelian scheme, and
uniqueness extends the corresponding rational section.
Lemma~110.74.3 makes the field group torsion by a finite
separable splitting extension and the norm construction
\cite{BStacksRegularTorsors}.
Therefore \(H^1(T,A)\) is torsion.  Its
\(\ell\)-primary subgroup is finite by
Theorem~\ref{thm:B-all-bases-surface-gate}, so
\(H^1(T,A)\otimes\Z_{(\ell)}\) is finite.
Localizing the defining exact sequence for \(K\) now proves
(B.24i) for reduced separated \(S\).

A nilpotent thickening changes the smooth commutative-group
cohomology through a finite succession of additive Lie-sheaf
kernels, killed by \(p\).  Its long exact sequences show
that \(H^1(S,A)\otimes\Z_{(\ell)}\) is unchanged.
Finally, take a finite affine open cover of a possibly
nonseparated \(S\).  Ordinary open-cover torsor descent
identifies the kernel of restriction in \(H^1\) with the
\v Cech \(H^1\) of the section groups.  Its localized terms
are finitely generated by the preceding section argument,
while the restriction target is a finite direct sum of the
finitely generated groups for the affine opens.  This proves
(B.24i) in its full stated scope.

For the curve family, put \(J=\operatorname{Pic}^0_{X/S}\)
and \(P=\operatorname{Pic}_{X/S}\).  The exact sequence
of \(\et\) sheaves
\[
 0\longrightarrow J\longrightarrow P
       \xrightarrow{\deg}\Z\longrightarrow0
\]
identifies the kernel of \(H^1(S,P)\to H^1(S,\Z)\)
with \(H^1(S,J)/\delta H^0(S,\Z)\).
The latter has finite \(\ell\)-primary torsion by (B.24i).
Moreover \(H^1(S,\Z)\) is torsion-free: for every positive
integer \(n\), the global-section map
\(H^0(S,\Z)\to H^0(S,\Z/n)\) is surjective on the
finitely many connected components, and the multiplication-by-
\(n\) exact sequence gives \(H^1(S,\Z)[n]=0\).
The \(\ell\)-primary torsion of \(H^1(S,P)\) therefore
lies in that finite kernel, proving the final assertion.
\end{proof}

This removes the infinite-period degree torsor as an obstruction
to finiteness of the \(\ell\)-primary \(H^1\) of the
relative Picard sheaf.  A full Brauer quotient on an arbitrary
singular base also contains the image of
\(H^2(S,\Gm)/\delta P(S)\).  Its torsion need not come
from \(\Br(S)\); the non-torsion saturation in this base
term is addressed for proper bases in
Proposition~\ref{prop:B-singular-proper-Brauer-gate}, and
for every finite-type base by the independent unit-cup
argument of Theorem~\ref{thm:B-all-bases-Brauer-units}.

\begin{lemma}[The integral degree-two kernel of a cover]
\label{lem:B-Gm-proper-kernel}
\status{Proved without a Tate hypothesis}
Let \(q:T\to S\) be a morphism of finite-type
\(\F_p\)-schemes whose perfection is a \(v\)-cover,
and put \(T_a=T\times_S\cdots\times_ST\)
with \(a+1\) factors.  Then
\[
 \ker\bigl(H^2_{\et}(S,\Gm)[1/p]
       \longrightarrow H^2_{\et}(T,\Gm)[1/p]\bigr)
\]
is a finitely generated \(\Z[1/p]\)-module.  Consequently,
for every prime \(\ell\ne p\), the same kernel after tensoring
with \(\Z_{(\ell)}\) is finitely generated over
\(\Z_{(\ell)}\).  These are assertions about the full
cohomology groups, including their possible non-torsion parts.
In particular the lemma applies to proper surjections and to
finite-type \(h\)-covers obtained by combining a finite open
cover with proper surjections.
\end{lemma}
\begin{proof}
Write \(U^{\mathrm{perf}}\) for the perfection of a
characteristic-\(p\) scheme.  The universal homeomorphism
\(U^{\mathrm{perf}}\to U\) identifies the small \(\et\)
topoi.  Under this identification its multiplicative sheaf is
\(\varinjlim(\Gm\xrightarrow{p}\Gm\xrightarrow{p}\cdots)\):
this follows on affine \(\et\) objects by taking units of the
perfected ring.  Cohomology of a quasi-compact quasi-separated
scheme commutes with filtered colimits.  Therefore, naturally
in every finite-type \(U\),
\[
 H^i_{\et}(U^{\mathrm{perf}},\Gm)
       =H^i_{\et}(U,\Gm)[1/p].                     \tag{B.24f}
\]
The two inputs are respectively the Stacks Project's
Theorem~59.45.2 and Proposition~59.45.4, and Lemma~69.5.1
\cite{BStacksPerfectEtale,BStacksCohomColimits}.

On the category of perfect quasi-compact quasi-separated
schemes, use the \(v\)-topology.  Bhatt--Scholze's
Theorem~4.1(ii), with \(n=1\), says that vector bundles form
a \(v\)-stack; hence so does the groupoid of line bundles.
Its automorphism sheaf at the trivial bundle is \(\Gm\),
and it is the classifying stack of \(\Gm\) because line
bundles are already Zariski locally trivial.  Thus
\[
 H^0_v(U^{\mathrm{perf}},\Gm)
       =\Gamma(U^{\mathrm{perf}},\mathcal O^\times),
 \qquad
 H^1_v(U^{\mathrm{perf}},\Gm)
       =\operatorname{Pic}(U^{\mathrm{perf}}).
\]
For the change-of-topology morphism
\(\epsilon:(U^{\mathrm{perf}})_v\to
(U^{\mathrm{perf}})_{\et}\), this gives
\(\epsilon_*\Gm=\Gm\) and \(R^1\epsilon_*\Gm=0\),
since line bundles become trivial on Zariski covers.
The low-degree Leray sequence consequently supplies a natural
injection
\[
 H^2_{\et}(U^{\mathrm{perf}},\Gm)
       \lhook\joinrel\longrightarrow
 H^2_v(U^{\mathrm{perf}},\Gm).                     \tag{B.24g}
\]
No comparison in all higher degrees is being presumed.
Perfection commutes with the fibre products here, and the
perfected morphism is a \(v\)-cover by hypothesis.
A proper surjection becomes a \(v\)-cover, as do open
covers and their composites; the proper-surjection case and
its use in descent are also recorded in
Bhatt--Scholze's Proposition~4.5 and its proof
\cite[Theorem~4.1(ii), Proposition~4.5, author-version
pp.~15--17]{BBhattScholze2017}.

The ordinary cohomological descent spectral sequence for that
cover is
\[
 E_1^{a,b}=H^b_v(T_a^{\mathrm{perf}},\Gm)
       \Longrightarrow H^{a+b}_v(S^{\mathrm{perf}},\Gm).
\]
The kernel of restriction in total degree two has a two-step
filtration whose associated groups are \(E_\infty^{1,1}\)
and \(E_\infty^{2,0}\).  They are subquotients of
\[
 \operatorname{Pic}(T_1^{\mathrm{perf}})
       =\operatorname{Pic}(T_1)[1/p],\qquad
 \Gamma(T_2^{\mathrm{perf}},\mathcal O^\times)
       =\Gamma(T_2,\mathcal O^\times)[1/p].
\]
Both are finitely generated over \(\Z[1/p]\).
For the first this is Jaffe's Theorem~10.11.  For the second,
his Proposition~10.8 gives finite generation on the reduction,
and the kernel of reduction on units is killed by a bounded
power of \(p\).  Even if the map on global units is not
surjective, its image is a subgroup of a finitely generated
abelian group, which suffices
\cite[Proposition~10.8 and Theorem~10.11,
preprint pp.~42--43]{BJaffe1994}.
The ring \(\Z[1/p]\) is noetherian, so the indicated
subquotients, their extension, and all its submodules are
finitely generated.  By the natural injection (B.24g), the
kernel in the statement embeds in this \(v\)-cohomological
kernel.  This proves the assertion.  Further localization is
exact and proves the \(\ell\)-local assertion.
\end{proof}

\begin{proposition}[Proper singular bases introduce no further prime-to-\(p\) gate]
\label{prop:B-singular-proper-Brauer-gate}
\status{Proved equivalence; the surface Tate assertion remains unresolved}
Fix \(p\ne\ell\).  Divisor Tate for all smooth projective
surfaces over finite extensions of \(\F_p\) is equivalent to
the following assertion:
\[
 \begin{gathered}
 H^2_{\et}(S,\Gm)\otimes\Z_{(\ell)}
 \quad\hbox{is a finitely generated }\Z_{(\ell)}\hbox{-module}\\
 \hbox{for every proper finite-type }\F_p\hbox{-scheme }S.
 \end{gathered}
                                                               \tag{B.24h}
\]
Under these equivalent hypotheses, \(\Br(S)[\ell^\infty]\)
is finite for every such \(S\), including singular and
nonreduced ones.  In particular, if \(S\) is proper and
\(f:X\to S\) is a proper curve family, the \(\ell\)-primary
cokernel of \(\Br(S)\to\Br(X)\) is finite, whether or not
there is a global positive-degree line bundle on that family.
\end{proposition}
\begin{proof}
Assume surface Tate.  By Chow's lemma followed by de Jong's
projective alterations, choose a proper surjection \(T\to S\)
with \(T\) a finite disjoint union of smooth projective
varieties over finite fields.  Nilpotents can first be removed,
and the finitely many irreducible components treated separately
\cite[Theorem~4.1 and Remark~4.2]{c_dejong1996}.
The Kahn--Morrow reduction used in the proof of
Theorem~\ref{thm:B-relative-surface-gate} gives divisor Tate
on each component of \(T\).  Properness makes the arithmetic
\(H^1\) term in (B.24b) vanish by its weight \(-1\);
Kummer and finite cotype therefore make
\(\Br(T)[\ell^\infty]\) finite.  Since \(T\) is regular,
\(H^2_{\et}(T,\Gm)\) is torsion
\cite[Proposition~1.4 and Corollary~1.8]{BGrothendieckBrauerII1966}.
Consequently
\(H^2_{\et}(T,\Gm)\otimes\Z_{(\ell)}\) is finite.
Lemma~\ref{lem:B-Gm-proper-kernel} makes the kernel of
restriction from \(S\) finitely generated, so its source is
an extension of a finite module by a finitely generated module.
This proves (B.24h).

The torsion submodule of a finitely generated
\(\Z_{(\ell)}\)-module is a finite \(\ell\)-group.
Localization identifies the original \(\ell\)-primary
torsion with that submodule.  This proves the stated Brauer
finiteness; applying it directly to the proper total scheme
\(X\) proves the relative consequence.  It avoids replacing
the full Leray base term by its torsion subgroup.
Conversely, apply (B.24h) to a smooth projective surface.
Its full \(H^2(\Gm)\) is torsion, so the localized
module is finite.  The proper rational Kummer argument in
Theorem~\ref{thm:B-relative-surface-gate} then gives divisor
Tate.  This proves the equivalence.
\end{proof}

Properness is essential to this particular argument.  For a
smooth nonproper alteration \(T\), the group
\(\Br(T)[\ell^\infty]\) can contain
\(\Q_\ell/\Z_\ell\).  An extension of that group by a
finitely generated lattice need not be finitely generated:
\(\Z_{(\ell)}\subset\Q\) has quotient
\(\Q_\ell/\Z_\ell\) as an abstract abelian group.
Thus the kernel lemma alone does not control the nonproper
case.  The unit-cup argument in
Theorem~\ref{thm:B-all-bases-Brauer-units} supplies the
missing control and extends the relative conclusion to
every finite-type base.

\begin{lemma}[Arithmetic unit cups generate the smooth Brauer group up to finite error]
\label{lem:B-unit-cup-generation}
\status{Proved conditional on divisor Tate for surfaces}
Fix a finite field \(k\) of characteristic \(p\) and a
prime \(\ell\ne p\).  Assume divisor Tate for all smooth
projective surfaces over finite extensions of \(k\).
For every smooth quasi-projective finite-type \(k\)-scheme \(T\),
the cyclic-cup map
\[
 c_T:\Gamma(T,\mathcal O_T^\times)\otimes
                  \Q_\ell/\Z_\ell
       \longrightarrow \Br(T)[\ell^\infty]         \tag{B.24j'}
\]
has finite cokernel.  Here \(u\otimes a/\ell^r\) maps to
the cyclic class \((\chi_{a/r},u)\), where the arithmetic
character \(\chi_{a/r}:\Gal(\bar k/k)\to
\Q_\ell/\Z_\ell\) sends arithmetic Frobenius to
\(a/\ell^r\).  The construction is natural for every
\(k\)-morphism.
\end{lemma}
\begin{proof}
The cyclic class can be defined as the cup of the
\(\Z/\ell^r\)-valued character with the Kummer class of
\(u\), followed by \(H^2(T,\mu_{\ell^r})\to\Br(T)\).
These definitions agree as \(r\) increases and are bilinear;
thus they define (B.24j').  Let
\[
 N_T=\bigl(\Gamma(\bar T,\mathcal O_{\bar T}^\times)
                     /\Gamma(\bar T,\mathcal O_{\bar T}^\times)_{\rm tors}
              \bigr)\otimes\Q_\ell.
\]
Geometric Kummer embeds \(N_T\) in
\(H^1(\bar T,\Q_\ell(1))\).  We prove that its quotient
is pure of weight \(-1\), while \(N_T\) has finite
Frobenius action.  Treating components after a finite constant
extension suffices for this assertion.

First suppose \(T=Y\setminus D\) for \(Y\) smooth
projective, where \(D\) denotes the union of its
codimension-one boundary components.  Boundary subsets of
codimension at least two do not affect \(H^1\) or units.
Remove also the singular loci and pairwise intersections of
those boundary components; this is a codimension-at-least-two
subset of \(Y\), so purity does not change its \(H^1\)
or \(H^2\).  The codimension-one localization sequence
therefore gives
\[
 0\longrightarrow H^1(\bar Y,\Q_\ell(1))
 \longrightarrow H^1(\bar T,\Q_\ell(1))
 \longrightarrow \bigoplus_{D_i}\Q_\ell[D_i]
 \xrightarrow{c_1}H^2(\bar Y,\Q_\ell(1)).         \tag{B.24k'}
\]
The kernel of \(c_1\) in the divisor space is exactly
the image of geometric units tensored with \(\Q_\ell\).
Indeed the N\'eron--Severi cycle map is injective, and the
kernel on integral boundary divisors is therefore, up to a
finite multiple, algebraic equivalence to zero.  Every point
of \(\operatorname{Pic}^0(\bar Y)\) is torsion, since it
is defined over some finite field.  A further finite multiple
of such a divisor is consequently principal.  Conversely a
unit's principal divisor is supported on the boundary and has
zero Chern class.  These are relations between finitely
generated integral divisor and N\'eron--Severi groups, so
tensoring their rational kernel with \(\Q_\ell\) gives
the full displayed kernel.  The residue of the Kummer class
of a unit is its vector of boundary orders.  Thus units split
this kernel inside (B.24k').  The remaining summand is
\(H^1(\bar Y,\Q_\ell(1))\), pure of weight \(-1\).
The boundary components, hence their relation lattice, have
finite Galois action.  This proves both assertions in this case.

For general smooth quasi-projective \(T\), compactify and apply a
projective smooth alteration to obtain a smooth projective
\(Y\), an open \(U\subset Y\), and a proper generically
finite morphism \(g:U\to T\) of degree \(d>0\)
\cite[Theorem~4.1 and Remark~4.2]{c_dejong1996}.
There are pullback and proper trace maps
\[
 H^1(\bar T,\Q_\ell(1))
  \xrightarrow{g^*}H^1(\bar U,\Q_\ell(1))
  \xrightarrow{\operatorname{Tr}_g}H^1(\bar T,\Q_\ell(1)),
 \qquad \operatorname{Tr}_g g^*=d.
\]
They are the purity/Gysin trace for a projective local complete
intersection map of virtual dimension zero, followed by proper
adjunction; the degree formula is checked on the finite-flat
generic locus
\cite[Corollary~3.1.2, Lemmas~3.5.8--3.5.10 and
Proposition~2.5.13]{c_riou2014}.
Both \(T\) and \(U\) have ample line bundles, so the
ampleness hypotheses in those cited statements are satisfied.
This trace sends a unit Kummer class to the Kummer class of
its field norm.  The norm is a unit on normal \(T\): its
divisor is the pushforward of the zero divisor of the unit
on \(U\).  The compatibility may be checked on a dense
finite-flat locus, where Kummer transfer is the norm
(for a separable extension this is the product of conjugates;
for a purely inseparable extension it is the corresponding
power).  Restriction on \(H^1\) from smooth \(T\) to
that dense open is injective by codimension-one purity.
Thus the compatibility holds on \(T\).
The weight-zero unit summand on \(U\) traces into the unit
summand on \(T\), and trace of the weight-\(-1\) part
generates the quotient on \(T\), because
\(d^{-1}\operatorname{Tr}_g g^*=1\).
That quotient is therefore pure of weight \(-1\).
The unit lattice on \(T\) has finite Frobenius action as
well, for example by its pullback into the finite boundary
lattice on \(U\).

Put \(\Gamma=\Gal(\bar k/k)\).
The weight-\(-1\) quotient has neither \(\Gamma\)-invariants
nor coinvariants.  Hence
\[
 H^1(\Gamma,N_T)\xrightarrow{\sim}
 H^1\bigl(\Gamma,H^1(\bar T,\Q_\ell(1))\bigr).
\]
Since Frobenius has finite order on \(N_T\), averaging
identifies its coinvariants with its invariant subspace.
Finite-field descent, or norms after a finite constant
extension, identifies that subspace with
\(\Gamma(T,\mathcal O_T^\times)\otimes\Q_\ell\).
On a component whose field of constants has degree \(e\)
over \(k\), the arithmetic character restricts with the
factor \(e\); this nonzero factor does not affect the
rational assertion or the eventual finite cokernel.
Cup with the arithmetic degree-one class is the resulting
map onto this \(H^1(\Gamma,-)\) term.
By (B.24c), whose divisor-Tate hypothesis follows from
Kahn--Morrow under the assumption of the lemma, this term
surjects onto \(V_\ell\Br(T)\).
It is exactly the rational Tate-module map of (B.24j').
Finally \(\Br(T)[\ell]\) is finite by finite-coefficient
cohomology, so \(\Br(T)[\ell^\infty]\) has finite cotype.
A homomorphism from a finite sum of
\(\Q_\ell/\Z_\ell\)'s which is surjective on rational
Tate modules has finite cokernel.  This proves the lemma.
\end{proof}

\begin{lemma}[Unit cups have finite kernel on every finite-type base]
\label{lem:B-unit-cup-kernel}
\status{Proved without a Tate hypothesis}
Let \(k\) be a finite field of characteristic \(p\), let
\(\ell\ne p\), and let \(U\) be any finite-type
\(k\)-scheme.  The natural cyclic-cup map
\[
 c_U:\Gamma(U,\mathcal O_U^\times)\otimes
                 \Q_\ell/\Z_\ell
       \longrightarrow\Br(U)[\ell^\infty]
\]
from Lemma~\ref{lem:B-unit-cup-generation} has finite kernel.
No separatedness, reducedness, or smoothness hypothesis on
\(U\) is required.
\end{lemma}
\begin{proof}
The nilpotent units are killed by a bounded power of \(p\)
and disappear after tensoring with \(\Q_\ell/\Z_\ell\).
On the reduction the units form a finitely generated group by
Jaffe's Proposition~10.8.  Thus, after discarding torsion,
the source is \(\Lambda\otimes\Q_\ell/\Z_\ell\) for
a finite free abelian group \(\Lambda\)
\cite[Proposition~10.8]{BJaffe1994}.

Restrict its units to the function fields \(K_i\) of the
finitely many reduced irreducible components.  Choose a normal
projective model of each \(K_i/k\).  A rational function
with zero order at every prime divisor on such a model is a
global unit and hence lies in its finite constant field.
Consequently the divisor map on all these models is injective
on \(\Lambda\).  A finite number of their divisorial
valuations \(v_1,\ldots,v_t\) already gives an injective
map
\[
 \Lambda\longrightarrow\Z^t,
 \qquad u\longmapsto\bigl(v_j(u)\bigr)_j,
\]
since only finitely many divisors occur in a fixed finite set
of generators.  A zero-dimensional component contributes only
torsion units and can be omitted.

For a cyclic class \((\chi,u)\), its residue at \(v_j\)
is, up to the fixed cup-order sign,
\[
       v_j(u)\,\chi|_{\kappa(v_j)}.
\]
This is the unramified-character residue formula in
\citet[section~1.4.1, Proposition~1.4.5 and the cup-residue
formula immediately following it, author-draft p.~37]
       {DCTSkorobogatov2019}.
If the algebraic closure of \(k\) in \(\kappa(v_j)\)
has degree \(d_j\), restriction of the arithmetic character
multiplies \(\Q_\ell/\Z_\ell\) by \(d_j\).
Thus the residues of \(c_U\) contain the homomorphism
induced by the integer matrix
\[
             u\longmapsto\bigl(d_jv_j(u)\bigr)_j.
\]
It has full column rank.  Its Smith normal form shows that
the induced map on \(\Lambda\otimes\Q_\ell/\Z_\ell\)
has finite kernel.  The kernel of \(c_U\) is a subgroup of
this finite kernel, proving the assertion.  Restriction to
generic fields and their residues is legitimate regardless
of singularities or separation of the original scheme.
\end{proof}

\begin{theorem}[The full arbitrary-base Brauer assertion has the surface-Tate gate]
\label{thm:B-all-bases-Brauer-units}
\status{Proved equivalence; divisor Tate for surfaces remains unresolved}
Fix \(p\ne\ell\), put \(R=\Z_{(\ell)}\) and
\(D=\Q_\ell/\Z_\ell\), and for a finite-type
\(\F_p\)-scheme \(S\) put
\[
 A_S=H^2_{\et}(S,\Gm)\otimes R,
 \qquad C_S=\operatorname{im}
       \bigl(c_S:\Gamma(S,\mathcal O_S^\times)\otimes D
                         \longrightarrow A_S\bigr).
\]
The following universal assertions are equivalent:
\begin{enumerate}[label=(\roman*)]
\item Divisor Tate holds for every smooth projective surface
      over every finite extension of \(\F_p\).
\item \(A_S/C_S\) is a finitely generated \(R\)-module
      for every finite-type \(\F_p\)-scheme \(S\).
\item Unit cups generate \(\Br(S)[\ell^\infty]\) up
      to a finite cokernel for every such \(S\).
\end{enumerate}
In particular, under these equivalent hypotheses,
\[
 \bigl(H^2_{\et}(S,\Gm)/H^2_{\et}(S,\Gm)_{\rm tors}\bigr)
                  \otimes R
\]
is finitely generated for every \(S\).  These statements
include nonproper, singular, nonreduced, and nonseparated bases.
\end{theorem}
\begin{proof}
Assume (i).  Choose a finite affine open cover of \(S\),
and choose a proper smooth alteration cover of each of its
reduced irreducible components.  Their disjoint union gives
an \(h\)-cover \(q:T\to S\) with \(T\) smooth and
quasi-projective over finite fields (choose projective
alterations over the affine opens).  It need not be proper over a
nonseparated \(S\), and no such properness is claimed.
Every term \(T_a\) of its \v Cech nerve is a finite-type
\(\F_p\)-scheme.  The proof of
Lemma~\ref{lem:B-Gm-proper-kernel} uses only that the
perfected cover is a \(v\)-cover, so it applies to this
cover as well.  Hence
\[
 K=\ker\bigl(A_S\longrightarrow A_T\bigr)
\]
is a finitely generated \(R\)-module: the comparison
\(H^2_{\et}\hookrightarrow H^2_v\) and the degree-two
\v Cech filtration bound it by the localized Picard group
of \(T_1\) and the localized units of \(T_2\).

Since \(T\) is regular, \(A_T=\Br(T)[\ell^\infty]\).
It has finite cotype by Kummer and finite-coefficient \(\et\)
cohomology.  Lemma~\ref{lem:B-unit-cup-generation} gives
a positive integer \(N\), which can be taken to be a
power of \(\ell\), annihilating \(A_T/C_T\).

Write \(U_0,U_1,U_2\) for the unit groups of
\(S,T,T_1\), respectively, modulo their torsion.
They are finite free abelian groups.  Multiplicative-sheaf
descent on perfections gives the rational equalizer
\[
 0\longrightarrow U_0\otimes\Q
   \longrightarrow U_1\otimes\Q
   \xrightarrow{\partial}U_2\otimes\Q.
\]
The associated \(D\)-coefficient middle homology
\[
 \frac{\ker(U_1\otimes D\to U_2\otimes D)}
             {\operatorname{im}(U_0\otimes D)}
\]
is finite.  This follows directly from Smith normal form;
equivalently, its only possible contribution is the finite
\(\ell\)-torsion in the cokernel of a map of finite
free lattices.  Fix an integer \(L>0\) annihilating it.
By Lemma~\ref{lem:B-unit-cup-kernel}, fix \(M>0\)
annihilating \(\ker(c_{T_1})\).

For any \(a\in A_S\), choose \(u\in U_1\otimes D\)
such that
\[
             Nq^*a=c_T(u).
\]
The two pullbacks of \(q^*a\) to \(T_1\) agree.
Naturality of cyclic cups therefore gives
\(c_{T_1}(\partial u)=0\), and consequently
\(M\partial u=0\).  By the choice of \(L\), there
is \(v\in U_0\otimes D\) whose pullback is \(LMu\).
It follows that
\begin{equation}
        LMN\,a-c_S(v)\in K.
 \label{eq:B-global-unit-finite-error}
\end{equation}
All three integers are fixed for the cover; they do not
depend on \(a\).

The group \(C_S\) is divisible, being the image of a
divisible group, and is contained in \((A_S)_{\rm tors}\).
Consider the induced map
\[
       A_S/C_S\longrightarrow A_T/q^*C_S.
\]
Its kernel is \((K+C_S)/C_S\), a quotient of the
finitely generated module \(K\).  By
\eqref{eq:B-global-unit-finite-error}, its image is killed
by \(h=LMN\).  The target has finite cotype, since it is
a quotient of the finite-cotype \(\ell\)-primary group
\(A_T\).  Its subgroup killed by \(h\) is therefore
finite.  Thus \(A_S/C_S\) is an extension of a finite
module by a finitely generated \(R\)-module and is
finitely generated.  This proves (ii).

Because \(C_S\) is torsion, the torsion of \(A_S/C_S\)
is exactly \((A_S)_{\rm tors}/C_S\): if a multiple of
\(a\) lies in \(C_S\), then that multiple is torsion,
so \(a\) is torsion as well.  The torsion submodule of
a finitely generated \(R\)-module is finite.  Since
\((A_S)_{\rm tors}=\Br(S)[\ell^\infty]\), (ii)
implies (iii).  It also proves the stated finite generation
modulo full torsion.

Finally, assume (iii) and let \(S\) be a smooth projective
surface over a finite extension of \(\F_p\).
Its global units form a finite group, so \(C_S=0\).
Thus \(\Br(S)[\ell^\infty]\) is finite, and the
proper Kummer criterion used in
Theorem~\ref{thm:B-relative-surface-gate} gives divisor Tate.
This proves (i) and completes the equivalence.
\end{proof}

\begin{corollary}[Relative Brauer finiteness over every finite-type base]
\label{cor:B-all-bases-relative-Brauer}
\status{Proved equivalence with the universal surface-Tate assertion}
Assume the equivalent assertions of
Theorem~\ref{thm:B-all-bases-Brauer-units}.
If \(f:X\to S\) is a morphism of finite-type
\(\F_p\)-schemes such that pullback identifies their
global unit groups, then
\[
 \operatorname{coker}\bigl(
       \Br(S)[\ell^\infty]\longrightarrow
       \Br(X)[\ell^\infty]\bigr)
\]
is finite.  In particular this holds for every smooth proper
morphism with geometrically connected fibres, including every
curve family in B.3 over an arbitrary finite-type base.
Conversely, the universal assertion for those curve families
implies divisor Tate for all smooth projective surfaces over
finite extensions of \(\F_p\).
\end{corollary}
\begin{proof}
Functoriality of cups and equality of units give
\(f^*C_S=C_X\).  Thus the displayed cokernel is a
quotient of the finite group
\(\Br(X)[\ell^\infty]/C_X\).
For a smooth proper morphism with geometrically connected
fibres, the Stacks Project's Lemma~36.32.6 gives
\(f_*\mathcal O_X=\mathcal O_S\), hence equality of
units, including on nonreduced bases
\cite[Lemma~36.32.6, tag 0E0L]{BStacksConnectedUnits}.
The converse already
follows from the smooth regular-base Lefschetz-pencil
families with section in
Theorem~\ref{thm:B-relative-surface-gate}.
This proof uses actual torsion classes throughout and does
not replace a full Leray \(H^2(\Gm)\) term by its
torsion subgroup.
\end{proof}

\subsubsection{One finite group for all primes different from the characteristic}

\begin{lemma}[Prime-uniform generation by arithmetic unit cups]
\label{lem:B-unit-cup-prime-uniform}
\status{Proved conditional on divisor Tate for surfaces}
Let \(k\) be a finite field of characteristic \(p\), and assume
divisor Tate for all smooth projective surfaces over finite
extensions of \(k\), at one fixed prime \(\ell_0\ne p\).
Put \(D=(\Q/\Z)[p']\).  For every smooth quasi-projective
\(k\)-scheme \(T\), arithmetic cyclic cups give a map
\[
 c_T:\Gamma(T,\mathcal O_T^\times)\otimes D
       \longrightarrow \Br(T)[p']
\]
with finite cokernel.  The cokernel is one finite group across
all primes different from \(p\); separate finiteness of its
primary components is not substituted for this assertion.
\end{lemma}
\begin{proof}
First suppose that \(U\) is a dense open of a smooth projective
geometrically connected \(Y/k\), and put
\(\Gamma=\operatorname{Gal}(\bar k/k)\).  Let \(B\) be the
free group on the geometric prime divisors in \(Y\setminus U\).
Regularity gives the exact sequence
\[
 0\longrightarrow\bar k^\times
 \longrightarrow E=\Gamma(\bar U,\mathcal O_{\bar U}^\times)
 \longrightarrow B\longrightarrow\operatorname{Pic}(\bar Y)
 \longrightarrow\operatorname{Pic}(\bar U)\longrightarrow0.
\]
Boundary subsets of codimension at least two do not change these
units or Picard groups.  The lattice
\(N=E/\bar k^\times\subset B\) is finite free, with finite
\(\Gamma\)-action.

All positive cohomology of \(\bar k^\times\) under
\(\Gamma\) vanishes: in degree one this is Hilbert~90, or
surjectivity of \(x\mapsto x^{\#k-1}\), and in higher degrees
it follows from torsion cohomological dimension one.  Hence
global units of \(U\) surject onto \(N^\Gamma\), and
\[
 H^2(\Gamma,E)\simeq H^2(\Gamma,N)
 \simeq H^1(\Gamma,N\otimes\Q/\Z)
 \simeq N_\Gamma\otimes\Q/\Z.
\]
For the middle isomorphism, higher continuous cohomology of the
discrete rational representation \(N\otimes\Q\) vanishes by
averaging over finite quotients.  The last isomorphism is the
procyclic torsion-cohomology formula.  Since \(N\) has finite
Galois action, \(N^\Gamma\to N_\Gamma\) has finite cokernel.
Tensoring with the divisible group \(\Q/\Z\) makes it
surjective.  The same holds on the prime-to-\(p\) summand.
The induced map is precisely the arithmetic cyclic cup: its
degree-two cocycle is the Bockstein of the arithmetic character
multiplied by the invariant unit.  Consequently the image of
\(H^2(\Gamma,E)[p']\) in \(H^2(U,\Gm)\) is contained
in \(C_U=\operatorname{im}(c_U)\).

The other term needed in total degree two is finite over all
primes as well.  Let \(L\) be the image of \(B\) in
\(\operatorname{Pic}(\bar Y)\), and let
\(P=\operatorname{Pic}^0_{Y/k,\mathrm{red}}\).
The group \(P(\bar k)\) is torsion, so
\(F=L\cap P(\bar k)\) is finite.  With
\(Q=\operatorname{NS}(\bar Y)/\operatorname{im}(L)\),
there is an exact sequence
\[
 0\longrightarrow P(\bar k)/F
 \longrightarrow\operatorname{Pic}(\bar U)
 \longrightarrow Q\longrightarrow0.
\]
Frobenius minus identity on \(P\) is an isogeny, since its
differential is minus identity.  It is therefore surjective
on geometric points and \(H^1(\Gamma,P(\bar k))=0\).
Together with \(H^2(\Gamma,F)=0\), this gives
\(H^1(\Gamma,P(\bar k)/F)=0\).  The group \(Q\) is
finitely generated with finite Galois action, and
\(H^1(\Gamma,Q)\) is finite: its finite torsion part is
harmless, and the free part reduces to finite-group cohomology,
because a profinite group has no nonzero continuous homomorphism
to a discrete torsion-free lattice.  Thus
\(H^1(\Gamma,\operatorname{Pic}(\bar U))\) is finite.

The Hochschild--Serre filtration on the kernel of
\[
 H^2(U,\Gm)\longrightarrow H^2(\bar U,\Gm)^\Gamma
\]
has a quotient of \(H^2(\Gamma,E)\) and a subquotient of
\(H^1(\Gamma,\operatorname{Pic}(\bar U))\) as its two
associated groups.  Every \(H^2(\Gm)\) group here is torsion
because the schemes are regular.  The preceding calculations
therefore show that the prime-to-\(p\) part of this kernel,
modulo \(C_U\), is finite.  Arithmetic cups vanish geometrically.
No degeneration or splitting of the spectral sequence is used.

The remaining geometric input is prime-uniform.  Surface Tate,
the Kahn--Morrow reduction, and the unconditional divisor
splitting of Lemma~\ref{lem:B-divisor-splitting} imply
\(\Br(\bar U)^\Gamma[\ell_0^\infty]\) is finite by
geometric Kummer and finite cotype.  Qin's Theorem~5.2 then
implies
\[
                 \Br(\bar U)^\Gamma[p']\quad\text{is finite}.
\]
This theorem applies to smooth geometrically connected opens;
it does not require an additional full Frobenius-semisimplicity
conjecture~\cite[Theorem~5.2 and its proof in
Lemmas~3.2--3.4, Propositions~4.2--4.3, and Theorem~4.4]
{BQin2021Uniform}.
For precision concerning versions, the splitting and
prime-independence results cited in that proof as an earlier
Qin manuscript are Proposition~2.6 and Theorem~2.7 in the
opened August~2025 Ertl--Keller--Qin version
\cite[Proposition~2.6 and Theorem~2.7]{BErtlKellerQin2025}.
The arithmetic kernel calculation now proves that
\(\Br(U)[p']/C_U\) is finite.

For general geometrically connected smooth quasi-projective
\(T\), choose a projective generically finite alteration
\(g:U\to T\) of degree \(d\), where \(U\) is open in
a smooth projective scheme.  Apply de~Jong's alteration theorem
to a projective compactification; source and target have ample
line bundles~\cite[Theorem~4.1 and Remark~4.2]{c_dejong1996}.
For every \(n\) prime to \(p\), the finite-coefficient Gysin
map gives
\[
 g_*:H^2(U,\mu_n)\longrightarrow H^2(T,\mu_n),
 \qquad g_*g^*=d.
\]
The scalar identity does not require \(d\) invertible modulo
\(n\): the composite is an endomorphism of the constant sheaf,
and its value on the dense finite-flat locus is multiplication
by \(d\).  Riou's finite-flat normalization includes inseparable
generic degree~\cite[Proposition~2.5.13, Remark~2.5.14,
Corollary~3.1.2, and Lemmas~3.5.8--3.5.10]{c_riou2014}.

These maps descend to compatible maps on \(\Br[n]\).
Indeed, a Kummer class from \(\operatorname{Pic}(U)/n\)
restricts to zero at the generic field.  Its trace has zero
generic Brauer class, hence zero Brauer class on regular \(T\),
by injectivity into the function-field Brauer group
\cite[Proposition~1.4 and Corollary~1.8]{BGrothendieckBrauerII1966}.
Thus the induced trace is independent of the Kummer lift; no
unstated integral derived trace on \(\Gm\) is needed.
The projection formula and the Kummer norm identity give
\[
 \operatorname{cor}_g(\chi|_U,u)
             =(\chi,\operatorname{Nm}_{k(U)/k(T)}u).
\]
The norm of a global unit on \(U\) is a global unit on \(T\).
Properness gives centers for every extension of every
codimension-one valuation of normal \(T\); all orders of the
unit are zero, so their norm valuation is zero.  The degree-one
Kummer trace identity can be checked on the dense finite-flat
locus and extended by the injectivity of
\(H^1(T,\mu_n)\to H^1(T^\circ,\mu_n)\) on normal schemes.
This argument includes purely inseparable factors.

Pullback and trace therefore act on
\(B_T=\Br(T)[p']/C_T\) and \(B_U=\Br(U)[p']/C_U\),
with composite \(d\).  The first case, applied over the finite
constant field of \(U\), makes \(B_U\) finite; the character
rescaling in the final paragraph identifies its cup image with
the one used here.  Thus one
integer \(h>0\) annihilates \(B_T\).  The subgroup \(C_T\)
is divisible, so \(B_T[h]\) is a quotient of
\(\Br(T)[h]\), after removing the \(p\)-power factor of
\(h\).  Kummer and finite-coefficient cohomology make the latter
finite.  Hence \(B_T=B_T[h]\) is finite.

Finally, if a connected component has finite constant field
\(k'/k\), the character normalized over \(k\) restricts to
\([k':k]\) times the one normalized over \(k'\).
Multiplication by this integer on \(D\) is surjective, so
the cup images coincide.  Finite disjoint unions complete the
proof.
\end{proof}

\begin{theorem}[Whole prime-to-characteristic Brauer finiteness and the exact gate]
\label{thm:B-Brauer-prime-uniform}
\status{Proved equivalence; the surface-Tate condition remains unresolved}
Fix \(p\), put \(R=\Z[1/p]\) and \(D=(\Q/\Z)[p']\),
and for every finite-type \(\F_p\)-scheme \(S\) define
\[
 A_S=H^2_{\et}(S,\Gm)\otimes R,
 \qquad C_S=\operatorname{im}
 \bigl(\Gamma(S,\mathcal O_S^\times)\otimes D\to A_S\bigr).
\]
The following conditions are equivalent:
\begin{enumerate}[label=(\roman*)]
\item Divisor Tate holds, at one prime different from \(p\),
      for all smooth projective surfaces over all finite
      extensions of \(\F_p\).
\item \(A_S/C_S\) is a finitely generated \(R\)-module
      for every finite-type \(\F_p\)-scheme \(S\).
\item \(\Br(S)[p']/C_S\) is a finite group for every such
      \(S\).
\item The whole prime-to-\(p\) Brauer cokernel is finite for
      every proper smooth geometrically connected curve family
      over every such \(S\).
\end{enumerate}
Under these conditions the relative Brauer conclusion also holds
for every finite-type morphism identifying global units,
in particular every smooth proper morphism with geometrically
connected fibres, in arbitrary relative dimension.
All statements include singular, nonreduced, nonproper, and
nonseparated bases.
\end{theorem}
\begin{proof}
Assume (i).  Repeat the proof of
Theorem~\ref{thm:B-all-bases-Brauer-units}, now over
\(R=\Z[1/p]\).  We give the points that make this a single
prime-uniform argument.  Choose the same finite-type
\(h\)-cover \(q:T\to S\), with \(T\) a finite disjoint
union of smooth quasi-projective schemes.  The kernel
\[
 K=\ker(A_S\to A_T)
\]
is finitely generated over \(\Z[1/p]\) itself by
Lemma~\ref{lem:B-Gm-proper-kernel}; its proof used the
localized Picard and unit groups, not a prime-dependent bound.
Since \(T\) is regular, \(A_T=\Br(T)[p']\).
Lemma~\ref{lem:B-unit-cup-prime-uniform} supplies one integer
\(N\) annihilating \(A_T/C_T\).

The unit-cup kernel in Lemma~\ref{lem:B-unit-cup-kernel}
is finite for \(D\) as well: the proof uses one full-rank
integer valuation matrix.  Its Smith invariant factors have
finite prime-to-\(p\) torsion, supported at finitely many
primes.  Likewise the rationally exact sequence of finite unit
lattices for \(S,T,T_1\) has finite middle homology after
tensoring with \(D\).  Choose fixed integers \(M,L\)
annihilating these two finite groups.  The exact same cup and
\v Cech calculation gives, for every \(a\in A_S\), a
\(v\in\Gamma(S,\mathcal O_S^\times)\otimes D\) with
\[
                      LMN\,a-c_S(v)\in K.
\]
Thus the image of \(A_S/C_S\to A_T/q^*C_S\) is killed
by the fixed integer \(h=LMN\), and its kernel is a quotient
of \(K\).

One does not need, and does not assume, a finite-cotype
description summed over all primes.  Instead, \(q^*C_S\)
is divisible.  If \(hb\in q^*C_S\), choose
\(c\in q^*C_S\) with \(hc=hb\); then \(b-c\in A_T[h]\).
Consequently
\[
              A_T[h]\longrightarrow(A_T/q^*C_S)[h]
              \quad\text{is surjective}.
\]
The source is finite by Kummer and finite-coefficient cohomology,
using the prime-to-\(p\) part of \(h\).
The image of \(A_S/C_S\) is therefore finite, and its kernel
is finitely generated, proving (ii).

Localization cannot turn a nontorsion element into a torsion
element.  Hence \((A_S)_{\rm tors}=\Br(S)[p']\).
Since \(C_S\) is torsion, the torsion submodule of \(A_S/C_S\)
is exactly \(\Br(S)[p']/C_S\).  A finitely generated torsion
\(\Z[1/p]\)-module is finite, giving (iii).
If pullback identifies units, functoriality gives
\(f^*C_S=C_X\); its relative Brauer cokernel is a quotient
of \(\Br(X)[p']/C_X\), hence finite.  For the smooth proper
families in (iv), equality of units follows from the
base-change-compatible identity \(f_*\mathcal O_X=\mathcal O_S\)
\cite[Lemma~36.32.6, tag 0E0L]{BStacksConnectedUnits}.
This proves (iv) and the stronger concluding assertion.

Finally (iv) implies the fixed-prime relative-curve assertion,
which implies (i) by the pencil argument of
Theorem~\ref{thm:B-relative-surface-gate}.
Alternatively (iii) applied to a smooth projective surface has
\(C_S=0\), because its units are finite, and the proper
Kummer criterion again gives (i).  This completes the
equivalences.  The characteristic-primary counterexample in
Theorem~\ref{thm:B-Brauer-characteristic} is unaffected.
\end{proof}

\begin{corollary}[Prime-uniform abelian torsors and an integral lattice]
\label{cor:B-torsors-prime-uniform}
\status{Proved equivalence with the same unresolved surface-Tate condition}
Under the equivalent conditions of
Theorem~\ref{thm:B-Brauer-prime-uniform}, every abelian scheme
\(A/S\) over a finite-type \(\F_p\)-scheme satisfies
\[
 H^1_{\et}(S,A)[p']\ \text{finite},
 \qquad
 H^1_{\et}(S,A)\otimes\Z[1/p]
       \ \text{finitely generated over }\Z[1/p].
\]
Conversely, the universal first assertion implies the surface-Tate
condition.  These are statements about one whole prime-to-\(p\)
group, including all primes simultaneously.  They include
nonreduced and nonseparated bases and abelian schemes without
a chosen global polarization.
\end{corollary}
\begin{proof}
Let \(f:B=A^\vee\to S\), let \(e\) be its zero section,
and write \(P=\operatorname{Pic}_{B/S}=R^1f_*\Gm\) on the
\(\et\) site.  Relative Picard duality identifies
\(P^0\) with \(A\).  Here \(P^0\) means fibrewise algebraic
equivalence to zero on every test scheme; it remains the dual
abelian scheme over a nonreduced base.  In particular we use the
arbitrary-abelian-scheme form, not just the projective form of
the theorem~\cite[section~2.3.2, Remark~2.3.6 and Theorem~2.3.7]
{BRobert2023AbelianSchemes}.
Set \(N=P/A\) as an \(\et\)-sheaf; no local constancy of
\(N\) is asserted or needed.  Inversion on \(B\) acts as
\(-1\) on \(A\), by duality, and as \(+1\) on \(N\).
For the latter assertion, \([-1]^*L\otimes L^{-1}\) is
algebraically trivial on every geometric fibre; this is exactly
the defining fibrewise criterion for \(P^0\).  The underlying
line-bundle identities are the theorem-of-the-square identities
in Lemma~7.16(ii) and Corollaries~7.23(ii) and~7.25 of the
opened duality chapter~\cite{BMoonenDuality2011}.

Naturality of the connecting homomorphism in
\[
 0\longrightarrow A\longrightarrow P\longrightarrow N
       \longrightarrow0
\]
therefore gives
\[
 -\delta(n)=\delta([-1]^*n)=\delta(n),
 \qquad 2\delta(n)=0\quad(n\in H^0(S,N)).
\]
Consequently the kernel of \(H^1(S,A)\to H^1(S,P)\)
is killed by two.

The zero section splits the unit map in the derived category:
\[
       Rf_*\Gm\simeq\Gm\oplus K,
       \qquad K\in D^{\ge1}(S_{\et}),
       \qquad\mathcal H^1(K)=P.
\]
The truncation triangle
\(P[-1]\to K\to\tau_{\ge2}K\) has
\(\mathbb H^1(S,\tau_{\ge2}K)=0\).  Hence it induces
an actual injection
\[
 H^1(S,P)\lhook\joinrel\longrightarrow
 H^2(S,K)\lhook\joinrel\longrightarrow H^2(B,\Gm).
\]
Its image is in the kernel of \(e^*\), and therefore meets
\(f^*H^2(S,\Gm)\) trivially.  In particular torsion inputs
give actual torsion Brauer classes; no lift from a torsion
quotient of a nontorsion group is being assumed.  We obtain
\[
 H^1(S,A)[p']\longrightarrow
 \operatorname{coker}\bigl(\Br(S)[p']\to\Br(B)[p']\bigr)
\]
with kernel killed by two.  If \(p=2\), that kernel is zero.
If \(p\ne2\), it is finite because Kummer gives a surjection
\(H^1(S,A[2])\to H^1(S,A)[2]\), and the finite \(\et\)
sheaf \(A[2]\) has finite cohomology over the finite-type
finite-field scheme \(S\).  A finite affine open cover proves
the same finiteness when \(S\) is nonseparated.
The target is finite by
Theorem~\ref{thm:B-Brauer-prime-uniform}, since \(B\to S\)
is smooth proper with geometrically connected fibres.  This
proves the first assertion.

For the second assertion, repeat the proof of
Proposition~\ref{prop:B-integral-torsor-lattice} over
\(\Z[1/p]\).  Its actual \v Cech-cocycle map already
embeds the proper-cover restriction kernel, after inverting
\(p\), in a subquotient of a finitely generated
\(\Z[1/p]\)-module of abelian-scheme sections.  On the regular
alteration the torsor group is torsion; its localization is
now finite by the first assertion, simultaneously at every
prime different from \(p\).  The same extension argument,
bounded \(p\)-primary nilpotent comparison, and finite open-cover
descent therefore prove the stated finite generation.
Finally, the universal first assertion implies its fixed-prime
version, whose surface-Tate equivalence is
Theorem~\ref{thm:B-all-bases-surface-gate}.
\end{proof}

\begin{lemma}[The prime in the divisor conjecture for a surface]
\label{lem:B-surface-Tate-prime}
\status{Proved}
Let \(Y\) be a smooth projective geometrically integral surface over a
finite field \(k\) of characteristic \(p\).  Surjectivity of
\[
 \operatorname{Pic}(Y)\otimes\Q_\ell\longrightarrow
 H^2(Y_{\bar k},\Q_\ell(1))^{\Gamma_k}
\]
for one \(\ell\ne p\) is equivalent to its surjectivity for every
\(\ell\ne p\).  It is also equivalent to the multiplicity of the
eigenvalue \(1\) of Frobenius on \(H^2(Y_{\bar k},\Q_\ell(1))\)
being \(\operatorname{rank}\operatorname{NS}(Y_{\bar k})^{\Gamma_k}\).
\end{lemma}

\begin{proof}
Put \(N=\operatorname{NS}(Y_{\bar k})\otimes\Q\).  Its intersection
pairing is nondegenerate by the Hodge index theorem.  The Galois action
on \(N\) factors through a finite group and preserves that pairing.
Averaging over that group shows that the restriction to \(N^{\Gamma_k}\)
is nondegenerate: an invariant vector orthogonal to all invariant
vectors is orthogonal to every vector.  Norms of divisors show that
the image of \(\operatorname{Pic}(Y)\otimes\Q\) is precisely
\(N^{\Gamma_k}\).  Divisor cycle classes embed this space into the
Frobenius invariants for every \(\ell\).

Suppose that surjectivity holds for one \(\ell\), and write \(V\) for
the corresponding degree-two twisted cohomology and \(F\) for
Frobenius.  The subspace \(V^F=N^{\Gamma_k}\otimes\Q_\ell\) is
nondegenerate for the cup product.  Its orthogonal complement is
\(F\)-stable and contains no fixed vector.  Thus \(F-1\) is invertible
on that complement; in particular the generalized \(1\)-eigenspace
equals \(V^F\).  The characteristic polynomial of Frobenius on
degree-two cohomology is independent of \(\ell\).  For every other
prime, the algebraic invariant subspace therefore already has the
full dimension of the generalized \(1\)-eigenspace.  This proves the
assertions, including the converse multiplicity criterion.  No
general Frobenius semisimplicity conjecture is used.
\end{proof}

\begin{theorem}[A smaller surface test class: vanishing coherent irregularity]
\label{thm:B-irregularity-zero-surface-gate}
\status{Proved equivalence; the restricted Tate conjecture remains unresolved}
Fix a finite field \(k\) and a prime \(\ell\ne\operatorname{char}k\).
The divisor Tate conjecture at \(\ell\) for every smooth projective
geometrically integral surface over \(k\) is equivalent to the same
assertion restricted to surfaces \(Y/k\) satisfying
\[
                         H^1(Y,\mathcal O_Y)=0.
\]
Consequently every assertion of
Theorem~\ref{thm:B-relative-surface-gate}, and the equivalent universal
relative Brauer assertions proved above, can be tested using only
these surfaces over the finite extensions of \(\F_p\).
This is a proper subclass restriction: for an elliptic curve \(E/k\),
the surface \(E\times\mathbb P^1_k\) has
\(\dim_kH^1(\mathcal O)=1\) and is excluded.  The restricted hypothesis
imposes only the displayed intrinsic vanishing; it does not retain
an arbitrary original surface as auxiliary data.
\end{theorem}

\begin{proof}
The implication to the restricted class is immediate.  For the
converse, Lemma~\ref{lem:B-surface-Tate-prime} lets us use the
prime-independent notation \(T^1\) of
\citet[Conjecture~1.5 and Theorem~1.6, pp.~1967--1968]{DYuan2021Hodge}.
We give the reduction, keeping its base field and its zero-dimensional
case explicit.  The full published proof is in that paper's
Section~4.2, pp.~1996--1997; its preparatory results are
Lemmas~4.2 and~4.4 and Theorems~4.3 and~4.5,
pp.~1994--1996.

Let \(X/k\) be any smooth projective geometrically integral surface.
The finite-field Lefschetz-pencil construction used by Yuan produces
a smooth blowup \(X'\) and a semistable fibration
\(X'\to\mathbb P^1_k\), over this same \(k\).
Write \(J/k(t)\) for the Jacobian of its smooth generic fibre.
Blowups add only algebraic divisor classes; the Artin--Tate
comparison then gives
\[
                  T^1(X)\ \Longleftrightarrow\ T^1(X')
                         \ \Longleftrightarrow\ \operatorname{BSD}(J),
\]
where \(\operatorname{BSD}\) here means equality of the analytic and
Mordell--Weil ranks, not the conjectural leading-coefficient formula.

By Yuan's Theorem~1.1, p.~1965, there is an isogeny over \(k(t)\)
\[
                         J\sim A\times C_{k(t)},
\]
where \(C/k\) is an abelian variety and the Hodge bundle of \(A\)
on \(\mathbb P^1_k\) is ample.  Semi-abelian reduction is preserved
under isogeny, so \(A\) has it everywhere.  The constant factor has
rank zero: a map \(\mathbb P^1_k\to C\) is constant, hence its
Mordell--Weil group is \(C(k)\).  Its analytic rank is also zero,
as the Frobenius eigenvalues of \(C\) have absolute value
\(\sqrt{\#k}\).  Isogeny invariance and additivity of the ranks,
together with the known nonnegativity of each analytic-rank defect,
give \(\operatorname{BSD}(J)\Longleftrightarrow
\operatorname{BSD}(A)\).  If \(A=0\), this already proves
\(T^1(X)\), and nothing further is required.

Assume \(\dim A\ge1\).  Yuan's Theorem~4.5 supplies a projective flat
regular model \(\psi:P\to\mathbb P^1_k\) of \(A^\vee\) with
\[
                    R^1\psi_*\mathcal O_P
                    \simeq\operatorname{Lie}(\mathcal A/\mathbb P^1_k),
\]
where \(\mathcal A\) is the N\'eron model of \(A\).  This bundle is
the dual of its ample Hodge bundle.  It has no global sections.
Since \(\psi_*\mathcal O_P=\mathcal O_{\mathbb P^1_k}\), the
coherent Leray sequence yields \(H^1(P,\mathcal O_P)=0\).
Regularity over the perfect field \(k\) makes \(P\) smooth.

If \(\dim P=2\), set \(Y=P\).  Otherwise successively take
sufficiently ample smooth hypersurface sections over \(k\) until
obtaining a surface \(Y\subset P\).  At each step choose a smooth
closed fibre over a point of sufficiently large residue degree.
Poonen's Bertini density theorem
\citep[Theorem~1.1, author-version p.~2]{DPoonen2004Bertini}
makes it possible to choose the
section simultaneously smooth on the total space and on that fibre:
the first density is fixed and positive, while the second tends to
one with that residue degree.  Smoothness of the generic fibre
follows.  Serre vanishing and duality give, at each cut,
\[
 H^1(Y,\mathcal O_Y)=0,
 \qquad
 H^1(P_\eta,\mathcal O_{P_\eta})
       \lhook\joinrel\longrightarrow H^1(Y_\eta,\mathcal O_{Y_\eta}),
\]
with the evident interpretation for the successive intermediate
sections.  For the first equality the required vanishing is
\(H^2(P,\mathcal O_P(-Y))=0\); for the injection on a generic
fibre it is \(H^1(P_\eta,\mathcal O_{P_\eta}(-Y_\eta))=0\).
The same ampleness ensures geometric connectedness.

Let \(B\) be the Jacobian of the resulting smooth generic curve
\(Y_\eta\).  Pullback of line bundles gives \(A\to B\), whose map
on tangent spaces is the displayed injection.  Its kernel is
finite, and Poincar\'e reducibility makes \(A\) an isogeny factor of
\(B\).  Thus
\[
 T^1(Y)\ \Longrightarrow\ \operatorname{BSD}(B)
       \ \Longrightarrow\ \operatorname{BSD}(A)
       \ \Longrightarrow\ T^1(X).
\]
Our hypothesis supplies the first assertion because
\(H^1(Y,\mathcal O_Y)=0\).  Every construction was over \(k\).
Applying this fixed-field result for every finite \(k/\F_p\)
handles all the constant fields in
Theorem~\ref{thm:B-relative-surface-gate}; viewing a geometrically
connected surface over its field of constants and using induction
of its Galois representation also handles the non-geometrically
connected surfaces appearing there.
\end{proof}

\begin{remark}[Sources and scope of the reduction]
\label{rem:B-irregularity-zero-source-scope}
\status{Verified source scope}
\citet[Remark~4.13, p.~56]{BErtlKellerQin2025} points to the
all-dimensional reduction and the additional restriction above.
The existing Kahn--Morrow reduction supplies the passage from
arbitrary dimension to surfaces; Yuan's Theorem~1.6 supplies the
further vanishing of coherent irregularity.  Yuan's cited
Nguyen thesis and K\"unnemann--Faltings--Chai inputs are used through
his exact published statements and proof; they are not represented
as independently opened here.  De Jong's unpublished reduction note
has not been obtained.  None of these reductions proves the Tate
conjecture for the remaining irregularity-zero surfaces.
\end{remark}

\begin{corollary}[Canonical surfaces suffice, with arbitrary numerical thresholds]
\label{cor:B-canonical-surface-gate}
\status{Proved strict test-class reduction; the restricted Tate conjecture remains unresolved}
Fix a finite field \(k\), a prime \(\ell\ne\operatorname{char}k\),
an integer \(N\ge1\), and a real number \(\varepsilon>0\).
Divisor Tate at \(\ell\) for every smooth projective geometrically
integral surface over \(k\) is equivalent to its restriction to
surfaces \(Y/k\) with
\[
 H^1(Y,\mathcal O_Y)=0,\qquad \omega_Y\text{ very ample},
 \qquad K_Y^2\ge N,\qquad p_g(Y)\ge N,
\]
and, if desired, the simultaneous extra conditions
\[
             c_2(Y)>0,\qquad
             \left|\frac{K_Y^2}{c_2(Y)}-1\right|<\varepsilon.
\]
Here \(c_2(Y)=\deg c_2(T_Y)\).
Thus the preceding universal Brauer and abelian-torsor gates can
be confined to an intrinsic class of canonically embedded surfaces
of general type.  This does not assert Tate for those surfaces.
\end{corollary}
\begin{proof}
Use the construction in the preceding theorem for an arbitrary
surface \(X/k\).  If its ample-Hodge factor \(A\) is zero,
\(T^1(X)\) is already proved.  Otherwise let
\(P\to\mathbb P^1_k\) be the smooth projective regular model
of \(A^\vee\) constructed there, with \(H^1(P,\mathcal O_P)=0\).
If \(\dim P=2\), replace \(P\) by \(P\times_k\mathbb P^1_k\),
using projection to the original base.  K\"unneth leaves
\(H^1(\mathcal O)\) unchanged, and its smooth generic fibre has
the same \(\operatorname{Pic}^0=A\).  Thus the new ambient
dimension is at least three in every case.

Perform the high-degree smooth cuts of Yuan's proof only until
reaching a smooth projective threefold \(W\) over \(k\).
They preserve \(H^1(W,\mathcal O_W)=0\), a smooth generic
fibre over \(\mathbb P^1_k\), and the injection from the
original generic-fibre \(H^1(\mathcal O)\).
Choose a very ample line bundle \(L\) on \(W\), with
\(H=c_1(L)\).  The final cut \(Y\in|L^d|\) may be taken
for arbitrarily large \(d\), simultaneously smooth on \(W\)
and on one fixed sufficiently high-degree smooth closed fibre,
by the positive-density argument already proved above.
Serre duality and vanishing give
\[
 H^2(W,L^{-d})=0,\qquad
 H^1(W_\eta,L^{-d}|_{W_\eta})=0
 \quad(d\gg1).
\]
Consequently \(H^1(Y,\mathcal O_Y)=0\), and the resulting
generic Jacobian still contains \(A\) up to isogeny.
The same high-degree choices ensure geometric connectedness.

There is one more threshold, compatible with all those above:
\(\omega_W\otimes L^d\) is very ample for every sufficiently
large \(d\)~\cite[Lemma~29.40.8, tag 0FVC]{BStacksCanonicalCuts}.
Taking determinants in the exact conormal sequence gives
\[
                  \omega_Y=(\omega_W\otimes L^d)|_Y,
\]
which is therefore very ample.  This determinant argument is
valid in every characteristic; the conormal sequence is short
exact because both \(W\) and \(Y\) are smooth
\cite[section~37.63, tag 06BB]{BStacksCanonicalCuts}.

All numerical conditions can be imposed in this same cut, not
as separate later reductions.  Put \(K=c_1(\omega_W)\).
Adjunction and the normal-bundle sequence give
\[
 \begin{aligned}
 K_Y^2&=d^3H^3+2d^2KH^2+dK^2H,\\
 c_2(Y)&=d^3H^3+d^2KH^2+d\,c_2(W)H.
 \end{aligned}
\]
Their common leading coefficient is \(H^3>0\), so
\(K_Y^2\to\infty\), \(c_2(Y)>0\) for large \(d\),
and \(K_Y^2/c_2(Y)\to1\).
The exact divisor sequence also gives
\[
 \chi(\mathcal O_Y)
   =\chi(\mathcal O_W)-\chi(L^{-d})
   =\frac{H^3}{6}d^3+O(d^2).
\]
Since \(Y\) is geometrically connected and \(H^1(\mathcal O_Y)=0\),
Serre duality identifies
\(p_g(Y)=\chi(\mathcal O_Y)-1\), which tends to infinity.
Increasing the one threshold for \(d\) therefore gives all
the displayed inequalities.  The proof of the preceding theorem
still gives \(T^1(Y)\Rightarrow T^1(X)\), over the original
fixed field \(k\).  This proves the equivalence.
\end{proof}

This restriction is intrinsic: it records no arbitrary original
surface as a blowup centre or other additional datum.
For example, \(\mathbb P^2_k\) has \(H^1(\mathcal O)=0\)
but is excluded by the canonical ampleness condition.  The
numerical refinements above are consequences of the same
high-degree-cut argument and are counted as part of this one
reduction, not as successive movements of the conjectural gate.

\begin{theorem}[A geometrically simply connected surface test]
\label{thm:B-simply-connected-surface-gate}
\status{Proved strict intrinsic test-class equivalence; restricted Tate unresolved}
Fix a finite field \(k\) of characteristic \(p\), a prime
\(\ell\ne p\), an integer \(N\geq1\), and a real number
\(\varepsilon>0\).  Divisor Tate for all smooth projective
geometrically integral surfaces over \(k\) is equivalent to
divisor Tate for just those surfaces \(Y/k\) satisfying
\[
 \pi_1^{\et}(Y_{\overline k})=1,
 \qquad H^1(Y,\mathcal O_Y)=0,
 \qquad \omega_Y\text{ very ample},
\]
and, simultaneously,
\[
 K_Y^2\geq N,\qquad p_g(Y)\geq N,\qquad c_2(Y)>0,
 \qquad \left|\frac{K_Y^2}{c_2(Y)}-1\right|<\varepsilon.
\]
No extension of the fixed field \(k\) is needed in the reduction.
Consequently the universal prime-to-characteristic torsor and
relative-Brauer finiteness assertions of
Theorem~\ref{thm:B-Brauer-prime-uniform} and
Corollary~\ref{cor:B-torsors-prime-uniform} have this smaller
intrinsic surface test class, when quantified over finite constant
extensions as before.  The numerical conditions are inherited from
Corollary~\ref{cor:B-canonical-surface-gate}, not independent
additional reductions.
\end{theorem}
\begin{proof}
The implication from all surfaces to the indicated subclass is
immediate.  For the converse start with a smooth projective
geometrically integral surface \(X/k\).  The already audited
proof of Yuan's Theorem~1.6 reduces its divisor Tate assertion,
over this same \(k\), to BSD for a semistable abelian variety
\(A/K\), where \(K=k(t)\), whose N\'eron Hodge bundle over
\(\mathbf P^1_k\) is ample.  If \(A=0\), the original Tate
assertion is already known.  Otherwise Yuan's Theorem~4.5,
applied to \(A^\vee\), gives a projective flat regular model
\(\psi:P\to\mathbf P^1_k\), with generic fibre \(A^\vee\),
and
\[
 R^1\psi_*\mathcal O_P\simeq
       \operatorname{Lie}(\mathcal A/\mathbf P^1_k),
 \qquad H^1(P,\mathcal O_P)=0.
\]
Here \(\mathcal A\) is the N\'eron model of \(A\); the
vanishing follows from ampleness of the dual bundle and Leray
over \(\mathbf P^1\)
\citep[Theorem~4.5 and Section~4.2, author pp.~33--36]{DYuan2021Hodge}.
The generic zero section extends over every point of the regular
base curve by properness, giving a global section and a
\(k\)-rational point.  Since \(k\) is perfect, \(P\) is smooth
over \(k\).

\emph{The geometric fundamental group of \(P\).}
Shrink to a nonempty open \(U\subset\mathbf P^1_k\) on which
the model is the abelian scheme extending its generic fibre.
The proper smooth homotopy sequence gives a right-exact sequence
\[
 \pi_1(P_{\overline\eta})\longrightarrow
 \pi_1(P_{U,\overline k})\longrightarrow
 \pi_1(U_{\overline k})\longrightarrow1,
\]
split on the right by the section
\citep[Expos\'e~X, Theorem~1.3 and Corollary~1.4,
pp.~202--203]{c_sga1_2003}.
No left injectivity is asserted.  Normality makes
\(\pi_1(P_{U,\overline k})\to\pi_1(P_{\overline k})\)
surjective.  The section subgroup dies in the latter group,
because it factors through \(\pi_1(\mathbf P^1_{\overline k})=1\).
Thus
\(G=\pi_1(P_{\overline k})\) is a quotient of the abelian
group \(\pi_1(P_{\overline\eta})\), and is itself abelian
\citep[Lemma~58.10.7]{BStacksFundamentalGroups};
\citep[Expos\'e~XI, Section~2 and Theorem~2.1,
pp.~221--222]{c_sga1_2003}.

The vanishing of \(H^1(P_{\overline k},\mathcal O)\) makes
the tangent space of the connected Picard scheme zero, so that
connected Picard scheme is trivial.  For every \(m\) prime to
\(p\), Kummer therefore gives
\[
 H^1(P_{\overline k},\mu_m)
       \simeq\operatorname{Pic}(P_{\overline k})[m]
       =\operatorname{NS}(P_{\overline k})[m].
\]
The torsion subgroup of N\'eron--Severi is one finite group;
hence the entire prime-to-\(p\) character group of \(G\) is
finite.  Artin--Schreier gives
\(H^1(P_{\overline k},\Z/p)=0\), since \(F-1\) is
surjective on \(\overline k\) and \(H^1(\mathcal O)=0\).
An abelian profinite group with a nonzero pro-\(p\) part has a
quotient \(\Z/p\).  That part is therefore zero.  Characters
separate points of an abelian profinite group, so \(G\) is
finite, of order prime to \(p\)
\citep[Expos\'e~XI, Proposition~6.4, Corollary~6.6 and its
following paragraph, pp.~234--235; Propositions~6.7--6.8 and
Corollary~6.10, pp.~236--237]{c_sga1_2003}.

\emph{A cover over the original field.}
Use the rational point on the zero section to split the arithmetic
short exact sequence
\[
 1\longrightarrow G\longrightarrow\pi_1(P)
      \longrightarrow\Gal(\overline k/k)\longrightarrow1.
\]
This arithmetic sequence, unlike the general homotopy sequence
above, is left exact as well
\citep[Expos\'e~IX, Theorem~6.1, p.~195]{c_sga1_2003};
\citep[Lemma~58.14.3]{BStacksFundamentalGroups}.
The splitting image is a closed subgroup of finite index
\(n=|G|\).  It defines a connected finite \'etale cover
\(P'\to P\) of degree \(n\), geometrically connected with
\(\pi_1(P'_{\overline k})=1\).  The whole section
\(\mathbf P^1_k\to P\) lifts: its fundamental-group image
is exactly this splitting subgroup.  In particular this
construction requires no finite extension of \(k\).

Its generic fibre \(B=P'_K\) is integral, being a localization
of the integral scheme \(P'\).  It is a finite \'etale cover
of \(A^\vee\) and has a \(K\)-point over zero.  Its field
of global functions is consequently \(K\), since it embeds
in \(K\) via that point.  Proper smooth base change of
\(H^0(\mathcal O)\) proves geometric connectedness.  The
pointed finite \'etale cover has a unique abelian-variety
structure making
\[
             B\longrightarrow A^\vee
\]
an isogeny of degree \(n\).  The geometric assertion is the
complete pointed-cover argument in SGA~1, Expos\'e~XI,
Section~2, pp.~221--222; uniqueness descends the group law
to \(K\).  Recall that \(p\nmid n\).

\emph{Rebuilding the model and preserving the Hodge bundle.}
One must not assume that \(H^1(P',\mathcal O_{P'})=0\).
Instead dualize the isogeny to obtain \(A\to B^\vee\).
There is a quasi-inverse whose two composites are \([n]\).
Both morphisms extend to their N\'eron models, and the induced
maps on invariant differentials compose to multiplication by
\(n\).  Because \(n\) is invertible on \(\mathbf P^1_k\),
the Hodge bundles of \(A\) and \(B^\vee\) are isomorphic.
The latter is therefore ample.  Semi-abelian reduction is
isogeny invariant
\citep[Section~1.2, Definition~1; Section~7.3, Lemma~5,
Proposition~6 and Corollary~7, pp.~180--181]{BBoschLutkebohmertRaynaud1990}.

Set \(A_{\mathrm{new}}=B^\vee\) and repeat Yuan's model
step, applying Theorem~4.5 to
\(A_{\mathrm{new}}^\vee=B\).  It constructs a model of
the identical generic abelian variety \(B\), over the same
\(K\), without another isogeny or a base extension.  Denote
this smooth projective model by \(P''\).  Its formula
\[
 R^1\psi''_*\mathcal O_{P''}
       \simeq\operatorname{Lie}(\mathcal B^\vee/\mathbf P^1_k)
\]
and the ample dual Hodge bundle give \(H^1(P'',\mathcal O)=0\).
The two smooth proper varieties \(P'\) and \(P''\) are
birational, since both model \(B/K\).  Purity and birational
invariance of the fundamental group therefore give
\[
                   \pi_1(P''_{\overline k})=1
\]
\citep[Expos\'e~X, Corollaries~3.3--3.4 and their proofs,
pp.~213--214]{c_sga1_2003}.

\emph{The surface cut.}
Apply the complete high-degree-cut construction in
Corollary~\ref{cor:B-canonical-surface-gate} to \(P''\).
When \(\dim P''=2\), first use \(P''\times\mathbf P^1\);
this preserves its geometric fundamental group, coherent
\(H^1\), and the generic Picard-zero variety \(B^\vee\).
At every subsequent smooth ample hypersurface cut, the ambient
dimension is at least three, so the map on geometric fundamental
groups is an isomorphism
\citep[Expos\'e~XII, Corollary~3.5, corrected p.~120,
original p.~152]{c_sga2_2005}.
The same Serre vanishings and finite-field Bertini conditions
preserve \(H^1(\mathcal O)=0\), injectivity on generic-fibre
\(H^1(\mathcal O)\), and smoothness of that generic fibre.
The final sufficiently high cut on a threefold makes the
canonical bundle very ample by adjunction and imposes all the
displayed numerical conditions.  Thus the resulting surface
\(Y/k\) satisfies every condition in the statement.

Let \(J\) be the Jacobian of \(Y_K\).  The tangent-space
injection makes the natural map \(B^\vee\to J\) have finite
kernel.  Hence \(B^\vee\), and therefore the isogenous
variety \(A\), is an isogeny factor of \(J\).  Tate for
\(Y\) implies BSD for \(J\), then BSD for \(A\), and then
Tate for \(X\), by the already established Yuan--Artin--Tate
chain.  This proves the equivalence over the unchanged field.
\end{proof}

The new intrinsic condition is a strict restriction of the
preceding test class in every characteristic, after a finite
constant extension if necessary.  Choose a prime \(r\geq5\)
different from \(p\), taking \(r=5\) unless \(p=5\), when
\(r=7\) suffices.  Over a finite field containing \(\mu_r\),
one may take \(k=\F_{p^{\operatorname{ord}_r(p)}}\).
The Fermat surface
\[
 F_r:\ x_0^r+x_1^r+x_2^r+x_3^r=0\quad\subset\mathbf P^3
\]
is smooth.  The cyclic group \(C_r\) acts diagonally with
weights \(0,1,2,3\).  Every nonidentity element has four
distinct eigenvalues, so its projective fixed locus is the four
coordinate points, none lying on \(F_r\).  The action is free;
its projective quotient \(Z=F_r/C_r\) is smooth and the
quotient map is finite \'etale of degree \(r\).
Lefschetz gives \(\pi_1(F_{r,\overline k})=1\), hence
\(\pi_1(Z_{\overline k})=C_r\).  The hypersurface sequence
gives \(H^1(F_r,\mathcal O)=0\), and the trace, with
\(r\) invertible, gives \(H^1(Z,\mathcal O)=0\).
Pad \(Z\) by \(\mathbf P^1\) and take a sufficiently high
smooth ample surface section.  Lefschetz preserves its
nontrivial group \(C_r\), while Serre vanishing and the
previous adjunction calculation give \(H^1(\mathcal O)=0\),
very ample canonical bundle and all the old numerical tails.
These surfaces belong to the old test class and are excluded by
the new geometric simple-connectedness condition.  This example
does not assert strictness over each individual smallest field;
the reduction theorem itself does hold over each fixed field.
The exact finite exponent checks are recorded in
\texttt{part\_b/code/fermat\_free\_action.py} and
\texttt{part\_b/output/fermat\_free\_action.txt}; the geometric
quotient and Lefschetz assertions are the proof above, not
inferences from those finite computations.


\begin{theorem}[The unqualified finite-field Brauer assertion is false]
\label{thm:B-Brauer-characteristic}
\status{Proved explicit counterexample; prime-to-characteristic finiteness remains separate}
Let
\[
 S=\Spec\F_2[t],\qquad
 E/\F_2:\quad Y^2Z+YZ^2=X^3,
 \qquad X=E\times_{\F_2}S,
\]
and let \(f:X\to S\) be the projection, with its origin section
\(s\).  Both schemes are regular, and \(f\) is smooth and proper
of relative dimension one.  There are explicit injections
\[
 \bigoplus_{r\geq0}\F_2\,[t^{2r+1}]
  \simeq \F_2[t]/\F_2[t^2]
  \lhook\joinrel\longrightarrow H^1(S,E_S)[2]
  \lhook\joinrel\longrightarrow
       \bigl(\Br(X)/f^*\Br(S)\bigr)[2].
\]
Consequently \(\Br(S)\to\Br(X)\) is not surjective modulo a
finite group.  This answers the literal unqualified question at
12:79 negatively, even for a constant genus-one relative curve with
a section over an affine line.  It does not refute the surrounding
manuscript's distinct \(\ell\ne p\) finiteness assertion.
\end{theorem}
\begin{proof}
The homogeneous gradient ideal of the displayed cubic is
\((X^2,Y^2,Z^2)\), so it is smooth, and its three \(\F_2\)-points
are \(O=[0:1:0]\), \([0:0:1]\), and \([0:1:1]\).
Let \(F:E\to E^{(2)}=E\) be relative Frobenius.  Its kernel is
\(\alpha_2\), with the following direct verification.  In the
chart \(Y=1\) around \(O\), put \(u=X/Y\), \(v=Z/Y\).
The equation is \(v+v^2=u^3\); imposing Frobenius kernel equations
\(u^2=v^2=0\) gives \(v=0\).  Its algebra is therefore
\(\F_2[u]/u^2\).  To identify the group law, write
\(\mathcal F(u,w)=u+w+cuw\) modulo \((u^2,w^2)\).
The invariant differential on \(E\) is \(dx\), where
\[
 v=u^3+u^6+u^{12}+\cdots,
 \qquad x=u/v=u^{-2}+u+u^4+\cdots.
\]
Thus \(dx=(1+0u+\cdots)du\).  Translation invariance says this
differential is \(du/\partial_w\mathcal F(u,0)\), whose linear
coefficient is \(c\).  Hence \(c=0\), proving the additive
Hopf law and \(\ker F=\alpha_2\).  Equivalently one may use
the Frobenius and Cartier-duality calculation in
Conrad's Examples~3.13, 3.16--3.17 and~4.4
\cite{BConrad2006}.

Every morphism \(\mathbf A^1\to E\) extends across infinity
by properness.  A nonconstant morphism \(\mathbf P^1\to E\)
is impossible: after factoring its purely inseparable part, the
separable map contradicts the genus formula.  Thus
\(E(S)=E(\F_2)\), and Frobenius is the identity on these three
points.  The fppf exact sequence
\[
 0\longrightarrow\alpha_{2,S}\longrightarrow E_S
   \xrightarrow{F}E_S\longrightarrow0
\]
therefore injects \(H^1_{\mathrm{fppf}}(S,\alpha_2)\)
into \(H^1(S,E_S)\).  The latter can be computed etale because
an \(E_S\)-torsor is smooth and consequently has etale-local
sections.  The identity \(VF=[2]\) makes the image two-torsion.
On the affine base, the sequence
\(0\to\alpha_2\to\mathbf G_a\xrightarrow{(-)^2}\mathbf G_a\to0\)
and vanishing of higher cohomology of a quasi-coherent sheaf give
\[
 H^1_{\mathrm{fppf}}(S,\alpha_2)
   =\F_2[t]/\F_2[t^2]
   =\bigoplus_{r\geq0}\F_2\,[t^{2r+1}].
\]
The basis class is the explicitly specified \(\alpha_2\)-torsor
\[
 T_r=\Spec\bigl(\F_2[t,w]/(w^2-t^{2r+1})\bigr),
 \qquad w\longmapsto w+a\quad(a^2=0).
\]
Its image \(\alpha_r\) is the contracted-product torsor
\(T_r\times^{\alpha_2}E_S\), using the kernel immersion
\(a\mapsto[a:1:0]\).

It remains to check that these classes give actual Brauer classes,
without making a prime-to-characteristic proper-base-change argument.
The section supplies a retraction of the derived unit map, so in
the derived category of etale sheaves on \(S\) one has
\[
 Rf_*\Gm\simeq\Gm\oplus C,\qquad C\in D^{\geq1},
 \qquad H^1(C)=\operatorname{Pic}_{X/S}=E_S\oplus\Z.
\]
Here the degree splitting of the relative Picard sheaf is given by
\(\mathcal O_E(O)\), and \(E=E^t\) by the origin polarization.
The truncation triangle
\[
 \operatorname{Pic}_{X/S}[-1]\longrightarrow C
       \longrightarrow\tau_{\geq2}C\longrightarrow
\]
injects \(H^1(S,\operatorname{Pic}_{X/S})\) into
\(H^2(S,C)\), since
\(H^1(S,\tau_{\geq2}C)=0\).
This injects \(H^1(S,E_S)\) into
\(H^2(X,\Gm)\), in the kernel of \(s^*\).
It is disjoint from the image of \(f^*H^2(S,\Gm)\), because
\(s^*f^*=1\).  In particular each \(\alpha_r\) yields a
Brauer class \(\beta_r\) of exact order two, and their images
are linearly independent modulo \(f^*\Br(S)\).
These are also Azumaya Brauer classes: \(X\) is a regular surface,
so Grothendieck's Corollary~2.2 identifies the two Brauer groups
\cite[Corollary~2.2]{BGrothendieckBrauerII1966}.
Equivalently \(\beta_r\) is the unique lift of \(\alpha_r\)
in the first Leray filtration whose pullback by \(s\) is zero.
This proof neither needs vanishing of \(R^2f_*\Gm\) nor passes
from separate primary finiteness to finiteness of a whole torsion group.

The smoothness, three rational points, Frobenius kernel algebra,
differential coefficient, and finite initial square-quotient matrices
are checked in
\path{part_b/code/characteristic_brauer_counterexample.py} and
\path{part_b/output/characteristic_brauer_counterexample.txt}.
The infinite independence itself follows from parity of polynomial
exponents, so it is not inferred from finitely many tests.
\end{proof}

The corrected conclusions remain explicitly scoped.  For \(\ell\ne p\),
Proposition~\ref{prop:B-relative-curve} gives the equivalence with
the corresponding Jacobian torsor finiteness, and
Theorem~\ref{thm:B-finite-chain-all-bases} proves the requested
Kummer chain on every finite-type base.  The subsequent
Theorem~\ref{thm:B-Brauer-prime-uniform} and
Corollary~\ref{cor:B-torsors-prime-uniform} strengthen the corrected
finiteness assertions to one whole prime-to-\(p\) group on every
finite-type base.  Their universal truth is equivalent to divisor
Tate for surfaces, and Theorem~\ref{thm:B-irregularity-zero-surface-gate}
reduces that condition to surfaces with \(H^1(\mathcal O)=0\).
That restricted conjecture remains unresolved.  The unqualified
assertion including the characteristic-primary part has the complete
counterexample above.

\begin{proposition}[A quantified obstruction to principalizing the normalized neighbourhood]
\label{prop:B-Kahn-principalization}
\status{Proved counterexample to the proposed automatic step; universal finiteness unchanged}
There is a fixed smooth affine surface \(U/\F_2\), with
\(\BPic(U)=0\), and smooth irreducible principal divisors
\(Z_r\subset U\), \(r\geq1\), with Nisnevich neighbourhoods
\(U_{1,r}\to U\) isomorphic over \(Z_r\), such that the
normalization \(\bar U_{1,r}\) of \(U\) in their function fields
has both of the following properties:
\begin{enumerate}
\item Restriction
\(H^1(\bar U_{1,r},\Q_5/\Z_5)\to
H^1(Z_r,\Q_5/\Z_5)\) is split surjective.
\item The selected divisor \(Z_r\) in \(\bar U_{1,r}\)
has exact Picard order \(5^r\), so is not principal.
\end{enumerate}
Moreover the geometric degree-one residue invariants of \(Z_r\)
have order \(5^r\).  Thus condition (1) of
Kahn's proposed argument in \S3.2.3 does not force condition (2),
even for this one fixed surface and prime.  The case \(r=1\)
has normalization of degree four, prime to five.
\end{proposition}
\begin{proof}
Take
\[
 E/\F_2:\quad y^2+y=x^3+x+1,
 \qquad C=E\setminus\{O\},\qquad U=C\times\Gm.
\]
The projective cubic is smooth: its homogeneous gradient ideal is
\((x^2,y^2,z^2)\).  Counting over \(\F_2\) gives
\(E(\F_2)=\{O\}\).  The divisor sequence therefore gives
\(\BPic(C)=\BPic(E)/\Z[O]=0\).  Its coordinate ring is
a Dedekind PID, and adjoining an invertible variable gives a UFD;
hence \(\BPic(U)=0\).  The smooth projective compactification
\(E\times\mathbf P^1\) has its Neron--Severi group generated
by the two factor classes, both defined over \(\F_2\).

Write \(\F_{16}=\F_2[v]/(v^4+v+1)\).  The point
\[
 P=(v^2+v,v)
\]
has order five and Frobenius orbit
\(P,\pi P,-P,-\pi P\).  The degree-four point \(z_1\)
that it defines on \(C\) is explicitly cut out by
\(x^2+x+1=0\): the curve equation then becomes
\(x=y^2+y\), and elimination gives \(y^4+y+1=0\).
The exact computation enumerates all 25 points of \(E(\F_{16})\)
and verifies that its 24 nonzero points have order five.

The Frobenius relation is \(\pi^2-2\pi+2=0\), hence
\(\pi^4=-4\).  Choose \(P_r\in E[5^r](\overline\F_2)\)
with \([5^{r-1}]P_r=P\).  It has exact order \(5^r\).
An exponent fixing it must be a multiple of four, and
\(\pi^{4k}P_r=(-4)^kP_r\).  The elementary valuation identity
\[
 v_5\bigl((1-5)^k-1\bigr)=1+v_5(k)
\]
therefore makes its orbit length
\(d_r=4\cdot5^{r-1}\).  Let \(z_r\) be this closed point of
\(C\), and set \(Z_r=z_r\times\Gm\).
It is smooth and irreducible; it is principal because
\(\BPic(C)=0\).

Set \(k_r=\F_{2^{d_r}}\) and take
\[
 \bar U_{1,r}=C_{k_r}\times\Gm,\qquad
 U_{1,r}=
 \left(E_{k_r}\setminus
 \bigl(\{O\}\cup\{\pi^aP_r:1\leq a<d_r\}\bigr)\right)
 \times\Gm.
\]
The constant extension \(\bar U_{1,r}\to U\) is finite etale.
The indicated affine open \(U_{1,r}\) is a Nisnevich neighbourhood:
the inverse image of \(Z_r\) is precisely
\(\{P_r\}\times\Gm\), and maps isomorphically to \(Z_r\).
The normal finite scheme \(\bar U_{1,r}\) has its function field,
so is the required normalization.
Projection to \(\Gm_{k_r}\) splits restriction along
\(\{P_r\}\times\Gm\), proving (1).

The selected Cartier divisor has class pulled back from
\([P_r-O]\in\BPic(C_{k_r})=E(k_r)\).
This pullback is injective, by the section \(t=1\), and the class
has exact order \(5^r\).  This proves (2).  Finally
\[
 H^1(\Gm_{\bar k_r},\Q_5/\Z_5)
       =\Q_5/\Z_5(-1),\qquad
 v_5(2^{d_r}-1)=v_5(16^{5^{r-1}}-1)=r,
\]
which computes the geometric invariant group claimed.
The finite case and modular valuation checks are reproduced in
\path{part_b/code/kahn_principalization_counterexample.py} and
\path{part_b/output/kahn_principalization_counterexample.txt}.
\end{proof}

The placement of the normalization is essential.  After shrinking
to \(U_{1,r}\), the selected class can be killed by the removed
divisors; for \(r=1\) the removed point \(-P\) already does so.
Kahn's condition (2) asks for a principal equation on the whole
\(\bar U_{1,r}\), also congruent to one on the other components.
The example refutes an automatic passage to that condition for the
chosen neighbourhood, and a uniform bound for this construction.
It does not exclude another choice of neighbourhood.  Indeed
\(E\times\mathbf P^1\) itself is of abelian type, so Tate holds
for it.  The remaining B.3 gate has not thereby moved
\cite[Corollary~3.3 and \S\S3.2.2--3.2.4]{BKahn2025}.


\subsection{B.4: the rational number-ring calculation and the integral correction}
\begin{theorem}[All rational twists on the compactified integer ring]\label{thm:B-SpecZ-Q}
\status{Proved}
In the Flach--Morin theory of \(\overline{\Spec\Z}\),
\[
H^i_{\BW}(\overline{\Spec\Z},\Q(n))=
\begin{cases}
\Q,&(n,i)=(0,0)\text{ or }(1,3),\\
\Q,&n\ge3\text{ odd},\ i=1,\\
\Q,&n\le-2\text{ even},\ i=2,\\
0,&\text{otherwise}.
\end{cases}\tag{B.5}
\]
The unit in degree zero pairs perfectly with the fundamental class in
degree three and twist one.  For every odd \(n\ge3\), the degree-one
regulator line in twist \(n\) pairs perfectly with the degree-two line
in twist \(1-n\).  The dualizing object is \(\Q(1)[3]\).
\end{theorem}
\begin{proof}
Flach--Morin's Corollary~3.12, in absolute dimension one, is
\[
\BRGamma_{\BW}(\bar X,\Q(n))\simeq
\BRGamma_{\BM}(X,\Q(n))\oplus
\BRHom(\BRGamma_{\BM}(X,\Q(1-n)),\Q[-3]).\tag{B.6}
\]
For \(X=\Spec\Z\), ordinary rational motivic cohomology is \(\Q\)
in degree zero and twist zero.  In positive twists its only possible
nonzero groups are \(H^1_{\BM}(X,\Q(n))\): twist one is
\(\Z^\times\otimes\Q=0\), while Borel's ranks give a one-dimensional
space for odd \(n\ge3\), and zero for even \(n\ge2\).
All negative-twist groups vanish.  Indeed, rational motivic cohomology
of this regular scheme is its Adams-graded rational \(K\)-theory;
the finite-field localization terms are torsion, the positive
even \(K\)-groups have rank zero, and the odd ranks are Borel's
Proposition~12.2~\cite{BBorel1974}.  Substitution in (B.6) proves
every case of (B.5).  Theorem~3.22 of Flach--Morin identifies the
pairing on (B.6) with the two cross-evaluation maps, so it is perfect
and has the claimed degree and twist~\cite[Corollary~3.12,
Theorem~3.22, and \S5.8.3]{BFlachMorin2018}.
\end{proof}

Theorem~\ref{thm:B-primitive-generator} below gives a terminating
construction of a primitive integral class in each odd positive line.
Its evaluation-dual gives the negative rational line.  These choices
are specified by finite searches; they are not canonical motivic
normalizations.

\begin{proposition}[The omitted order-two groups]\label{prop:B-integral-errors}
\status{Refuted; corrected statement proved}
The two integral tables asserted in 15:173 are false in Flach--Morin
Weil--etale cohomology.  Their correct values are
\[
H^i_{\BW}(\overline{\Spec\Z},\Z(1))=
\begin{cases}\Z/2,&i=1,\\\Z,&i=3,\\0,&\text{otherwise},\end{cases}
\qquad
H^i_{\BW}(\overline{\Spec\Z},\Z)=
\begin{cases}\Z,&i=0,\\\Z/2,&i=3,\\0,&\text{otherwise}.\end{cases}\tag{B.7}
\]
\end{proposition}
\begin{proof}
The number-ring table in Flach--Morin \S5.8.3 gives, in twist one,
the units, the class group and \(\Z\) in degrees one, two and three.
For \(\Z\) these are \(\{\pm1\}\), zero and \(\Z\).
In twist zero it gives \(\Z\) in degree zero,
\(\Hom(\mathcal O_K^\times,\Z)\oplus\BCl(\mathcal O_K)^D\)
in degree two, and \(((\mathcal O_K^\times)_{\rm tors})^D\) in degree
three, where \(D=\Hom(-,\Q/\Z)\).  Specializing to \(K=\Q\)
proves (B.7).  The order-two classes are dual under the integral derived
pairing; tensoring with \(\Q\) removes both.  This checks, and repairs,
the exact R5 step used in Theorem~\ref{thm:B-SpecZ-Q}.
\end{proof}

\paragraph{Effective completion of the integral calculation.}
The source's structural table alone leaves finite arithmetic groups
and primitive free classes unspecified.  The continuation below fills
both requirements: Theorem~\ref{thm:B-effective-odd-torsion} computes
all odd-primary invariant factors with cohomological generators,
Proposition~\ref{prop:B-two-primary} supplies every two-primary
factor, and Theorem~\ref{thm:B-primitive-generator} constructs
primitive integral free classes by terminating finite searches.
This is an effective presentation for each requested twist, with no
claim of an efficient bound or of uniform Vandiver vanishing.

\begin{proposition}[The complete structural table]
\label{prop:B-integral-structure}
\status{Proved}
For \(m\ge2\), put
\(C_m=H^2_{\et}(\Spec\Z,\Z(m))\), a finite group, and
\(a_m=1\) for odd \(m\ge3\), \(a_m=0\) for even \(m\).
Let \(w_m\) be the order of
\(H^0(\Q,\Q/\Z(m))\).  Explicitly,
\[
 w_m=\begin{cases}
 2,&m\text{ odd},\\
 2^{2+v_2(m)}\displaystyle\prod_{\substack{p\text{ odd prime}\\p-1\mid m}}
 p^{1+v_p(m)},&m\text{ even}.
 \end{cases}
\]
Define \(\epsilon_{i,n}\) to be one if
\(i\equiv n\pmod2\) and either \(1\le i\le n\) or
\(n<i<0\), and zero otherwise.  In addition to (B.7), all the
integral groups have the following structure:
\[
H^i_{\BW}(\overline{\Spec\Z},\Z(n))\simeq
\begin{cases}
\Z^{a_n}\oplus\Z/w_n,&n>1,\ i=1,\\
C_n,&n>1,\ i=2,\\
(\Z/2)^{\epsilon_{i,n}},&n>1,\ i\ge3,\\
(\Z/2)^{\epsilon_{i-1,n}},&n<0,\ i<0,\\
\Z^{a_{1-n}}\oplus C_{1-n}^{D},&n<0,\ i=2,\\
\Z/w_{1-n},&n<0,\ i=3,\\
0,&\text{all remaining cases with }n\notin\{0,1\}.
\end{cases}\tag{B.22}
\]
The direct-sum splittings and generators of the free groups in this
display are not canonical.  The arithmetic computation denoted by \(C_m\) is explicitly
supplied by the two following results.
\end{proposition}
\begin{proof}
The table in Flach--Morin \S5.8.3 expresses the answer through
\[
 H^{1,m}=H^1_{\et}(\Spec\Z,\Z(m)),
\]
the groups \(C_m\), their
integral and Pontryagin duals, and the displayed real-place terms.
Its finite generation and rank calculation give
\(H^{1,m}\simeq\Z^{a_m}\oplus(H^{1,m})_{\rm tors}\).
The rational-to-torsion coefficient sequence, together with the
positive-twist localization sequence, identifies the torsion subgroup
with \(H^0(\Q,\Q/\Z(m))\).  For each odd prime \(p\), the
cyclotomic character has image \(\Z_p^\times\).  Its \(m\)-th
power acts trivially on a cyclic subgroup of order \(p^b\) exactly
when \(p-1\mid m\) and \(b\le1+v_p(m)\).  At two, use
\(\Z_2^\times=\{\pm1\}\times(1+4\Z_2)\): an odd twist
has invariants of order two, and an even twist has invariants of
order \(2^{2+v_2(m)}\).  This proves the formula for \(w_m\).
Substitution in the source's integral table proves (B.22), including
the bounded negative-degree terms.  Rationalization recovers
Theorem~\ref{thm:B-SpecZ-Q}.
\end{proof}

\begin{proposition}[Every two-primary entry and its generator]
\label{prop:B-two-primary}
\status{Proved}
For the finite group \(C_m\) in (B.22),
\[
 C_m[2^\infty]=
 \begin{cases}\Z/2,&m\ge2\text{ even},\\0,&m\ge3\text{ odd}.
 \end{cases}
\]
For even \(m\), a generator is the inverse image, under the
canonical localization and completion isomorphism, of
\[
 \delta\bigl(\iota_m([-1])\bigr)
 \in H^2(\Spec\Z[1/2],\Z_2(m)).
\]
Here \([-1]\in H^1(\Spec\Z[1/2],\Z/2)\) is the
Kummer class of \(-1\), \(\iota_m:\Z/2\to\Q_2/\Z_2(m)\)
is the order-two inclusion, and \(\delta\) is the connecting map
for \(0\to\Z_2(m)\to\Q_2(m)\to\Q_2/\Z_2(m)\to0\).
\end{proposition}
\begin{proof}
Rognes--Weibel's Lemma~5.5 identifies
\(H^1(\Z[1/2],\Q_2/\Z_2(m))\) with its real restriction
\(\Z/2\) for nonzero even \(m\); its proof gives precisely
\(\iota_m([-1])\) as generator.  For positive even \(m\),
Borel--Soule gives \(H^1(\Z[1/2],\Q_2(m))=0\) and
\(H^2(\Z[1/2],\Q_2(m))=0\).  Hence \(\delta\) is an
isomorphism.  For odd \(m\ge3\), Proposition~6.13 of the same
source says that the finite group
\(H^2(\Z[1/2],\Z_2(m))\) has two-rank \(s+t-1\),
where \(s\) counts primes above two and \(t\) is the
class group's two-rank.  For \(\Q\), \(s=1,t=0\), so this
group vanishes~\cite[Lemma~5.5, Propositions~6.12--6.13]{BRognesWeibel2000}.
Finally, the localization contribution at two is in the positive
finite-field twist \(m-1\), of odd order \(2^{m-1}-1\).
Thus localization and etale motivic comparison identify the just
computed groups with \(C_m[2^\infty]\).  In particular these
are the unmodified motivic groups; the separate real-place terms in
(B.22) must still be retained.
\end{proof}

\subsubsection{An effective finite presentation for the odd-primary terms}
\begin{theorem}[Finite arithmetic computation of each odd-primary term]
\label{thm:B-effective-odd-torsion}
\status{Proved as a terminating effective presentation}
Fix an integer \(n\ge2\).  There is an unconditional terminating
algorithm, specified below using integer polynomial enumeration and
finite matrices, which returns the invariant factors and cohomological
generators of the odd-primary subgroup of
\(C_n=H^2_{\et}(\Spec\Z,\Z(n))\).  It requires neither Vandiver's
conjecture nor a class-group vanishing assertion.  This is an effective
presentation, not a claim that the resulting invariant factors obey
one presently proved closed vanishing pattern.
\end{theorem}
\begin{proof}
We give the complete finite procedure and its termination argument.
Write \(m=2n-2\), \(N=2m+1\), and put
\[
 A_N=2^{4N}(N+1)^{2N^2},\qquad
 w_N=N(N+1)/2-N,
\]
\[
 K_n=A_N^{(N+1)\binom{A_N}{N+w_N-m-1}},\qquad
 B_n=\max(4n-2,K_n).
\tag{B.A1}
\]
These are explicit positive integers.  Only odd primes \(\ell\le B_n\)
need be considered.  Indeed, Soul\'e's Theorem~3 gives
\[
 \prod_{\ell\ge4n-1}
 \#H^2(\Spec\Z[1/\ell],\Z_\ell(n))\le k(2n-2).
\]
In his Proposition~3 and Theorems~1--2, the bound for \(\Q\) is
\(k(m)=a(N)^{(N+1)\binom{a(N)}{N+w_N-m-1}}\).
His Proposition~2 and the calculation in \S4.4 give
\(a(N)=2^{4N}(\prod_iT_i)^{1/2}\), with
\(T_i=\gamma^N\prod_{j\ne i}B_j^2\),
\(B_j=(j+1)/2\), and \(\gamma\le1+N/4\).
Thus \(a(N)\le A_N\), and the binomial estimate is monotone in this
range, giving \(k(m)\le K_n\).  A nonzero \(\ell\)-group has order
at least \(\ell\); hence every prime greater than \(B_n\) contributes
zero.  We use only the explicit bound, not an asymptotic or a
conjectural size estimate~\cite[Proposition~3, Theorems~1--3,
and \S4.4]{Soule2003Bound}.

Fix one of these primes, an integer \(r\ge1\), and set
\(q=\ell^r\), \(D=q\varphi(q)\), \(U=\Spec\Z[1/\ell]\).
Let \(\mathcal G=\operatorname{Gal}(\Q_{\{\ell,\infty\}}/\Q)\)
be the Galois group of the maximal extension unramified at every
finite prime other than \(\ell\).  On the cyclic module
\(M_q=\Z/q(1-n)\), let
\(a_g=\chi_q(g)^{1-n}\in(\Z/q)^\times\).
An inhomogeneous continuous crossed cocycle
\(z:\mathcal G\to M_q\) defines the homomorphism
\[
 g\longmapsto(z(g),\chi_q(g)):
 \mathcal G\longrightarrow
 (\Z/q)\rtimes(\Z/q)^\times,
\tag{B.A2}
\]
where the action on the first factor is the \((1-n)\)-th power.
Its image has order at most \(D\).  Consequently every such cocycle
factors through a Galois extension of \(\Q\) of degree at most
\(D\), ramified at no finite prime except \(\ell\).

Here is a finite, explicitly bounded enumeration of all those fields.
Put
\[
 \Delta=\ell^{D(1+\lfloor\log_\ell D\rfloor)}.
\tag{B.A3}
\]
A Galois field \(E/\Q\) of degree \(d\le D\), ramified only at
\(\ell\), has \(|\operatorname{disc}E|\le\Delta\).
To verify this, at \(\ell\) write the ramification and residue
degrees as \(e,f\).  Over the maximal unramified subextension an
Eisenstein polynomial for a uniformizer has degree \(e\).  The
valuations of its nonzero derivative terms have distinct residues
modulo \(e\).  Its different exponent is therefore at most
\(e-1+e v_\ell(e)\), the valuation of the leading derivative term.
Taking the norm of the different at all primes gives
\[
 v_\ell(\operatorname{disc}E)
 \le d(1+v_\ell(e))\le D(1+\lfloor\log_\ell D\rfloor).
\]
There are no other finite discriminant factors.

For each \(2\le d\le D\), enumerate all monic integer
polynomials of degree \(d\) whose coefficients have absolute value
at most
\[
 H(d,\Delta)=2^d\bigl(d^{d+3}\Delta\bigr)^d.
\tag{B.A4}
\]
Every degree-\(d\) field with discriminant bounded by \(\Delta\)
occurs.  In fact, in its standard real Minkowski embedding the
integer lattice has covolume at most \(\sqrt\Delta\), and every
nonzero vector has Euclidean norm at least one: an algebraic integer
of nonzero norm has at least one conjugate of absolute value at
least one.  The unit ball contains the cube of side \(2/\sqrt d\).
Minkowski's second theorem therefore bounds the last successive
minimum by \(d^{d/2}\sqrt\Delta\).  Choose \(d\) linearly
independent integers \(b_i\) of norm at most
\(2d^d\sqrt\Delta\).  For the \(t=d(d-1)/2\) pairs of
distinct embeddings, each equality
\(\sigma(\sum c_i b_i)=\tau(\sum c_i b_i)\) excludes at most
\((t+1)^{d-1}\) elements of the grid
\(0\le c_i\le t\).  Since
\(t(t+1)^{d-1}<(t+1)^d\), some grid point gives a primitive
integer \(\alpha\).  Every conjugate of \(\alpha\) has absolute
value at most \(d^{d+3}\Delta\).  The elementary symmetric
formulae for its minimal polynomial give (B.A4).

For each polynomial, decide irreducibility, compute its integer ring
and discriminant, and retain it exactly when its field is Galois
and has no finite ramification outside \(\ell\).  All these are
finite exact operations.  For example, the integer ring can be
computed by enumerating the finitely many intermediate lattices
between \(\Z[\alpha]\) and its trace-dual lattice, checking
multiplication and integrality, and taking the maximal integral
order.  Polynomial factorization over number fields decides normality
and constructs composita and automorphisms.  These operations,
including their unconditional algorithmic form, are described by
Lenstra in \S\S2--4~\cite{Lenstra1992Algorithms}; no polynomial-time
claim is needed here.  Include \(\Q\) separately.

Let \(L\) be the compositum of all the retained fields and
\(\Q(\mu_q)\), and compute \(G=\operatorname{Gal}(L/\Q)\).
Choose a prime of \(L\) above \(\ell\) and let \(D_\ell\subset G\)
be its decomposition group.  It is computable by testing which
automorphisms preserve that prime ideal.  The action \(a_g\) is
read off from the action on the specified \(q\)-th root of unity.
For variables \((z_g)_{g\in G},b\) in \(\Z/q\), form the
following finite system of linear congruences:
\[
 \begin{aligned}
 z_{gh}-z_g-a_gz_h&=0 &&(g,h\in G),\\
 z_d-(a_d-1)b&=0 &&(d\in D_\ell).
 \end{aligned}
\tag{B.A5}
\]
Let \(Z_q\) be its solution module and quotient it by the cyclic
submodule generated by
\[
 \bigl((a_g-1)_{g\in G},1\bigr).
\tag{B.A6}
\]
The resulting finite abelian group is exactly
\(H^1_c(U,\Z/q(1-n))\).  Indeed, the mapping cone defining
compact support represents a degree-one cocycle as a global
crossed cocycle together with a local degree-zero trivialization;
its cocycle and boundary equations are (B.A5)--(B.A6).  At the
real place the modified Tate complex is acyclic because \(\ell\)
is odd.  Formula (B.A2) proves that every global cocycle occurs on
\(G\).  The local trivialization condition depends only on its
image \(D_\ell\), so no additional local cocycles have been
discarded.  Changing the chosen prime conjugates the presentation
and gives the canonical isomorphic compact-support group.

Finite-coefficient Artin--Verdier duality now gives
\[
 H^2(U,\Z/q(n))
 \simeq\operatorname{Hom}\bigl(Z_q/\langle\text{(B.A6)}\rangle,
                                  \Q/\Z\bigr).
\tag{B.A7}
\]
Smith normal form computes both cyclic factors and generators on
the right.  More explicitly, if \(x_j\) are the returned generators
of orders \(d_j\), the characters with value \(1/d_j\) on
\(x_j\) and zero on the other generators are the requested
dual cohomological generators.

Finally write \(C_{n,\ell}=H^2(U,\Z_\ell(n))\), which is finite
for \(n\ge2\).  The odd-primary Galois group here has
\(\ell\)-cohomological dimension two.  The finite coefficient
groups satisfy the Mittag--Leffler condition; hence continuous
cohomology in degree three is zero.  The coefficient exact sequence
therefore identifies the group in (B.A7) with
\(C_{n,\ell}/\ell^r C_{n,\ell}\).
Compute successively for \(r=1,2,\ldots\), and stop once two
successive orders agree.  If
\(C_{n,\ell}=\bigoplus_j\Z/\ell^{e_j}\), the orders have
exponents \(\sum_j\min(r,e_j)\).  Equality at \(r,r+1\)
holds exactly when every \(e_j\le r\).  Thus the stopping
criterion is sufficient and, by finiteness, is eventually reached.
At that stage (B.A7) is the whole group with its generators.

Positive-twist localization, or equivalently the prime-to-\(\ell\)
order \(\ell^{n-1}-1\) of the relevant finite-field term, identifies
\(C_n[\ell^\infty]\) with \(C_{n,\ell}\).  The finite prime
bound (B.A1) and the just-proved termination for each prime finish
the algorithm.  No enumeration in this proof is asserted to be
computationally practical, and the finite verification accompanying
it tests the cocycle presentation, not all of its astronomically
large number-field searches.
\end{proof}


\subsubsection{A terminating construction of primitive integral regulator classes}

\begin{theorem}[An integral generator specified by finite searches]
\label{thm:B-primitive-generator}
\status{Proved as an effective construction}
For each odd integer \(m\ge3\), the following procedure specifies an
higher Chow cycle \(\gamma_m\) over \(\Q\) with integer
coefficients whose canonical integral-base class generates
\[
 H^1_{\et}(\Spec\Z,\Z(m))/\mathrm{tors}
 \simeq\Z.
\]
Every search terminates.  The procedure requires the finite arithmetic
computation of \(C_m\) in Theorem~\ref{thm:B-effective-odd-torsion}
and Proposition~\ref{prop:B-two-primary}, but requires no uniform
vanishing conjecture for these groups.  Its output is a cycle given by
rational polynomial equations and integer multiplicities, together with
an explicitly found boundary certifying the saturation step.
\end{theorem}
\begin{proof}
We first specify exactly the finite tests used by our searches.  Put
\(\square=\mathbf P^1\setminus\{1\}\), and write
\(z^m(\Q,r)\) for the integral cubical higher Chow complex, modulo
degenerate cycles.  At stage \(H=1,2,\ldots\), enumerate at most \(H\) polynomial
generators per ideal, coordinate degrees at most \(H\), rational
coefficients with numerator of absolute value at most \(H\) and
denominator between one and \(H\), at most \(H\) cycle
components, and multiplicities with numerator and denominator bounded
by \(H\).  Order each finite stage lexicographically.  Every
finite rational cycle description eventually occurs.  Multihomogeneous
ideals in \((\mathbf P^1)^r\), with components contained in a
coordinate \(1\) removed, describe all the required cycles.  Ideal
factorization, dimension, saturation, and intersection multiplicities
compute their irreducible components and their restrictions to every
face \(z_i=0,\infty\).  These are finite operations over \(\Q\).
For ideal membership, dimension, primary decomposition and associated
primes, respectively, one can use Seidenberg's numbered statements
5, 6, 36 and 46; the remark following statement 46 verifies the
required effectiveness hypotheses over \(\Q\)~\cite{BSeidenberg1974}.
Intersection multiplicities are the lengths of the zero-dimensional
quotients at the computed generic points, computed by finite basis
linear algebra over their function fields.
The admissibility test is the finite set of dimension inequalities for
all faces.  Equality of two cycles is tested component by component;
degeneracy is tested by elimination in each coordinate and equality
with the inverse image of the resulting ideal.  Thus the differential
\(\partial\), the assertion \(\partial Z=0\), and any specified
boundary identity are decidable for each candidate.  The procedure
never asks whether an arbitrary cycle is a boundary without searching
for a particular bounding chain.

We also use rational coefficients during the first search.  A cycle is
called real-proper with exactly the finite family of conditions used
by the regulator complex.  For \(r=2m-1\), for each
\(J\subset\{1,\ldots,r\}\), each choice of a zero or infinity
face on \(J\), and each \(0\le k\le r\), put
\(I=\{1,\ldots,k\}\setminus J\).  Require proper real
intersection with the chosen faces and the cuts
\(z_i\in\R_{<0}\) for \(i\in I\).  These are precisely the
conditions specified in Kerr--Li, footnote~4~\cite{BKerrLi2020}.  After writing complex coordinates in real and
imaginary parts, these are finite dimension questions for semialgebraic
sets over \(\Q\).  Cylindrical algebraic decomposition decides them:
take a decomposition adapted to the finitely many polynomial equalities
and inequalities and read off the dimensions of its cells.  Collins's
Section 3, algorithms DECOMP and ELIM and Theorem 8, gives the
required finite procedure~\cite{BCollins1975}.  No
approximation to a period occurs in this test.

Here is the first search.  Set \(t_i=z_i/(z_i-1)\) and let \(Z_m\)
be the closure of the following graph in \(\square^{2m-1}\):
\[
 \left(
 \frac{t_1}{t_1-1},\ldots,\frac{t_{m-1}}{t_{m-1}-1},
 1+t_1\cdots t_{m-1},t_1^2,\ldots,t_{m-1}^2
 \right).
\]
For \(1\le i\le m-1\), let \(p_i\) project to the first
\(m-i\) coordinates, and put
\[
 f_a=1+\prod_{j=1}^a\frac{z_j}{z_j-1}.
\]
Enumerate tuples of rational admissible real-proper cycles
\((W_1,\ldots,W_{m-1})\) in \(z^m(\Q,2m-1)_\Q\)
until
\[
 \operatorname{supp}(W_i)\subset p_i^{-1}V(f_{m-i}),
 \qquad
 \partial\left(Z_m+\sum_iW_i\right)=0.
\tag{B.G1}
\]
Use the zero divisor in the open cube in this formula, so clearing
denominators does not add a component in a removed face.  These support
conditions are again ideal-membership tests.  Kerr--Yang's
Theorem~3.3 and Remark~3.4, specialized to \(N=2\), prove that a
tuple satisfying (B.G1) exists.  Their existence argument precedes and
does not use the explicit formulae in their \S4: Proposition~2.13
annihilates the residue class, and its bounding chain is decomposed
according to the faces of the relative cycle complex.  The pieces are
the \(W_i\), with exactly the support conditions (B.G1).
Consequently the first search terminates~\cite[Theorem~3.3 and Remark~3.4]{BKerrYang2018}.  Write
\(\widetilde Z_m=Z_m+\sum W_i\), let \(D_m\) be a common
denominator of its actual coefficients, and set
\[
 d_m=(m-1)!D_m,\qquad \xi_m=d_m\widetilde Z_m.
\tag{B.G2}
\]
This is an integral closed cycle.  Indeed, the cubical cycle groups
modulo degeneracies are free abelian on the nondegenerate irreducible
cycles, so a rationally zero integral boundary is already zero.

We check the regulator normalization before saturating.  In the
unscaled real coordinate of
\(H^1_D(\Spec\Q/\R,\R(m))\simeq\R\), Kerr--Yang's
Theorem~3.8 gives
\[
 \rho(\xi_m)
 =-d_m2^{m-1}\operatorname{Li}_m(-1)
 =d_m(2^{m-1}-1)\zeta(m)>0.
\tag{B.G3}
\]
One can verify the vanishing of the correction terms directly: on
their support the product of the corresponding \(t_j\)'s is
\(-1\), whereas on the simultaneous negative-real cuts all these
\(t_j\) lie in \([0,1]\).  The graph term is the iterated
integral in the proof of their Theorem~3.8.  Taking real parts removes
the ambiguity by \((2\pi i)^m\Q\), because \(m\) is odd.
We use this equation, not the inconsistent \(N=2\) sentence following
their Corollary~3.9~\cite[Theorem~3.8]{BKerrYang2018}.

Put \(c_m=|C_m|\), now an explicitly computed positive integer.
Let \(g\) be a positive primitive class in
\(H^1_{\et}(\Z,\Z(m))/\mathrm{tors}\).
Flach--Morin's regulator in \S5.8.3 uses the target lattice
\((2\pi i)^{m-1}\Z\), so its covolume is
\(R_m=\rho(g)/(2\pi)^{m-1}\).  Both maps are the Beilinson
regulator from the same motivic complex; passage from the complex
Deligne quotient to this real coordinate is the real-part projection,
with no factorial or additional factor two.  The factor
\(2/(2\pi i)^{m-1}\) in Kerr--Yang Lemma 2.8 converts their
Beilinson regulator to the Borel regulator; Flach--Morin uses the
Beilinson regulator itself, so that conversion is not made here.
Since \(w_m=2\),
their equation (88), proved for \(\Q\) by Proposition~5.35,
and the functional equation of \(\zeta\) give
\[
 |\zeta'(1-m)|
 =\frac{c_m\rho(g)}{2(2\pi)^{m-1}}
 =\frac{(m-1)!\zeta(m)}{2(2\pi)^{m-1}}.
\]
It follows that the index of \(\xi_m\) in the free lattice is
the explicitly known positive integer
\[
 h_m=\frac{d_m(2^{m-1}-1)c_m}{(m-1)!}
     =D_m(2^{m-1}-1)c_m.
\tag{B.G4}
\]
All factors and lattices here are those specified in the opened
source~\cite[\S5.8.3, equation~(88), Proposition~5.35]{BFlachMorin2018}.
The integer in (B.G4) is determined by exact arithmetic; no numerical
recognition of a regulator ratio is used.

For completeness the cycles may be searched over \(\Q\).
Positive-twist localization identifies
\[
 H^1_M(\Spec\Z,\Z(m))
 \xrightarrow{\ \sim\ }
 H^1_M(\Spec\Q,\Z(m))=CH^m(\Q,2m-1),
\tag{B.G5}
\]
because the adjacent finite-field groups are
\(H^{-1}_M(\F_p,\Z(m-1))\) and
\(H^0_M(\F_p,\Z(m-1))\), both zero for \(m-1>0\).
Geisser's Theorem~3.2 gives the higher-Chow localization sequence,
and his Theorem~1.2(2) gives the motivic-to-etale comparison through
degree \(m+1\), under the norm-residue theorem.  That theorem is
Voevodsky's Theorem~6.16, with the integral field comparison in
Theorem~6.18.  Thus the comparison invoked in Flach--Morin
\S5.8.3 is unconditional in the degree used here
\cite[Theorems~1.2(2), 3.2]{BGeisserDedekind2004}
\cite[Theorems~6.16, 6.18]{BVoevodsky2011}.
The torsion of the group in (B.G5) has order \(w_m=2\).

Finally enumerate pairs of integral admissible cycles
\[
 \gamma\in z^m(\Q,2m-1),\qquad
 \eta\in z^m(\Q,2m)
\]
until the finite tests give
\[
 \partial\gamma=0,
 \qquad \partial\eta=2(h_m\gamma-\xi_m).
\tag{B.G6}
\]
Equation (B.G4) proves existence: any integral representative of the
positive primitive free class satisfies
\(h_m[\gamma]-[\xi_m]\in\mathrm{tors}\), and twice this
difference is a boundary by the definition of higher Chow homology.
Hence this search also terminates.  Conversely (B.G6) implies that
\(h_m[\gamma]=[\xi_m]\) in the free quotient.  Since the latter
class has index exactly \(h_m\), \([\gamma]\) is primitive.
We take the first pair in the specified enumeration and define
\(\gamma_m=\gamma\).  The returned \(\eta\) certifies the
integral saturation.  This proves the claimed construction.
\end{proof}

\begin{lemma}[The degeneracy in the displayed uniform correction]
\label{lem:B-KY-degenerate}
\status{Proved source-formula defect; no general boundary formula asserted}
For every weight \(m\ge4\), the displayed \(W_2^{(i)}\)
terms in Kerr--Yang \S4D are cubically degenerate, or zero,
as rational higher Chow cycles.  Their existence theorem used above
does not depend on those displayed terms.
\end{lemma}
\begin{proof}
Every coordinate of a displayed \(W_2^{(i)}\), other than the
last coordinate \(v-1\), is homogeneous of degree zero in
\((u,v)\).  In the nontrivial ratio for \(i\le m-2\), the
numerator has degree \(m-3\) and the denominator degree
\(1+(m-4)=m-3\).  For the last term the denominator has
degree \(-1+(m-2)=m-3\).  The remaining coordinates are
ratios such as \(v t_j^N/u\), \(v/u\), or
\(u/(v\prod t_j^N)\), hence have the same homogeneity.
On the dense chart \(v\ne0\), substitute \(u=sv\).
All but the final coordinate now depend only on \(s,t_j\),
while \(v-1\) independently varies through a dense open subset
of the final cube.  The image closure is consequently the pullback
of its projection with that coordinate omitted, which is a
degenerate cycle.  If one of the other coordinates is identically
one, the cycle is already empty in the open cube.

The symbolic check for weights four, five and six, with \(N=2\),
is preserved in \texttt{source\_and\_graph\_checks.py}; the
homogeneity argument proves the assertion in every stated weight.
Kerr--Li independently records that the first weight-four formula
is not closed, because its proposed \(W_2\) correction is
degenerate~\cite[p.~548, footnote~10]{BKerrLi2020}.
We do not infer a complete uniform boundary computation from this
single defect.  The construction in
Theorem~\ref{thm:B-primitive-generator} instead uses the independent
existence theorem and tests the boundary of its actual output.
\end{proof}



\begin{proposition}[An exact Vandiver gate inside B.4]
\label{prop:B-Vandiver}
\status{Proved}
Fix an odd prime \(p\), set \(L=\Q(\zeta_p)\),
\(\Delta=\Gal(L/\Q)\), and let \(A\) be the \(p\)-primary
part of \(\BCl(\mathcal O_L)\).  Write \(A[j]\) for its
\(\omega^j\)-eigenspace, where \(\omega\) is the Teichmuller
character.  For every \(m\ge2\),
\[
 H^2_{\Bcts}(\Z[1/p],\Z_p(m))/p
 \simeq (A/p\otimes\mu_p^{\otimes(m-1)})^\Delta
 \simeq A[1-m]/p.\tag{B.23}
\]
Consequently the simultaneous assertion
\[
 C_m[p^\infty]=0\quad\text{for every odd }m\ge3
\]
is equivalent to the \(p\)-instance of Vandiver's conjecture,
\(p\nmid h(\Q(\zeta_p)^+)\).
\end{proposition}
\begin{proof}
There is exactly one prime of \(L\) over \(p\), and it is
principal.  Thus removing it does not change \(A\).  Since \(L\)
has no real places, the invariant-sum description of the Brauer group
gives \(\Br(\mathcal O_L[1/p])[p^\infty]=0\): a single
local invariant has zero sum only when it vanishes.  Kummer theory
therefore identifies
\(H^2(\mathcal O_L[1/p],\mu_p)\) with \(A/p\).
Tensoring by \(\mu_p^{\otimes(m-1)}\) and using the exactness
of \(\Delta\)-invariants, since \(|\Delta|=p-1\), gives the
middle term of (B.23).  The coefficient sequence for multiplication
by \(p\) has no obstruction in degree three: the \(p\)-cohomological
dimension of this number-ring Galois group is two.  This proves
(B.23).  The group on its left before reduction modulo \(p\) is
finite for \(m\ge2\); Nakayama's lemma therefore detects its
vanishing by (B.23).

Etale motivic comparison and localization identify its finite
\(p\)-primary group with \(C_m[p^\infty]\).  The local term at
\(p\) is zero here, because the finite-field positive-twist groups
in twist \(m-1\ge1\) have order prime to \(p\).  As odd
\(m\ge3\) runs through residue classes modulo \(p-1\),
\(1-m\) runs through every even residue class.  Their sum is the
plus part of \(A\); its vanishing is exactly the indicated
class-number condition, since the quadratic extension to the real
subfield has degree prime to \(p\).  This also rederives the precise
equivalence recorded in Kurihara's Corollary~1.5~\cite{BKurihara1992}.
\end{proof}

This reduction does not logically say that every possible algorithm
for individual \(C_m\) must first prove Vandiver.  It says precisely
that the often-proposed uniform answer ``zero in all these odd
twists'' is equivalent to Vandiver, and that an all-twist table must
settle or retain these eigenspaces.  The 2026 paper \emph{Almost all
primes are partially regular}, Theorem~1.1 and its partial Vandiver
corollary, treats a growing but restricted range of eigenspaces for
a density-one set of primes; its statement does not assert the
universal vanishing in Proposition~\ref{prop:B-Vandiver}
\cite[Theorem~1.1 and Corollary~1.4]{BChen2026}.

\begin{example}[A certified irregular-prime instance]
\label{ex:B-Vandiver37}
\status{Proved by a complete arithmetic certificate}
Let \(L=\Q(\zeta_{37}+\zeta_{37}^{-1})\), with
\(\zeta_{37}\) a primitive thirty-seventh root of unity.
The preserved PARI/GP computation constructs its degree-eighteen
real cyclotomic polynomial, returns the trivial computed class
group, and returns \(1\) from \texttt{bnfcertify(B,1)}.
The documented meaning of that flag is that the actual class
group is a quotient of the computed group; since the latter is
trivial, the actual class group is trivial
\cite{BPariClass2026}.  Consequently
\(C_m[37^\infty]=0\) for every odd \(m\ge3\), by
Proposition~\ref{prop:B-Vandiver}.

The same exact computation gives
\[
 \frac{B_{32}}{32}=-\frac{7709321041217}{16320},\qquad
 v_{37}(B_{32}/32)=1,
\]
so this is an irregular prime instance.  The executable input is
\path{part_b/code/vandiver37_quotient.gp}; the full polynomial,
certificate output, PARI version \(2.19.1\), and elapsed time
\(47\) milliseconds are in
\path{part_b/output/vandiver37_quotient.txt}.
The interrupted stronger full-unit certification is retained as
unused output and plays no role in this proof.
\end{example}
\paragraph{Localization and the zeta order.}
For a finite set \(S\) of rational primes, the rational absolute
localization triangle has support term
\[
 \bigoplus_{v\in S}\BRGamma_{\BW}(\F_v,\Q(n-1))[-2].
\]
By (B.1) this is zero if \(n\ne1\).  This proves the localization
isomorphism asserted in 15:173, with exactly the corrected exclusion
\(n\ne1\), for the specified rational absolute theory.  In twist
one, its two support groups give instead the unit and invariant-sum
sequence (B.17).  An identification with ordinary motivic cohomology
would destroy the degree-three fundamental class, so localization
must be taken within the same theory throughout.

The ranks in (B.5) do not by themselves give the order of
\(\zeta(s)\) by the naive weighted Euler characteristic: for
example the isolated line in twist three would give \(-1\), while
\(\zeta(3)\ne0\).  The needed archimedean and regulator contribution
is explicit in the compactly supported Weil--Arakelov theory.  For
\(X=\Spec\Z\), Flach--Morin's computation in \S5.8.3 gives
\[
 H^i_{{\rm ar},c}(X,\widetilde{\R}(n))=
 \begin{cases}
 \R,&n\le-2\text{ even},\ i=1,2,\\
 \R,&n=1,\ i=2,3,\\
 0,&\text{otherwise}.
 \end{cases}\tag{B.24}
\]
It follows, with order negative for a pole, that
\[
\sum_i(-1)^i i\dim_\R H^i_{{\rm ar},c}
=\operatorname{ord}_{s=n}\zeta(s)
=\begin{cases}1,&n\le-2\text{ even},\\-1,&n=1,\\0,&\text{otherwise}.
\end{cases}\tag{B.25}
\]
Here is why no unproved leading-value formula is being used: the
Borel regulator is an isomorphism after tensoring with \(\R\),
so the positive-twist regulator lines cancel against their real
targets in the defining cone.  For negative twists the dual lines
and the real-place term remain in the two degrees shown; at twist
one the degree map leaves the two fundamental real lines.  These
are the cone computations displayed on pp.~1505--1506 of the
source.  Independently, the functional equation, the simple pole at
one, and the nonvanishing Euler product for real \(s>1\) give
exactly the right-hand side of (B.25); \(\zeta(0)=-1/2\) covers
the remaining edge case.  Lead B2 is therefore correct about the
dual regulator lines and the negative even zeros, with this
archimedean completion required for the full order formula.

\subsection{B.5: finite duality and the rational statement that survives}

\begin{theorem}[Finite absolute duality, with the coefficient category specified]
\label{thm:B-finite-duality}
\status{Proved}
Let \(X\) be regular, separated, and of finite type over \(\Spec\Z\),
of pure absolute dimension \(d\).  Let \(F\) be a bounded
constructible complex of \(\Z/m\)-sheaves on \(X_{\et}\).
There is a perfect pairing of finite groups
\[
 \widehat H_c^i(X,F)\ \times\
\operatorname{Ext}^{2d+1-i}_{X,\Z/m}
       (F,\Z/m(d))\longrightarrow\Q/\Z.\tag{B.26}
\]
Here \(\Z/m(d)\) is the etale motivic complex, the Ext is in the
derived category of \(\Z/m\)-sheaves, and the hat denotes the
Tate modification at real places.  If \(m\) is invertible on \(X\),
\(\Z/m(d)\simeq\mu_m^{\otimes d}\).  Thus, for a locally
constant finite free \(F\), the second factor is the cohomology of
the usual twisted dual sheaf.  The real-place modification disappears
for odd \(m\), for schemes with no real fibre, and for schemes over
a finite field; properness over \(\Spec\Z\) alone does not remove it.
\end{theorem}
\begin{proof}
Geisser--Schmidt's Theorem~4.6, applied over the global field \(\Q\),
gives
\[
\widehat H_c^i(X,F)\times
\operatorname{Ext}^{2-i}_{X}(F,\Z_X^c(0))\longrightarrow\Q/\Z.
\]
Its hypotheses allow any bounded constructible torsion complex here;
the exclusion of characteristic-primary torsion in their positive
characteristic global-field formulation is irrelevant when the base
is \(\Spec\Z\).  For our regular pure-dimensional \(X\), the
cycle-complex indexing gives
\(\Z_X^c(0)/m=\Z/m(d)[2d]\).
Their Lemma~3.7 gives
\[
\BRHom_X(F,\Z_X^c(0))\simeq
\BRHom_{X,\Z/m}(F,\Z_X^c(0)/m)[-1].
\]
The two shifts change degree \(2-i\) to \(2d+1-i\), proving
(B.26).  Their Theorem~3.9 identifies the cycle complex with
\(\mu_m^{\otimes d}[2d]\) when \(m\) is invertible.  In the
smooth case this recovers precisely Milne's Theorem~II.7.6 and
Corollary~II.7.7, whose relative dimension \(d-1\) gives total
degree \(2(d-1)+3=2d+1\)~\cite{BGeisserSchmidt2018,AMilneADT}.
Milne's updated addendum explicitly directs the regular flat extension
to this result~\cite{BMilneErrata2025}.

The correction at infinity is part of the compact-support functor,
not of the algebraic properness assumption.  It is the cone of
restriction to Tate cohomology for \(\Gal(\C/\R)\), as in
Geisser--Schmidt's equations~(2.12)--(2.13).  Tate cohomology of that
group vanishes after inverting two, but for the trivial module
\(\Z/2\) it is \(\Z/2\) in every degree.  This proves both
the stated exceptions and the need to retain the correction in general.
\end{proof}

There is a concrete characteristic-primary check on the notation.
For a smooth complete curve \(C/\F_q\), \(q\) a power of \(p\),
one has
\(\Z/p^a(1)=\nu_a^1[-1]\), where \(\nu_a^1\) is the sheaf
of logarithmic de Rham--Witt one-forms.  Hence the dualizing complex
is \(\nu_a^1[2]\).  Milne's Theorem~II.7.11 gives the perfect
pairing
\[
 H^i(C,\Z/p^a)\times H^{2-i}(C,\nu_a^1)
 \longrightarrow\Z/p^a.
\]
It is exactly (B.26) after the coefficient shift.  Replacing
\(\Z/p^a(1)\) by the etale sheaf \(\mu_{p^a}\), which is zero
on a smooth reduced characteristic-\(p\) scheme, would instead give
a false assertion.  This also checks the edge case \(p=2\).

\begin{proposition}[Finite duality without separatedness]
\label{prop:B-finite-duality-nonsep}
\status{Proved for compact support defined by open-cover descent}
The finite pairing (B.26) extends to every regular finite-type
\(\Z\)-scheme of pure dimension \(d\), when compact support
on a possibly nonseparated scheme is defined by the following
finite open-cover construction.  This definition agrees with the
usual Tate-modified compact support when the scheme is separated.
\end{proposition}
\begin{proof}
Choose a finite affine open cover \(X=\bigcup_{j=1}^N U_j\).
Every nonempty intersection \(U_I=\bigcap_{j\in I}U_j\) is
quasicompact and separated: it is an open subscheme of any one
of its affine members.  It is therefore of finite type over
\(\Spec\Z\), and Theorem~\ref{thm:B-finite-duality} applies
to it.  Use functorial cochain models, including the real-place
cones, as in Geisser--Schmidt \S2; totalization is therefore of
an actual diagram of complexes~\cite{BGeisserSchmidt2018}.  Define
\[
 \begin{gathered}
 \widehat R\Gamma_c(X,F):=\operatorname{Tot} C_c^\bullet(\mathcal U,F),\\
 C_c^{-r}(\mathcal U,F):=
 \bigoplus_{|I|=r+1}\widehat R\Gamma_c(U_I,F),\qquad 0\le r<N.
 \end{gathered}
\]
where the notation means the finite alternating Cech total complex,
with the \(|I|=r+1\) term in Cech degree \(-r\), and the
structure maps are extension by zero for open immersions.
The real-place Tate modification is included in
each term; its maps are natural for these open immersions.

This construction is independent of the cover.  For two covers
use their intersection refinement.  On each separated intersection,
the usual compact-support Mayer--Vietoris sequence identifies the
refinement totalization with the original term.  Totalizing the
finite double Cech complex in either order proves both refinement
isomorphisms and their compatibility.  The same argument for a
separated \(X\) identifies the definition with its usual compact
support.  No compactification of a nonseparated scheme is assumed.

For the duality, put \(D_X=\Z/m(d)\).  The finite alternating
resolution of \(F\) by the sheaves \(j_{I!}(F|_{U_I})\)
is exact, as is checked on each geometric stalk by the augmented
simplex of the members containing that point.  Open restriction
of the motivic cycle complex gives \(D_X|_{U_I}=D_{U_I}\),
and adjunction gives
\[
 R\operatorname{Hom}_{X,\Z/m}
 (j_{I!}(F|_{U_I}),D_X)
 \simeq R\operatorname{Hom}_{U_I,\Z/m}(F|_{U_I},D_{U_I}).
\]
Thus the dual side is the corresponding finite Cech cochain
totalization.  The local trace pairings of (B.26) are compatible
with the open-immersion maps, so their duality isomorphisms
identify these two totalizations with the shift \(2d+1\).
All totalizations here use finitely many terms.  Their groups in
each degree are finite, and \(\operatorname{Hom}(-,\Q/\Z)\)
is exact.  Taking cohomology therefore gives the perfect pairing
(B.26) on \(X\), with precisely the declared compact-support
convention.
\end{proof}

\begin{theorem}[Rational duality for the Tate cases proved here]
\label{thm:B-rational-duality}
\status{Proved}
For \(X=\Spec\mathcal O_K\), the Flach--Morin compactified
complexes have perfect rational pairings
\[
 H^i_{\BW}(\bar X,\Q(n))\times
 H^{3-i}_{\BW}(\bar X,\Q(1-n))\longrightarrow\Q.
 \tag{B.27}
\]
For \(U=\Spec\mathcal O_K[1/S]\), the same assertion pairs
the localized rational absolute complex with the compact-support
complex, with the archimedean completion inherited from \(\bar X\).
For a smooth projective geometrically connected curve \(C/\F_q\),
the corresponding Weil--etale pairing is
\[
 H^i_{\BW}(C,\Q(n))\times H^{3-i}_{\BW}(C,\Q(1-n))
 \longrightarrow\Q.
 \tag{B.28}
\]
Both statements hold for every integral twist \(n\), including
negative twists, and for finite direct sums of these Tate objects.
\end{theorem}
\begin{proof}
For the proper number ring, all the finite-generation and duality
conditions used in Flach--Morin's construction hold in dimension
one, as verified in their \S5.8.3.  Their Theorem~3.22 therefore
applies, and Corollary~3.12 identifies its rational pairing with
cross-evaluation between ordinary motivic cohomology and the shifted
dual of the complementary twist.  This proves (B.27) and explains
why ordinary motivic cohomology alone lacks its degree-three partner.

For the open scheme put
\(D_S(n)=\bigoplus_{v\in S}\BRGamma_{\BW}(\F_v,\Q(n))\).
Compact support is the fibre of restriction
\(\BRGamma_{\BW}(\bar X,\Q(n))\to D_S(n)\).
Localization is the cone of the Gysin map
\(D_S(n-1)[-2]\to\BRGamma_{\BW}(\bar X,\Q(n))\).
The point calculation (B.1) gives a perfect trace pairing between
\(D_S(n)\) and \(D_S(-n)\) of total degree one.  Local
restriction and the Gysin map are adjoint under that pairing and
the global trace: this follows already for finite coefficients from
(B.26), and on the rational point summands is the equality between
degree and the invariant-sum map.  Dualizing the compact-support
triangle in complementary twist therefore gives exactly the
localization triangle.  Since evaluation is perfect on its two
known terms, the induced cone pairing is perfect as well.  This is
the asserted open-scheme duality, and specifies which compact
support is meant.

For the projective curve, Lichtenbaum's Theorem~5.4 gives
\[
 \BRGamma_{\BW}(C,\Gm)\simeq
 \BRHom(\BRGamma_{\BW}(C,\Z),\Z[-2]).
\]
Since \(\Z(1)=\Gm[-1]\), this is the integral source of
the degree-three pairing in (B.28) for twists zero and one.
Their rational groups are \(\Q\) in degrees \(0,1\) at twist
zero and \(\Q\) in degrees \(2,3\) at twist one, by
Lichtenbaum's Theorem~6.1.  For every other twist, the unconditional
curve comparison in Geisser's Theorem~9.5 and the Frobenius weights
make both sides zero.  This proves every case of (B.28).
\end{proof}

\begin{corollary}[Finite-monodromy coefficient motives]
\label{cor:B-Artin-Tate-duality}
\status{Proved}
Let \(X\) be either an open number ring or a smooth projective
curve over a finite field.  Let \(W\) be a finite-dimensional
rational representation of a finite quotient of
\(\pi_1^{\et}(X)\).  Then the rational absolute or Weil--etale
theory just specified has a common rational structure for all
\(W(n)\), every integer \(n\), and perfect pairings
\[
 H^i_{\rm abs}(X,W(n))\ \times\
 H^{3-i}_{\rm abs,c}(X,W^\vee(1-n))\longrightarrow\Q.
\]
Compact support includes the stated archimedean completion on a
number ring; it is unnecessary for a proper finite-field curve.
The degree-one realization is the Bloch--Kato finite-condition
group with the same geometric conditions at the omitted places.
\end{corollary}
\begin{proof}
Choose a finite etale Galois cover \(Y\to X\), with group \(G\),
trivializing \(W\).  Apply the preceding Tate calculations to
each connected component of \(Y\), over its number field or full
constant finite field.  Descent defines the rational complex by
\[
 R\Gamma_{\rm abs}(X,W(n))=
  \bigl(R\Gamma_{\rm abs}(Y,\Q(n))\otimes_\Q W\bigr)^G,
\]
and similarly with compact support.  The invariants functor is
exact, with projector \(|G|^{-1}\sum_{g\in G}g\), so it commutes
with scalar extension to every \(\Q_\ell\).  Restriction and
corestriction preserve the unramified, crystalline and de Rham
local kernels defining \(f\) and \(g\); their composite is
the extension degree.  At a retained place above \(\ell\),
the finite etale cover is unramified, so crystalline descent
applies.  At a removed place above \(\ell\), the condition is
\(g\), for which de Rham descent allows a ramified extension.
Thus the degree-one descent is exactly
the claimed Bloch--Kato group.  Higher comparison splittings
can be averaged over \(G\), giving the same graded assertion.

Combine the perfect pairing on \(Y\) with evaluation
\(W\otimes W^\vee\to\Q\) and divide the trace by \(|G|\).
Its restriction to invariant subspaces is perfect: for an invariant
nonzero vector, average any vector pairing nontrivially with it.
This proves the assertion and its compatibility with localization.
\end{proof}

\begin{proposition}[The finite-field coefficient-duality criterion]
\label{prop:B-finite-field-coefficient-gate}
\status{Proved equivalence for all Chow-heart coefficients; the universal hypothesis remains conjectural.}
Fix a finite field \(k\), and let \(\mathcal H\) be the rational Chow
heart in \(\mathcal D=DM_{\rm gm}(k,\Q)\).  For a smooth projective
\(X/k\) of dimension \(d\), use the conventions
\[
 h_{\rm coh}(X)=M(X)^\vee=M(X)(-d)[-2d],\qquad
 P_{X,n}=h_{\rm coh}(X)(n)[2n].
\]
Thus \(\mathcal H\) is generated by the \(P_{X,n}\) under finite sums
and direct summands, and is closed under duality.  Use the prescribed
rational Weil extension of ordinary geometric motives, normalized by
a generator \(e\) of degree one at \(\Spec k\).  The following are
equivalent.
\begin{enumerate}
\item For every \(P\in\mathcal H\) and every integer \(i\), the
evaluation--trace pairing
\[
 W^i(P)\times W^{1-i}(P^\vee)\longrightarrow\Q
\]
is perfect between finite-dimensional spaces.
\item Beilinson's rational--numerical equality and Parshin's rational
vanishing hold for every smooth projective variety over \(k\).
\item For every \(P\in\mathcal H\), the prescribed constant coefficient
on \(\mathbf P^1_k\) has perfect finite-dimensional pairings
\[
 W^i(\mathbf P^1,\pi^*P)\times
 W^{3-i}(\mathbf P^1,\pi^*(P^\vee(1)))\longrightarrow\Q
 \quad\text{for every }i,
\]
with the usual projective-bundle formula, cup product and Gysin trace.
\end{enumerate}
Here ``Chow heart'' does not mean a single homogeneous cohomological
weight.  To infer the universal criterion from an assertion restricted
to homogeneous pure Chow summands, one must justify extension to the
whole Chow heart; an actual finite Chow--K\"unneth splitting supplies
such a justification.
\end{proposition}
\begin{proof}
Put \(A^i(P)=\Hom_{\mathcal D}(\mathbf1,P[i])\).  Our conventions give
\[
 P_{X,n}^\vee=P_{X,d-n},\qquad
 A^i(P_{X,n})=H_M^{2n+i}(X,\Q(n))
             =\operatorname{CH}^n(X,-i)_\Q.
\]
Since the higher-Chow complex has no negative simplicial degrees,
\(A^i(P)=0\) for \(i>0\), first on these generators and then on every
\(P\in\mathcal H\).  The same applies to \(P^\vee\).

The rational Weil splitting, together with its multiplicative
projection formula, gives a model
\[
 C_W(P)=E\otimes_\Q R\Hom_{\mathcal D}(\mathbf1,P),\qquad
 E=\Q\oplus\Q e,\quad |e|=1,\quad e^2=0,
\]
and hence
\[
 W^i(P)=A^i(P)\oplus e A^{i-1}(P).
\]
These are the rational coefficient formulas in Kahn's
Remark~10.5.2(2), Theorem~11.2.1 and Corollary~11.2.2(c)
\cite[pp.~47--49]{BKahn2012FullFaithfulness}.  For clarity, place
\(e\) on the left.  If \(a+eb\) and \(c+ef\) have total degrees
\(i\) and \(1-i\), the trace is exactly
\[
 b_P^i(a+eb,c+ef)=(-1)^i\operatorname{ev}(a,f)
                         +\operatorname{ev}(b,c).
\]
The sign comes from moving the degree-\(i\) element \(a\) past \(e\).
Both evaluations on the right have ordinary degree zero.

Suppose (1) holds.  For \(i<0\), negativity gives
\(W^{1-i}(P^\vee)=0\).  Perfection forces \(W^i(P)=0\), and therefore
its direct summand \(A^i(P)\) vanishes.  Together with negativity,
\(A^i(P)=0\) for every \(i\ne0\).  Taking \(P=P_{X,n}\) is precisely
Parshin's vanishing in every twist.  In degrees zero and one the
pairing is now ordinary evaluation
\[
 \operatorname{CH}^n(X)_\Q\times
 \operatorname{CH}^{d-n}(X)_\Q
 \longrightarrow\Q,\qquad
 (\alpha,\beta)\longmapsto\deg(\alpha\beta).
\]
Its kernel on either side is the numerically trivial subgroup, so
perfection gives Beilinson's equality.

Conversely assume (2).  Parshin and negativity make \(A^i\) vanish
outside degree zero on every generator.  Beilinson makes the displayed
intersection pairing nondegenerate.  Numerical cycle groups are
finite dimensional: any independent finite family of numerical classes
admits a full-rank rational intersection matrix, whose rank is bounded
by the corresponding finite Weil-cohomology dimension.  The pairing
is consequently perfect.  This persists under finite sums and direct
summands, since the dual idempotent is adjoint to the original
idempotent under evaluation.  Thus only
\(W^0(P)=A^0(P)\) and \(W^1(P)=eA^0(P)\) remain, and the two cross
pairings are perfect.  This proves (1).  It is the Chow-heart version
of Kahn's Proposition~10.5.1 and Remark~10.5.2(2), whose argument
also uses exactly the negative higher-Chow degrees and numerical
intersection pairing \cite[pp.~46--47]{BKahn2012FullFaithfulness}.

For the curve statement put \(h=c_1(\mathcal O(1))\), with degree two
and twist one.  The projective-bundle formula reads
\[
\begin{aligned}
 W^i(\mathbf P^1,\pi^*P)
   &=W^i(P)\oplus h W^{i-2}(P(-1)),\\
 W^{3-i}(\mathbf P^1,\pi^*(P^\vee(1)))
   &=W^{3-i}(P^\vee(1))\oplus hW^{1-i}(P^\vee).
\end{aligned}
\]
Since \(\pi_*h=1\) and \(h^2=0\), the trace pairing is block
off-diagonal.  On the two pairs of components it is
\[
 ((x,u),(y,v))\longmapsto
 b_P^i(x,v)+b_{P(-1)}^{i-2}(u,y).
\]
There is no extra sign from \(h\), whose degree is even.  The
unit--unit component is killed by the Gysin projection, and the
\(h\)--\(h\) component is zero.  Thus (3) implies (1) on the first
cross-block, including finite dimensionality.  Conversely
\(P'=P(-1)[-2]\) belongs to \(\mathcal H\), with
\[
 W^i(P')=W^{i-2}(P(-1)),\qquad
 W^{1-i}((P')^\vee)=W^{3-i}(P^\vee(1)).
\]
Condition (1) for \(P\) and \(P'\) proves perfection of both blocks,
and hence (3).

Finally, the argument uses the prescribed ordinary-motive Weil
extension.  It does not replace those motivic groups by arbitrary
abstract dual completions.  Nor does it assume that the image of
ordinary motives is itself a triangulated subcategory of the larger
Weil category; Kahn's Remark~11.3.2 explicitly explains this failure.
A cohomological K\"unneth projector does not by itself provide an
actual Chow--K\"unneth splitting.  This proves the stated limitation
for a formulation involving only homogeneous pure coefficients.
\end{proof}

Under B1's standing universal Tate and semisimplicity hypotheses, the
verified Geisser--Kahn implication supplies Parshin once universal
Beilinson holds.  In that setting the preceding criterion retains the
same Chow-phantom gate.  Without those standing hypotheses it is the
joint Beilinson--Parshin condition.  This finite-field dependency is
separate from the number-ring f-admissibility condition below; neither
follows from finite-coefficient duality or the proved Artin--Tate cases.


\paragraph{The scope still to be proved.}
For a general pure motive \(M\), a rational analogue of (B.26)
must specify an absolute or Weil--etale motivic category, finite
local conditions, a trace, and a compact-support operation including
infinity.  The proved finite theorem does not create those rational
objects: tensoring a finite torsion pairing with \(\Q\) gives zero.
For a dualizable coefficient motive on a regular base of absolute
dimension \(d\), the corrected derived assertion is
\[
 R\Gamma_{\rm abs}(X,M)\simeq
 R\operatorname{Hom}_{\Q}
 \bigl(R\Gamma_{\rm abs,c}(X,M^\vee(d)),\Q\bigr)[-2d-1],
\]
with the compactified arithmetic convention just proved in dimension
one.  A compatible rational trace and realization into a perfect
local/global duality imply that this pairing is perfect, since
extension of scalars from \(\Q\) to any \(\Q_\ell\) is
faithfully flat.  The finite theorem alone supplies neither the
rational trace nor that realization.
Even on a number ring the equality of the expected degree-one
cohomology of \(V_\ell A\) with \(A(K)\otimes\Q_\ell\)
contains the Sha-finiteness assertion in (B.21).  Thus the arbitrary
pure-motive portion of B.5 is not proved by the all-twist Tate
calculation, and remains explicitly recorded as a separate task.

\begin{proposition}[The exact number-ring rational gate for a dual pair]
\label{prop:B-general-rational-gate}
\status{Proved equivalence for dual completions; the realization hypothesis is conjectural}
Fix a pure coefficient motive \(M/K\), and write
\(M^*=M^\vee(1)\).  Suppose rational motivic spaces
\(A^a(M)\) and \(A^a(M^*)\), for \(a=0,1\), are specified
with their realization maps.  They admit a finite-dimensional
rational dual completion having the graded Selmer comparisons
in every degree if and only if the degree-zero and degree-one
maps for both motives satisfy Fontaine's f-admissibility
conditions \(c^0_\ell,c^1_\ell\) for every \(\ell\).
For the prescribed motivic spaces, this is a genuine realization
conjecture, not the assertion that some abstract vector spaces of
the same dimensions can be chosen.
\end{proposition}
\begin{proof}
Under those conditions put
\[
\begin{aligned}
 A^2(M)&=A^1(M^*)^\vee,&
 A^3(M)&=A^0(M^*)^\vee,\\
 R A(M)&=\bigoplus_{i=0}^3 A^i(M)[-i].
\end{aligned}
\]
Evaluation, using \((M^*)^*=M\), gives the perfect pairing
\(RA(M)\otimes RA(M^*)\to\Q[-3]\).
The degree-zero and degree-one realizations are the assumed
maps.  The Poitou--Tate calculation (B.16), with the splittings
already specified, gives the degree-two and degree-three
comparisons.  Hence this is a rational dual completion with all
the required graded comparisons on the proper number ring.
Conversely any such completion, with the prescribed degree-zero
and degree-one maps, implies exactly those f-admissibility
conditions upon taking those two cohomology groups.  No existence
of a six-operation motivic category or canonical product on all
these completions has been deduced.  The proposition identifies
the exact remaining comparison gate for this specified graded
construction, rather than concealing it in higher-degree notation.
\end{proof}
}{}
\IfFileExists{sections/part_c.tex}{\section{Correspondences, tautological classes, and extension problems}
\label{sec:part-c}
\providecommand{\Lie}{\operatorname{Lie}}
\providecommand{\tr}{\operatorname{tr}}
\providecommand{\Alt}{\operatorname{Alt}}
\providecommand{\sgn}{\operatorname{sgn}}
\providecommand{\Gr}{\operatorname{Gr}}
\providecommand{\id}{\operatorname{id}}

\subsection{C.1: a finite presentation for every cohomological tautological ring}

We use a Weil cohomology with coefficient field $L$ of characteristic zero, the usual Koszul commutativity constraint, and cohomological grading. Let $X$ be smooth, projective, geometrically connected, and of dimension $n>0$ over an algebraically closed field. Put $D=2n$, $A=H^*(X)$, and $\kappa(a)=\int_Xa$. We choose the top class $y$ with $\kappa(y)=1$. All the generators under consideration have degree divisible by $D$; we write $T_k[d]=T_k\cap H^{Dd}(X^k)$. For $n=0$, $X$ is a point and $T_k=L$ in degree zero, which we henceforth set aside.

\begin{definition}[Pointed partitions]
\label{c:def-pointed}
\status{Definition, with degrees verified}
Let $\mathcal P_k$ be the set of partitions of $\{0,1,\ldots,k\}$. The block containing $0$ is the distinguished block. For $P\in\mathcal P_k$, put $J=P(0)\setminus\{0\}$ and define
\[
 \varepsilon_P=\prod_{i\in J}y_i\prod_{B\in P,\,0\notin B}\delta_B,
 \qquad d(P)=k-|P|+1.
\]
Here $\delta_B$ is the pullback of the small diagonal in $X^B$ and $\delta_{\{i\}}=1$. Thus $\varepsilon_P\in T_k[d(P)]$. This describes the pointed and unpointed partial diagonals of manuscript 11:55 without duplicate notation.
\end{definition}

\begin{theorem}[C.1(b): the complete graph product rule]
\label{c:graph-rule}
\status{Proved}
Choose a tree on each block of $P$ and of $Q$, using the edges $\{0,i\}$ in a distinguished block, and take the union as a multigraph on $\{0,\ldots,k\}$. Write $b(C)=|E(C)|-|V(C)|+1$ for the cycle rank of a connected component. A component not containing $0$ contributes
\[
 \begin{cases}
 \delta_C,&b(C)=0,\\
 \chi(X)\prod_{i\in C}y_i,&b(C)=1,\\
 0,&b(C)\geq2,
 \end{cases}
\]
whereas the component containing $0$ contributes $\prod_{i\in C\setminus\{0\}}y_i$ if $b(C)=0$, and zero otherwise. Multiplying these contributions gives $\varepsilon_P\varepsilon_Q$, independently of the chosen trees. In particular, the $\varepsilon_P$ span $T_k$.
\end{theorem}

\begin{proof}
A tree of pairwise equalities imposes a small diagonal, since its edges impose independent diagonal conditions. In the algebra $A^{\otimes k}$ this also follows by composing the coevaluation for the perfect pairing $\kappa(ab)$. After the tree conditions have been imposed, each further edge in the same component is the self-intersection of the diagonal. Its pullback is the Euler element
\[
 e_A=\mu(\Delta_A)=\chi(X)y,
\]
where $\Delta_A\in A\otimes A$ is the coevaluation and $\mu$ is multiplication. To verify the last equality directly, choose homogeneous dual bases; the coefficient of $y$ is the supertrace of the identity, namely $\sum_r(-1)^r\dim H^r(X)$. Since $n>0$, $y^2=0$. Consequently a second excess edge gives zero. If the component already contains a point condition, an excess edge gives $y e_A=0$; the tree itself makes every vertex in that component equal to the point. The coevaluation identities show that changing the tree leaves the same diagonal class. Every product of generators is therefore either zero or a scalar times a pointed-partition class. These arguments establish the previously unchecked products of 11:55--57, including the cycle and point cases.
\end{proof}

We next give an actual finite system of generators for all relations, valid even when the restricted intersection pairing on $T_k$ is degenerate. Using only that restricted pairing would not suffice.

\begin{theorem}[C.1(a): all relations from the multilinear forms]
\label{c:matrix-presentation}
\status{Proved}
Choose a homogeneous basis $(e_a)$ of $A$, containing $e_0=1$. For fixed $d$, form a matrix $C_{k,d}$ whose columns are the partitions $P$ with $d(P)=d$ and whose rows are the basis tensors $e_{a_1}\otimes\cdots\otimes e_{a_k}$ of total degree $D(k-d)$. Its entry is
\begin{equation}
\label{c:contraction-matrix}
 (C_{k,d})_{\mathbf a,P}
 =\epsilon(\mathbf a,P)
 \prod_{B\in P,\,0\notin B}\kappa\Bigl(\prod_{i\in B}e_{a_i}\Bigr)
 \prod_{i\in P(0)\setminus\{0\}}[a_i=0].
\end{equation}
In each block the indices are in increasing order; $\epsilon$ is the Koszul sign of collecting the factors into those blocks. Then
\[
 \dim_L T_k[d]=\operatorname{rank}_L C_{k,d},\qquad
 \sum_Pc_P\varepsilon_P=0\ \Longleftrightarrow\ (c_P)_P\in\ker C_{k,d}.
\]
A finite presentation of $T_k$ consists of one variable $z_P$ for each $P\in\mathcal P_k$, the unit relation, all product relations of Theorem~\ref{c:graph-rule}, and one linear relation for each vector in a basis of each $\ker C_{k,d}$. Thus Gaussian elimination gives generators for the entire ideal of relations. The input is exactly the graded vector space, its unit and trace, and the multilinear forms $\kappa(u_1\cdots u_m)$ on positive-degree elements.
\end{theorem}

\begin{proof}
The projection formula for each diagonal gives
\[
 \int_{X^k}\varepsilon_P(e_{a_1}\otimes\cdots\otimes e_{a_k})
 =(C_{k,d})_{\mathbf a,P}.
\]
For a pinned coordinate, $\kappa(ye_a)=[a=0]$; for an unpinned block it is the integral of the product of its inputs. These are all factors, with precisely the displayed Koszul permutation sign. The full Poincar\'e pairing on $A^{\otimes k}$ is perfect, so its contraction with all rows detects the zero vector. This proves both the kernel and rank assertions. The graph relations reduce any polynomial to a linear combination of $z_P$; the matrix kernel is exactly the kernel remaining after this reduction. No nondegeneracy of the restricted pairing on the span of the $\varepsilon_P$ is used. Finally, multilinear forms involving a unit are obtained by deleting that unit, so the stated positive-degree data suffice.
\end{proof}

\begin{proposition}[The modified diagonal and the correction to R8]
\label{c:modified-diagonal}
\status{Proved; proposed generation by this relation alone refuted}
The class
\[
 M=\delta_{123}-\delta_{12}y_3-\delta_{23}y_1-\delta_{13}y_2
       +y_1y_2+y_2y_3+y_1y_3
 \in H^{4n}(X^3)
\]
vanishes if and only if $\kappa(uvw)=0$ for every three positive-degree classes $u,v,w$. The degree $H^{2n}$ stated in R8 is incorrect when $n$ denotes the geometric dimension. Neither $M$ nor its higher-arity analogues, without further dimension relations, generate all relations in every $T_k$.
\end{proposition}

\begin{proof}
Pair $M$ with $u\otimes v\otimes w$. If all three inputs have positive degree, all terms except the small diagonal vanish and the pairing is $\kappa(uvw)$. If at least one input is $1$, the terms cancel by inclusion and exclusion; multilinearity reduces every input to these cases. Poincar\'e duality proves the equivalence. For the generation assertion take $X=\mathbb P^1$ and $k=2$. There is the linear relation $\delta_{12}=y_1+y_2$; relations with at least three distinct indices are absent in this ring, while the universal self-intersection and point relations have polynomial degree at least two. Their homogeneous ideal cannot contain this nonzero linear polynomial. More generally, the dimension relations in Theorem~\ref{c:pair-ranks} below are additional indispensable relations. We distinguish this correction from the valid triple-product criterion in R8.
\end{proof}

\begin{theorem}[Dimensions for curves and for surfaces with $b_1=0$]
\label{c:pair-ranks}
\status{Proved}
For a curve of genus $g$, put $b=-2g$; for a smooth projective surface with $b_1=0$, put $b=b_2$. Define $q_{ij}=\delta_{ij}-y_i-y_j$. All relations are generated by
\begin{align}
 y_i^2&=0,&y_iq_{ij}&=0,&q_{ij}^2&=b\,y_iy_j,&q_{ij}q_{ik}&=y_iq_{jk},\label{c:pair-local}
\end{align}
and the following dimension relations, with all displayed indices distinct. For a surface they are
\begin{equation}
\label{c:orth-det}
 \sum_{\sigma\in S_{b+1}}\sgn(\sigma)
   \prod_{a=1}^{b+1}q_{i_a,j_{\sigma(a)}}=0;
\end{equation}
for a curve they are
\begin{equation}
\label{c:sympl-pf}
 \sum_{\text{perfect matchings }\mathcal M\text{ on }2g+2\text{ slots}}
       \prod_{\{i,j\}\in\mathcal M}q_{ij}=0.
\end{equation}
There is no sign in \eqref{c:sympl-pf}, because the symplectic signs are already included by the odd degree of $H^1$ and the Koszul constraint.

For a partition $\lambda$, write $f^\lambda$ for the number of standard Young tableaux of shape $\lambda$, explicitly $|\lambda|!$ divided by the product of the hook lengths. Put
\[
 r_b(m)=
 \begin{cases}
 \displaystyle\sum_{\lambda\vdash m,\,\ell(\lambda)\leq b} f^{2\lambda},&b\geq0\quad\text{(surface)},\\[4pt]
 \displaystyle\sum_{\lambda\vdash m,\,\lambda_1\leq g} f^{2\lambda},&b=-2g\quad\text{(curve)},
 \end{cases}
 \qquad r_b(0)=1.
\]
Then the answer in every degree is
\begin{equation}
\label{c:dimension-pair}
 \dim T_k[d]=
 \sum_{\substack{m\geq0,\ a=d-m\geq0\\a+2m\leq k}}
 \frac{k!}{a!(2m)!(k-a-2m)!}\,r_b(m).
\end{equation}
The curve formula includes $g=0$, when $r_0(m)=0$ for $m>0$ and the answer is $\binom{k}{d}$.
\end{theorem}

\begin{proof}
In both cases $A=L1\oplus W\oplus Ly$, with $W=H^1$ for a curve and $W=H^2$ for a surface. The product on $W$ is a perfect alternating, respectively symmetric, pairing with values in $Ly$. Direct contraction gives \eqref{c:pair-local}; the superdimension of $W$ gives $b$. In particular the positive triple products vanish. Repeated use of those relations reduces a monomial to $y_I\prod_{\{i,j\}\in\mathcal M}q_{ij}$, with disjoint indices. Factorwise cohomological degrees separate different point sets $I$ and different sets of paired vertices. For a fixed set of $2m$ paired vertices, the remaining problem is exactly the span of the pairing tensors in $W^{\otimes2m}$.

For completeness, the matching representation of $S_{2m}$ decomposes without multiplicity into the Specht modules indexed by the even-row shapes $2\lambda$. Its contraction matrix has nonzero eigenvalue exactly for $\ell(\lambda)\leq b$ in the orthogonal case and $\lambda_1\leq g$ in the symplectic case. The orthogonal calculation and the argument that every kernel Specht module is generated by the first determinant are stated and proved in \citet[Proposition 3.5 and \S3.6, pp.~508--510]{c_yin2015}; the symplectic calculation, including the all-plus convention and generation by \eqref{c:sympl-pf}, is worked out in \citet[\S6, equation (2), and the argument preceding Corollary 6.4]{c_tavakol2018}. These calculations concern a matrix of pairings and therefore apply to the present $W$ independently of whether it comes from a K3 surface or a hyperelliptic curve. To see the generation step, a kernel Young diagram has either a column of length $b+1$ or a row of length $2g+2$; antisymmetrizing that column, respectively symmetrizing that row with the Koszul signs, supplies \eqref{c:orth-det}, respectively \eqref{c:sympl-pf}. Decomposing the larger stabilizer into cosets reduces every larger kernel generator to its pullbacks and products. The hook formula gives $r_b(m)$. Finally, choose the $a$ point slots and the $2m$ pairing slots; the number of choices is the multinomial coefficient in \eqref{c:dimension-pair}. Summing gives the formula and proves completeness of the relations.
\end{proof}

\begin{proposition}[Explicit dimension algorithms for projective spaces and abelian varieties]
\label{c:other-dimensions}
\status{Proved}
The following finite integer matrices determine every dimension by ordinary exact rational rank; in particular they give a uniform computable answer for every $n,g,k,d$, with no geometric unknown in their entries.

For $X=\mathbb P^n$, the columns are $P\in\mathcal P_k$ with $d(P)=d$, and the rows are $(a_1,\ldots,a_k)\in\{0,\ldots,n\}^k$ with $\sum a_i=nd$. The entry is $1$ precisely when
\[
 a_i=n\quad(i\in P(0)\setminus\{0\}),\qquad
 \sum_{i\in B}a_i=n(|B|-1)\quad(0\notin B),
\]
and is zero otherwise. Its rank is $\dim T_k[d]$.

For an abelian variety of dimension $g>0$, use the exterior algebra on the $2gk$ symbols $e_{i,a}$, ordered lexicographically. In this algebra put
\[
 y_i=\prod_{a=1}^{2g}e_{i,a},\qquad
 \delta_{ij}=\prod_{a=1}^{2g}(e_{i,a}-e_{j,a}).
\]
For each $P$ multiply these expressions along any of its trees and its pinned vertices. The coefficient matrix in the exterior-monomial basis of degree $2gd$ has entries in $\Z$ and rank $\dim T_k[d]$. These ranks depend only on $g,k,d$.
\end{proposition}

\begin{proof}
For projective space, $H^*(X)=L[h]/(h^{n+1})$, $\int h^n=1$, and the small diagonal on a block $B$ is the sum of the monomials satisfying the displayed condition, each with coefficient one. This follows either by pairing with monomials or by induction along a tree. For an abelian variety, $H^*(X)=\bigwedge H^1(X)$. Choose a volume-normalized basis; the diagonal formula follows from the difference map $(x,y)\mapsto x-y$, whose pullback on $H^1$ is $e\mapsto e\otimes1-1\otimes e$, and whose inverse image of the origin is the diagonal. Expansion in the basis is therefore an explicit expression for every $\varepsilon_P$. In both cases the coefficient matrix has precisely these classes as columns, so its rank is their dimension. These matrices are distinct from an unevaluated Chow-theoretic condition: their entries require only finite integer arithmetic.
\end{proof}

The exact-arithmetic program \texttt{code/part\_c/check\_tautological.py} implements the Frobenius algebra construction independently of the display formulas. It verified 2,286 graph products and the curve/surface rank formula in the finite ranges recorded in its JSON output. For example, the dimensions for $k=4$, in normalized degrees $d=0,\ldots,4$, are
\[
\begin{array}{c|ccccc}
 X&0&1&2&3&4\\\hline
 \mathbb P^1&1&4&6&4&1\\
 \mathbb P^2&1&10&19&10&1\\
 \mathbb P^3&1&10&25&14&1\\
 \text{curve }g=1&1&10&20&10&1\\
 \text{curve }g=2&1&10&21&10&1
\end{array}
\]
The non-palindromic row for $\mathbb P^3$ is one reason not to replace the full pairing in Theorem~\ref{c:matrix-presentation} by an assumed perfect tautological pairing.

\subsection{C.1(c): Chow lifting, and the exact K3 obstruction}

Write $T_k^{\CH}(X,o)$ for the rational Chow subring generated by the same diagonals and a degree-one zero-cycle $o$. The requested lifting of rational cohomological relations is exactly the injectivity of the cycle class map on this subring. For a general Weil coefficient field, this rational statement is distinguished from injectivity after extending Chow coefficients to that field; the latter coefficient convention is addressed separately in Corollary~\ref{c:finite-weil-criterion}. In Chow groups the graph proof has Euler class $c_n(T_X)$ in place of $\chi(X)o$; this replacement itself is a condition to verify, not a consequence of the cohomological computation.

\begin{theorem}[The abelian and hyperelliptic cases]
\label{c:chow-positive}
\status{Known, with the required implication proved}
For an abelian variety with $o$ its origin, every cohomological relation in the rings of C.1(a) holds in $T_k^{\CH}(X,o)$ for every $k$. The same conclusion holds for a hyperelliptic curve with $o$ a Weierstrass point (and for an elliptic curve with its origin and for $\mathbb P^1$).
\end{theorem}

\begin{proof}
For an abelian variety, \citet[Theorem 6.2.5 and Corollary 6.2.6]{c_osullivan2011} identifies the algebra of symmetrically distinguished cycles with the numerical Chow algebra, compatibly with products and with pullbacks and pushforwards by homomorphisms. The origin is the pushforward of $1$ by the zero homomorphism from the zero-dimensional abelian variety. Each diagonal is the pushforward of $1$ by a diagonal homomorphism. Their products and pullbacks are consequently symmetrically distinguished. A homologically zero cycle is numerically zero; the stated isomorphism then forces a distinguished such cycle to vanish in rational Chow groups. This establishes the assertion for precisely the requested generators, without assuming the Hodge or Tate conjecture.

For a hyperelliptic curve of genus $g\geq2$, $K_X=(2g-2)o$ in $\Pic(X)\otimes\Q$. The modified diagonal vanishes by \citet[Proposition 4.8 and its relation with the cycle on $X^3$ in \S4]{c_grossschoen1995}; equivalently one may use the explicit divisor calculation in that proposition followed by the finite symmetric-power push--pull. The remaining relations are precisely \eqref{c:pair-local} and \eqref{c:sympl-pf}. \citet[\S6, equation (2), and Corollary 6.4]{c_tavakol2018} proves them and the perfectness of every tautological intersection pairing on the fixed curve. A homologically zero tautological cycle pairs to zero with every complementary tautological cycle, hence is zero by that perfectness. The identification of Tavakol's class $a_i=K_i/(2g-2)$ with $o_i$ verifies the change of generators. Genus one is covered both by the abelian argument and by Tavakol's Example 6.2. Genus zero is the polynomial Chow presentation of projective space.
\end{proof}

\begin{theorem}[Classification of the allowed zero-cycles on an abelian variety]
\label{c:abelian-point-classification}
\status{Proved over every algebraically closed field}
Let $A/F$ be an abelian variety over an algebraically closed field of any characteristic, and let $o\in\CH_0(A)_\Q$ have degree one. The following are equivalent:
\begin{enumerate}[label=(\roman*)]
\item all the cohomological point-and-diagonal relations hold in $T_k^{\CH}(A,o)$ for every $k$;
\item $\Delta_*o=o\boxtimes o$ in $\CH_0(A\times A)_\Q$;
\item $o=[a]$ in $\CH_0(A)_\Q$ for a point $a\in A(F)$.
\end{enumerate}
The point is unique modulo torsion in $A(F)$. In particular, specifying the origin is sufficient but is not necessary for the requested lifting.
\end{theorem}

\begin{proof}
The equality in (ii) is the cohomological relation $\delta_{12}y_1=y_1y_2$, since both classes are degree-one zero-cycles on $A^2$. Thus (i) implies (ii). If (iii) holds, simultaneous translation by $a$ on every factor preserves all partial diagonals and carries the origin class to $[a]$. Theorem~\ref{c:chow-positive} proves (i).

We prove the remaining implication algebraically, so that no characteristic-zero restriction is hidden. Put $g=\dim A$ and write $\widehat A$ for the dual abelian variety. For the normalized Poincar\'e bundle $\mathcal P$ on $A\times\widehat A$, define the rational Chow Fourier transform by
\[
 \mathcal F_A(z)=(p_2)_*\bigl(p_1^*z\,\exp(c_1(\mathcal P))\bigr).
\]
It is an isomorphism and $\mathcal F_{\widehat A}\mathcal F_A=(-1)^g[-1]^*$. These identities hold over the arbitrary algebraically closed field of \citet[Theorem 2.2, p.~156, with the field convention on p.~154]{c_mukai1981}: apply Grothendieck--Riemann--Roch to his Fourier--Mukai identity, using the trivial tangent bundle of an abelian variety. The Chern character identifies rational $K$-theory with rational Chow groups. We will also use the directly checked Poincar\'e identities below; over $\C$ they are formulas F1--F3 and Proposition 1 of \citet[pp.~647--648]{c_beauville1986}.

Set $z=\mathcal F_A(o)$. Its codimension-zero part is $\deg(o)=1$. Fourier transformation of (ii), using the defining Poincar\'e bundle and the projection formula, gives
\[
 m_{\widehat A}^*z=p_1^*z\,p_2^*z.
\]
Indeed Fourier transformation changes $\Delta_*$ into pullback by the dual homomorphism $m_{\widehat A}$, and changes an exterior product into the exterior product of the transforms. The ideal of positive-codimension classes is nilpotent, so the finite series $D=\log z$ is defined and satisfies
\[
 m_{\widehat A}^*D=p_1^*D+p_2^*D,
 \qquad [n]^*D=nD\quad(n\geq1).
\]

Write $D=\sum_{r=1}^gD_r$ by codimension. We verify that $D_r=0$ for $r\geq2$ without assuming the conjectural vanishing of any negative Beauville piece. Choose an integer $n\geq2$ prime to the characteristic. If $w_q$ is the codimension-$q$ part of $\mathcal F_{\widehat A}(D_r)$, its defining formula uses the power $c_1(\mathcal P)^{g-r+q}$. The identity $(1\times[n])^*\mathcal P=\mathcal P^{\otimes n}$ and flat base change therefore give
\[
 [n]^*w_q=n^{g-r+q}w_q.
\]
Fourier functoriality for the isogeny $[n]$, applied to $[n]^*D_r=nD_r$, gives $[n]_*w_q=nw_q$. Since $[n]_*[n]^*=n^{2g}$, a nonzero $w_q$ requires
\[
 n^{2g}=n^{g-r+q+1},\qquad q=g+r-1.
\]
The inequality $q\leq g$ forces $r\leq1$. Fourier inversion now gives $D=D_1$. On the N\'eron--Severi group $[n]^*$ acts by $n^2$, whereas it acts on $D_1$ by $n$. Hence $D_1\in\Pic^0(\widehat A)\otimes\Q$.

Biduality identifies $\Pic^0(\widehat A)$ with $A(F)$. This group is divisible: multiplication by every positive integer is a surjective isogeny and is surjective on points over an algebraically closed field, also when the integer is divisible by the characteristic. Thus there is $a\in A(F)$ with $D_1=c_1(\mathcal P_a)$ in rational Chow groups. Directly from its definition, $\mathcal F_A([a])=\exp(c_1(\mathcal P_a))=z$. Fourier inversion proves $o=[a]$. Finally $[a]=[b]$ in rational Chow groups if and only if their Albanese classes differ by torsion: the forward implication follows from Albanese, and the reverse from the identical Fourier exponentials of torsion-equivalent points. This proves the uniqueness assertion and the theorem.
\end{proof}

\begin{proposition}[A finite criterion for curves]
\label{c:curve-chow-criterion}
\status{Proved over every algebraically closed field}
Let $C$ be a smooth projective connected curve and $o\in\CH_0(C)_\Q$ have degree one. All point-and-diagonal cohomological relations hold in all its powers if and only if $M_o=0$ in $\CH^2(C^3)_\Q$. If $g\geq2$, this forces
\[
 o=\frac{K_C}{2g-2}\quad\text{in }\CH^1(C)_\Q.
\]
For $g=0$ or $1$ every degree-one rational zero-cycle has the required property. Over $\C$, the genus-at-least-two criterion is also the multiplicativity criterion in \citet[Propositions 3.1 and 3.5]{c_fulaterveervial2021}.
\end{proposition}

\begin{proof}
Necessity follows because $M_o$ is always cohomologically zero for a curve. Suppose it vanishes. Intersecting with $p_1^*o$ and pushing to the last two factors gives
\[
 \Delta_*o=o\boxtimes o.
\]
Intersecting instead with $\delta_{12}$ and pushing to the third factor gives
\[
 c_1(T_C)=(2-2g)o.
\]
In the second calculation, the small-diagonal term contributes $c_1(T_C)$, the term $\delta_{12}o_3$ contributes $(2-2g)o$, and the four nonzero remaining point terms cancel in pairs. These identities and the expansion of $M_o=0$ give all the local relations \eqref{c:pair-local} for $q_{ij}=\delta_{ij}-o_i-o_j$, with parameter $b=-2g$.

We check the remaining matching relation, including for a rational degree-one zero-cycle. Choose a base point $e$ and the Abel--Jacobi map $C\to J(C)$. The Jacobian's principal autoduality and the theorem of the square identify the pullback of its normalized symmetric Poincar\'e class, up to a common sign, with $\delta_{12}-e_1-e_2$. Translation of the Abel--Jacobi map by $z\in J(C)$ adds the corresponding algebraically trivial divisor to each factor. Since $J(C)(k)$ is divisible, we may choose $z$ so that the new pullback is $q_{12}$ in rational Chow groups, up to the same sign. Thus every matching monomial in the $q_{ij}$ is the pullback of a matching monomial in Poincar\'e divisors on a power of $J(C)$.

Let $\theta$ be a symmetric principal polarization on $J(C)$ and use integer parameters $n_1,\ldots,n_{2g+2}$. Pull back $\theta^{g+1}=0$ along the homomorphism $(x_i)\mapsto\sum n_ix_i$. The theorem of the square gives a polynomial identity whose quadratic form is
\[
 \sum_i n_i^2\theta_i+\sum_{i<j}n_in_jP_{ij}.
\]
The coefficient of $n_1\cdots n_{2g+2}$ in its $(g+1)$st power is $(g+1)!$ times the sum over all perfect matchings of the products of $P_{ij}$. A rational polynomial vanishing at every integer tuple has every coefficient zero. Pulling back to $C^{2g+2}$ proves precisely \eqref{c:sympl-pf}. Theorem~\ref{c:pair-ranks} now proves sufficiency in every power. These are algebraic arguments valid also in characteristic two.

For genus zero, $\CH_0(C)_\Q=\Q$. For genus one, choose an origin. The degree-zero Chow group is its Jacobian group, so divisibility shows that every degree-one rational zero-cycle equals a point class. Apply Theorem~\ref{c:abelian-point-classification}. The formula for $o$ when $g\geq2$ follows from the second displayed identity.
\end{proof}

\begin{proposition}[The exact additional obstruction for regular surfaces]
\label{c:surface-chow-criterion}
\status{Proved equivalence; does not assume finite-dimensionality}
Let $S/\C$ be a smooth projective surface with $b_1(S)=0$, and let $o$ have degree one. All the point-and-diagonal cohomological relations lift in all powers if and only if $M_o=0$ and $h(S)$ is Kimura finite-dimensional. If they hold, $o$ is uniquely determined: for every ample divisor $H$,
\[
 o=\frac{H^2}{\deg H^2}.
\]
\end{proposition}

\begin{proof}
The projector $\pi_2=\Delta-o\times S-S\times o$ has realization $H^2(S)$; the other two summands are Tate. Write $b=b_2(S)$. The alternating cross-matching determinant on $2(b+1)$ factors is $(b+1)!$ times the alternating projector on $h^2(S)^{\otimes(b+1)}$. Therefore the lifting of all relations implies $\bigwedge^{b+1}h^2(S)=0$, hence finite-dimensionality. It also implies $M_o=0$, since three positive-degree classes have zero product on such a surface.

Conversely, let $M_o=0$. Acting on $o$, on the diagonal, and on $H\boxtimes H$ gives respectively
\[
 \Delta_*o=o\boxtimes o,\qquad c_2(T_S)=(b+2)o,
 \qquad H^2=(\deg H^2)o.
\]
The same expansions as for the curve give all local relations \eqref{c:pair-local}, now with $b=b_2(S)$. Finite-dimensionality supplies the remaining alternating determinant: the realization is purely even of dimension $b$, so its Kimura exterior dimension is $b$. Theorem~\ref{c:pair-ranks} says that the local identities and these determinants generate every cohomological relation. This proves sufficiency and uniqueness. This formal argument is consistent with the larger divisor-generated criterion in \citet[Proposition 4.5]{c_fulaterveervial2021}; it explicitly applies to the narrower ring asked for here.
\end{proof}

\begin{theorem}[A finite cycle criterion when the cohomology is even]
\label{c:finite-even-criterion}
\status{Proved in the pure-even case}
Let $X/\C$ be smooth projective and connected, with $H^{\rm odd}(X,\Q)=0$,
let $n=\dim X>0$, and let $o\in\CH^n(X)_\Q$ have degree $1$.
Put $d=\dim_\Q H^*(X,\Q)=\chi(X)$.  There is a finite list of rational
tensor expressions $P_1,\ldots,P_s$, constructed by the procedure below
from the rational cup-product table and trace of $X$, for which the
following criterion is necessary and sufficient for all the relations
of Theorem~\ref{c:matrix-presentation} to hold in rational Chow groups
on every power:
\begin{enumerate}[label=(\roman*)]
\item $c_n(T_X)=d\,o$ and $\Delta_*o=o\boxtimes o$;
\item $\bigwedge^{d+1}h(X)=0$ in rational Chow motives;
\item $P_a(h(X),\mu,1,\operatorname{tr},o,\beta^{-1})=0$ for
      $1\leq a\leq s$.
\end{enumerate}
Here $\mu$ is the small-diagonal multiplication,
$\operatorname{tr}:h(X)\to\Q(-n)$ is the trace,
$\beta=\operatorname{tr}\circ\mu$, and
$\beta^{-1}:\Q(-n)\to h(X)\otimes h(X)$ is its inverse pairing.
Every condition is the vanishing of a specified finite linear
combination of correspondence diagrams.  Subject to (i), each can
equivalently be written as a finite list of cycles in the requested
point-and-diagonal rings.  The construction uses finite-dimensional
polynomial algebra, so is a finite criterion, rather than one separate
test for each power of $X$.
\end{theorem}

\begin{proof}
Write $A=H^*(X,\Q)$ as an ordinary vector space; the pure-even
hypothesis makes its tensor symmetry the ordinary one.  Retain the
Tate degree by a one-dimensional auxiliary representation $L$, and
put $G=\operatorname{GL}(A)\times\mathbb G_m$, with the second factor
acting on $L$.  The data
\[
 t=(\mu,1,\operatorname{tr},o,\beta^{-1})
\]
form one rational point of the finite-dimensional rational
$G$-representation
\[
 W=(A\otimes A^\vee\otimes A^\vee)\oplus A
   \oplus(A^\vee\otimes L^n)
   \oplus(A\otimes L^{-n})
   \oplus(A\otimes A\otimes L^{-n}).
\]
In the point $t$, the coordinate denoted by $o$ is the cohomology class
of the degree-one zero-cycle.  Let $I\subset\Q[W]$ be the ideal of
the Zariski closure of $G\cdot t$.  We describe all steps producing
the finite list.

First compute generators of $I$ by elimination: write the orbit map
in matrix entries of $g\in\operatorname{GL}(A)$ and a nonzero scalar
$z$, introduce inverses with equations $u\det(g)=1$ and $vz=1$, and
eliminate $g,u,z,v$ from the graph ideal.  The resulting elimination
ideal is the ideal of the image closure.  Hilbert's basis theorem
makes the resulting list finite.  Enlarge its span to a finite-dimensional
$G$-stable subspace $E\subset I$ generating $I$ as an ideal.
This enlargement is finite because the action on the polynomial ring
is locally finite and preserves the finite degree bound of the chosen
generators.  It can be computed by closing their span under the matrix
units of the Lie algebra and the second torus action; $G$ is connected
in characteristic zero.

The inclusion $E\hookrightarrow\Q[W]$ is equivalently a polynomial
covariant
\[
 p_E\in(E^\vee\otimes\Q[W])^G.
\]
Choose a split $G$-embedding of $E^\vee$ into a finite direct sum $T$
of mixed tensor powers of $A,A^\vee,L,L^\vee$.  Such a split embedding
exists because rational representations of $G$ are semisimple and
are generated by its defining representations and their duals.
For an explicit terminating search, enumerate these finite direct
sums and compute the rational intertwiner spaces.  Enumerate their
rational elements until one has full rank on $E^\vee$, and then
solve the linear intertwiner equations for a retraction.  The
existence of a split embedding proves termination.  The components of the image of
$p_E$ in $T\otimes\Q[W]$ are the covariants we denote by $P_a$.

We recall the two representation-theoretic facts needed here.
Over $\C$, averaging a Hermitian form over
$U(d)\times U(1)$ gives invariant complements.  This compact group
is Zariski dense in $G_\C$, so the complements are $G_\C$-stable.
The intertwiner and splitting equations have rational coefficients;
Gaussian elimination descends their solvability to $\Q$.
Every finite-dimensional rational representation embeds, by its
matrix coefficients, into finitely many copies of $\Q[G]$.
The relevant coefficient functions have bounded polynomial degree
and bounded powers of $\det^{-1}$ and the second torus coordinate.
They therefore occur in mixed tensor powers, using symmetric and
exterior projectors.  Complete reducibility makes the embedding
split.  This proves the existence statements just used.  It also
proves exactness of invariants for the locally finite rational
modules below: place a preimage and its translates in a
finite-dimensional submodule and split the resulting surjection.
The exact references are \citet[Theorem~4.14, pp.~88--89;
Theorem~22.42, pp.~478--479; Remark~12.53(a), p.~248,
and the Reynolds operator on p.~249]{c_milneag2022}.

We explain why these covariants are tensor expressions of the promised
kind.  Polarize a homogeneous polynomial covariant in the finitely
many tensor coordinates of $W$.  Its coefficients are invariants in
mixed tensor powers of $A$ and $A^\vee$.  The first fundamental
theorem says that these are sums of complete contractions, with any
uncontracted indices supplying the target tensor.  The only linear
relations among such contractions are generated by the
$(d+1)$-fold antisymmetrizer and by the scalar loop value $d$.
These are precisely \citet[Theorems~1.2--1.5, pp.~4--5]{c_procesi2021tensor};
polarization and restitution apply because the characteristic is zero.
This use is of the classical tensor theorems, and does not apply a
matrix-only statement to arbitrary tensor variables without that step.
The extra $\mathbb G_m$ factor merely selects the matching Tate
degrees.  Thus each $P_a$ is a finite rational linear combination of
the prescribed diagrams, and every universal identity in their
ordinary $d$-dimensional interpretation follows from the stated
antisymmetrizer and loop relation.

Now let $F$ be any polynomial covariant with target a mixed tensor
space $T'$ that vanishes at $t$.  Equivariance implies that it
vanishes on $G\cdot t$ and its closure.  Hence
\[
 F\in(T'\otimes I)^G.
\]
The multiplication map $E\otimes\Q[W]\twoheadrightarrow I$ is
$G$-equivariant.  Since $G$ is reductive over $\Q$, taking invariants
is exact, and $F$ lifts to
\[
 \widetilde F\in(T'\otimes E\otimes\Q[W])^G.
\]
Contracting its $E$ index against $p_E$ gives $F$ exactly.  Extend
that contraction through the chosen split embedding into $T$.
The first fundamental theorem again writes $\widetilde F$ as a
diagram.  Consequently the tensor ideal of all relations at $t$
is generated by the finite covariants $P_a$, together with the
$(d+1)$-fold antisymmetrizer and the loop relation.  No reductivity
of the stabilizer of $t$ has been assumed or used.

Evaluate this presentation in Chow motives.  The scalar loop value
is $\chi(X)=d$ by the categorical trace formula.  Condition (ii)
annihilates the dimension relation, and condition (iii) annihilates
the finite orbit-ideal generators.  Therefore every cohomological
diagram relation vanishes in Chow, in particular every relation in
the requested rings.  This proves sufficiency.
The additional identities for the auxiliary line,
$\bigwedge^2 L=0$ and $\operatorname{tr}(\id_L)=1$, already hold
for the fixed even invertible Tate object $L=\Q(-1)$; they are
not additional conditions on $(X,o)$.

Conversely suppose all the requested relations hold in Chow.
The cohomological relations
$\Delta\cdot o_1=o_1o_2$ and
$\Delta^2=d\,o_1o_2$ hold there.  The first is the second assertion
of (i); pushing the second to $X$ gives its first assertion.
The alternating projector in (ii) is a rational linear combination
of permutation graphs on $X^{d+1}$, each a product of pairwise
diagonals in $X^{2(d+1)}$.  Its cohomology is zero, so (ii) follows.

Finally the correspondence diagrams used above do not enlarge the
vanishing test after (i).  The Frobenius identities reduce each
connected diagram to a small diagonal with some point insertions
and some handles.  A handle inserts $c_n(T_X)$; a point insertion
inserts $o$.  By (i) both reduce to scalar multiples of $o$, two
such insertions vanish for dimension reasons, and a single insertion
on a small diagonal is the external product of the point classes.
Closed components give their degrees.  Thus every diagram reduces
to the pointed partitions of Definition~\ref{c:def-pointed}.
This is the same graph reduction as Theorem~\ref{c:graph-rule},
now including contractions and using the two verified Chow
identities.  Each $P_a$ has zero cohomology and reduces to a relation
of the requested kind; it consequently vanishes in Chow.  This
proves necessity and the finite criterion.
\end{proof}

\begin{remark}[Scope of the finite criterion]
\label{c:finite-criterion-scope}
\status{Proved scope clarification; the arbitrary-pair classification remains unresolved}
Theorem~\ref{c:finite-even-criterion} applies to arbitrary cup-product
algebras with pure-even cohomology, not only to surfaces.  It gives
a finite algebraic construction of the cycle conditions, without
claiming to decide their rational equivalence by computation.
For mixed cohomology the direct use of
$\operatorname{GL}_{a|b}$ would require exactness of its invariant
functor, which fails.  Theorem~\ref{c:finite-super-criterion}
replaces that step by projective lifting and an external algebraic
parity correspondence.  Corollary~\ref{c:finite-weil-criterion}
also separates finite existence over an arbitrary Weil coefficient
field from effective construction with specified coefficient data.
Theorem~\ref{c:finite-unconditional-criterion} subsequently removes
the sign hypothesis by first passing to scalar forms in the
covariant polynomial category.  These finite criteria do not
decide the resulting Chow vanishings.
The unrestricted classification of all $(X,o)$ therefore remains
distinct from the specific K3 obstruction proved below.
\end{remark}


\begin{proposition}[Products and the failure of the tautological parity converse]
\label{c:product-parity}
\status{Product assertion proved; the proposed parity converse refuted}
Let $(X,o_X)$ and $(Y,o_Y)$ be pointed smooth projective complex
varieties for which all the requested point-and-diagonal relations
hold in rational Chow groups on every power.  Then the same is
true of $(X\times Y,o_X\boxtimes o_Y)$.
Nevertheless, this property does not imply that cohomological
parity is represented by a correspondence in the tensor category
generated by the marked Frobenius diagrams.  A counterexample is
\[
 X=E\times\mathbb P^1,\qquad
 E:\ y^2z=x^3-xz^2,\qquad
 o=[(0:1:0)]\boxtimes[\infty].
\]
Its parity correspondence is algebraic, but lies outside that
diagram category.
\end{proposition}

\begin{proof}
Under $(X\times Y)^k=X^k\times Y^k$, each marked diagonal is
the external product of the two corresponding marked diagonals.
Let $A_k$ and $B_k$ be the rational Chow tautological rings of the
two factors.  The two injective cycle maps and the rational Betti
K\"unneth isomorphism make
\[
 A_k\otimes_\Q B_k\longrightarrow
 H^*(X^k,\Q)\otimes_\Q H^*(Y^k,\Q)
\]
injective.  Its factorization through the external-product map
to $\CH^*(X^k\times Y^k)_\Q$ shows that the cycle map is
injective on the image of that map.  The requested product ring
is contained in this image, proving the first assertion.

For the displayed pair, the elliptic and projective-line cases
were proved in Theorem~\ref{c:chow-positive}.
Also $c_2(T_X)=0=\chi(X)o$, because $T_E$ is trivial, and
$\Delta_*o=o\boxtimes o$ because $o$ is an actual point.
The Frobenius graph reduction therefore applies to every
contracted diagram, and its degree-zero endomorphisms of $h(X)$
are contained in
\[
 T^{\CH}_2(X,o)\cap\CH^2(X^2)_\Q
   =\operatorname{span}_\Q\{\Delta_X,o\times X,X\times o\}.
\]
The latter two correspondences act only on $H^0$ and $H^4$,
respectively.  Thus any correspondence in this span acts by the
same scalar on all of $H^1,H^2,H^3$.  Both $H^1$ and $H^2$ are
nonzero, whereas parity must act as $-1$ on the former and $+1$
on the latter.  This is impossible.  Explicitly, the three diagonal
matrices of this span have rank three as vectors of endomorphism
coefficients, and adjoining parity raises the rank to four; the
exact computation is recorded with the accompanying code.

On the other hand, the graph of
$[-1]_E\times\id_{\mathbb P^1}$ acts as parity on cohomology,
so parity is algebraic in this example.  The counterexample
refutes only the proposed implication that parity belongs to the
\emph{narrow diagram category}.  It neither refutes algebraicity
of parity nor rules out a different finite mixed-parity criterion.
\end{proof}


\begin{lemma}[Rational projective mixed tensors]
\label{c:super-projective-tensors}
\status{Proved}
Let $A=\Q^{a|b}$ be a finite supervector space, let $L$ be an even
line, and put $G=\operatorname{GL}(A)\times\mathbb G_m$.
We use even morphisms and the category $\mathcal R$ of rational
$G$-supermodules on which
$\epsilon=(\operatorname{diag}(1_a,-1_b),1)$ acts as intrinsic parity.
Every finite object $E$ of $\mathcal R$ is a quotient of a finite
projective-injective object $P$ of $\mathcal R$.  Such $P$, its dual,
and a fixed nonzero projective object $P_0$ are direct summands of
finite sums of mixed tensor powers of $A,A^\vee,L,L^\vee$.
The modules, quotient maps, mixed tensor embeddings and retractions
can be constructed over $\Q$ by terminating finite procedures.
Projectivity also applies to surjections of locally finite rational
modules, including the polynomial modules used below.
\end{lemma}

\begin{proof}
Write $G_0=\operatorname{GL}_a\times\operatorname{GL}_b\times
\mathbb G_m$, $\mathfrak g=\operatorname{Lie}G$, and
$B=U(\mathfrak g_0)$.  We take
\[
 P=U(\mathfrak g)\otimes_B E|_{G_0},\qquad
 P\longrightarrow E,\quad v\otimes e\longmapsto ve.
\]
The odd part of $\mathfrak g$ is
$\operatorname{Hom}(\Q^a,\Q^b)\oplus
\operatorname{Hom}(\Q^b,\Q^a)$.
The PBW theorem makes $P$ finite, of total dimension
$2^{2ab}\dim E$.  Its natural rational $G_0$-action, together with
the compatible $\mathfrak g$-action, integrates to an algebraic
$G$-action over $\Q$.  Here we use
\citet[Corollary~3.6, p.~9; Proposition~5.4 and its proof,
pp.~23--24; Theorem~5.8 and Corollary~5.10, p.~26]{c_masuokashibata2017super}.
The identity $\operatorname{hy}(G)=U(\mathfrak g)$ in characteristic
zero is stated on p.~22 there.  In the proof of Proposition~5.4,
PBW splits off the finite exterior factor, the rational even action
supplies its coefficient functions, and injectivity of the pairing
with the hyperalgebra proves the coassociativity of the resulting
coaction.  Thus the statement applies over $\Q$, not only after
extension to $\C$.  The induced module still satisfies the
$\epsilon$-parity condition.

Induction is left adjoint to restriction to $G_0$.  The latter has
semisimple finite rational representations, so $P$ is projective.
In this case induction and coinduction coincide, with no character
or parity twist.  We give the reason to check this point over $\Q$:
\[
 \det(\mathfrak g_1)
 =\det(\Q^a)^{-b}\det(\Q^b)^a
  \det(\Q^a)^b\det(\Q^b)^{-a}=1,
 \qquad \dim\mathfrak g_1=2ab.
\]
Choose an ordered odd PBW basis.  Projection to its top exterior
coefficient gives a $B$-bimodule map $U(\mathfrak g)\to B$:
the displayed determinant identity is precisely the vanishing of
the possible even-character twist, and the top degree is even.
Multiplication followed by this map is a nondegenerate relative
pairing.  Complementary exterior monomials pair to $\pm1$, while
the commutator correction terms are triangular in exterior degree.
It identifies $U(\mathfrak g)$ with its $B$-linear dual with the
induction/coinduction bimodule structures.  Tensoring with $E$
gives $\operatorname{Ind}(E)\simeq\operatorname{Coind}(E)$.
Right adjunction and semisimplicity over $G_0$ now make $P$
injective as well.  We may take $P_0=\operatorname{Ind}(\Q)$.

For completeness, the tensor-generation assertion can be checked
without assuming $a\geq b$.  Matrix coefficients and the coaction
embedding express every finite rational superrepresentation as a
subquotient of finite sums of mixed tensors and their parity shifts.
Indeed $\Q[G]$ is generated by the entries of the standard
supermatrix, its inverse, and the two powers of the last torus
coordinate.  Its usual localization by the determinants of the two
even blocks has the same generators: an inverse diagonal block is
recovered from the corresponding block of the full inverse by a
finite geometric series in the odd, hence nilpotent, correction.
Every finite set of coefficient functions therefore comes from
bounded mixed tensor degrees.  To retain the correct category,
the operator consisting of intrinsic parity times $\rho(\epsilon)$
splits the full superrepresentation category into $\mathcal R$ and
its parity shift.  It is $+1$ on the mixed tensors in question and
$-1$ on their parity shifts.  Taking its $+1$ summand shows that an
object of $\mathcal R$ is already a subquotient of the unshifted
mixed tensors.  This also explains the category convention in
\citet[section~2, pp.~3--4]{c_heidersdorf2017mixed};
his Theorem~1.1, p.~2, includes the assertion that every projective
is a mixed tensor.

Write $P=U/V$ with $U\subset T$, where $T$ is a finite sum of mixed
tensors.  Projectivity splits $U\twoheadrightarrow P$, so $P$
embeds in $T$; injectivity retracts that embedding.  The dual is
then a summand of $T^\vee$.  These operations are rational.  For an
explicit terminating search, we enumerate finite mixed tensor sums,
compute their rational intertwiner spaces using the even and odd
Lie generators and the torus weights, find an intertwiner of full
rank, and solve the linear equations for its retraction.  The
existence just proved ensures termination.  All the linear
equations have rational coefficients.

Finally $\operatorname{Hom}_{G}(P_0,-)=
\operatorname{Hom}_{G_0}(\Q,-)$ is exact on locally finite rational
modules.  More generally the same assertion holds for the induced
$P$ with $E$ in place of $\Q$.  Exactness follows by placing
finitely many desired preimages in a finite $G_0$-submodule and
using its semisimplicity.  Alternatively, a finite-dimensional
image in a locally finite $G$-module has preimages in a finite
$G$-submodule, so finite projectivity suffices.  This proves the
last assertion without treating an infinite polynomial module as
a finite representation.
\end{proof}

\begin{theorem}[A finite criterion under the sign conjecture]
\label{c:finite-super-criterion}
\status{Proved under the explicitly stated hypothesis $C^+(X)$}
Let $X/\C$ be smooth projective and connected of dimension $n>0$,
let $o\in\CH^n(X)_\Q$ have degree $1$, and put
\[
 a=\dim H^{\rm even}(X,\Q),\qquad
 b=\dim H^{\rm odd}(X,\Q),\qquad
 N=(a+1)(b+1).
\]
Assume $C^+(X)$: a rational algebraic correspondence on $X$ realizes
the even cohomological projector.  Equivalently, some
$u\in\operatorname{End}(h(X))$ realizes cohomological parity
$(-1)^i$ on $H^i(X,\Q)$.  Neither idempotence of its Chow lift nor
membership in the point-and-diagonal ring is assumed.
There is a finite list of rational marked-Frobenius tensor
expressions $P_1,\ldots,P_s$, constructed from the rational cup
table by the procedure below, such that all the cohomological
relations in the requested rings on every power of $X$ hold in
rational Chow if and only if the following finite conditions hold:
\begin{enumerate}[label=(\roman*)]
\item $c_n(T_X)=(a-b)o$ and $\Delta_*o=o\boxtimes o$;
\item $y_\lambda(h(X)^{\otimes N})=0$, where
      $\lambda=((b+1)^{a+1})$ and $y_\lambda$ is its rectangular
      Young symmetrizer;
\item $P_s(h(X),\mu,1,\operatorname{tr},o,\beta^{-1})=0$ for
      every member of the constructed finite list.
\end{enumerate}
The Young symmetrizer is the row symmetrizer followed by the column
antisymmetrizer, with the tensor symmetry of Chow motives; no
choice of its nonzero normalization affects (ii).
Subject to (i), conditions (ii)--(iii) are finite lists of cycles
in the specified point-and-diagonal rings.  No assertion that
rational equivalence itself is algorithmically decidable is made.
\end{theorem}

\begin{proof}
We spell out the superalgebra construction, because exactness of
the invariant functor is unavailable here.  The hypothesis $C^+$
has exactly the usual meaning in
\citet[section~1, p.~1389]{c_tabuada2018sign}; if $v$ realizes the
even projector, $u=2v-\id$ realizes parity.
Put $A=H^*(X,\Q)$ with its parity, $L=\Q(-1)$ as an even line,
and $G=\operatorname{GL}(a|b)\times\mathbb G_m$.
All tensor representations used in the covariant construction,
including $W$, $R$ and $I$ below, belong to the category $\mathcal R$
of Lemma~\ref{c:super-projective-tensors}.  The entire regular
representation $\mathcal O(G)$ need not belong to that category.
The same marked tensors as in
Theorem~\ref{c:finite-even-criterion} give an even rational point
$t=(\mu,1,\operatorname{tr},o,\beta^{-1})$ in
\[
 W=(A\otimes A^\vee\otimes A^\vee)\oplus A
   \oplus(A^\vee\otimes L^n)
   \oplus(A\otimes L^{-n})
   \oplus(A\otimes A\otimes L^{-n}).
\]
Write $R=\mathcal O(W)$ and define
\[
 I=\ker\bigl(R\longrightarrow\mathcal O(G)\bigr),
 \qquad f\longmapsto\bigl(g\longmapsto f(g\cdot t)\bigr).
\]
Thus $I$ is the ideal of the schematic super-orbit closure.  We
retain its odd directions and nilpotents; the ideal of the orbit
of ordinary complex points would not suffice.

The supercommutative algebra
$R=\operatorname{Sym}(W_{\bar0}^\vee)\otimes
\bigwedge(W_{\bar1}^\vee)$ is finite over its ordinary even
polynomial subalgebra and is noetherian.  We can compute generators
of $I$ by orbit-graph elimination.  Coordinates on $G$ consist of
the two even matrix blocks, the odd off-diagonal blocks, and the
inverses of the two even determinants and of the last torus
coordinate.  Introduce ordinary variables for those inverses and
their defining equations.  Since the odd variables have a finite
exterior basis, elimination is a finite-module polynomial
elimination with explicitly known signed multiplication matrices.
The usual noetherian termination of module Gr\"obner elimination
therefore applies.  Enlarge the resulting finite generating span
to a finite $G$-stable submodule $E\subset I$.  The action preserves
polynomial degree, so this is a finite linear computation; one can
close the span under the rational coaction or the even and odd Lie
generators together with the torus weights.

Use Lemma~\ref{c:super-projective-tensors} to construct a finite
projective-injective $P\twoheadrightarrow E$ and a split mixed
tensor embedding of $P^\vee$.  The composite $P\to E\hookrightarrow R$
corresponds to the polynomial covariant
\[
 p_P\in(P^\vee\otimes R)^G.
\]
Its components in the chosen mixed tensor sum are our finite
expressions $P_1,\ldots,P_s$.  We also fix the nonzero projective
$P_0=\operatorname{Ind}_{G_0}^{G}(\Q)$ and a mixed tensor projector
whose image is $P_0$.

We verify both the diagram interpretation and its exact relations.
\citet[Theorem~2.1, Corollary~2.2, Theorem~2.3 and
Corollary~2.5]{c_lehrerzhang2021super} state the first and second
fundamental theorems for the general linear supergroup in this
scope: mixed invariants are contractions, and their kernel is the
tensor ideal generated by the Young symmetrizer with $a+1$ rows
and $b+1$ columns.  Although the article's title concerns an
orthosymplectic supergroup, these numbered statements are explicitly
its preliminary general linear theorems.  To pass from the
endomorphism version to all mixed tensors we bend dual indices;
the central scalar in $\operatorname{GL}(A)$ forces the two
remaining tensor degrees to be equal.  The stated kernel in
$\Q[\mathfrak S_r]$ is the ideal generated by the rectangular
symmetrizer in the first $N$ tensor positions.  Bending back gives
exactly its tensor ideal in all mixed degrees.  The scalar loop
has value $a-b$.

For polynomial covariants, the supersymmetric power in each finite
degree is the image of the averaging idempotent for the signed
symmetric-group action on the corresponding tensor power of
$W^\vee$.  That idempotent provides a split inclusion before taking
invariants.  Polarization and the preceding tensor theorem
therefore express the covariant as a finite contraction diagram;
no exactness of supergroup invariants is needed.  This is also the
superpolarization proved in Appendix~A, Proposition~A.1 and
Corollary~A.2 of the same source.  The source states these theorems
over $\C$, but all invariant spaces, maps and symmetrizers here
are defined by rational linear equations.  The equality of their
images and kernels descends along $\Q\subset\C$.
The extra torus just enforces the Tate degree.  The fixed Tate
line already satisfies $\bigwedge^2L=0$ and
$\operatorname{tr}(\id_L)=1$.

We now replace the invalid reductive invariant argument.  Let
$F\in(T'\otimes R)^G$ be any polynomial diagram covariant that
vanishes at $t$, where $T'$ is a mixed tensor target.  Equivariance
on superalgebra-valued points implies that its coefficients lie
in $I$, so $F\in(T'\otimes I)^G$.  Multiplication gives an
equivariant surjection $P\otimes R\twoheadrightarrow I$.
Regard $F\otimes\id_{P_0}$ as a $G$-map
\[
 P_0\longrightarrow T'\otimes I\otimes P_0.
\]
Projectivity lifts it through
\[
 T'\otimes P\otimes R\otimes P_0
     \twoheadrightarrow T'\otimes I\otimes P_0.
\]
The lift has bounded polynomial degree: its domain is finite and
the target is locally finite.  The preceding invariant theorem
writes it as a finite diagram between the chosen mixed tensor
summands.  Composing with $p_P$ and multiplying polynomial factors
recovers $F\otimes\id_{P_0}$.  Consequently every relation,
after tensoring with this one fixed $P_0$, lies in the tensor ideal
generated by the finite list $P_s$ and the rectangular universal
relation.  This is a projective lifting argument, not an assertion
that $(\cdot)^G$ is exact.

Evaluate in rational Chow motives, assuming (i)--(iii).  The loop
value is already $\chi(X)=a-b$.  Condition (ii) kills all universal
mixed tensor relations.  In particular, a diagram lift $e$ of the
projector defining $P_0$ is now an actual Chow idempotent: the
universal identity $e^2-e$ is among those relations.  Thus
$P_0(h(X))=\operatorname{im}e$ is defined.  Condition (iii) gives
\[
 F(h(X))\otimes\id_{P_0(h(X))}=0.
\]
It remains to cancel the projective factor, whose categorical
dimension can itself be zero.

Let $u$ be the algebraic correspondence realizing parity supplied
by $C^+(X)$.  On the mixed tensor sum $T(h(X))$ containing $P_0(h(X))$
we tensor $u$ and its dual, and use the identity on the even Tate
line; denote the resulting endomorphism by $u_T$.
Set $u_0=e u_T e$ on $P_0(h(X))$.  Its Betti realization is
parity on $P_0(A)$, because the even projector $e$ commutes with
parity in supervector spaces.  Hence
\[
 \operatorname{tr}(u_0)
  =\operatorname{str}\bigl((-1)^{\rm parity}|P_0(A)\bigr)
  =\dim P_0(A)_{\bar0}+\dim P_0(A)_{\bar1}>0.
\]
The first equality is an equality of rational scalars: Betti
realization is injective on $\operatorname{End}(\mathbf1)=\Q$.
Compose the preceding vanishing with $\id\otimes u_0$ and take
the partial categorical trace.  We obtain
$\operatorname{tr}(u_0)F(h(X))=0$, whence $F(h(X))=0$.
No equality $u^2=1$ in Chow, and no commutation of $u_T$ with $e$
in Chow, was used.  This proves sufficiency for every diagram
relation and therefore for every requested relation.

Conversely, all requested relations imply (i) by
$\Delta\cdot o_1=o_1o_2$ and
$\Delta^2=(a-b)o_1o_2$ and their pushforwards.  The Young
symmetrizer in (ii) is a finite rational combination of permutation
graphs on $X^N$, each a product of the allowed diagonals on
$X^{2N}$, and its Betti class vanishes.  Thus (ii) follows as well.
After (i), the Frobenius graph normalization in the proof of
Theorem~\ref{c:finite-even-criterion} reduces every dualized or
contracted marked diagram to a rational combination of pointed
partitions: each handle inserts $c_n(T_X)$, each marking inserts
$o$, and the two identities in (i) remove these insertions or turn
them into external point classes.  Closed components give their
degrees.  Every $P_s$ vanishes at $t$, so its normalized cycle is
one of the requested cohomological relations.  It therefore
vanishes in Chow.  This proves necessity and the finite iff.
\end{proof}

\begin{corollary}[Finite existence over an arbitrary Weil coefficient field]
\label{c:finite-weil-criterion}
\status{Proved under $C^+_H(X)$; effectiveness requires the stated coefficient data}
Let $X/k$ be smooth projective and connected over an arbitrary
algebraically closed field and of positive dimension, let $H$ be a
Weil cohomology with coefficient field $K$ of characteristic zero,
and let $o$ be a
rational zero-cycle of degree one.  Suppose $C^+_H(X)$, meaning
that a rational algebraic correspondence realizes parity on
$H^*(X)$.  The finite criterion (i)--(iii) of
Theorem~\ref{c:finite-super-criterion} still exists, with a finite
list of \emph{rational} marked diagrams, and is equivalent to
injectivity of
\[
 T_k^{\CH}(X,o)\longrightarrow H^*(X^k)
 \quad\hbox{for every }k.
\]
Here injectivity refers to rational linear relations, as required
for cycles in $\CH_\Q$.  No rational form of the whole cohomology
algebra is assumed.  The finite list is effectively constructible
if a finite presentation of the $\Q$-algebra generated by its
chosen structure constants is supplied.  For arbitrary elements
of an unspecified coefficient field, only finite existence is
asserted.  In dimension zero, $X=\operatorname{Spec}(k)$ and
all the requested relations hold immediately.
\end{corollary}

\begin{proof}
Choose a homogeneous $K$-basis of $H^*(X)$ and let
$A_\Q=\Q^{a|b}$ be a formal supervector space of the same
dimensions.  It is not asserted to be a rational form of the cup
algebra.  Build $G_\Q$, $W_\Q$ and $R_\Q=\mathcal O(W_\Q)$
from this formal space exactly as above.  The actual structure
tensors give a $K$-point $t$ of $W_\Q$.  Replace the orbit ideal by
\[
 I_\Q=\ker\left(R_\Q\longrightarrow
        \mathcal O(G_\Q)\otimes_\Q K\right),
 \qquad f\longmapsto\bigl(g\longmapsto f(g\cdot t)\bigr).
\]
The map is a $G_\Q$-comodule map, with the regular action on the
first target factor and trivial action on the coefficient field.
Its kernel is consequently a rational $G_\Q$-stable ideal.
Equivalently, the kernel assertion follows by tensoring its
underlying exact sequence with $\mathcal O(G_\Q)$ and using the
coaction identity.  A rational covariant $F$ satisfies $F(t)=0$
if and only if its coordinates lie in $I_\Q$: in one direction
equivariance gives $F(g\cdot t)=gF(t)$ on all superalgebra-valued
points, and in the other we evaluate at the identity.

Noetherianity of $R_\Q$ gives finite ideal generators; closing
their span in a bounded polynomial degree gives a finite rational
stable $E\subset I_\Q$.  Lemma~\ref{c:super-projective-tensors}
and every subsequent projective lift take place over $\Q$.
Evaluation now takes place in $K$-cohomology.  The trace cancellation
still gives the positive integer $\dim_\Q P_0$, since
$\operatorname{End}(\mathbf1)=\Q$ embeds in $K$.
The same necessity argument uses rational permutation graphs and
the Frobenius normalization in Chow.  Thus the finite iff survives
unchanged.

For the effectiveness assertion, let $c_1,\ldots,c_m\in K$ be
all the coordinates of $t$, including the inverse pairing, and
write $B=\Q[c_1,\ldots,c_m]$.  If a finite presentation
$B=\Q[z_1,\ldots,z_m]/J$ is given, the orbit map lands already
in $\mathcal O(G_\Q)\otimes B$, which injects into
$\mathcal O(G_\Q)\otimes K$.  Orbit-graph elimination including
the equations of $J$ therefore computes $I_\Q$.  Without a
presentation of $B$, no such computational claim follows merely
from the abstract knowledge of the elements $c_i$.

If instead one asks for every $K$-linear relation to vanish in
$\CH_\Q\otimes_\Q K$, the same proof is performed entirely over
$K$, using the full $K$-orbit ideal.  That is a separate
coefficient convention; it is not deduced by tensoring an
injective rational cycle map with $K$ and identifying its target
with the original cohomology space.
\end{proof}

\begin{corollary}[The sign hypothesis over an algebraic closure of a finite field]
\label{c:finite-field-sign}
\status{Proved for the stated cohomology theories}
For $X/\overline{\F}_p$, the sign hypothesis in
Corollary~\ref{c:finite-weil-criterion} holds for $\ell$-adic
cohomology, $\ell\ne p$, and crystalline cohomology with
characteristic-zero coefficients.  More generally it holds for
every Weil theory satisfying the Poincar\'e duality, weak
Lefschetz and zeta-function hypotheses in section~II of
\citet{c_katzmessing1974}.  Thus in this setting the finite
criterion has no additional sign-conjecture hypothesis.
\end{corollary}

\begin{proof}
Descend $X$ to a smooth projective geometrically connected variety
over a finite field $\F_q$.  The finitely many points and rational
coefficients representing $o$ also descend after enlarging
$\F_q$; the parity construction itself does not depend on this
choice of $o$.  Put
$Q_i(T)=\det(T-F\mid H^i(X))$.
\citet[Theorem~1, pp.~74--75; Theorem~2(1), pp.~76--77]{c_katzmessing1974}
proves that $Q_i\in\Q[T]$ and that all its roots have absolute
value $q^{i/2}$ in every complex embedding, with the same
polynomials for the cohomology theories just specified.
Consequently the nonconstant $Q_i$ are pairwise coprime.
By the Chinese remainder theorem there are polynomials
$r_i\in\Q[T]$ with $r_i\equiv1\pmod{Q_i}$ and
$r_i\equiv0\pmod{Q_j}$ for $j\ne i$.
Cayley--Hamilton shows that $r_i(F)$ acts as the projector onto
$H^i(X)$.  Each is a rational combination of graphs of iterates
of the Frobenius morphism, and hence an algebraic correspondence.
Their sum over even $i$ gives the required sign projector after
base extension to $\overline{\F}_p$.  This is the explicit
polynomial argument in Theorem~2(1), so it does not require
Frobenius semisimplicity.
\end{proof}


\begin{remark}[What the additional finite criterion settles]
\label{c:super-criterion-scope}
\status{Proved scope clarification; unrestricted classification remains unresolved}
Theorem~\ref{c:finite-super-criterion} supplies a finite cycle
criterion for mixed cohomology whenever $C^+(X)$ is known.  The
parity correspondence may be completely external to the requested
ring, as it is for the positive pair of
Proposition~\ref{c:product-parity}.  Thus that example's obstruction
to a tautological parity projector does not obstruct this proof.
We do not infer $C^+(X)$ from the point-and-diagonal relations.
Theorem~\ref{c:finite-unconditional-criterion} below removes that
hypothesis from the finite criterion by a different scalar-form
construction, which needs no cancellation step.
Corollary~\ref{c:finite-weil-criterion} removes the rational Betti
restriction for finite existence; its effectiveness qualification
is essential.
The criterion does not prove the finite Chow vanishing conditions
for an arbitrary input variety.  The general classification in
C.1(c), and in particular the K3 vanishing below, remain open
items of this task rather than consequences of finite generation.

Condition (ii) retains a precise finite-dimensionality obstruction.
Under $C^+_H(X)$ it is equivalent to Kimura finite-dimensionality,
by \citet[Theorem~3.2, p.~286]{c_delpadronemazza2005}.
To check the input, their Proposition~3.1 applies the partial-trace
identity of Theorem~2.1 to an algebraic even projector: all its
power traces equal $a$, which lies outside the forbidden contents
for the minimal $(a+1)$-by-$(b+1)$ rectangle.  The numerical ideal
is therefore nilpotent.  The same rectangle identity propagates
to tensor powers by the super tensor relations, so their
numerical ideals are nilpotent as well.  The even and odd
cohomological projectors lift across these nilpotent ideals, and
the cohomologically zero exterior and symmetric projectors on
their tensor powers are nilpotent idempotents and hence zero.
This gives the Kimura decomposition.  Conversely its even and
odd dimensions give the rectangular Schur vanishing.  In
particular, arbitrary Schur-finiteness with an unspecified larger
partition is not being substituted for this exact condition.
\end{remark}


\begin{lemma}[Exact lifting for covariant polynomial tensors]
\label{c:covariant-polynomial-lifting}
\status{Proved}
Let $V_0=\Q^{u|v}$, let $L_0$ be an even line, and set
$G=\operatorname{GL}(V_0)\times\mathbb G_m$.  Fix $q\geq1$ and put
\[
 W=\bigoplus_{m=1}^{q}(V_0^\vee)^{\otimes m}\otimes L_0,
 \qquad R=\mathcal O(W)
 =\operatorname{Sym}_{\mathrm{super}}
       \left(\bigoplus_{m=1}^{q}V_0^{\otimes m}\otimes L_0^{-1}\right).
\]
Grade $R$ by total $V_0$-degree $r$ and number $d$ of form factors.
Every $R_{r,d}$ is finite and is a direct summand of a finite sum
of copies of $V_0^{\otimes r}\otimes L_0^{-d}$.
If $I\subset R$ is a $G$-stable ideal and $E\subset I$ is a finite
homogeneous $G$-submodule generating $I$, then
\begin{equation}
 \Hom_G\bigl(V_0^{\otimes r}L_0^{-d},(E\otimes R)_{r,d}\bigr)
 \longrightarrow
 \Hom_G\bigl(V_0^{\otimes r}L_0^{-d},I_{r,d}\bigr)
 \label{eq:c-polynomial-lift}
\end{equation}
is surjective.  The same conclusion holds after a finite direct
sum in the source.  All the splittings and lifts can be found by
rational linear algebra.
\end{lemma}

\begin{proof}
The degree-$r$ polynomial representations of
$\operatorname{GL}(u|v)$ are modules for its Schur superalgebra.
The coefficient-coalgebra description and the explicit
coefficient-basis proof in
\citet[section~5, (5.2), Lemma~5.1 and Theorem~5.2,
author pp.~13--15]{c_brundankujawa2003}
identify that algebra, over $\Q$, with
\[
 S(u|v,r)=\End_{\Q[\mathfrak S_r]}(V_0^{\otimes r}),
\]
where the symmetric group acts by superpermutations.
Maschke's theorem makes this a semisimple algebra: decompose the
finite $\Q[\mathfrak S_r]$-module into its simple isotypical
parts and take the product of their matrix endomorphism algebras.
This is also the characteristic-zero argument in
\citet[Theorem~1.3.1(a)(i), proof of Theorem~4.2.6(a),
author pp.~3 and~12]{c_hemmerkujawanakano2006}.
An ungraded splitting can be made even by replacing a map $f$
by $(f+\Pi f\Pi)/2$, where $\Pi$ denotes intrinsic parity on the
source and target.  Conjugation by $\Pi$ preserves the module
intertwiner equations, so this averaging still gives a splitting.
Thus the category with even morphisms is semisimple as well.
Tensoring by a fixed character of the independent last torus does
not change this conclusion.

The degree-$d$ supersymmetric power is the image of the averaging
idempotent on the $d$-fold tensor power of its generators.  After
fixing total degree $r$, the latter is a finite sum of copies of
$V_0^{\otimes r}L_0^{-d}$.  Finiteness follows also directly from
$d\leq r\leq qd$, except for $R_{0,0}=\Q$.
The two central scalar actions split every stable ideal into
these homogeneous components.  The multiplication surjection
$(E\otimes R)_{r,d}\twoheadrightarrow I_{r,d}$ is therefore a
surjection of finite polynomial representations of the same
degree and Tate character.  Its source, target and the source of
the displayed Hom functor belong to the semisimple category just
proved.  It splits, proving the assertion.  Intertwining and
splitting impose finite rational linear equations.  We have used
exactness in the covariant polynomial category, not exactness of
invariants on arbitrary mixed rational superrepresentations.
\end{proof}

\begin{theorem}[A finite criterion without a sign hypothesis]
\label{c:finite-unconditional-criterion}
\status{Proved; finite existence is unconditional, and effectiveness has the stated input condition}
Let $X/k$ be smooth projective and connected, of dimension $n>0$
over an algebraically closed field, let $H$ be a Weil cohomology
with coefficient field $K$, and let
$o\in\CH^n(X)_\Q$ have degree one.  Set
\[
 a=\dim_K H^{\mathrm{even}}(X),\qquad
 b=\dim_K H^{\mathrm{odd}}(X),\qquad
 N=a(b+1),\qquad \lambda=((b+1)^a).
\]
In Chow motives put $A=h(X)$, $L=\Q(-n)$,
\[
 \epsilon(x)=\operatorname{tr}(o\cdot x),\qquad
 \pi=\id_A-\eta\epsilon,\qquad V=(A,\pi),
 \qquad
 \kappa_m=\operatorname{tr}\circ\mu^{(m)}\circ\pi^{\otimes m}
       :V^{\otimes m}\longrightarrow L.
\]
Here $\eta$ is the unit and $\mu^{(m)}$ is iterated cup product;
the displayed notation for $\epsilon$ includes its usual Tate
cancellation.  Since $\epsilon\eta=1$, $\pi$ is an actual Chow
projector, independently of any standard conjecture.  There is
a finite list of rational expressions $P_1,\ldots,P_s$ in the
forms $\kappa_1,\ldots,\kappa_{2n}$, using tensor products,
permutations and rational linear combinations, such that the
following are equivalent:
\begin{enumerate}[label=(\Alph*)]
\item All rational cohomological relations among the specified
      point and partial-diagonal generators on every power of $X$
      hold in rational Chow.
\item The following finite list of Chow identities holds:
      \begin{enumerate}[label=(\roman*)]
      \item $c_n(T_X)=(a-b)o$ and $\Delta_*o=o\boxtimes o$;
      \item $\kappa_{2n+1}=0$;
      \item the rectangular Young symmetrizer
            $y_\lambda:V^{\otimes N}\to V^{\otimes N}$ is zero;
      \item $P_\nu=0$ for $1\leq\nu\leq s$.
      \end{enumerate}
\end{enumerate}
Each identity in (ii)--(iv), after Poincar\'e-duality bending,
is a finite rational combination of cycles in the requested
point-and-diagonal rings.  The list can be constructed effectively
from a rational cohomology table, and more generally from a finite
presentation of the $\Q$-algebra generated by its structure
constants.  With an arbitrary unspecified Weil coefficient
field, finite existence is asserted, without an algorithm for
testing its constants.  No decidability assertion for rational
equivalence is made.  No algebraic parity projector is assumed.
More explicitly, the bent form $\kappa_m$ is the modified diagonal
\begin{equation}
 \Gamma_m(X,o)=
 \sum_{\varnothing\ne J\subseteq\{1,\ldots,m\}}
       (-1)^{m-|J|}\Delta_J\prod_{i\notin J}p_i^*o
       \ \in\CH^{n(m-1)}(X^m)_\Q,
 \label{eq:c-explicit-higher-modified-diagonal}
\end{equation}
where a singleton diagonal means the unit on that coordinate.
Thus condition (ii) is the explicit identity
$\Gamma_{2n+1}(X,o)=0$.
\end{theorem}

\begin{proof}
We first normalize the Frobenius algebra.  Poincar\'e duality
identifies the equality $\Delta_*o=o\boxtimes o$ with
\[
 \epsilon\mu=\epsilon\otimes\epsilon.
\]
Indeed the two sides pair with $x\otimes y$ as, respectively,
$\operatorname{tr}(oxy)$ and
$\operatorname{tr}(ox)\operatorname{tr}(oy)$.  This calculation
is an identity of Chow correspondences, since the pairing and
its inverse are the diagonal evaluation and coevaluation.
Consequently $V=\ker\epsilon$ is an algebra ideal, with a
restricted multiplication $\mu_V$.  For $m\geq1$ we have
\begin{equation}
 \kappa_{m+1}
   =\kappa_m\circ(\mu_V\otimes\id_V^{\otimes(m-1)}).
 \label{eq:c-high-form-recursion}
\end{equation}
Thus (ii) kills all $\kappa_m$ with $m>2n$ in Chow.  In
cohomology $H(V)=H^{>0}(X)$, so these same forms vanish by
degree for every $m>2n$.  This explains why a single additional
modified-diagonal identity suffices to truncate the forms.

Here is the finite construction.  Let $V_0=\Q^{a-1|b}$ and let
$L_0$ be a formal even line.  We do not assert that $V_0$ is a
rational form of the given cohomology algebra.  Apply
Lemma~\ref{c:covariant-polynomial-lifting} with $q=2n$ and write
$G=\operatorname{GL}(V_0)\times\mathbb G_m$.
The cohomological forms give an even $K$-point
$t=(\kappa_1,\ldots,\kappa_{2n})$ of $W$.  Define
\begin{equation}
 I_\Q=\ker\left(R\longrightarrow\mathcal O(G)\otimes_\Q K\right),
 \qquad f\longmapsto\bigl(g\longmapsto f(g\cdot t)\bigr).
 \label{eq:c-pure-form-orbit}
\end{equation}
This is the kernel of a rational comodule map, with the regular
action on $\mathcal O(G)$ and trivial action on $K$.  It is
therefore a $G$-stable ideal.  We use the schematic super-orbit:
odd directions and nilpotents are retained.  Taking only the
ordinary-point orbit would be a different and insufficient
construction.  The two central scalar actions make $I_\Q$
homogeneous in $(r,d)$.

The supercommutative algebra $R$ is a finite module over the
polynomial algebra on its even generators, and hence is
noetherian.  Choose finite ideal generators of $I_\Q$ and close
their homogeneous span under $G$.  This yields a finite
homogeneous stable submodule $E\subset I_\Q$ that still generates
the ideal.  Each $E_{r,d}$ lies in the finite semisimple module
$R_{r,d}$ and hence is a split summand of a finite sum of copies
of $V_0^{\otimes r}L_0^{-d}$.  Choose a finite such sum $T$ and
maps
\[
 s:E\longrightarrow T,\qquad p:T\longrightarrow E,
 \qquad ps=\id_E.
\]
For each tensor summand $T_\nu=V_0^{\otimes r_\nu}L_0^{-d_\nu}$,
the composite
\begin{equation}
 T_\nu\longrightarrow T\xrightarrow{p}E\hookrightarrow R
 \label{eq:c-finite-form-relations}
\end{equation}
is a polynomial multilinear form expression.  After evaluation
at $t$ and cancellation of the indicated line twist it is a map
$V_0^{\otimes r_\nu}\otimes K\to L_0^{d_\nu}\otimes K$.
Its Chow expression is $P_\nu$.

We justify both the expression and every equality among such
expressions.  In each bidegree the symmetric-power averaging
idempotents split $R_{r,d}$ into a finite sum of pure tensor
powers with its fixed line character.  The first fundamental
theorem in \citet[Theorem~2.1]{c_lehrerzhang2021super} says that
the intertwiners between equal pure powers are generated by
permutations.  Hence \eqref{eq:c-finite-form-relations} is a rational combination of products
of the universal forms, evaluated on disjoint groups of inputs.
The second fundamental theorem, their Theorem~2.3, says that all
relations between those permutation maps are generated by the
Young symmetrizer with $(a-1)+1=a$ rows and $b+1$ columns.
The central scalar excludes unequal pure degrees.  Inserting
the averaging idempotents merely gives direct summands of this
same description.  These statements descend from $\C$ to $\Q$
because the maps and their kernel equations are rational.
The fixed Tate line is even and invertible, and its usual
one-dimensional identities already hold in Chow motives.
Thus (iii) makes every representation-theoretic equality used
in this construction a valid Chow tensor identity.

Now let a rational universal form
\[
 F:V_0^{\otimes r}L_0^{-d}\longrightarrow R
\]
be equivariant and satisfy $F(t)=0$.  Equivariance on all
superalgebra-valued points gives $F(g\cdot t)=gF(t)=0$;
equivalently its image is contained in $I_\Q$.  Conversely,
membership in $I_\Q$ implies zero evaluation at the identity.
The multiplication map $E\otimes R\twoheadrightarrow I_\Q$
and Lemma~\ref{c:covariant-polynomial-lifting} give a lift
\[
 \widetilde F:V_0^{\otimes r}L_0^{-d}
                 \longrightarrow(E\otimes R)_{r,d}.
\]
Compose with $s\otimes\id$ to obtain a lift with target
$(T\otimes R)_{r,d}$.  Expanding this lift with the pure-tensor
first fundamental theorem expresses $F$ as a sum of expressions
that apply one of the maps \eqref{eq:c-finite-form-relations} to some inputs and other
universal forms to the remaining inputs.  All intervening
objects are covariant polynomial representations of total
$V_0$-degree $r$ and line character $-d$; no exactness assertion
for an arbitrary mixed invariant functor has entered.
The equality is valid as a formal tensor identity modulo the
rectangular relation by the preceding second fundamental
theorem.  Conditions (iii)--(iv) therefore imply
$F(\kappa_1,\ldots,\kappa_{2n})=0$ in Chow.  There is no
projective auxiliary tensor to cancel and no use of $C^+(X)$.

We next return from forms to the requested cycles.  By (i) and
the graph product rule, every monomial in the generators reduces
to a scalar times a pointed-partition cycle.  A homogeneous
relation between these cycles, with $k$ coordinates and $d$
unpointed blocks, bends to a map
$A^{\otimes k}\to L^d$.  For an unpointed block $B$ its factor
is $\operatorname{tr}(\prod_{i\in B}x_i)$; for a pinned
coordinate it is $\epsilon(x_i)$.  Split every input using the
actual decomposition $A=\mathbf1\oplus V$.  A pinned $V$ input
annihilates its term.  Unit inputs in an unpointed block are
deleted; if no input remains, its value is
$\operatorname{tr}(1)=0$, since $n>0$.  Every other block is one
of the forms $\kappa_m$.  Condition (ii) removes blocks with
$m>2n$.  Each remaining restriction is exactly a rational form
$F$ of the preceding paragraph.  If the original relation is
cohomologically zero, each such restriction vanishes at $t$,
and hence vanishes in Chow.  The direct-sum decomposition of
$A^{\otimes k}$ and Poincar\'e duality prove the original cycle
relation.  This proves (B)$\Rightarrow$(A).

For necessity, the cohomological relations
$\Delta_*o-o\boxtimes o=0$ and
$\Delta_X^2-(a-b)o\boxtimes o=0$ belong to the requested ring
on $X^2$.  Pushing the second along a projection and using
self-intersection gives $c_n(T_X)=(a-b)o$.  The form
$\kappa_{2n+1}$ is cohomologically zero by positive degrees.
Its bent Chow cycle is obtained from the small diagonal by
inserting $\pi$ in each input.  Expanding
$\pi=\id-\eta\epsilon$ gives pointed partial-diagonal cycles.
The same expansion shows that $y_\lambda(V)$, viewed as a
correspondence on $X^N$ and hence a cycle on $X^{2N}$, is a
rational combination of products of pair diagonals and point
classes.  It is cohomologically zero by the rectangular theorem.
Finally every $P_\nu$ is a product of the forms just described,
with permuted inputs, and is cohomologically zero because
$E\subset I_\Q$.  Bending gives the claimed cycles in the
requested rings.  Thus (A) implies (ii)--(iv) as well as (i).
For the displayed formula for $\Gamma_m$, the adjoint of
$\pi=1-\eta\epsilon$ under Poincar\'e duality is
$1-o\operatorname{tr}$.  Expanding its action on the small
diagonal gives the terms indexed by nonempty $J$ in
\eqref{eq:c-explicit-higher-modified-diagonal}.  The omitted empty
term integrates every coordinate and is
$\operatorname{tr}(1)=0$.

For effectiveness, the orbit map is computed using two even
matrix blocks, odd off-diagonal entries, inverses of the two
even determinants, and the last torus coordinate and its inverse.
If the coordinates of $t$ generate a presented algebra
$B=\Q[c_1,\ldots,c_h]$, compute the kernel into
$\mathcal O(G)\otimes B$ by orbit-graph elimination, including
the defining equations of $B$.  This is finite-module
polynomial elimination over the even variables: the odd
variables have a finite exterior basis, with known signed
multiplication.  Noetherian module Gr\"obner elimination
terminates.  The kernel agrees with \eqref{eq:c-pure-form-orbit}
because $B\hookrightarrow K$.  Stable spans, the splittings,
permutation expressions and all lifts then use finite rational
linear algebra.  Without such a presentation only the
noetherian existence argument is used.  This completes the proof.
\end{proof}

\begin{corollary}[A finite power detects all the relations]
\label{c:finite-power-cutoff}
\status{Proved; the cut-off is effective with the same coefficient input}
In Theorem~\ref{c:finite-unconditional-criterion}, let $r_\nu$
be the input arity of $P_\nu$ and put
\[
 B=\max\{2,\ 2n+1,\ 2a(b+1),\ r_1,\ldots,r_s\}.
\]
Then all rational cohomological relations hold in Chow on every
power of $X$ if and only if the cycle map is injective on
$T_B^{\CH}(X,o)$.  In particular such a finite cut-off exists
without $C^+(X)$ or a rational form of the Weil cohomology.
Its effective construction has exactly the coefficient-data
qualification in the theorem.
\end{corollary}

\begin{proof}
All finite tests in the theorem are bent cycles on powers of
arity at most $B$.  The two normalization tests are detected by
the indicated relations on $X^2$, the high form by $X^{2n+1}$,
the rectangle by $X^{2N}$, and $P_\nu$ by $X^{r_\nu}$.
Injectivity on $T_B^{\CH}$ implies injectivity on every
$T_k^{\CH}$ with $k\leq B$: pull back along $p:X^B\to X^k$,
and retract that pullback by
\[
 z\longmapsto p_*\left(z\prod_{i=k+1}^{B}p_i^*o\right).
\]
The composite on $\CH^*(X^k)_\Q$ is the identity because
$\deg o=1$.  Thus injectivity on the single specified power
implies all the finite tests, and the theorem applies.  The
converse is immediate.  The cut-off depends only on the
cohomological form data and the choices in its finite
construction, not on an unproved Chow vanishing.
\end{proof}

\begin{lemma}[A coefficient-independent recursive elimination bound]
\label{c:recursive-elimination-bound}
\status{Proved}
For integers $v,d\geq1$, let $\mathcal B(v,d)$ be the tree of
finite sequences $(\alpha_0,\ldots,\alpha_{r-1})$ in $\Z_{\geq0}^v$
satisfying
\[
 |\alpha_i|\leq d2^i,
 \qquad \alpha_j\not\leq\alpha_i\quad(0\leq j<i<r),
\]
where $\leq$ is coordinatewise.  This is a finite tree.  Its
maximum sequence length $\mathcal L(v,d)$ is a total computable
integer-valued function.  Put
\[
 \mathcal D(v,d)=d\,2^{\mathcal L(v+1,d)}.
\]
Over any field, an ideal in $v$ ordinary polynomial variables
generated in degrees at most $d$ has, for every monomial order,
a Gr\"obner basis in degrees at most $\mathcal D(v,d)$.
In particular this bound applies to an elimination order.
\end{lemma}

\begin{proof}
We give the elementary argument, so no closed Gr\"obner-degree
estimate is needed.  Dickson's assertion that an infinite
sequence in $\Z_{\geq0}^v$ contains a comparable pair follows
by induction on $v$.  An infinite sequence of nonnegative
integers has an infinite nondecreasing subsequence: either some
value repeats infinitely or one successively selects larger
values.  Apply this observation to the first coordinate and
then apply induction to the remaining coordinates.  Thus
$\mathcal B(v,d)$ has no infinite branch.  Each node has only
finitely many children, by its displayed degree bound.
If the tree were infinite, repeatedly choosing a child with
infinitely many descendants would construct an infinite branch.
It is consequently finite.  Enumerating its successive finite
levels until the next level is empty computes $\mathcal L(v,d)$.

Homogenize the given generators with one new variable $h$.
Extend the desired monomial order by comparing total degree
first and, at equal degree, comparing the monomials after
$h=1$ in the desired order.  Process the homogeneous generators
one at a time, reducing each by the previous nonzero remainders,
and then run the usual pair-reduction algorithm.  Every accepted
nonzero remainder has a leading monomial not divisible by any
previously accepted leading monomial.  Homogeneous division
does not increase degree, and the least common multiple of
two leading monomials has degree at most the sum of their
degrees.  If the accepted polynomials are numbered from zero,
their degrees are therefore at most $d2^i$.  Their leading
exponent vectors form a sequence in $\mathcal B(v+1,d)$.
At most $\mathcal L(v+1,d)$ polynomials can be accepted, and
all their degrees are at most $\mathcal D(v,d)$.
There are then only finitely many pairs to process.

For completeness, the pair criterion here says that if every
pair polynomial reduces to zero, every element of the generated
ideal has leading monomial divisible by one of the recorded
leading monomials.  Indeed write such an element as a sum of
monomial multiples of the recorded polynomials.  If the greatest
displayed leading monomial cancels, subtract the corresponding
pair relations to lower that greatest monomial.  A representation
minimizing it, and then the number of terms attaining it, cannot
have such a cancellation.  Its greatest monomial is therefore
the leading monomial of the element.  This proves the criterion
and that the resulting list is a Gr\"obner basis.

Dehomogenization gives a Gr\"obner basis for the original ideal.
To check this even though the homogenized generator ideal need
not be saturated, if $f$ belongs to the original ideal, then
$h^e f^{\mathrm h}$ belongs to the homogenized generator ideal
for some $e\geq0$, by homogenizing an expression for $f$ in
the generators to a common degree.  A homogeneous leading
monomial of our basis divides that of $h^e f^{\mathrm h}$.
After setting $h=1$, it divides the leading monomial of $f$
in the desired order.  Degrees do not increase.  This proves
the bound for arbitrary, including elimination, orders.
\end{proof}

\begin{corollary}[A uniform arity bound for rational cohomology tables]
\label{c:uniform-arity-bound}
\status{Proved in the stated rational-table scope}
Suppose the graded Weil cohomology algebra in
Theorem~\ref{c:finite-unconditional-criterion} has a rational
table.  This includes rational Betti cohomology over $\C$.
Put $q=2n$, $u=a-1$, $v=b$, $s=u+v$, and
\[
 M_0=\frac12\sum_{m=1}^{q}\bigl(s^m+(u-v)^m\bigr),\qquad
 M_1=\frac12\sum_{m=1}^{q}\bigl(s^m-(u-v)^m\bigr).
\]
Define the following integers using the function of the preceding
lemma:
\[
 O=M_1+4uv,\qquad
 E=M_0+2(u^2+v^2)+2,\qquad
 J=2^O,\qquad P=E+J,
\]
\[
 D=\mathcal D(P,q+2),\qquad
 B_{n,a,b}=\max\{2,\ q+1,\ 2a(b+1),\ q(D-1+M_1)\}.
\]
Then all the point-and-diagonal cohomological relations hold in
rational Chow on all powers if and only if the cycle map is
injective on $T_{B_{n,a,b}}^{\CH}(X,o)$.
The bound depends only on $(n,a,b)$, and is a total computable
integer.  It is not a practical or minimal bound, and does not
decide the required Chow injectivity.
\end{corollary}

\begin{proof}
We bound the input arities of the finite form relations in
Theorem~\ref{c:finite-unconditional-criterion}.  Its parameter
superspace has $M_0$ even and $M_1$ odd coordinates.
Present $\operatorname{GL}(u|v)$ by two independent even
supermatrices $G,H$ with relations $GH=HG=1$; $H$ is the inverse
of $G$.  This presentation has $2(u^2+v^2)$ even and $4uv$ odd
matrix entries.  Present the independent line torus by $z,w$
with $zw=1$.  The orbit-graph equation for a form of arity $m$
is a coordinate of the parameter tensor minus a coordinate of
$z(H^\vee)^{\otimes m}t_m$, with the fixed dual-action signs.
It has total degree at most $m+1\leq q+1$; all group equations
have degree two.  Because the table is rational, its structure
constants add no variables or degrees.

There are $E$ even and $O$ odd variables in this graph
presentation.  The free supercommutative algebra is a free module
of rank $J=2^O$ over the polynomial ring $A_0$ on its even
variables, with exterior-monomial basis $e_1,\ldots,e_J$.
If $f_1,\ldots,f_t$ are the graph equations, their ideal is the
$A_0$-submodule generated by all $e_j f_i$.  The polynomial
coordinates of these module generators have even-variable
degree at most $q+1$.  Encode a vector $(p_1,\ldots,p_J)$ by
the ordinary polynomial $\sum p_jY_j$, and form the ideal in
$A_0[Y_1,\ldots,Y_J]$ generated by the encoded module generators
and all $Y_iY_j$.  Its part of total $Y$-degree one is exactly
the original module.  This ordinary ideal has $P=E+J$ variables
and generators of degree at most $q+2$.

Choose an elimination order that eliminates the even group
variables and every $Y_j$ whose exterior monomial involves an
odd group variable.  Retain the even parameter variables and
the remaining $Y_j$.  The ordinary ideal is homogeneous in
total $Y$-degree.  The homogeneous pair-reduction construction in
Lemma~\ref{c:recursive-elimination-bound} preserves this additional
grading, both before and after dehomogenization, and therefore
gives a bounded Gr\"obner basis homogeneous in $Y$-degree.  The elimination basis elements of
$Y$-degree one give module generators for the intersection with
the original parameter superalgebra, namely the schematic
super-orbit ideal.  Quadratic $Y$ terms cannot contribute to
this degree-one part.

Lemma~\ref{c:recursive-elimination-bound} bounds these ordinary
basis elements by $D$.  Their coefficients in the retained even
variables have degree at most $D-1$.  Returning to exterior
monomials adds at most $M_1$ odd parameter factors.  Every
parameter coordinate has input arity at most $q$.  The orbit
ideal therefore has generators of input arity at most
$q(D-1+M_1)$.  Taking homogeneous components and closing their
span under the supergroup preserves input arity.  The split
pure tensor modules used to define $P_\nu$ can consequently
all be taken within this same bound.  The assertion now follows
from Corollary~\ref{c:finite-power-cutoff}.
\end{proof}

\begin{remark}[The coefficient qualification in the uniform bound]
\label{c:uniform-arity-scope}
\status{Proved scope clarification}
The finite criterion and finite cut-off of
Theorem~\ref{c:finite-unconditional-criterion} and
Corollary~\ref{c:finite-power-cutoff} apply to an arbitrary
Weil field.  The displayed bound depending only on $(n,a,b)$
has the additional rational-table hypothesis.  For a presented
coefficient algebra whose $h$ named generators include every
coordinate of $t=(\kappa_1,\ldots,\kappa_q)$, with its exact
defining ideal generated in degrees at most $e$, the same proof gives a bound by adding
$h$ to $E$ and replacing the generator-degree input $q+2$
by $\max(q+3,e+1)$.  The extra one accounts for a structure
constant occurring as a variable in an orbit equation.  These
$h$ coefficient variables are eliminated together with the group
variables, so the result is the rational contraction $I_\Q$.
We make no uniform rational-contraction claim for unspecified
elements of an arbitrary coefficient field.

The accompanying exact certificate exhausts only
$\mathcal L(1,d)=d+1$ for $1\leq d\leq4$ and
$\mathcal L(2,1)=4$, and checks the square-zero elimination
mechanism for the explicit super map
$z=x$, $\theta_1=\eta$, $\theta_2=x\eta$.
It does not evaluate the enormous geometric cut-off above.
No closed Dub\'e degree formula is used in its proof.
\end{remark}


\begin{remark}[What the finite criterion settles]
\label{c:finite-unconditional-scope}
\status{Proved scope statement; the actual general Chow vanishing remains unresolved}
Theorem~\ref{c:finite-unconditional-criterion} removes the sign
hypothesis from the finite criterion, including for an arbitrary
Weil coefficient field.  It supplies finite generators for the
all-power obstruction and does not assert that these cycles
vanish for every variety.  It also does not make rational
equivalence decidable.  Its list is not claimed minimal; in
particular the more precise K3 obstruction identified above is
unchanged.  The explicitly requested all-power K3 assertion and
the unrestricted evaluation of the criterion for every $(X,o)$
therefore remain separate unresolved parts of C.1(c).
If $n=0$, connectedness over the algebraically closed field gives
$X=\operatorname{Spec}(k)$, and every relation holds immediately.
\end{remark}


\begin{theorem}[The narrow K3 ring already detects finite-dimensionality]
\label{c:k3-gate}
\status{Proved equivalence; the general vanishing is unresolved}
Let $S/\C$ be a projective K3 surface and let $o$ be its Beauville--Voisin class. The following assertions are equivalent:
\begin{enumerate}[label=(\roman*)]
\item for every $k$, all the cohomological relations among the point and diagonal classes hold in $\CH^*(S^k)_\Q$;
\item with $\pi_2=\Delta_S-o\times S-S\times o$, the cycle
\begin{equation}
\label{c:k3-determinant}
 D_{23}=\sum_{\sigma\in S_{23}}\sgn(\sigma)
          \prod_{a=1}^{23}(\pi_2)_{a,23+\sigma(a)}
 \quad\text{vanishes in }\CH^{46}(S^{46})_\Q;
\end{equation}
\item $\bigwedge^{23}h^2(S)=0$ in rational Chow motives;
\item $h(S)$ is finite-dimensional in the sense of Kimura--O'Sullivan.
\end{enumerate}
Thus the request for every K3 surface is not supplied by the Beauville--Voisin small-diagonal theorem. The first missing universal determinant for this narrower ring occurs on $46$ factors.
\end{theorem}

\begin{proof}
The projectors $o\times S$, $\pi_2$, and $S\times o$ are mutually orthogonal Chow projectors, giving $h(S)=\mathbf1\oplus h^2(S)\oplus\Q(-2)$ (with the contravariant/cohomological convention). The cycle $\pi_2$ belongs to the requested ring. As a correspondence on $S^{23}$, $D_{23}$ is $23!$ times the alternating projector on $h^2(S)^{\otimes23}$. Consequently its vanishing is exactly (iii), because an object in a pseudo-abelian category is zero precisely when its defining idempotent is zero. Its cohomology class is zero since $\dim H^2(S)=22$, so (i) implies (ii). Assertion (iii) makes $h^2(S)$ evenly finite-dimensional, and the other two summands are Tate motives; hence it implies (iv). Conversely, finite-dimensionality makes the even summand have exterior dimension $22$, giving (iii); alternatively the dimension assertion is the finite-dimensionality statement recalled in \citet[\S2.4, p.~506]{c_yin2015}. Finally, \citet[unnumbered main Theorem, p.~503, and \S\S2.5, 3.1--3.6]{c_yin2015} proves that finite-dimensionality is equivalent to injectivity for the larger ring generated by all divisor classes and diagonals. Its restriction gives (i).

We can also see sufficiency directly in the smaller ring. \citet[Theorem 1 and Proposition 3.2]{c_beauvillevoisin2004} gives $c_2(T_S)=24o$ and the small-diagonal relation. These imply every local relation in \eqref{c:pair-local} with $b=22$. The only additional generators supplied by Theorem~\ref{c:pair-ranks} are the permutations and products of \eqref{c:k3-determinant}. This proves that no larger unlisted family of relations is hidden in the equivalence.
\end{proof}

The statement just proved is a reduction, not a proof of (i) for general K3 surfaces. The current-source check includes \citet[Theorem L, explicitly Theorem 9.7 and Corollary 9.8]{c_floccarifu2026}. That result proves abelian Chow motives for K3 surfaces whose transcendental quadratic space embeds in $U_\Q^{\oplus2}\oplus\langle-1\rangle\oplus\langle-b\rangle$, with $b>0$ and $b\equiv3\pmod4$; it supplies new four-dimensional families of generic Picard rank $16$. It does not assert the vanishing of \eqref{c:k3-determinant} for every projective K3 surface. The general classification of all pairs $(X,o)$ requested in C.1(c) also remains a separate open row; the finite presentation above gives its exact cycle conditions but does not turn them into a geometric classification.

\begin{corollary}[Removing the algebraic part of the K3 gate]
\label{c:k3-transcendental-gate}
\status{Proved smaller equivalent vanishing test; general vanishing unresolved}
Let $S/\C$ be a projective K3 surface of Picard rank $\rho$, let
$D_1,\ldots,D_\rho$ be a rational basis of $\operatorname{NS}(S)_\Q$,
and let $B=(\deg D_aD_b)_{a,b}$.  Define
\[
 \pi_{\rm alg}=\sum_{a,b}(B^{-1})_{ab}D_a\times D_b,
 \qquad \pi_{\rm tr}=\pi_2-\pi_{\rm alg},\qquad
 t_2(S)=(S,\pi_{\rm tr},0).
\]
Every assertion of Theorem~\ref{c:k3-gate} is also equivalent to
\[
 \bigwedge^{23-\rho}t_2(S)=0.
\]
Equivalently, the alternating sum of the products
$(\pi_{\rm tr})_{a,\,23-\rho+\sigma(a)}$, over
$\sigma\in S_{23-\rho}$, vanishes on $S^{2(23-\rho)}$.
For Picard rank $1$ this is a test on $44$ factors.  This smaller
correspondence uses divisor classes outside the narrow requested
ring; its factor count therefore does not contradict the unconditional
narrow-ring result through $45$ factors.
\end{corollary}

\begin{proof}
The intersection matrix $B$ is nonsingular by the Hodge index
theorem.  Direct composition of the product correspondences shows
that $\pi_{\rm alg}$ is an idempotent with image $\Q(-1)^{\oplus\rho}$,
and that it is orthogonal to the degree-zero and degree-four
projectors.  Thus
\[
 h^2(S)=\Q(-1)^{\oplus\rho}\oplus t_2(S).
\]
For a direct sum in a rational pseudo-abelian symmetric category,
the exterior-power decomposition gives
\[
 \bigwedge^{23}h^2(S)
 =\bigoplus_{a=0}^{\rho}
   \bigwedge^a\bigl(\Q(-1)^{\oplus\rho}\bigr)
       \otimes\bigwedge^{23-a}t_2(S).
\]
The $a=\rho$ summand is
$\Q(-\rho)\otimes\bigwedge^{23-\rho}t_2(S)$.
Its vanishing follows from that of the whole direct sum; tensoring
with the inverse Tate object proves one implication.
Conversely, every higher exterior power is a direct summand of
the tensor product of the $(23-\rho)$th exterior power with a
further tensor power.  Hence the stated vanishing kills every
summand in the display.  The correspondence form is exactly the
alternating-projector description.  This also checks directly the
transcendental reduction in \citet[\S2.5, pp.~506--507, equation~(7)]{c_yin2015}
while explaining its relation to the smaller ring in the question.
\end{proof}


\begin{corollary}[The unconditional K3 range]
\label{c:k3-stable-range}
\status{Proved}
For every projective complex K3 surface with its Beauville--Voisin class, the requested relations hold in every power $S^k$ with $k\leq45$. More generally they hold in normalized degree $d$ of every $S^k$ whenever
\[
 \min(d,k-d)\leq22,
\]
where normalized degree $d$ means Chow codimension $2d$. In this range no finite-dimensionality assumption is required.
\end{corollary}

\begin{proof}
The Beauville--Voisin identities give the local relations, so every Chow monomial reduces to a linear combination of disjoint matching monomials
$\prod_{i\in I}o_i\prod_{\{a,b\}\in\mathcal M}q_{ab}$.
If the matching has $m$ edges, degree $d$ means $|I|+m=d$, and the number of used factors is $|I|+2m$. Consequently $m\leq\min(d,k-d)$. For $m\leq22$, the orthogonal matching tensors in $(H^2(S))^{\otimes2m}$ are linearly independent: the first determinantal relation requires $23$ edges, as proved in Theorem~\ref{c:pair-ranks}. Monomials with different point and matching supports lie in different K\"unneth multidegrees. Thus their cohomological classes are independent, and no homologically zero Chow combination remains. For $k\leq45$, every matching has at most $22$ edges. This proves both assertions.
\end{proof}

The remaining relation cannot be obtained merely by asserting that all lower-power relations hold. It first appears in the middle normalized degree of $S^{46}$ and is exactly \eqref{c:k3-determinant}. We also tested three possible ways to remove this obstruction. A Kuga--Satake correspondence giving a splitting in homological motives would produce a composite $1+\eta$ in Chow, with $\eta$ homologically trivial; promoting it to a Chow splitting still requires nilpotence or another argument making $1+\eta$ invertible. Specializing to dense loci with finite-dimensional motives does not prove the generic Chow vanishing, because the rational-equivalence locus is only controlled as a countable union of closed loci without a uniform family of equivalences. Finally, the known small-diagonal identities are precisely the local matching reductions and do not supply the first orthogonal determinant. These are failed proof steps, not counterexamples to finite-dimensionality.

The further nilpotence test has a precise limitation.  Put
$d=22-\rho(S)$ and let $a_{d+1}$ be the normalized alternating
idempotent on $t_2(S)^{\otimes(d+1)}$.  Since $a_{d+1}^2=a_{d+1}$,
composition-nilpotence of this particular correspondence is equivalent
to its vanishing.  Also, the symmetric-category identity
\[
 \operatorname{ptr}_{m}(a_m)=\frac{d-m+1}{m}a_{m-1}
\]
shows that its partial trace is zero at $m=d+1$; this operation does
not produce a smaller nonzero obstruction.  In the formula, permutations
fixing the last factor contribute $d/m$, and deleting that factor from
a nontrivial cycle contributes $-(m-1)/m$.  Nilpotence only in
$\End(h(S))$ has not been shown to propagate to this higher tensor
power.  The all-powers quantifier in
\citet[Theorem~0.2, author p.~2]{c_jannsen2006finite} cannot be dropped.

There is a valid conditional route in the stated K3 families of
\citet[Theorems~3.5--3.6, with Proposition~1.3]{c_voisin2014bloch}:
her relative argument supplies the needed Chow lifting once the general
Kuga--Satake correspondence is algebraic.  This does not prove that
remaining algebraicity premise.  Nor does the codimension-two argument
apply directly to the higher-codimension alternating projector.
The detailed method audit records these distinctions; none changes the
unresolved vanishing status.

The current-source search also checked the official withdrawal notice of the October 2026 paper claiming the Hodge conjecture for all K3 products \citep{c_k3withdrawal2026}. We exclude that withdrawn proof. Even a valid homological Kuga--Satake splitting alone would leave the Chow nilpotence step just identified.

\begin{example}[An explicit genus-three failure]
\label{c:fermat-counterexample}
\status{Proved counterexample}
Let $C/\C$ be the plane quartic $X^4+Y^4=Z^4$, and take $o=(0:1:1)$. Its modified diagonal $M_o$ vanishes in every Weil cohomology, but
\[
 M_o\ne0\quad\hbox{in }\CH^2(C^3)_\Q.
\]
Thus the cohomological relation in Proposition~\ref{c:modified-diagonal} fails in rational Chow groups for this completely specified genus-three pointed curve.
\end{example}

\begin{proof}
The three partial derivatives $4X^3,4Y^3,-4Z^3$ have no common projective zero, so $C$ is smooth. Adjunction for a smooth plane quartic gives $g=(4-1)(4-2)/2=3$. The displayed point lies on $C$. Every product of three positive-degree cohomology classes on a curve is zero; Proposition~\ref{c:modified-diagonal} therefore gives the cohomological vanishing.

Put $J=\operatorname{Jac}(C)$ and $i:C\to J$, $x\mapsto[x-o]$. We use the precise nontorsion assertion of \citet[Theorem 4.1, printed p.~104]{c_bloch1984}: the Ceresa cycle $\Sigma=i_*[C]-[-1]_*i_*[C]$ for this Fermat quartic has infinite order modulo algebraic equivalence. We checked the detected class in its proof: the projected $5$-adic Abel--Jacobi class is nonzero by Theorem 3.4; a hypothetical torsion class would come from invariants of a $\Q_5/\Z_5$ module on which Frobenius at $11$ acts as $-11$, and $-11-1$ is a $5$-adic unit. Restriction--corestriction prevents any finite extension from killing the nontorsion class. The assertion persists over $\C$: finite rational- or algebraic-equivalence data over $\C$ spread over a finitely generated $\overline\Q$-algebra and specialize to a closed point, giving such data over a finite extension, which would contradict the theorem.

We make the passage to $M_o$ directly in rational Chow groups. For this three-dimensional Jacobian, the decomposition theorem and Proposition 3(a) of \citet[pp.~647--649]{c_beauville1986} give
\[
 \CH^2(J)_\Q=\CH^2(J)_{(0)}\oplus\CH^2(J)_{(1)}\oplus\CH^2(J)_{(2)}.
\]
The odd cycle $\Sigma$ lies entirely in grade $1$. Proposition 1(iv) of the same source therefore gives $[m]_*\Sigma=m^3\Sigma$.

For $s:C^3\to J$, $(x,y,z)\mapsto i(x)+i(y)+i(z)$, direct pushforward of the seven terms gives
\[
 s_*M_o=[3]_*i_*[C]-3[2]_*i_*[C]+3i_*[C].
\]
This is also the identity in \citet[Proposition 5.3]{c_grossschoen1995}. Taking its odd part under $[-1]_*$ gives
\[
 \frac{1-[-1]_*}{2}s_*M_o
 =\frac12\bigl([3]_* -3[2]_*+3\bigr)\Sigma.
\]
The identity just proved in the Chow group makes this odd part exactly
\[
 \tfrac12(3^3-3\cdot2^3+3)\Sigma=3\Sigma\ne0.
\]
Hence $s_*M_o$ and therefore $M_o$ are nonzero in rational Chow groups, as claimed.

The polynomial smoothness certificate, the point substitution, the genus value, and the coefficient $3^3-3\cdot2^3+3=6$ are recorded by \texttt{check\_explicit\_bad\_reduction.py}. The nontorsion input is the opened theorem and its proof, not a claim that this finite polynomial certificate computes a Chow group.
\end{proof}

\subsection{C.2: the special bivariant systems}

The relevant category must retain the two maps of manuscript 11:34. Requiring an ordinary algebra homomorphism to preserve both the unit and the Frobenius counit would change that category and would exclude the geometric maps under discussion.

\begin{definition}[A gauged special bivariant system]
\label{c:bivariant-def}
\status{Definition, with the corrected exchange axiom}
Over a characteristic-zero field $L$, a special bivariant system assigns to every finite set $I$ a finite-dimensional graded-commutative algebra $A(I)$, with $A(\varnothing)=L$, a trace $\kappa_I$ of degree $-D|I|$, and a perfect trace pairing. For a map $\alpha:I\to J$ it assigns
\[
 \alpha^*:A(I)\longrightarrow A(J),\qquad
 \alpha_*:A(J)\longrightarrow A(I).
\]
The first is a unital graded algebra homomorphism; the second has degree $D(|I|-|J|)$. Pullbacks compose covariantly in maps of indexing sets, pushforwards contravariantly, and
\[
 \kappa_I(\alpha_*x\,a)=\kappa_J(x\,\alpha^*a),\qquad
 \alpha_*(x\alpha^*a)=\alpha_*x\,a.
\]
For disjoint unions the exterior-product maps $A(I)\otimes A(J)\to A(I\amalg J)$ are isomorphisms compatible with associativity, units, traces, and the Koszul symmetry. These maps are compatible with pullback and pushforward. The exchange axiom is base change for product squares (a projection pulled back by a map on the other factor); arbitrary self-intersection squares are not included. The gauge is the chosen trace, not a quotient up to multiplication by a scalar.

A morphism in the bivariant category from $A$ to $B$ is a pair $(f^*,f_*)$, with $f^*:B\to A$ a unital graded algebra map and $f_*:A\to B$ its trace adjoint. On powers these maps are tensor powers and satisfy the corresponding pullback and pushforward compatibilities. The adjoint need not preserve the unit.
\end{definition}

\begin{theorem}[C.2: equivalence and the exchange axiom]
\label{c:bivariant-equivalence}
\status{Proved, with the category specified above}
Special gauged bivariant systems are equivalent to finite-dimensional graded-commutative Frobenius algebras of even top degree, with bivariant adjoint morphisms. The functor taking a system to its one-factor algebra and diagonal coevaluation is fully faithful. For a smooth projective connected variety, $I\mapsto H^*(X^I)$ satisfies these axioms. Correspondence composition defined by these axioms is associative and has the diagonal as identity.
\end{theorem}

\begin{proof}
From a Frobenius algebra $(A,\mu,1,\kappa)$, set $A(I)=A^{\otimes I}$. For $\alpha:I\to J$, let $\alpha^*$ multiply the inputs along each fibre $\alpha^{-1}(j)$, inserting $1$ for an empty fibre, with the Koszul permutation signs. Associativity and commutativity of $\mu$ give functoriality. Define $\alpha_*$ as its unique trace adjoint. Perfectness gives existence and uniqueness; taking adjoints proves pushforward functoriality. Pairing both sides against an arbitrary element proves the projection formula. Product-square base change follows by separating the tensor factors: contraction in the projected factor commutes with the map on the other factor. This also verifies all unit and exterior-product compatibilities.

Conversely, specialness identifies $A(I)$ with $A(\{1\})^{\otimes I}$. A fold of two indices is the product, the empty fibre gives its unit, and the trace gives a perfect pairing. Every map of finite sets is a composite of folds, injections, and bijections. Its pullback is therefore forced by this algebra; its pushforward is forced by adjointness. The coevaluation $\Delta$ is the unique tensor whose contraction is the identity. This proves reconstruction and full faithfulness. A bivariant pair on one factor extends to tensor powers, and the required compatibilities follow respectively from multiplicativity and its adjoint. We do not impose the generally false extra identity $(f^*\otimes f^*)\Delta_B=\Delta_A$.

For Weil cohomology, K\"unneth, Poincar\'e duality, and the projection formula give exactly these maps. A map $\alpha:I\to J$ induces the geometric map $X^J\to X^I$, so all variances agree with the displayed formulas. Product base change is again K\"unneth and trace in the omitted factor. In contrast, pulling a diagonal back through itself produces its Euler self-intersection, so unrestricted Cartesian base change without an excess term is false. Refined Gysin base change may be included on transverse squares, with the transversality hypothesis retained.

Finally, composition of kernels is tensor contraction in the middle factor. Two successive contractions on different factors commute by associativity of tensor product and the Koszul rule. The resulting two contractions are exactly the two parenthesizations of a triple composition. Contraction with the coevaluation is the identity by duality. This proves the associativity requested in 11:48. It also supplies the precise exchange axiom behind ``axiome cart\'esien oubli\'e !'' (11:46), and answers ``il convient de pr\'eciser les axiomes'' (11:42) on the special objects actually used here.
\end{proof}

\begin{remark}[The R7 hypotheses used in this Part]
\label{c:r7-audit}
\status{Used and rechecked, with its internal-Hom counterexample confirmed}
We checked the proof in the supplied reading, \S\S2.1--2.3, against 12:150--158. The conditional construction there is not a construction of a faithful realization on all Chow motives: it first quotients by the chosen homological equivalence, assumes its equality with numerical equivalence and the stated positivity, and then takes the semisimple envelope. Our Chow assertions above do not use this conditional passage. In particular it cannot be used to kill the homologically zero K3 determinant.

The sign correction is forced by the cross-product formula: interchanging a class of degree $a$ with one of degree $b$ multiplies it by $(-1)^{ab}$. Thus the original realization is symmetric into graded vector spaces with that commutativity constraint. This is the convention used in our Frobenius equivalence. A level filtration by thick subcategories yields a direct-sum level decomposition only when the ambient category is semisimple; then we group the simple constituents according to their least allowed twist. Thickness alone gives no splitting of extensions. The geometric level definition of C.4 therefore does not assert a canonical splitting or identify itself with coniveau.

The internal-Hom error of 12:156 already occurs in the semisimple effective Tate category with objects finite sums of $\Q(-a)$, $a\geq0$. For every effective $Z$,
\[
 \Hom(Z\otimes\Q(-1),\Q)=0.
\]
Hence the effective internal Hom from $\Q(-1)$ to $\Q$ is the zero object. After adjoining all Tate twists the corresponding internal Hom is $\Q(1)$, which is nonzero. Thus effective internal Homs need not be preserved by Tate inversion. In a rigid tensor category the valid identity is
\[
 \underline{\Hom}(M,N(-i))
 =M^\vee\otimes N(-i)
 =M^*\otimes N,
 \qquad M^*=M^\vee(-i).
\]
Evaluation and coevaluation give these isomorphisms and their inverse maps. Merely postulating the effective object $M^*$ without the duality identities does not prove them. These checks confirm precisely the restrictions of R7 needed here.
\end{remark}

\subsection{C.3: the algebraic part of the Lie object}

We separate the Lie algebra of an algebraic tannakian group from the analytic Lie algebra of a compact $\ell$-adic image. They are different objects. In a neutral tannakian category over $\Q$, the algebraic part of an adjoint Lie object means its maximal trivial subobject, computed as $\Hom(\mathbf1,\Lie G)$. The same definition is made after neutralizing a gerbe, with its descent action retained.

\begin{proposition}[The invariant-centre identity]
\label{c:lie-centre}
\status{Proved}
For a pro-reductive group $G$ over a characteristic-zero field,
\[
 \Hom_G(\mathbf1,\Lie G)=(\Lie G)^G=\Lie Z(G).
\]
If $G$ is connected, the map from its centre to its abelianization induces an isomorphism on Lie algebras. If $G^\circ$ has protorus abelianization $S$ and $\Gamma=G/G^\circ$, then
\[
 (\Lie G)^G=(\Lie S)^\Gamma.
\]
The invariant condition under $\Gamma$ cannot be discarded.
\end{proposition}

\begin{proof}
An infinitesimal element fixed by every conjugation is exactly an infinitesimal central element: the equations defining these two closed subgroup functors give the same kernel after differentiation. For a connected reductive group, $Z(G)^\circ\times G^{\mathrm{der}}\to G$ is a central isogeny. Its differential is an isomorphism, and the semisimple derived Lie algebra has no adjoint invariants. This proves the connected statement. The connected component acts trivially on its centre; the remaining action is precisely that of the component group. Passing to the inverse limit of finite-dimensional reductive quotients gives the pro-reductive assertions, since both Lie algebras and fixed-vector conditions are computed compatibly in those quotients.
\end{proof}

\begin{theorem}[Lead C3 with its arithmetic action]
\label{c:serre-centre}
\status{Proved under the stated Serre-quotient hypothesis}
Let $K$ be finitely generated over $\Q$, embedded in $\C$, and let $K_0$ be the algebraic closure of $\Q$ in $K$. Suppose that the geometric connected motivic group is pro-reductive, that its maximal protorus quotient is the full Serre protorus $S$, and that the action of $\Gal(\overline K/K)$ on this quotient is the usual action through $\Gal(\overline\Q/K_0)$. If $E$ is the largest CM subfield of $K_0$, then
\[
 \dim_\Q (\Lie G)^G=
 \begin{cases}1+[E:\Q]/2,&E\text{ exists},\\1,&K_0\text{ has no CM subfield}.
 \end{cases}
\]
These hypotheses hold for the full category of motivated motives of abelian type, including Artin motives. The conclusion concerns the full arithmetic adjoint invariants; the centre of the geometric connected group alone is generally an infinite-dimensional pro-vector space.
\end{theorem}

\begin{proof}
We first prove the character calculation rather than infer it from the dimension of a geometric torus. Write $\Q^{\mathrm{CM}}$ for the union of finite Galois CM extensions of $\Q$. Complex conjugation is central in its Galois group. The rational character space of the Serre protorus is the union of the spaces of locally constant rational-valued functions on the finite Galois quotients satisfying
\[
 f(\sigma)+f(c\sigma)=w(f),\qquad w(f)\text{ constant}.
\]
It decomposes as the one-dimensional constant space plus the space of $c$-odd functions. Taking coinvariants under $\Gamma_{K_0}$ is exact at each finite level, because averaging over the finite image identifies coinvariants with invariants in characteristic zero. At a finite Galois CM field $L$, invariant functions are functions on the embeddings of $L\cap K_0$. If this intersection is totally real, conjugation fixes every embedding and no odd function remains. If it is CM, conjugation pairs its embeddings without fixed points and the odd space has dimension $[L\cap K_0:\Q]/2$.

Every subfield of a Galois CM field is either totally real or CM: central conjugation descends to it, and the fixed field is totally real. The CM subfields of the finite number field $K_0$ have a compositum which is again CM, hence a largest one $E$ whenever any exists. Taking $L$ containing its normal closure makes the preceding dimensions stabilize at the asserted value. Thus the coinvariant character space is finite-dimensional of that dimension. Since $\Lie S=\Hom_\Q(X^*(S)\otimes\Q,\Q)$ as a pro-vector space, its fixed part is dual to those coinvariants. The preceding proposition proves the formula.

Restriction $\Gamma_K\to\Gamma_{K_0}$ is surjective: $K/K_0$ is regular, hence linearly disjoint from every finite extension of $K_0$, and the compatible extension property for finite Galois extensions gives surjectivity on inverse limits. Therefore no extra constants enter the invariant calculation.

For motivated motives of abelian type, \citet[Theorem 0.6.2 and the remark following its proof in \S6]{c_andre1996} proves that all Hodge classes on abelian varieties are motivated, including the absolute-Hodge classes over an arbitrary characteristic-zero base. Consequently the geometric Hodge realization is fully faithful on this category; compare \citet[Theorem 8.25 and Proposition 8.23]{c_delignehodge1982}. Its maximal subcategory with torus group is the category of CM Hodge structures, generated by CM abelian varieties. The Serre quotient and its arithmetic action are identified by \citet[unnumbered Theorem, p.~262, and Lemmas 1--2]{c_delignetaniyama1982}. A CM object defined over a finite extension is brought into the category over $K$ by restriction of scalars and then recovered after that extension as a direct summand; thus this assertion uses the full abelian-type category, not merely the CM varieties already having all endomorphisms over $K$. These facts establish precisely the hypotheses used in our character calculation. They also explain why ``no CM subfield'' does not mean that the geometric category contains no CM objects.
\end{proof}

\begin{example}[The finite-field analytic/algebraic distinction]
\label{c:finite-field-lie}
\status{Proved counterexample to the one-dimensional algebraic reading}
Let $E/\F_5$ have equation $y^2=x^3-x$. It is smooth since its discriminant $64$ is nonzero modulo $5$. It has eight rational points, so geometric Frobenius has polynomial $T^2+2T+5$ on $H^1$ and eigenvalues $\pi=-1+2i$, $\bar\pi=-1-2i$. Its algebraic Frobenius torus is
\[
 T=\operatorname{Res}_{\Q(i)/\Q}\Gm,
 \qquad\dim T=2.
\]
Indeed the two primes of $\Q(i)$ above $5$ give valuations $(1,0)$ and $(0,1)$ on $(\pi,\bar\pi)$; hence $\pi^a\bar\pi^b=1$ implies $a=b=0$. There is consequently no character cutting out a proper algebraic subgroup containing Frobenius. The adjoint action of $T$ is trivial, so its whole Lie algebra is algebraic, but it is two-dimensional.

On the other hand, for $\ell=3$ a suitable positive power of $\pi$ belongs to $1+3\Z_3[i]$. The analytic closure of the cyclic Frobenius subgroup has Lie algebra $\Q_3\log(\pi^N)$, of dimension one. The logarithm is nonzero because $\pi$ is not a root of unity. This is the distinction made by Serre in his letter of 8 September 1964, printed p.~202 of \citet{c_grothendieckserre}. Thus ``tout'' on 12:59 is valid for the adjoint Lie algebra of a commutative algebraic motivic band, as in the finite-field category under the Tate hypotheses of Part A; the accompanying identification with a single $\Q(0)$ does not describe that algebraic Lie object. Treating an analytic one-dimensional image as the motivic Lie algebra is an error, not an unconditional proof of a finite-field motivic assertion.
\end{example}

\subsection{C.4: weights, levels, and the overstrengthened bound}

\begin{definition}[The level convention of 12:111--115]
\label{c:level-def}
\status{Definition, checked against 12:111--115}
We fix $\ell$ invertible on the base and use $\Q_\ell(1)$ of weight $-2$. In the dimension-normalized convention, a special lisse sheaf of weight $r\geq0$ is a lisse subquotient of $R^rg_*\Q_\ell$ for a smooth projective morphism $g$ of relative dimension $r$. This includes the inserted dimension clause on 12:111; the transcription marks its following qualifier as uncertain. ``Strictly special'' means a direct summand, as clarified in 12:113. We explicitly call the version without a restriction on the relative dimension the broader convention. Over a field the two conventions agree, by the cut-and-pairing argument in Lemma~\ref{c:proper-level}; Theorem~\ref{c:general-level-bound} first uses the dimension-normalized convention, and Proposition~\ref{c:broader-coefficients} removes that clause for the general polygon and for the requested margin bound. Effective admissibility means that, after a finite stratification and the finite extensions allowed in 12:111, there is a finite filtration by extensions by zero of such sheaves. A sheaf has levels at most $\nu$ if these pieces may be chosen as $V(-a)$ with $a\geq0$, where $V$ is special of weight $r\leq\nu$. The simultaneous condition of weights in $\Phi$ requires the same pieces to satisfy $r+2a\in\Phi$. This is the geometric generation definition in 12:112; it is not simply a definition in terms of the eigenvalue weights, and it is not automatically the coniveau filtration.
\end{definition}

\begin{proposition}[C.4: the general $n-1$ bound is false]
\label{c:level-counterexample}
\status{Refuted, with the precise manuscript correction}
In the generality requested in the prompt, the bound
\[
 \operatorname{level}(R^if_*F)_{\bar y}
 \leq\min\{i,d_y(f)-1\}+j
\]
is false. Take $f:E\to\Spec\F_5$ for the elliptic curve in Example~\ref{c:finite-field-lie}, $F=\Q_\ell$, $j=0$, $i=1$, and any $\ell\ne5$. Then $d_y(f)=1$ and the right side is zero, whereas $H^1(E_{\overline\F_5},\Q_\ell)$ has level exactly one.
\end{proposition}

\begin{proof}
The sheaf $F$ is strictly special of weight zero and $E$ is regular, smooth and projective. Its nonzero $H^1$ is special of weight one, so its level is at most one. A level-zero object, by Definition~\ref{c:level-def}, has successive constituents $V(-a)$ with $V$ of weight zero and $a\geq0$; all its weights are even. The characteristic polynomial $T^2+2T+5$ has two roots of absolute value $\sqrt5$, so this $H^1$ is pure of odd weight one and cannot have level zero. The program records all eight points and the polynomial, so the counterexample does not depend on a catalogue lookup.

The supplied transcription of the final general-case box of 12:120 reads $\min(i,n)+j$, not $\min(i,n-1)+j$. The margin asks ``peut on remplacer par $n-1$ si $f$ de dim $\leq n-1$ ?'' The elliptic structural morphism has relative dimension $n$, so it does not refute that additional margin hypothesis. We therefore refute the strengthened prompt, not an assertion that the manuscript actually made without the condition.
\end{proof}

\begin{proposition}[Sharp elementary cases]
\label{c:level-special-cases}
\status{Proved in the stated cases}
If $f:X\to\Spec k$ is smooth projective of dimension $n$ and $F=\Q_\ell$, then $H^i(X)$ is pure of weight $i$ over a finite field and has level at most $\min(i,2n-i)$. Over a smooth pair $(Y,D)$ with a strict normal-crossing divisor and open immersion $j:Y\setminus D\hookrightarrow Y$, if $r$ branches meet at a geometric point $\bar y$, then
\[
 (R^ij_*\Q_\ell)_{\bar y}\simeq
 \bigwedge^i\bigl(\Q_\ell(-1)^r\bigr).
\]
It vanishes for $i>r$ and, when nonzero, has weight $2i$ and level zero. If a strictly special lisse sheaf extends across this pair, the projection formula tensors this description with its stalk.
\end{proposition}

\begin{proof}
For $i\leq n$, the cohomology itself is special of weight $i$, hence has level at most $i$. For $i>n$, hard Lefschetz identifies it with $H^{2n-i}(X)(-(i-n))$, giving the second bound. Products of ordinary elliptic curves over finite fields attain these geometric level bounds; the Frobenius-slope proof is given in Proposition~\ref{c:level-optimality}. The normal-crossing formula is local: in strict henselian coordinates each missing coordinate hyperplane supplies its Kummer class in $\Q_\ell(-1)$, and the products of distinct Kummer classes form the exterior algebra. Purity and the localization sequence prove this by induction on the number of coordinate divisors. Frobenius acts with weight two on each Kummer class. For an extending lisse sheaf, the stalkwise projection formula applies and gives the last assertion.
\end{proof}

\begin{definition}[The local dimension on 12:120]
\label{c:local-dimension}
\status{Definition, with the punctured-trait correction verified}
For a finite-type morphism $f:X\to Y$ of universally catenary noetherian schemes we use
\[
 d_y(f)=\sup_{\substack{x\in X\\ f(x)\text{ generalizes }y}}
 \left\{\operatorname{trdeg}_{k(f(x))}k(x)
       +\dim\overline{\{f(x)\}}^{\,\Spec\mathcal O_{Y,y}}\right\}.
\]
The closure in the second summand is taken in the local base. This is the dimension function, normalized to be zero at $y$, pulled back to $X$. It is the formula visible in the supplied transcription, whose notation for the local spectrum is uncertain. It measures local source dimension, not maximum fibre dimension: for $\Spec K\hookrightarrow\Spec R$ with $R$ a discrete valuation ring and $y$ its closed point it equals $1$, although all nonempty fibres have dimension zero. After an equidimensional localization of the source and a compactification, its maximum is the local dimension of that compactification. Regular alterations preserve this maximum. The dimension-function convention and the corresponding local cohomological interpretation are made precise in \citet[\S1.1.1 and Corollary 1.2.2]{c_pillonistroh2014}.
\end{definition}

\begin{lemma}[Proper cohomology with simultaneous geometric bounds]
\label{c:proper-level}
\status{Proved}
Let $Z$ be proper of dimension at most $d$ over a field $k$, and let $a\geq0$. After the finite field extensions allowed in the admissibility definition, $H^a(Z_{\bar k},\Q_\ell)$ has a finite filtration by subquotients of special objects $V_r(-t)$ satisfying
\[
 t\geq\max(0,a-d),\qquad r\geq0,\qquad r+2t\leq a.
\]
In particular its weights are at most $a$ and its geometric levels are at most $\min(a,2d-a)$. A negative upper level bound means the object is zero. These are geometric generation assertions, not deductions from a Hodge polygon.
\end{lemma}

\begin{proof}
For a smooth projective variety of dimension $d$, hard Lefschetz identifies $H^a$ with $H^{2d-a}(-(a-d))$ when $a>d$; for $a\leq d$ use $H^a$ itself. The input is \citet[Theorem 4.1.1]{c_deligneweil1980}. If a special object of degree $r$ is required to be presented by a variety of dimension exactly $r$, take a smooth complete intersection of dimension $r$ in the smooth projective variety. Weak Lefschetz injects $H^r$ into its middle cohomology. The restriction pairing is nondegenerate by hard Lefschetz on the original variety, so this injection even has a Galois-equivariant left inverse. A finite extension of $k$ suffices to choose the hyperplanes. Thus both conventions of specialness give the same objects over a field.

For an integral proper $Z$, \citet[Theorem 7.3 and Corollary 7.4]{c_dejong1996}, applied over an algebraic closure and then descended to a finite extension, give a smooth projective $W$, a finite group $G$ acting on it, and a proper map $V=W/G\to Z$ which is a universal homeomorphism over a dense open $U\subset Z$. The construction and the finite group action descend because they involve finitely many algebraic coefficients. Passing through finite purely inseparable extensions does not change the \(\ell\)-adic topos. Let $S=Z\setminus U$ and $E=V\times_ZS$. Both have dimension at most $d-1$. Proper base change identifies
\[
 H^a(V,\Q_\ell)=H^a(W,\Q_\ell)^G:
\]
the finite quotient has no higher direct images of the constant sheaf, and the invariant functions on every geometric orbit are the constants. Averaging is exact over $\Q_\ell$, even if $\ell$ divides $|G|$.

The localization triangles for $U\subset Z$ and its inverse image, together with proper base change, give the abstract blowup sequence
\[
 H^{a-1}(E)\longrightarrow H^a(Z)
 \longrightarrow H^a(V)\oplus H^a(S).
\]
The induction hypothesis for $E$ gives $t\geq\max(0,(a-1)-(d-1))$ and weight at most $a-1$; for $S$ it gives an even stronger lower bound on $t$. The smooth-projective assertion for $W$ handles $V$. Subquotients and extensions preserve all three inequalities. Reducible proper schemes follow by the same localization argument for their components; nilpotents do not change the \(\ell\)-adic topos. This completes the dimension induction.

Finally finite-extension descent does not enlarge the bounds. For a finite separable extension $k'/k$, a representation is a direct summand of the induction of its restriction, by the trace divided by $[k':k]$. A smooth projective $k'$-scheme viewed as a $k$-scheme is still smooth projective of the same dimension, and realizes that induced cohomology. Purely inseparable extensions have already been handled. This justifies the descent used above.
\end{proof}

\begin{theorem}[Corrected bounds and the margin question]
\label{c:general-level-bound}
\status{Proved in the stated arithmetic finite-type setting}
Let $f:X\to Y$ be separated of finite type, with $X$ regular and $X,Y$ of finite type over $\Z[1/\ell]$ or over a field of characteristic different from $\ell$. Fix $y\in Y$ and put $n=d_y(f)$. Let $F$ be strictly special of weight $j$, in the dimension-normalized sense of the insertion on 12:111: it is a direct summand of $R^jg_*\Q_\ell$ for a smooth projective $g$ of relative dimension $j$. Set $\epsilon=1$ if no component of local dimension $n$ maps to $y$, and set $\epsilon=0$ otherwise. Put
\[
 L=\max(0,i-n+\epsilon),\qquad U=\min(i,n)+j.
\]
The stalk $(R^if_*F)_{\bar y}$ is effectively admissible, with a filtration into pieces of weights $w$ and levels bounded simultaneously by
\begin{equation}
\label{c:weight-level-polygon}
 2L\leq w\leq2U,
 \qquad \operatorname{level}\leq\min\{w-2L,\,2U-w\}.
\end{equation}
Consequently
\[
 \operatorname{weight}\leq2\min(i,n)+2j,
 \qquad
 \operatorname{level}\leq
 \begin{cases}
 \min(i,2n-i)+j,&\epsilon=0,\\
 \min(i,2n-1-i)+j,&\epsilon=1.
 \end{cases}
\]
For an open immersion at a boundary point, $\epsilon=1$. The condition that $f$ have relative dimension at most $n-1$ near $y$ also implies $\epsilon=1$. Hence the margin's replacement by $n-1$ is valid with its stated extra condition, since $\min(i,2n-1-i)\leq\min(i,n-1)$. Without that condition the corrected general bound is at most $\min(i,n)+j$, and Proposition~\ref{c:level-counterexample} prevents replacing $n$ by $n-1$ universally.
\end{theorem}

\begin{proof}
We first take $F=\Q_\ell$ and $j=0$. Work locally on $Y$ around $y$, compactify $X$ over $Y$, and retain only the closure of $X$. Applied to that compactification and its boundary, the alteration theorems of \citet[Theorem 4.1 for varieties over fields; Theorem 8.2 for arithmetic schemes]{c_dejong1996} give a proper regular alteration $h:\widetilde X\to Y$ with a strict normal-crossing boundary $D$, and an open $v:X'\hookrightarrow\widetilde X$ above $X$. The induced $\pi:X'\to X$ is proper and generically finite. The construction can be made component by component. Localization to $\Spec\mathcal O_{Y,y}$ preserves regularity, boundary codimensions, and the local dimension function.

There are maps
\[
 \Q_{\ell,X}\longrightarrow R\pi_*\Q_{\ell,X'}
 \xrightarrow{\operatorname{Tr}_\pi}\Q_{\ell,X}
\]
whose composite is multiplication by the generic degree on each component. Here the trace is not an unjustified trace for a nonfinite map: a projective generically finite map between regular schemes is a smoothable local complete-intersection morphism of virtual dimension zero. Absolute purity identifies its Gysin map with $\Q_\ell\to\pi^!\Q_\ell$, and proper adjunction gives the displayed trace. The degree computation is checked on the generic finite-flat locus, since an endomorphism of the constant sheaf on a connected scheme is a scalar. The precise construction and calculation are \citet[Corollary 3.1.2, Lemmas 3.5.8--3.5.10, and Proposition 2.5.13]{c_riou2014}; the equality also covers inseparable generic degree by the finite-flat degree formula. Dividing by the degree proves that $Rf_*\Q_\ell$ is a direct summand of $Rh_*Rv_*\Q_\ell$.

Let $D^{(q)}$ denote the disjoint union of the intersections of $q$ distinct boundary components, and put $D^{(0)}=\widetilde X$. Absolute purity gives
\[
 R^qv_*\Q_\ell=\Q_{\ell,D^{(q)}}(-q).
\]
This is \citet[Corollary 3.1.4]{c_riou2014}, equivalently \citet[Theorem 2.1.2]{c_pillonistroh2014}; it remains true for regular arithmetic schemes. The stalkwise Leray sequence and proper base change now give
\[
 E_2^{a,q}=H^a(F_q,\Q_\ell)(-q)
 \Longrightarrow (R^{a+q}(f\pi)_*\Q_\ell)_{\bar y},
 \qquad F_q=D^{(q)}\times_Y\bar y.
\]
Every $F_q$ is proper. By the definition of $n$ and regular codimension, $\dim F_q\leq n-q$ for $q\geq1$, while $\dim F_0\leq n-\epsilon$. This step does not assume that intersecting a special fibre with a vertical divisor lowers its dimension: the inequality already holds for the entire local scheme $D^{(q)}$, of codimension $q$ in a regular local source of dimension at most $n$. The extra drop for $F_0$ follows because a proper closed fibre containing a component of maximal local dimension would make that component map to $y$. Alteration does not change the image of its generic point.

Write $a=i-q$. Lemma~\ref{c:proper-level} says that each constituent of $E_2^{a,q}$ is $V_r(-(t+q))$, with
\[
 r+2t\leq i-q,
 \qquad t\geq\max(0,i-q-\dim F_q).
\]
For $q=0$ this gives $t\geq L$. For $q\geq1$ it gives $t+q\geq q+\max(0,i-n)\geq L$. Its total weight $w=r+2(t+q)$ satisfies $w\leq i+q$, and nonvanishing requires $q\leq\min(i,2n-i)$ (as well as $q\leq n$). It follows that $w\leq2\min(i,n)$ and $r\leq w-2L$. The other inequality $r\leq2\min(i,n)-w$ follows as follows. If $i\leq n$, use $w\leq i+q$ to obtain $r=w-2(t+q)\leq2i-w$. If $i>n$, use $t\geq i-n$ for $q\geq1$, and combine $q\geq w-i$ with the same formula, obtaining $r\leq2n-w$; the $q=0$ term obeys it because $w\leq i$ and $t\geq i-n+\epsilon$. Thus every term lies in \eqref{c:weight-level-polygon}. A spectral sequence gives a filtration by subquotients of its terms; the inequalities survive extensions and direct summands. No degeneration of this Leray sequence is needed.

For the coefficient sheaf, relative hard Lefschetz and \citet[Theorem 1.5 and Proposition 2.4]{c_deligne1968} give a derived splitting
\[
 Rg_*\Q_\ell\simeq\bigoplus_bR^bg_*\Q_\ell[-b].
\]
Accordingly $R^if_*F$ is a direct summand of $R^{i+j}(fg)_*\Q_\ell$. The source of $g$ is regular, its local dimension is $n+j$, and the value of $\epsilon$ is unchanged. Applying the constant-coefficient result with $(i,n)$ replaced by $(i+j,n+j)$ gives exactly \eqref{c:weight-level-polygon} and the displayed bounds. This proves the corrected statement, including the margin condition.
\end{proof}

\begin{proposition}[Removing the dimension clause from the coefficient]
\label{c:broader-coefficients}
\status{Proved}
Keep the finite-type hypotheses of Theorem~\ref{c:general-level-bound}, but suppose only that $F$ is a direct summand of $R^jg_*\Q_\ell$ for a smooth projective $g:Z\to X$ of arbitrary relative dimension. Put $n=d_y(f)$ and $L_0=\max(0,i-n)$, $U=\min(i,n)+j$. Then the simultaneous bounds
\[
 2L_0\leq w\leq2U,\qquad
 \operatorname{level}\leq\min(w-2L_0,2U-w)
\]
still hold. In particular the general level bound is $\min(i,2n-i)+j$.
If no component of local dimension $n$ maps to $y$, there is also the bound
\[
 \operatorname{level}\leq\min(i,n-1)+j.
\]
Consequently the requested weight bound and the manuscript's margin bound hold in this broader convention as well. The stronger margin polygon with $L_0$ replaced by $\max(0,i-n+1)$ remains covered by Theorem~\ref{c:general-level-bound} in its dimension-normalized convention.
\end{proposition}

\begin{proof}
We first explain the relative dimension reduction on a dense open set. We work componentwise, so the relative dimension $r$ of $g$ is constant. If $r\leq j$, replace $Z$ by $Z\times\mathbb P^{j-r}$; K\"unneth realizes $R^jg_*\Q_\ell$ as a direct summand of the degree-$j$ cohomology of this dimension-$j$ family. If $r>j$, choose over the generic point of $X$ a smooth complete intersection of relative dimension $j$ in a projective embedding of $Z$. It spreads to a smooth projective family $h:W\to U$ over a dense open $U\subset X$. The generic field is infinite unless the source component is zero-dimensional over a finite field; in that case a finite field extension permits Bertini and restriction/corestriction gives the same direct-summand conclusion over the original field.

Weak Lefschetz gives an injection $R^jg_*\Q_\ell|_U\to R^jh_*\Q_\ell$. The pairing induced on the source by middle-dimensional Poincar\'e duality on $W/U$ is cup product with the $(r-j)$th power of the hyperplane class followed by trace. Hard Lefschetz makes it nondegenerate. Taking its adjoint and inverse gives a morphism of lisse sheaves which retracts the injection. Thus $F|_U$ has the dimension-$j$ presentation required by the preceding theorem.

We prove the general polygon by induction on the inherited local dimension. Choose the dense open $U$ just constructed and write $D=X\setminus U$. Since $X$ is regular and excellent, $D$ has a finite stratification by regular locally closed schemes $Z_\alpha$, each of pure codimension $c_\alpha\geq1$ in its source component, with a filtration by closed supports. The dimension formula gives
\[
 d_y(f|_{Z_\alpha})\leq n-c_\alpha.
\]
For nested closed supports $D'\subset D''$, the localization triangle for $R\Gamma_{D''}F$ gives the quotient term obtained by pushing the supported complex on $D''\setminus D'$ from $X\setminus D'$ to $X$. Absolute purity therefore gives, for our regular strata, successive terms
\[
 R(i_\alpha)_*(F|_{Z_\alpha})(-c_\alpha)[-2c_\alpha],
 \qquad i_\alpha:Z_\alpha\longrightarrow X.
\]
Here $i_\alpha$ is the locally closed immersion, and the pushforward is derived ordinary pushforward. Purity applies after the open restriction in which that immersion is closed; its exact form is \citet[Corollary 3.1.2]{c_riou2014}. Composing with $Rf_*$, then using the induction hypothesis on $Z_\alpha$, bounds the supported terms.

Explicitly write $c=c_\alpha$, $n'=d_y(f|_{Z_\alpha})$ and $m=i-2c$. If $m<0$ there is no term. Otherwise the induction hypothesis supplies pieces $V_a(-t_0)$ with
\[
 t_0\geq\max(0,m-n'),\qquad a+t_0\leq\min(m,n')+j.
\]
The additional twist replaces $t_0$ by $t=t_0+c$. Since $n'\leq n-c$,
\[
 t\geq\max(c,i-n)\geq L_0,
 \qquad a+t\leq\min(i-c,n)+j\leq U.
\]
These are exactly the asserted simultaneous polygon inequalities. The open term $R(f|_U)_*(F|_U)$ satisfies them by Theorem~\ref{c:general-level-bound}. The localization triangle $R\Gamma_DF\to F\to Rv_*F|_U$, with $v:U\hookrightarrow X$, now proves them for $Rf_*F$: in each cohomological degree the middle group is an extension of a quotient of the first group by a subobject of the third. Finite extensions, subquotients and filtrations preserve the two inequalities. The zero-dimensional starting case was the field reduction above. This proves the general polygon.

For the extra margin bound we use a different descent that does not introduce vertical closed strata. The smooth-complete-intersection parameter space is a nonempty open in a product of relative projective spaces and is smooth and surjective over $X$. A smooth surjection has sections after an \emph{\'etale} covering. Quasi-compactness gives a finite-type surjective \'etale map $v:V\to X$ over which the coefficient has a dimension-$j$ presentation. Denote the terms of its \v Cech nerve by $v_a:V_a\to X$. They are \'etale over $X$, hence regular, and their inherited local dimensions are at most $n$. When that dimension is $n$, no component maps to $y$; when it is smaller, it is at most $n-1$. The normalized theorem consequently gives, for every $a,b$,
\[
 \operatorname{level}\bigl(R^b(fv_a)_*v_a^*F\bigr)_{\bar y}
       \leq\min(b,n-1)+j.
\]
The \'etale \v Cech descent spectral sequence is
\[
 E_1^{a,b}=\bigl(R^b(fv_a)_*v_a^*F\bigr)_{\bar y}
       \ \Longrightarrow\ (R^{a+b}f_*F)_{\bar y}.
\]
For torsion sheaves this is \citet[Proposition 96.21.1, Tag 072G; see also Theorem 59.19.2, Tag 03OW]{c_stacksdescent}. For a lisse $\Q_\ell$-sheaf choose a stable lattice, apply it modulo $\ell^s$, and pass to the inverse limit and then invert $\ell$. Constructibility makes the finite-level stalks finite, so their inverse systems are Mittag--Leffler. Each total degree has only finitely many relevant first-quadrant terms, which justifies this passage and convergence. Since $b\leq a+b=i$, every contributing piece has level at most $\min(i,n-1)+j$. This proves the margin assertion. No properness of the \'etale cover was needed, and no stronger reflected margin inequality is inferred from this spectral sequence.
\end{proof}

\begin{proposition}[Optimality and the limit of the word sharp]
\label{c:level-optimality}
\status{Proved optimality in the stated senses}
For constant coefficients the preceding upper weight and level bounds are attained for every nonnegative degree in the indicated nonzero ranges. The general and margin level bounds remain optimal after adding $j$. The coefficient $2j$ in the universal weight envelope cannot be replaced by $j$: for every $j\geq0$ there is a strictly special weight-$j$ example with $n=i=1$ whose $R^1$ boundary stalk is pure of weight $2j+2$.
\end{proposition}

\begin{proof}
For the general level bound, take the structural map of $E^n$ over $\F_5$, where $E$ is the ordinary elliptic curve of Example~\ref{c:finite-field-lie}. Its Frobenius eigenvalues have $5$-adic slopes $0,1$. In $H^i(E^n)$ the smallest slope is $\max(0,i-n)$ and the weight is $i$. Any special constituent $V_r(-t)$ has nonnegative slopes before twisting, since eigenvalues on smooth projective cohomology are algebraic integers \citep[Expos\'e XXI, Theorem~5.2.2 and Corollary~5.5.3, pp.~385--386 and 393--395]{c_katzdeligne1973niveau}. Hence such a lowest-slope eigenvalue requires $t\leq\max(0,i-n)$, and therefore $r=i-2t\geq\min(i,2n-i)$. The upper bound is attained. Tensoring with the constant strictly special object $H^j(E^j)$, which has a slope-zero eigenvalue, raises the required level by exactly $j$.

For the margin case let $B=E^{n-1}$ and choose a sufficiently ample projectively normal embedding. Let $Y$ be its affine cone and $X=Y\setminus\{y\}$, where $y$ is the vertex. Then $X$ is the complement of the zero section of a line bundle on the smooth $B$, so it is regular of dimension $n$ and $X\hookrightarrow Y$ is an open immersion. The blowup of the vertex gives the line-bundle resolution. Purity and proper base change identify its local cohomology with the Gysin sequence for that line bundle:
\[
 \operatorname{coker}\bigl(L:H^{i-2}(B)(-1)\to H^i(B)\bigr)
 \longrightarrow (R^if_*\Q_\ell)_{\bar y}
 \longrightarrow\ker\bigl(L:H^{i-1}(B)(-1)\to H^{i+1}(B)\bigr).
\]
The two outer terms give a short exact sequence. For $i\leq n-1$, a slope-zero class in $H^i(B)$ cannot be in the image of $L$, whose eigenvalues have slopes at least one. It gives level $i$. For $n\leq i\leq2n-1$, put $r=2n-i-1$. The primitive degree-$r$ cohomology contains a slope-zero class for the same reason. Hard Lefschetz carries it to the right-hand kernel after $i-n$ powers of $L$, and the boundary twist adds one further Tate twist. This has weight $i+1$ and slope $i-n+1$, forcing level at least $2n-i-1$. Tensoring again with $H^j(E^j)$ proves sharpness with $+j$. For $n=1$ the same construction is simply $\Gm\hookrightarrow\mathbb A^1$.

For weights with $j=0$, if $i\leq n$ use the origin-boundary stalk of
$\Gm^i\times\mathbb A^{n-i}\hookrightarrow\mathbb A^n$, whose degree-$i$ Kummer class has weight $2i$. If $n<i\leq2n$, put $q=2n-i$ and $r=i-n$. The product of the degree-$q$ class of $\Gm^q$ with the top class of $\mathbb P^r$ has degree $i$ and weight $2(q+r)=2n$. Take the product open immersion followed by projection to $\mathbb A^q$; if $q=0$ this is the projective structural map. For the margin range $i\leq2n-1$ we have $q\geq1$, so the same example satisfies the extra condition.

Finally let $U=\mathbb A^1_{\F_5}\setminus\{0,1\}$ and let $\pi:\mathcal E\to U$ be the Legendre family
\[
 y^2=x(x-1)(x-t).
\]
At $t=0$ its discriminant $16t^2(1-t)^2$ has order two and $c_4=16(1-t+t^2)$ is a unit. The node has tangent slopes whose squares are $-1$, which is a square in $\F_5$, so the reduction is split multiplicative. Thus on $H^1$ inertia acts by a nonzero rank-one unipotent operator, with invariant line $\Q_\ell$ and coinvariant line $\Q_\ell(-1)$. This is the Tate-curve calculation underlying SGA~7~I, IX, Proposition 2.2.5. Take $F_j=\operatorname{Sym}^j R^1\pi_*\Q_\ell$. The K\"unneth and Schur idempotents exhibit it as a direct summand of $R^j(\mathcal E^j/U)_*\Q_\ell$, so it has the required dimension-$j$ presentation. Its inertia coinvariants are $\Q_\ell(-j)$, and tame cohomology supplies one further twist:
\[
 (R^1u_*F_j)_{\overline0}=\Q_\ell(-j-1),
 \qquad u:U\hookrightarrow\mathbb A^1.
\]
This proves the claimed weight $2j+2$. Products with the Kummer and projective examples above attain $2\min(i,n)+2j$ for every $1\leq i\leq2n-1$. The finite matrix calculation of the symmetric-power monodromy and the displayed optimization inequalities are recorded in the accompanying program.
\end{proof}

\begin{proposition}[The top degree and its sharper weight]
\label{c:top-direct-image}
\status{Proved}
In the finite-type setting of Theorem~\ref{c:general-level-bound}, let $F$ be any lisse $\Q_\ell$-sheaf. Then $(R^if_*F)_{\bar y}=0$ for $i>2n$. If no component of local dimension $n$ maps to $y$, it vanishes already for $i>2n-1$. If $k(y)$ is finite and $F$ is pure of weight $j$, the degree-$2n$ stalk is pure of weight $2n+j$ when nonzero. This endpoint is attained by $\mathbb P^n\to\Spec\F_q$ with constant coefficient $H^j(E^j)$ for an elliptic curve $E/\F_q$. These assertions do not require the dimension-normalized presentation of $F$.
\end{proposition}

\begin{proof}
We first treat an integral source component of inherited local dimension $n$ whose image is not $y$. Compactify it as $v:X\hookrightarrow\bar X\xrightarrow{h}Y$, and pass to the strict localization at $y$. For a geometric point over $x\in\bar X_y$, let $c$ be the dimension of the strict local ring of $\bar X$ at $x$. The dimension formula gives $c+\operatorname{trdeg}_{k(y)}k(x)\leq n$. For $q>0$, Gabber's bound on the cohomological dimension of an open subset of a strictly local noetherian $c$-dimensional scheme gives
\[
 (R^qv_*F)_{\bar x}\ne0\ \Longrightarrow\ q\leq2c-1;
 \quad
 \dim\operatorname{Supp}\bigl((R^qv_*F)|_{\bar X_y}\bigr)
 \leq n-\Bigl\lceil\frac{q+1}{2}\Bigr\rceil.
\]
The exact input is \citet[Theorem 1.1, formula (1.1.1), p.~455]{c_illusie2014cd}. Its torsion statement passes to lisse $\Q_\ell$-sheaves by a stable lattice and the finite constructible stalks, whose inverse systems satisfy Mittag--Leffler. For $q=0$ the support is contained in $\bar X_y$, of dimension at most $n-1$: a fibre component of dimension $n$ would be a whole irreducible component mapping to $y$.

Proper base change and the composition spectral sequence express the desired stalk through
\[
 H^a\bigl(\bar X_{\bar y},(R^qv_*F)|_{\bar X_{\bar y}}\bigr).
\]
The cohomology of a constructible sheaf supported in dimension $s$ over an algebraically closed field vanishes above $2s$. For $q>0$ every nonzero term therefore has
\[
 a+q\leq2n-2\Bigl\lceil\frac{q+1}{2}\Bigr\rceil+q\leq2n-1;
\]
for $q=0$ it has $a\leq2n-2$. This proves the asserted vanishing on this component. Components of smaller local dimension obey the same stronger bound with their own dimension.

The irreducible components of regular $X$ are disjoint. The only possible contributions in degree $2n$ thus come from its $n$-dimensional components mapping to $y$. Such a component is a smooth variety over the finite field $k(y)$ in the weight assertion. Poincar\'e duality identifies its top cohomology with $H^0_c(X_{\bar y},F^\vee)^\vee(-n)$. On a connected nonproper component, this compactly supported section space is zero: a nonzero section of a lisse sheaf has support the entire component. On a proper component it is a subspace of a fibre of $F^\vee$, hence is pure of weight $-j$. Duality gives weight $2n+j$. Components are permuted through a finite group, which does not change weights. The case $n=0$ is the same finite-fibre argument. The projective-space example gives the stated optimality.
\end{proof}

\begin{proposition}[An improved bound in degree zero]
\label{c:improved-zero-weight}
\status{Proved; optimality for general coefficients of weight at least three remains unresolved}
Let $f:X\to Y$ be separated and of finite type over $\Z[1/\ell]$
or over a finite field of characteristic different from $\ell$, with
$X$ regular, and let $y\in Y$ have finite residue field.  Let $F$ be
strictly special of weight $j\geq0$, in either convention of
Definition~\ref{c:level-def}.  Put
\[
 b(0)=0,\qquad b(1)=1,\qquad b(j)=2j-2\quad(j\geq2).
\]
Then the eigenvalues of geometric Frobenius on
$(f_*F)_{\bar y}$ have weights at most $b(j)$.
In particular, the expected bound $j$ holds for $j\leq2$.
The assertion about weights does not replace the geometric level
definition or assert a corresponding smaller level.
\end{proposition}

\begin{proof}
We first prove the local-field input, including the two extreme
monodromy maps.  Let $K$ be a finite extension of $\Q_p$ or of
$\F_p((t))$, and let $Z/K$ be smooth projective of dimension $d$.
After a finite extension and a projective semistable alteration,
the pullback on cohomology is split injective by trace divided by
the generic degree.  Thus we may use a projective strictly
semistable model.  This reduction is
\citet[Theorem~6.5, pp.~83--86]{c_dejong1996}; the assertion
descends to the original inertia invariants because Frobenius over
the extension acts there as a power of the original Frobenius.

Let $Y^{(r)}$ be the disjoint union of the $(r+1)$-fold intersections
of the special-fibre components.  We use the weight spectral sequence
\[
 E_1^{a,b}=\bigoplus_{r\geq\max(0,-a)}
 H^{b-2r}(Y^{(a+2r)}_{\bar k},\Q_\ell)(-r)
 \Longrightarrow H^{a+b}(Z_{\bar K},\Q_\ell).
\]
The $E_1^{a,b}$ are pure of weight $b$, the sequence degenerates at
$E_2$, and $N$ maps $(a,b)$ to $(a+2,b-2)$, with the matching
Tate twist.  These are
\citet[Satz~2.10 and (2.11), pp.~40--41]{c_rapoportzink1982}.
For $H^d$, all weights lie in $[0,2d]$.

The extreme map
\[
 N^d:E_2^{-d,2d}\longrightarrow E_2^{d,0}(-d)
\]
is an isomorphism.  Suppressing twists, it is the natural map
\[
 \ker\bigl(H^0(Y^{(d)})\xrightarrow{\mathrm{Gys}}
 H^2(Y^{(d-1)})\bigr)
 \longrightarrow
 \operatorname{coker}\bigl(H^0(Y^{(d-1)})
 \xrightarrow{\mathrm{Res}}H^0(Y^{(d)})\bigr).
\]
Restriction and Gysin are adjoint.  On the rational vector space
of functions on the finite set $Y^{(d)}$, the pairing is a positive
sum of products.  Its restriction to the rational image of
restriction is nondegenerate, and remains so over $\Q_\ell$.
The displayed kernel is its orthogonal complement and maps
isomorphically to the quotient.  This is the argument of
\citet[Proposition~1.5, p.~9]{c_yamamura2026weights}.

When $d\geq2$, the next extreme map
\[
 N^{d-1}:E_2^{-d+1,2d-1}\longrightarrow
 E_2^{d-1,1}(-(d-1))
\]
is also an isomorphism.  Put $U=H^1(Y^{(d-1)})$, so the components
in this expression are smooth projective curves, and put
$R=\operatorname{im}[H^1(Y^{(d-2)})\to U]$.
The signed restriction map is induced by the corresponding
homomorphism from the product of the reduced Picard varieties of
the surface strata to the product of the curve Jacobians.
Consequently $R$, with the usual Tate twist understood, is the
Tate realization of an abelian subvariety of that product.
The product principal polarization restricts to a polarization on
this abelian subvariety.  Thus the cup pairing on $U$ is
nondegenerate on $R$.  Since Gysin is the adjoint of restriction,
its kernel is $R^\perp$, and $R^\perp\to U/R$ is an isomorphism.
This proves the next extreme map, with all signs included.
We have used the Picard-polarization argument in the proof of
\citet[Satz~2.13, pp.~41--43, especially p.~42]{c_rapoportzink1982},
not an application of its surface theorem to a higher-dimensional
variety.  The resulting general-dimensional statement is
\citet[Proposition~1.7, p.~10]{c_yamamura2026weights}.

Inertia is unipotent in this model, so its invariants are $\ker N$.
An element in the kernel cannot have a nonzero weight-$2d$ component,
and for $d\geq2$ it cannot have a weight-$(2d-1)$ component either:
on each of these components a positive power of $N$ is injective.
The weight filtration is strict here, since it is the decomposition
into disjoint generalized Frobenius eigenspaces by their weights.
Hence $H^d(Z)^{I_K}$ has weights at most $2d-2$ for $d\geq2$, and
at most $1$ for $d=1$.  Degree zero has weight zero.  This recovers
the middle-degree consequences of
\citet[Theorem~1.6 and Proposition~1.8, pp.~9--11]{c_yamamura2026weights}.
The proof just given works in both local-field characteristics.
For a direct summand of $H^j$ of a smooth projective variety of
dimension at least $j$, the cut-and-pairing construction of
Lemma~\ref{c:proper-level} realizes it as a direct summand in degree
$j$ of a smooth projective $j$-fold, after a finite extension.
If the dimension is $d<j$, replace $Z$ by
$Z\times\mathbb P^{j-d}$: the $H^0$ projector on the projective
factor splits off the original $H^j$, and the product has dimension
$j$.  When $j>2d$ the original group is already zero.
Thus the same bound $b(j)$ applies.
We use only these independently proved middle-degree consequences.
The displayed higher-degree hard-Lefschetz isomorphism on p.~10
of the cited preprint prints the twist $i-d$; in our convention
the correct twist is $d-i$.  Its subsequent numerical weight bounds
use the latter sign.  This printed slip is not used in our proof.

We now reduce the degree-zero stalk to these local fields.
Write $S=\Spec\mathcal O_{Y,\bar y}^{\mathrm{sh}}$ and
\[
 V=(f_*F)_{\bar y}=H^0(X\times_Y S,F).
\]
There are finitely many connected components; they are integral
because the source is regular.  Descend these finitely many
open-and-closed components and a compactification to a finite
\'etale neighbourhood of $y$.  We do not assert that every
$\ell$-adic section descends to one such neighbourhood.
After a finite residue extension the relevant components are fixed.
The closure of each component in a proper compactification meets
the closed fibre.  Choose a closed point $c$ there.  A component
lying over $y$ can be evaluated at a finite closed point directly,
giving only weight $j$.

For the other components we construct a curve approaching $c$ from
the chosen dense open.  In the local domain $A$ of its closure at
$c$, choose $0\ne g\in\mathfrak m_A$ with $D(g)$ contained in
that open; in the horizontal mixed-characteristic case include $p$
as a factor of $g$.  Since $A$ is an excellent local domain,
the equidimensional dimension formula gives
$\dim A/(g)=\dim A-1$; thus $g$ extends to a system of parameters
$g,a_1,\ldots,a_{d-1}$ of $A$.  A minimal prime $P$ of
$(a_1,\ldots,a_{d-1})$ has height at most $d-1$ and cannot
contain $g$, since that would make $P=\mathfrak m_A$.
Moreover $g$ is primary for the maximal ideal of $A/P$, so
$\dim(A/P)=1$.  The resulting integral curve has generic point
in the required open.  Normalize, localize at a point above $c$,
and complete.  Excellence gives a discrete valuation ring with
finite residue field, whose fraction field is a finite extension
of $\Q_p$ or of $\F_p((t))$.  The entire curve construction is
performed after the preceding finite descent, over a scheme of
finite type over $\Z$ or a finite field.

Let $T$ be this trait.  The map $T^{\mathrm{sh}}\to S$, and its
generic-point map to the chosen component, pull back a section in
$V$ to an inertia-invariant vector of the restricted sheaf.
This map is equivariant for the residue Galois group after the
finite residue extension.  Evaluation at one geometric point of
a connected component is injective for a lisse sheaf.  Taking
the finite sum of these evaluations therefore embeds $V$ into
the local-field invariant spaces just considered and the pure
weight-$j$ evaluation spaces.  Since $j\leq b(j)$ for every
$j\geq0$, their weights are all at most $b(j)$.
Restriction of a strictly special sheaf retains its presenting
smooth projective family and its projector, so the local-field
input applies.  Finally the bound descends through the finite
residue extension by taking powers of Frobenius.
\end{proof}

\begin{corollary}[The first remaining lower-endpoint case]
\label{c:threefold-weight-gate}
\status{Proved equivalence for this case; the stated monodromy map remains unresolved in general}
Let $Z$ be a smooth projective threefold over a mixed-characteristic
local field with a projective strictly semistable model, and
$\ell\ne p$.  The following are equivalent:
\begin{enumerate}[label=(\roman*)]
\item $H^3(Z_{\bar K},\Q_\ell)^{I_K}$ has weights at most $3$;
\item $N:\operatorname{Gr}^W_4H^3\to
       \operatorname{Gr}^W_2H^3(-1)$ is injective;
\item the weight--monodromy conjecture holds for $H^3(Z_{\bar K},\Q_\ell)$.
\end{enumerate}
Here the weight filtration $W$ is supplied by the weight spectral
sequence, and the weight--monodromy statement means that
$N^r:\operatorname{Gr}^W_{3+r}\to
\operatorname{Gr}^W_{3-r}(-r)$ is an isomorphism for every $r\geq0$.
\end{corollary}

\begin{proof}
The preceding proof already makes $N^3$ from weight $6$ to weight
$0$, and $N^2$ from weight $5$ to weight $1$, isomorphisms.
Consequently the only possible weight above $3$ in $\ker N$ is
weight $4$.  Strictness with respect to weights identifies its
weight-$4$ part with the kernel in (ii), proving (i)$\Leftrightarrow$(ii).
Poincar\'e duality on $H^3$, with values in $\Q_\ell(-3)$, pairs
weight $4$ perfectly with weight $2$.  Thus the source and target
in (ii) have equal dimension; injectivity is equivalent to an
isomorphism.  Together with the already verified $r=2,3$ cases,
and the identity for $r=0$, this is exactly (iii).
\end{proof}


\begin{proposition}[The exact surface-pairing obstruction in dimension three]
\label{c:threefold-pairing-kernel}
\status{Proved reduction; the remaining primitive pairing is unresolved}
Let $\mathcal X$ be a projective strictly semistable model of a smooth
projective threefold $X$ over a mixed-characteristic local field with
finite residue field $k$, and let $\ell\ne\operatorname{char}k$.
Write $Y^{(r)}$ for the disjoint union, over $\bar k$, of the
$r$-fold intersections of distinct components of its special fibre.
All cohomology below has coefficients $\Q_\ell$.  Put
\[
 A=H^2(Y^{(1)}),\quad U=H^2(Y^{(2)}),\quad
 B=H^0(Y^{(3)})(-1),\quad Z=H^0(Y^{(4)})(-1).
\]
Use the signed restriction and Gysin maps of the weight spectral
sequence, and define
\[
 \begin{split}
 R&=\operatorname{im}\bigl(\rho_A:A\longrightarrow U\bigr),\\
 K_0&=R\cap\gamma_B
       \ker\bigl(\rho_B:B\longrightarrow Z\bigr),\\
 b(u,v)&=\sum_{S\subset Y^{(2)}}\operatorname{Tr}_S(u_Sv_S)
           \ \in\Q_\ell(-2).
 \end{split}
\]
Then $K_0\subset\operatorname{rad}(b|_R)$, and there is a
Galois-equivariant isomorphism
\begin{equation}
\label{eq:c-threefold-radical-kernel}
 \ker\!\left(
 N:\operatorname{Gr}^W_4H^3(X_{\bar K})(1)
       \longrightarrow\operatorname{Gr}^W_2H^3(X_{\bar K})
 \right)
 \simeq
 \frac{\operatorname{rad}(b|_R)}{K_0}.
\end{equation}
Consequently the first remaining C.4 endpoint assertion is equivalent
to nondegeneracy of the induced form on $R/K_0$.

Choose a relatively ample class $L$.  The same assertion is
equivalent to nondegeneracy of the surface intersection form on
\begin{equation}
\label{eq:c-threefold-primitive-image}
 R_{\rm prim}
 =R\cap
   \ker\bigl(
       L:H^2(Y^{(2)})\longrightarrow H^4(Y^{(2)})(1)
       \bigr).
\end{equation}
In particular the remaining condition involves only a primitive
restriction image in the degree-two cohomology of the
double-intersection surfaces.  It does not require the unknown
nondegeneracy of an additional Gysin image.
\end{proposition}

\begin{proof}
Let
\[
 C=H^4(Y^{(1)})(1),\qquad D=H^2(Y^{(3)}).
\]
The weight spectral sequence of
\citet[Satz~2.10, formulas~(2.10.1)--(2.11)]{c_rapoportzink1982},
with the same normalization as Theorem~\ref{c:threefold-weight-gate},
gives
\[
 \begin{split}
 \operatorname{Gr}^W_4H^3(1)
  &=
  \frac{\ker\{U\oplus Z\longrightarrow C\oplus D,\ 
      (u,z)\longmapsto
      (\gamma_Uu,\rho_Uu+\gamma_Zz)\}}
       {\{(\gamma_Bb,\rho_Bb):b\in B\}},\\
 \operatorname{Gr}^W_2H^3
  &=
  \frac{\ker\{U\oplus Z\longrightarrow D,\
      (u,z)\longmapsto\rho_Uu+\gamma_Zz\}}
       {\{(\rho_Aa+\gamma_Bb,\rho_Bb):
                    a\in A,\ b\in B\}}.
 \end{split}
\]
Here $\gamma_U:U\to C$ and $\gamma_Z:Z\to D$ are the
signed codimension-one Gysin maps.
These displayed identifications use degeneration at $E_2$ by
the different Frobenius weights.  The monodromy map is induced
by the identity of $U\oplus Z$.  We choose the common
restriction/Gysin sign convention in both rows, as is possible
because the differential commutes with $N$.

If a class lies in its kernel, its representative has the form
$(\rho_Aa+\gamma_Bb,\rho_Bb)$.  Subtracting the boundary
$(\gamma_Bb,\rho_Bb)$ leaves $(r,0)$ with $r=\rho_Aa\in R$.
The remaining cocycle conditions are $\gamma_Ur=0$ and
$\rho_Ur=0$.  The latter is automatic from $\rho_U\rho_A=0$.
Two such representatives differ by a domain boundary precisely
when their difference is $\gamma_Bb$ with $\rho_Bb=0$.
Thus the kernel is
\[
 (R\cap\ker\gamma_U)/
       (R\cap\gamma_B\ker\rho_B).
 \]
Poincare duality identifies $\gamma_U$ with the adjoint of
$\rho_A$ under the perfect pairing between $A$ and $C$ and
the pairing $b$ on $U$.  Hence
$\ker\gamma_U=R^\perp$.
The relation $\gamma_U\gamma_B=0$ also shows that the denominator
lies in the radical.  This proves
\eqref{eq:c-threefold-radical-kernel}, including its Galois
equivariance and all its twists.
The equivalence with the C.4 endpoint follows from
Theorem~\ref{c:threefold-weight-gate}.

For the primitive formulation we apply
\citet[Theorem~0.2, complete proofs in \S\S1.1--2.5,
and the dimension-three converse in \S2.7,
author pp.~2--9]{c_msaito2003positivity}.
His dimension-three criterion consists of
\eqref{eq:c-threefold-primitive-image} and the primitive Gysin
image in $H^3(Y^{(1)})$.
The latter condition is automatic: \S2.7 identifies it with
the degree-two monodromy isomorphism, already proved by
restriction to a general hyperplane surface and
weight--monodromy for surfaces.
The other degree-zero and degree-one pairing conditions are
settled by rational positivity and Picard-variety
polarizations, as in the proof of
Theorem~\ref{c:improved-zero-weight}.
The remaining canonical pairing is exactly the intersection
form on the primitive degree-two part of each surface,
up to a common nonzero sign.  This proves the last equivalence.
\end{proof}

\begin{corollary}[Removing the rational divisor restriction image]
\label{c:threefold-divisor-split}
\status{Proved unconditional splitting; the spectral identification uses surface Tate}
Retain the notation of Proposition~\ref{c:threefold-pairing-kernel}.
Let $N_r=\operatorname{NS}(Y^{(r)})_\Q$ for $r=1,2$,
and let $\rho_{\rm NS}:N_1\to N_2$ be signed restriction.
Define the rational space
\[
 D_{\rm prim}=
 \operatorname{im}(\rho_{\rm NS})\cap
 \ker\!\left(
 N_2\longrightarrow
 \bigoplus_{S\subset Y^{(2)}}\Q,\quad
 (d_S)_S\longmapsto(\deg(L_Sd_S))_S
 \right).
\]
We regard $D_{\rm prim}\otimes\Q_\ell(-1)$ as a subspace of
$R_{\rm prim}$ by the divisor cycle map.  Then
\[
 R_{\rm prim}
 =\bigl(D_{\rm prim}\otimes\Q_\ell(-1)\bigr)
   \mathbin{\perp}R_{\rm rem},
 \qquad
 R_{\rm rem}=
 R_{\rm prim}\cap
       \bigl(D_{\rm prim}\otimes\Q_\ell(-1)\bigr)^\perp .
\]
The first summand has a nondegenerate negative-definite
rational intersection form.  Consequently the remaining
threefold assertion is equivalent to nondegeneracy of the
form on $R_{\rm rem}$.

If the divisor Tate conjecture holds for every smooth
projective surface over finite fields of the residue
characteristic, let
$T_r\subset H^2(Y^{(r)})$ be the Frobenius spectral summand
whose eigenvalues are not $|k|$ times roots of unity.
Then
\[
 R_{\rm rem}=\rho_A(T_1)\subset T_2.
\]
Thus, under that separate surface-Tate hypothesis, the
remaining pairing is entirely on the image of the
transcendental degree-two spectral summand.
\end{corollary}

\begin{proof}
All the divisor classes and restriction maps descend to a
finite extension of $k$.  The displayed primitive condition
is a rational matrix of divisor intersection numbers.
Kernels and images commute with extension from $\Q$ to
$\Q_\ell$, so its realized space is exactly the primitive
part of the image of rational divisor classes.
The surface Hodge index theorem makes the intersection
form negative definite on this rational primitive space,
component by component; the same rationality step is
spelled out in
\citet[Claim~4.6 and its full proof, author pp.~12--13]
{c_ito2004threefolds}.
In particular its Gram determinant is nonzero in $\Q$ and
remains nonzero in $\Q_\ell$.
Orthogonal projection onto this subspace, using the inverse
Gram matrix, preserves $R_{\rm prim}$ because the subspace
lies in $R_{\rm prim}$.  Its kernel is $R_{\rm rem}$, which
proves the asserted orthogonal direct sum.  This argument
does not identify the subspace with every algebraic vector
of $R_{\rm prim}$.

For the conditional assertion, the surface-Tate hypothesis
implies the divisor Tate conjecture on all the smooth
projective strata and all their finite constant-field
extensions by
\citet[Theorem~8, pp.~89--90]{BKahn2023Divisors};
the same reduction is used in
Theorem~\ref{thm:B-relative-surface-gate}.
The divisor subspace of $H^2$ has a Galois-stable orthogonal
complement: use the ample Lefschetz intersection form,
whose restriction to divisors is nondegenerate.
The divisor subspace has eigenvalues $|k|$ times roots of
unity.  Its complement has none: otherwise after a finite
constant extension it would have a nonzero eigenvector of
eigenvalue $|k'|$, contradicting Tate surjectivity and the
direct sum.  This also excludes a generalized eigenspace
of that type, since every nonzero generalized eigenspace
contains an eigenvector.  No additional full Frobenius
semisimplicity is used.

Restriction therefore preserves the two spectral blocks.
On the surface target the second block is orthogonal to
every divisor class and hence is primitive, because the
ample class is a divisor.  We obtain
\[
 R_{\rm prim}
   =\bigl(D_{\rm prim}\otimes\Q_\ell(-1)\bigr)
        \mathbin{\perp}\rho_A(T_1).
\]
Comparing with the unconditional orthogonal decomposition
proves the formula for $R_{\rm rem}$.
\end{proof}

\begin{remark}[The geometric content not removed by the pairing criterion]
\label{c:threefold-pairing-scope}
\status{Proved scope clarification; the primitive nondegeneracy is open}
The radical formula is an exact identification of the existing
kernel.  The primitive criterion additionally removes the
nonprimitive Lefschetz pieces and the Gysin-image condition.
It does not prove the resulting nondegeneracy.
In particular, containment of a $\Q_\ell$-subspace in the span
of divisor classes does not allow one to apply rational
Hodge-index positivity to that subspace without establishing
its rational origin.  For example, over $\Q_5$, choose
$\iota^2=-1$ with $\iota\equiv2\pmod5$.
The line $\Q_5(1,\iota)$ is isotropic for $-x^2-y^2$ although
this form is negative definite over $\Q$.
It is stable under the scalar Frobenius $2I$.
This is a counterexample to that linear-algebra inference,
not a geometric counterexample to weight--monodromy.

For the same reason, an isomorphism of a stratum motive with
an abelian motive would not alone identify its geometric
primitive form with a positive polarization.  The conditional
image-projector and positivity inputs in
\citet[Theorem~0.4 and \S\S4.1--4.3]{c_msaito2003positivity}
remain separate hypotheses.  We have not silently supplied
either input from Frobenius purity or from the 2026 theorem
for CM abelian varieties.
The corresponding warning is explicit in
\citet[submitted version, Remark~4.10 and footnote~3,
p.~14]{DMilne2002Polarizations}: equality of motive categories
under a Tate hypothesis does not by itself transfer the
positivity of the geometric forms.  The published-version
Theorem~3.3 and Corollary~3.4 concern abelian varieties.
\end{remark}

\begin{proposition}[A target-only sufficient positivity hypothesis]
\label{c:threefold-target-positivity}
\status{Proved conditional refinement; the general primitive pairing remains unresolved}
Retain the strictly semistable threefold setting and fixed
$\ell\ne\operatorname{char}k$ of
Proposition~\ref{c:threefold-pairing-kernel}, and put
$V=Y^{(1)}$, $S=Y^{(2)}$.
After a finite extension of the finite residue field, we may suppose
that the geometric components are individually defined.
Assume the Hodge standard conjecture
$I(S\times S,L\boxplus L)$ and equality of $\ell$-homological
and numerical equivalence for rational cycles on $V\times S$.
Here the first assumption includes the products of distinct surface
components. Then the primitive restriction image
$R_{\rm prim}$ has nondegenerate intersection pairing.
Consequently weight--monodromy holds on $H^3$ of the generic
threefold. No Lefschetz standard conjecture on $V$ is required
in this sufficient package.

Alternatively, suppose that for each component $V_i$ an inverse to
\[
 L_i:H^2(V_i)\longrightarrow H^4(V_i)(1)
\]
is induced by an algebraic correspondence.
Choose smooth ample surfaces $j_i:T_i\hookrightarrow V_i$ with
classes $m_iL_i$, and put $T=\coprod_iT_i$.
Then the same conclusion follows from
$I(S\times S,L\boxplus L)$ and
homological equals numerical equivalence on $T\times S$.
Only the relevant homogeneous degree-two mixed trace pairings
are needed in either formulation; full equality of the two
equivalences on the displayed products is a sufficient hypothesis.
\end{proposition}

\begin{proof}
Rational algebraic degree projectors over a finite field are available
from Frobenius and
\citet[Theorem~2(1), pp.~76--77]{c_katzmessing1974}.
We therefore work with the homogeneous degree-two homological
objects $H_V,H_S$ and the algebraic restriction morphism
$f=\rho:H_V\to H_S$.
The primitive projector $p$ on the surface target is algebraic:
on each component it is the degree-two projector minus
\[
 x\longmapsto
 \frac{\operatorname{Tr}(xL)}{\operatorname{Tr}(L^2)}\,L.
\]

We first justify the target algebra and its precise pairing.
The Lefschetz standard conjecture for surfaces is known and is
stable under products. For the latter assertion, the normalized
lowering operator on a product is
$\Lambda_1\otimes1+1\otimes\Lambda_2$.
The degree operator and the $\mathfrak{sl}_2$ Casimir have finitely
many rational eigenvalues; polynomial interpolation recovers
each primitive-string level and hence the usual unnormalized
lowering operator. Thus the product lowering is algebraic.
It follows that $B(S\times S)$, and hence $A(S\times S)$, holds.
Combined with the assumed $I(S\times S)$, this gives $D(S\times S)$
as in \citet[\S4.1, pp.~14--15]{c_msaito2003positivity}.
Consequently
\[
 A=\End(H_S)
\]
identifies with the homogeneous corner of the numerical
correspondence algebra. It is finite-dimensional over $\Q$ and
semisimple by \citet[Theorem~1, pp.~448--449]
{c_jannsen1992semisimple}.
Finite-dimensionality over $\Q$ is obtained this way, not inferred
from the dimension of its $\Q_\ell$ realization.

On all of $H^2(S)$ we use the Lefschetz-modified pairing
$\Phi(x,y)=\operatorname{Tr}(x\cup *y)$, with Saito's signed
Lefschetz involution. Its adjoint $a\mapsto a^\dagger$ is algebraic,
and the assumed product positivity gives
\begin{equation}
\label{eq:c-target-positive-trace}
 \operatorname{Tr}(a^\dagger a)>0\qquad(0\ne a\in A),
\end{equation}
by \citet[equation~(4.1.3), p.~15]{c_msaito2003positivity}.
On primitive $H^2$ one has $*=-1$, so $\Phi$ is the negative
of the raw intersection pairing. They have the same radical
there. The raw intersection pairing on all of $H^2(S)$ is
not asserted to have positive adjoint trace.

For algebraic homological correspondences, put
\[
 M=\Hom(H_V,H_S),\qquad N=\Hom(H_S,H_V).
\]
The pairing
\[
 M\times N\longrightarrow\Q,\qquad
 (u,v)\longmapsto\operatorname{Tr}(uv)
\]
is nondegenerate in its first variable under $D(V\times S)$.
Indeed, in the covariant cohomological convention the two
classes have complementary codimensions $3$ and $2$ on
$V\times S$. Their composition trace is their intersection
number. Projecting a complementary class onto its homogeneous
degree-two block does not change its pairing with $u$.
Thus vanishing against every reverse block makes $u$
numerically zero, and the hypothesis makes its cohomological
class zero.

The subspace $J=fN\subset A$ is a right ideal.
Since $A$ is semisimple, $J=eA$ for an idempotent $e\in A$.
There is $g_0\in N$ with $e=fg_0$, so
$\operatorname{im}e\subset\operatorname{im}f$.
For every $g\in N$ one has $(1-e)fg=0$.
The mixed trace pairing therefore gives $(1-e)f=0$, whence
$\operatorname{im}f=\operatorname{im}e$.
This argument requires no semisimplicity of $\End(H_V)$.

The finite additive Karoubi category generated by $H_S$
is semisimple, and its faithful realization is exact because
its exact sequences split. The kernel of
\[
 (1-p)e:\operatorname{im}e\longrightarrow H_S
\]
therefore has an algebraic idempotent projector $r$ and
realizes exactly
$\operatorname{im}f\cap\operatorname{im}p=R_{\rm prim}$.
We have taken the actual intersection, not replaced it
by $p(\operatorname{im}f)$.
Its radical for $\Phi$ is
$\operatorname{im}r\cap\ker r^\dagger$, another algebraic
subobject. If nonzero, it has a nonzero idempotent projector $q$.
The pairing vanishes on $\operatorname{im}q$ with itself,
so $q^\dagger q=0$, contradicting
\eqref{eq:c-target-positive-trace}.
This proves the required nondegeneracy. The dimension-three
criterion of Proposition~\ref{c:threefold-pairing-kernel}
now proves the asserted weight--monodromy statement.

For the alternative package, the algebraic maps
\[
 r_i=m_i^{-1}L_i^{-1}j_{i*}:H^2(T_i)\longrightarrow H^2(V_i)
\]
satisfy $r_i j_i^*=1$ by the projection formula.
In this composition $j_{i*}$ has target $H^4(V_i)(1)$,
which is exactly the source of the supplied inverse.
Thus $r=\bigoplus_i r_i$ is surjective and
$\rho r$ has exactly the image of $\rho$.
Apply the preceding proof to the two surface objects $H_T,H_S$
using the assumed mixed pairing on $T\times S$.
A finite constant-field extension for the components or smooth
sections is harmless: the geometric spaces and maps are the same,
and the Frobenius eigenvalues are replaced by their powers,
which preserves all their weights.
\end{proof}

This sufficient package isolates an input that the known results
for an individual surface do not supply.
$I(S)$ controls divisor classes, whereas the required
positive adjoint trace involves codimension-two correspondence
classes on the fourfold $S\times S$.
Likewise $D(S)$ for divisors does not supply $D(T\times S)$.
A supplied degree-two inverse Hard Lefschetz on a threefold
reduces the mixed product to products of surfaces; it does not
prove either retained product hypothesis.




\begin{proposition}[A simple-spectrum substitute for target positivity]
\label{c:threefold-simple-spectrum}
\status{Proved sufficient condition; the universal threefold gate remains open}
We keep the threefold semistable notation of
Proposition~\ref{c:threefold-pairing-kernel}. After fixing a finite
residue field $\F_q$ over which all the components and the chosen ample
classes are defined, put
\[
 V=Y^{(1)},\qquad S=Y^{(2)},\qquad
 f=\rho:H^2(V,\Q_\ell)\longrightarrow H^2(S,\Q_\ell),
 \qquad P=H^2(S,\Q_\ell)_{\rm prim}.
\]
Here primitivity is imposed separately on each surface component, using
the restrictions of the fixed ample class. Assume that numerical and
$\ell$-adic homological equivalence agree on $V\times S$, and that
the characteristic polynomial of geometric Frobenius on the entire
space $P$ is squarefree. Then the cup-product pairing on
\[
 \operatorname{im}(f)\cap P
\]
is nondegenerate. Consequently the first threefold monodromy condition
in Proposition~\ref{c:threefold-pairing-kernel} holds.

The squarefreeness condition concerns the union of all surface
components, including possible coincidences between their eigenvalues.
It is not enough to impose squarefreeness separately on each component.
We do not assume the standard conjecture of Hodge type on $S\times S$,
or a semisimplicity statement for the source motive. The proof uses
algebraic projectors that are idempotent on cohomology; it does not
require idempotence modulo rational equivalence.
\end{proposition}

\begin{proof}
All algebraic correspondences in this proof have rational coefficients
and are defined over the fixed field $\F_q$. We quotient their actions
by $\ell$-adic homological equivalence. The finite-field degree
projectors are algebraic by
\citet[Theorem~2(1), pp.~76--77]{c_katzmessing1974}; hence every map
below can be represented by an algebraic correspondence homogeneous
in degree two. On a surface component $S_i$, let $h_i$ be the chosen
ample divisor class and let $d_i=\deg(h_i^2)>0$. The operator
\[
 x\longmapsto \frac{\langle x,h_i\rangle}{d_i}h_i
\]
is algebraic and is the orthogonal projector onto its hyperplane line.
Summing these projectors gives an algebraic orthogonal projector
$1-p$ with image
\[
 W=\bigoplus_i\Q_\ell h_i\simeq\Q_\ell(-1)^{\oplus c},
 \qquad\operatorname{im}(p)=P.
\]
Its rational algebraic endomorphism algebra is
$\operatorname{Mat}_c(\Q)$: products of the divisor classes provide
the matrix units. Rational intersection numbers give the converse
inclusion. The zero-surface and zero-primitive cases are immediate;
we treat the nonzero spaces below.

We first obtain the actual intersection with $P$. This step is needed
because $\operatorname{im}(pf)$ need not equal
$\operatorname{im}(f)\cap P$. Put $H_V=H^2(V,\Q_\ell)$ and
\[
 k=(1-p)f:H_V\longrightarrow W.
\]
The mixed numerical-equals-homological hypothesis implies that
\[
 \Hom_{\rm alg}(H_V,W)\times\Hom_{\rm alg}(W,H_V)
 \longrightarrow\Q,\qquad (a,b)\longmapsto\operatorname{Tr}(ab)
\]
is nondegenerate in its first variable. Indeed, embed $W$ into $H^2(S)$
and project back to $W$ by the displayed algebraic maps; a mixed
numerical test for the resulting correspondence on $V\times S$
then restricts to a test in this pairing. The same argument will
apply with $P$ in place of $W$. Only these homogeneous mixed corners
of numerical-equals-homological equivalence are used.

Let $A_W=\End_{\rm alg}(W)=\operatorname{Mat}_c(\Q)$ and form the
right ideal
\[
 J=k\Hom_{\rm alg}(W,H_V)\subseteq A_W.
\]
Because $A_W$ is semisimple, $J=eA_W$ for an idempotent $e$. Since
$e\in J$, we can choose an algebraic map $h:W\to H_V$ such that
$kh=e$. For every algebraic $b:W\to H_V$ we have
$(1-e)kb=0$, and the preceding trace nondegeneracy implies $ek=k$.
It follows that the algebraic operator
\[
 u=1-hk\in\End_{\rm alg}(H_V)
\]
satisfies $u^2=u$, $ku=0$, and $u|_{\ker k}=1$. Thus
$\operatorname{im}(u)=\ker k$ on cohomology. In particular
\[
 a=fu:H_V\longrightarrow P
 \quad\text{has image}\quad
 \operatorname{im}(a)=\operatorname{im}(f)\cap P.
\]
No algebraicity of an inverse Lefschetz map or source image splitting
has been assumed in constructing $u$.

We now determine the rational algebraic endomorphism algebra of $P$.
Write $F$ for its geometric Frobenius and let $d=\dim P$. Its
characteristic polynomial $Q(T)$ belongs to $\Q[T]$: it is obtained
from the rational characteristic polynomial on $H^2(S)$ by removing
the $c$ hyperplane eigenvalues $q$. The squarefreeness hypothesis
implies that $F$ is semisimple with distinct eigenvalues over
$\overline{\Q}_\ell$. Hence
\[
 \End_{\Q_\ell[F]}(P)=\Q_\ell[F],\qquad
 \dim_{\Q_\ell}\Q_\ell[F]=d.
\]
Every base-field algebraic endomorphism $b$ of $P$ commutes with $F$,
so write $b=\sum_{i=0}^{d-1}b_iF^i$ with $b_i\in\Q_\ell$.
For each $0\le j<d$, the trace $\operatorname{Tr}(bF^j)$ is rational:
the algebraic projector $p$ isolates degree two, so the ordinary trace
equals the alternating correspondence trace, which is the rational
intersection number with the diagonal. The matrix
\[
 G_{ij}=\operatorname{Tr}(F^{i+j}),\qquad 0\le i,j<d,
\]
has rational entries and is invertible. In fact it is the trace-pairing
matrix of the finite étale algebra $\Q[T]/Q(T)$, or equivalently the
Gram matrix of the Vandermonde matrix in the distinct Frobenius
eigenvalues. Solving this rational linear system gives every $b_i\in\Q$.
Conversely each rational polynomial in the algebraic operator $F$ is
algebraic. We have proved the equality
\[
 A_P:=\End_{\rm alg}(P)=\Q[F]\simeq\Q[T]/Q(T).
\]
This argument does not assume numerical-equals-homological equivalence
on $S\times S$.

The cup-product pairing on $P$ is nondegenerate, and its adjoint
involution satisfies
\[
 F^\dagger=q^2F^{-1}.
\]
For each complex embedding of a number-field factor of $A_P$, the
image $\alpha$ of $F$ has absolute value $q$, by purity of the smooth
proper surface cohomology
\citep[Corollary~3.3.9, p.~207]{c_deligneweil1980}.
Therefore $q^2/\alpha=\overline\alpha$.
The adjoint involution preserves each number-field factor and is
complex conjugation there. In particular every idempotent of $A_P$
is fixed by the adjoint. More strongly, the ordinary trace satisfies
\[
 \operatorname{Tr}(b^\dagger b)
   =\sum_{\sigma:A_P\to\C}|\sigma(b)|^2>0
 \qquad(b\ne0),
\]
where $P$ has multiplicity one for each embedding. This positivity
is proved for this particular rational algebra; it is not an assertion
that an arbitrary $\Q_\ell$-subspace of a polarized space is
nondegenerate.

Finally apply the right-ideal argument again, now to
\[
 a\Hom_{\rm alg}(P,H_V)\subseteq A_P.
\]
The same mixed trace nondegeneracy produces an idempotent $e'\in A_P$
and an algebraic $b:P\to H_V$ with $e'=ab$ and $e'a=a$.
Thus $\operatorname{im}(a)=\operatorname{im}(e')$. Since
$(e')^\dagger=e'$, its image is orthogonal to its kernel and the
ambient nondegenerate pairing restricts nondegenerately to the image.
This proves the required assertion about
$\operatorname{im}(f)\cap P$. We conclude using the already checked
primitive restriction-image criterion in
Proposition~\ref{c:threefold-pairing-kernel}.

If the mixed equivalence hypothesis is supplied over
$\overline{\F}_q$, it still gives the base-field trace tests used
above. A witnessing rational correspondence is defined over some
finite extension; average its finitely many Frobenius conjugates.
The first member of the trace pairing is defined over $\F_q$, so
the trace is unchanged by this averaging. No extension of the chosen
Frobenius field, and hence no unproved preservation of squarefreeness
under taking Frobenius powers, enters the argument.
\end{proof}

\begin{remark}[Limits of the simple-spectrum criterion]
\label{c:threefold-simple-spectrum-scope}
\status{Proved scope clarification; no unconditional reduction of the universal gate}
The proposition replaces the full target positivity hypothesis of
Proposition~\ref{c:threefold-target-positivity} by one directly testable
condition on a Frobenius characteristic polynomial. It retains the
mixed numerical-equals-homological hypothesis needed to obtain the
algebraic kernel and image projectors. Neither that hypothesis nor the
squarefreeness of the primitive target spectrum has been proved for
all semistable threefolds. Eigenvalues on different components can
coincide, and distinct eigenvalues can become equal after a finite
residue-field extension. We have not produced an alteration or other
transfer that removes these conditions. Thus this is a sufficient
subcase, not movement on the general C.4 assertion.
\end{remark}


\begin{lemma}[A finite exterior-power test for monodromy purity]
\label{c:exterior-weight-criterion}
\status{Proved}
Let $V$ be a finite-dimensional mixed Frobenius module over
$\Q_\ell$, with integral weights and a nilpotent operator
$N:V\to V(-1)$.  Fix an integer $m$.  We use geometric Frobenius $\phi$ and
$\phi N\phi^{-1}=q^{-1}N$.  Suppose
\[
 \sum_w w\dim V_w=m\dim V,
\]
where $V=\bigoplus_wV_w$ is the decomposition by Frobenius
weights; Frobenius need not be semisimple.  Let $h$ be the
largest length of a Jordan string of $N$, with $h=0$ if $V=0$.
Then the following conditions are equivalent:
\begin{enumerate}
\item $V$ is pure of weight $m$ in the weight--monodromy sense;
\item for every $1\leq r\leq h$, all weights of
$\ker(N:\bigwedge^r V\to\bigwedge^r V(-1))$ are at most $rm$.
\end{enumerate}
Here the exterior power is the ordinary exterior power of vector
spaces, and $N$ acts on it as a derivation.  The determinant
hypothesis holds in particular when $V$ has a nondegenerate
Frobenius-equivariant pairing
$V\otimes V\to\Q_\ell(-m)$.
\end{lemma}

\begin{proof}
The generalized Frobenius eigenspaces of a given weight have a
sum defined over $\Q_\ell$: algebraic conjugation preserves the
weight of a Weil number.  They give the asserted direct
decomposition, and the Frobenius relation implies
$N(V_w)\subset V_{w-2}$.  None of these statements requires
Frobenius semisimplicity.

We first construct homogeneous Jordan strings.  Put
$K_L=\ker N\cap\operatorname{im}N^{L-1}$.  These are graded
subspaces with $K_{L+1}\subset K_L$.  Choose a homogeneous basis
$B_L$ of a complement of $K_{L+1}$ in $K_L$, and for every
$b\in B_L$ choose a homogeneous lift $v_b$ with
$N^{L-1}v_b=b$.  The strings
\[
 v_b,Nv_b,\ldots,N^{L-1}v_b
\]
are linearly independent together.  Indeed, in a relation,
apply the largest power of $N$ equal to the distance of a
remaining term from its bottom.  The surviving bottom vectors
belong to the independent union of the $B_L$, so their
coefficients vanish; descending induction finishes the argument.
They have the correct total dimension, since
\[
 \dim K_L=\operatorname{rank}N^{L-1}
               -\operatorname{rank}N^L,\qquad
 \sum_L L(\dim K_L-\dim K_{L+1})=\dim V.
\]
Thus they form a homogeneous Jordan basis.  The individual
strings need not be Frobenius-stable.

A string of length $L$, whose top vector has weight $b$, has
weights $b,b-2,\ldots,b-2(L-1)$ and centre $c=b-L+1$.
The monodromy filtration assigns these positions the indices
$L-1,L-3,\ldots,-L+1$.  Hence weight--monodromy purity of
weight $m$ is equivalent to every string centre being $m$.

The full wedge of the string,
\[
 v_b\wedge Nv_b\wedge\cdots\wedge N^{L-1}v_b,
\]
is nonzero, is killed by $N$, and has weight
\[
 Lb-2(0+\cdots+L-1)=L(b-L+1)=Lc.
\]
Each term in its derivative has a repeated vector or is zero.
Although the wedge line need not be Frobenius-stable, the
nonzero intersection of $\ker N$ with the weight-$Lc$ subspace
is Frobenius-stable.  Condition (2), applied with $r=L$, therefore
gives $c\leq m$ for every string.  The determinant hypothesis is
$\sum_{\text{strings}}L(c-m)=0$, so every centre is $m$.

Conversely, if every centre is $m$, the homogeneous strings
carry the usual $\mathfrak{sl}_2$ structure whose lowering
operator is $N$ and whose Cartan weights are the Frobenius
weights minus $m$.  On $\bigwedge^rV$ the Cartan weights are
the Frobenius weights minus $rm$.  A vector $z$ of Cartan
weight $a$ killed by the lowering operator has $a\leq0$:
if $t$ is maximal with $e^tz\ne0$, the relation
$f e^kz=-k(a+k-1)e^{k-1}z$ at $k=t+1$ gives
$a=-t$.  Here $e$ is nilpotent because it raises the finite
weight grading.  This proves (2), in fact for every exterior power.
Finally, a perfect pairing of the indicated weight pairs
$V_w$ perfectly with $V_{2m-w}$, and gives the determinant
hypothesis by summing dimensions.

This is a direct graded-string argument.  The related
determinant criterion and invariance of purity under finite
extension are recorded, with their proofs, in
\citet[Lemma~1.4(1), (2), (5), author pp.~6--7]{c_tayloryoshida2007}.
\end{proof}

\begin{theorem}[The universal trait endpoint and weight--monodromy]
\label{c:universal-trait-weight-equivalence}
\status{Proved equivalence; the universal assertions remain unresolved}
Fix a residue characteristic $p$ and a prime $\ell\ne p$.
We consider finite extensions $K/\Q_p$, with finite residue
field, and geometric $\ell$-adic cohomology.  The following
universal assertions are equivalent:
\begin{enumerate}
\item Weight--monodromy holds for $H^j(Z_{\bar K},\Q_\ell)$
for every smooth projective $Z/K$ and every $j\geq0$.
\item For every smooth projective $j$-fold $Z/K$,
$H^j(Z_{\bar K},\Q_\ell)^{I_K}$ has weights at most $j$.
\item The assertion in (2) holds for the generic fibres of
projective strictly semistable models over $\mathcal O_K$.
\end{enumerate}
In (2) the stated space is the closed geometric stalk of
$u_*H^j(Z_{\bar K},\Q_\ell)$ for the open immersion
$u:\Spec K\hookrightarrow\Spec\mathcal O_K$.  Thus it is exactly
the dimension-normalized, degree-zero trait endpoint.

For an individual $j$-fold, after a finite extension making its
inertia action unipotent, it suffices to test the bound
$rj$ on $\ker N$ in
$\bigwedge^r H^j(Z_{\bar K},\Q_\ell)$ for $1\leq r\leq j+1$.
Each such test is a direct-summand test in the middle cohomology
$H^{rj}(Z^r_{\bar K},\Q_\ell)$.  A failure of weight--monodromy
therefore yields a trait-endpoint failure on one of these powers,
or on a projective semistable alteration of that power.
\end{theorem}

\begin{proof}
We give the transfers explicitly.  Finite extension multiplies
$N$ by a nonzero scalar if the tame character is renormalized
and replaces Frobenius weights by the same weights relative to
the larger residue cardinality.  For the latter assertion one
may pass to a common normal extension killing the finite
Weil--Deligne inertia, and choose a positive power of a
Frobenius lift lying in its Weil group.  A further power of the
new Frobenius lift differs from it by inertia acting trivially
in the Weil--Deligne representation.  Taking roots of the
corresponding absolute values proves the assertion.
Consequently weight--monodromy is invariant under finite
extension; this also verifies directly
\citet[Lemma~1.4(2)]{c_tayloryoshida2007}.

For a projective semistable model, the spectral sequence used
in Proposition~\ref{c:improved-zero-weight} has
\[
 E_1^{a,b}=
 \bigoplus_{s\geq\max(0,-a)}
 H^{b-2s}(Y^{(a+2s)}_{\bar k},\Q_\ell)(-s).
\]
Here $Y^{(t)}$ denotes the disjoint union of the $(t+1)$-fold
intersections.  On total degree $j=a+b$, a nonzero summand
requires
$b\geq2s\geq\max(0,2b-2j)$, hence $0\leq b\leq2j$.
The terms have weight $b$, so all abutment weights lie in
$[0,2j]$.  Since $N$ lowers weights by two,
$N^{j+1}=0$.  The spectral-sequence and monodromy statements
are precisely \citet[Satz~2.10 and (2.11), pp.~40--41]{c_rapoportzink1982}.
For arbitrary smooth projective $Z$, the same conclusions
follow after a projective semistable alteration and finite
extension: pullback is split by pushforward divided by the
generic degree, so its cohomology is a Frobenius- and
$N$-stable direct summand.  This is the alteration and trace
reduction already checked in Proposition~\ref{c:improved-zero-weight},
using \citet[Theorem~6.5, pp.~83--86]{c_dejong1996}.

For a $j$-fold, Poincare duality pairs $H^j$ perfectly into
$\Q_\ell(-j)$.  Lemma~\ref{c:exterior-weight-criterion}
therefore applies with largest string length at most $j+1$.
The required ordinary exterior power is a Galois direct
summand of the Kunneth factor
$H^j(Z)^{\otimes r}\subset H^{rj}(Z^r)$.  Explicitly its
projector on that factor is
\[
 \frac1{r!}\sum_{\sigma\in\mathfrak S_r}
       (\operatorname{sgn}\sigma)^{j+1}\sigma^*.
\]
The geometric permutation $\sigma^*$ already contributes
the Koszul factor $(\operatorname{sgn}\sigma)^j$.
Thus we use the geometric antisymmetrizer for even $j$ and
the geometric symmetrizer for odd $j$.  These are direct
summands of Galois representations; no algebraic
Chow--Kunneth projector is assumed.  The variety $Z^r$ has
dimension $rj$, exactly as required by the normalization.
The rational denominator $r!$ causes no exception at $\ell=2$.

Assume (3).  Given a $j$-fold $Z$, pass first to a finite
extension on which its inertia action is unipotent.  For each
$1\leq r\leq j+1$, choose, after a further finite extension,
a projective strictly semistable alteration
$\pi_r:W_r\to Z^r$ on generic fibres over some $L_r$.
We allow the fields $L_r$ to depend on $r$; no ramified
base change of a semistable integral model is presumed to
remain regular.
The maps
\[
 H^{rj}(Z^r)\xrightarrow{\pi_r^*}H^{rj}(W_r)
 \xrightarrow{(\deg\pi_r)^{-1}\pi_{r*}}H^{rj}(Z^r)
\]
have composite the identity and commute with inertia and
monodromy.  On a semistable representation, inertia invariants
are exactly $\ker N$, because its inertia action is
$\exp(t_\ell(\tau)N)$ and $t_\ell$ has infinite image.
Assumption (3) consequently bounds the weights of
$\ker N$ on each indicated exterior summand over $L_r$ by $rj$.
The original inertia action was already unipotent, so passage
to $L_r$ leaves this kernel unchanged.  Frobenius-power
comparison descends the same bound separately for every $r$.
The lemma proves weight--monodromy for $H^j(Z)$.
Finite extension does not change that conclusion.

We next reduce arbitrary cohomological degrees to this middle
case, retaining actual split summands.  Let $\dim X=n$.
If $n\geq j$, choose a smooth complete intersection
$i:Z\hookrightarrow X$ of dimension $j$, for a very ample
class $L$.  Weak Lefschetz and hard Lefschetz give the
Galois-equivariant retraction
\[
 H^j(Z)\xrightarrow{i_*}
 H^{2n-j}(X)(n-j)
 \xrightarrow{(L^{n-j})^{-1}}H^j(X),
\]
whose composite with $i^*$ is the identity.  The inverse here
is an inverse of cohomological Galois modules; its algebraicity
is unnecessary.  If $n<j$, replace $X$ by
$X\times\mathbf P^{j-n}$ and take the $H^0$ Kunneth factor
of projective space, again an actual split summand.
Pure weight--monodromy modules are closed under split
submodules: the canonical monodromy filtration is additive
on a direct sum, so each summand inherits the pure graded
pieces.  Thus (3) implies (1).

If (1) holds, a monodromy string of centre $j$ has bottom
weight $j-L+1\leq j$.  After a semistable extension this proves
the bound on $\ker N$.  The original $I_K$-invariants form
a subspace of these invariants; comparison by Frobenius
powers, on which the original inertia acts trivially, gives
the same upper bound over $K$.  Hence (1) implies (2).
Finally (2) implies (3) by restriction of its range.
\end{proof}

\begin{example}[One exterior test can miss a failure]
\label{c:single-endpoint-insufficient}
\status{Proved formal counterexample; no geometric realization asserted}
Take $K=\Q_2$, $\ell=3$, $q=2$, and an unramified Weil
representation over $\Q_3$ with basis
$u_6,v_4,x_4,y_2$, subscripts denoting weights.  Put
\[
 \phi=\operatorname{diag}(8,4,4,2),\qquad
 N(u_6)=v_4,\quad N(x_4)=y_2,\quad N(v_4)=N(y_2)=0.
\]
Its two string centres are $5$ and $3$, so it is not
weight--monodromy pure of weight $4$.  Nevertheless
its determinant has weight $16=4\dim V$ and
$\ker N=\langle v_4,y_2\rangle$ has weights at most $4$.
It even has the nondegenerate symmetric pairing into
$\Q_\ell(-4)$ specified by
\[
 \langle u_6,y_2\rangle=1,\qquad
 \langle v_4,x_4\rangle=-1
\]
and symmetric counterparts, all other entries being zero.
The identities $\phi N\phi^{-1}=2^{-1}N$,
$N^{\mathsf T}Q+QN=0$, and
$\phi^{\mathsf T}Q\phi=2^4Q$ are immediate matrix identities.
In $\bigwedge^2V$, the nonzero vector $u_6\wedge v_4$
is killed by $N$ and has weight $10>8$.
The second exterior test detects the failure.
The exact rational-matrix certificate
\path{code/part_c/check_c4_exterior_witness.py}
and its companion
\path{code/part_c/check_c4_exterior_witness_output.json}
verify these identities and the exterior-power calculation.
This is a counterexample only to the proposed linear-algebra
implication from the single invariant bound and self-duality.
It is not a smooth projective fourfold or a counterexample
to geometric weight--monodromy.
\end{example}

\begin{remark}[Scope of the trait criterion]
\label{c:trait-exterior-scope}
\status{Proved scope clarification; universal C.4 remains incomplete}
The theorem identifies the universal middle-degree trait
endpoint with weight--monodromy and gives a finite power
witness to a failure.  It does not prove either universal
assertion.  Its bound $r\leq j+1$ uses the actual nonnegative
cohomological degree $j$; it is not a bound in terms of the
weight of an arbitrary Tate-twisted motive.

The finite inertia quotient can remove a vector from the
invariants at the original place, which is why the witness
is taken after a finite semistable extension.  All local fields
are quantified in the universal assertion, so this extension
is permitted.  The argument concerns the stated traits.
For a general morphism in C.4, the component and curve-selection
argument in Proposition~\ref{c:improved-zero-weight} is an
additional required transfer; no general sheaf bound is
inferred merely by naming the trait equivalence.
\end{remark}

\begin{proposition}[The lower-endpoint gate over a trait]
\label{c:lic-endpoint}
\status{Proved exact equivalence; the general mixed-characteristic vanishing is not asserted}
Let $R$ be a complete discrete valuation ring with finite residue field $k$, fraction field $K$, and $\ell\ne\operatorname{char}(k)$. Let $\mathcal Z/R$ be regular, proper, flat and generically smooth, with generic fibre $Z$. For $u:\Spec K\hookrightarrow\Spec R$, the assertion that
\[
 (u_*H^j(Z_{\bar K},\Q_\ell))_{\bar k}
       =H^j(Z_{\bar K},\Q_\ell)^{I_K}
 \quad\text{has weights at most }j
\]
is equivalent to surjectivity of the specialization map
\[
 H^j(\mathcal Z_{\bar k},\Q_\ell)
       \longrightarrow H^j(Z_{\bar K},\Q_\ell)^{I_K}.
\]
Thus, when $\dim Z=j$, improving the universal $i=0$ coefficient bound from $2j$ to $j$ includes exactly this local invariant cycle assertion.
\end{proposition}

\begin{proof}
\citet[Theorem 10.1(a), pp.~77--80 of the opened author version]{c_flachmorin2012topos} proves that the image of specialization is precisely the weight-at-most-$j$ part of the invariant space. The source has weights at most $j$, since the special fibre is proper. To check the step in that theorem, localization and Hochschild--Serre reduce surjectivity onto that weight part to the absence of weights at most $j$ in $H^{j+1}_{\mathcal Z_{\bar k}}(\mathcal Z_{R^{\rm sh}},\Q_\ell)$. Absolute purity and duality identify this supported group with $H^{2d+1-j}(\mathcal Z_{\bar k},\Q_\ell)^\vee(-d-1)$, where $d=\dim Z$; properness bounds its weights below by $j+1$. This proves the cited image assertion and both implications. The equivalence does not prove the surjectivity in general mixed characteristic; the known low-degree and equal-characteristic cases are recorded in \citet[\S3.1, p.~494]{c_esnaultwittenberg2016}.
\end{proof}

\begin{proposition}[The lower endpoint for abelian-type coefficients]
\label{c:abelian-invariant-weights}
\status{Proved}
Let $K$ be a complete discretely valued field with finite residue field of
characteristic $p$, and let $\ell\ne p$.  Suppose that $V$ is the
$\ell$-adic realization of a pure weight-$j$ motive which, after a finite
extension of $K$, is a direct summand of tensor constructions on the
$H^1$ of abelian varieties, their duals and Tate twists.  Finite-image
Artin factors are allowed.  Then the eigenvalues of geometric Frobenius
on $V^{I_K}$ have weights at most $j$.  In particular this holds for the
local realizations of pure absolute-Hodge or motivated motives of abelian
type over number fields.  The bound is attained for good reduction.
\end{proposition}

\begin{proof}
We give the weight--monodromy argument, including the tensor step.
After a finite extension we may suppose that all the finitely many
abelian generators have semistable reduction and all the Artin factors
are trivial.  Semistable reduction is SGA~7~I, IX, Theorem~3.6, printed
p.~351.  For one such abelian variety write
$U=T_\ell(A)\otimes\Q_\ell$, $V_0=U^{I_K}$, and let $W$ be the
toric part of its Tate module.  The N\'eron special fibre gives
\[
 0\longrightarrow W\longrightarrow V_0
       \longrightarrow T_\ell(B)\otimes\Q_\ell\longrightarrow0,
\]
where $B$ is the abelian quotient of the identity component.  The exact
sequence and the identification of the fixed part are IX, formula
(2.1.8), p.~328, and Proposition~2.2.5, pp.~332--333.  A torus has Tate
module of weight $-2$, and $T_\ell(B)$ has weight $-1$.  A polarization
and the orthogonality theorem IX, Theorem~2.4, p.~338, give
$W=V_0^\perp$ in the semistable case, as stated in Proposition~3.5,
p.~350.  Consequently $U/V_0$ is dual to $W$ with values in
$\Q_\ell(1)$ and has weight $0$.

Corollary~3.5.2, p.~351, says that inertia factors through the tame
$\Z_\ell(1)$ quotient and that $(g-1)^2=0$ for a generator $g$.
Thus its logarithm $N$ satisfies $N^2=0$ and $\ker N=V_0$.
The polarization is inertia invariant, so $N$ is skew adjoint and
$\operatorname{im}N=(\ker N)^\perp=W$.  It induces an isomorphism
\[
 N:U/V_0\xrightarrow{\sim}W(-1).
\]
All these statements are rational statements; no assumption on an
integral monodromy pairing is required.  The filtration
$0\subset W\subset V_0\subset U$ is therefore the monodromy
filtration centred at weight $-1$, with graded weights $-2,-1,0$.
Dualizing proves weight--monodromy, centred at $1$, for $H^1(A)$.
This derives exactly the form we use from the cited original pages
of SGA~7 \citep{c_sga7}.

For completeness, pure Weil--Deligne representations are stable under
the indicated tensor operations.  The monodromy filtration is
characterized by $NM_a\subset M_{a-2}(-1)$ and the isomorphisms
$N^a:\operatorname{gr}_a^M\xrightarrow{\sim}
\operatorname{gr}_{-a}^M(-a)$ for $a\geq0$.  On a tensor product
the convolution filtration has these properties: decompose the
nilpotent modules into $\mathfrak{sl}_2$ strings and apply the
Clebsch--Gordan decomposition to their tensor products.  Its graded
terms have the sums of the two graded weights.  Thus the centre
weights add.  Duality reverses the indices and weights, and a Tate
twist shifts the centre weight by an even integer.  An idempotent
commuting with the Weil action and $N$ preserves the unique
monodromy filtration; its image on every graded term is pure of the
same weight.  This proves stability under direct summands as well.
We apply these facts to the weight-$j$ part of the tensor presentation
of the given pure motive.  Its Weil--Deligne representation is pure
of weight $j$.

For a monodromy string centred at $j$, the vector killed by $N$ is
its lowest-weight vector.  Equivalently, $N$ is injective on every
positive-index graded term of its monodromy filtration.  Hence
$\ker N$ has weights at most $j$.  The original inertia invariants
are a Frobenius-stable subspace of $\ker N$ after the finite
extension.  If that extension has residue degree $f$, any lift of
its geometric Frobenius differs from the $f$th power of the original
Frobenius by an element of $I_K$.  These two operators therefore
agree on $V^{I_K}$.  The weight inequality is unchanged
on taking these powers, and therefore descends to $K$.
For a pure motive with good reduction the invariant space is the
whole pure weight-$j$ representation, proving the stated sharpness.
\end{proof}


The preceding results prove the requested upper-weight inequality, the stronger displayed uniform geometric-level bounds, and the manuscript's margin question in the explicit finite-type setting. They also settle the top coefficient endpoint and, for abelian-type coefficients over a trait, the lower endpoint. Proposition~\ref{c:lic-endpoint} identifies the remaining lower-endpoint problem rather than deriving it from a theorem about ordinary direct-image weights. Proposition~\ref{c:broader-coefficients} also removes the coefficient-dimension ambiguity for the general polygon and the manuscript's stated margin bound. Only the stronger reflected margin polygon retains the dimension-normalized hypothesis in our proof.

\subsection{C.5: good reduction in codimension one}

The quantifier on primes is the one in 12:140--142: set $S_\ell=\Spec R[1/\ell]$ and test $T_\ell$ at its height-one points. It does not require the ordinary $p$-adic Tate module to be unramified at a point of residue characteristic $p$.

\begin{theorem}[The positive purity range]
\label{c:vz-positive}
\status{Known, with its application proved}
Let $(R,\mathfrak m)$ be a two-dimensional regular local ring of mixed characteristic $(0,p)$ and let $U=\Spec R\setminus\{\mathfrak m\}$. Every abelian scheme over $U$ extends uniquely over $R$ if $p\notin\mathfrak m^p$. In particular the sufficient absolute ramification bound is $e\leq p-1$, including equality. More generally the same conclusion holds under the complete faithfully flat quotient condition of Vasiu--Zink Theorem 3.

If an abelian variety over $K=\operatorname{Frac}(R)$ has good reduction at every height-one prime, these conclusions apply to it. Equivalently, at each such prime one may use the prime-to-residue-characteristic unramifiedness test of the N\'eron--Ogg--Shafarevich criterion.
\end{theorem}

\begin{proof}
The first assertion is \citet[Corollary 4]{c_vasiuzink2010}. The more general condition is the existence of a faithfully flat complete regular local $R$-algebra $\widehat R$, of the same dimension and with algebraically closed residue field $k$, admitting a surjection to
\[
 W(k)[[T_1,T_2]]/(p-h),\qquad
 h\notin(p,T_1^p,T_2^p,T_1^{p-1}T_2^{p-1}).
\]
Their Theorem 3 gives extension and uniqueness under this condition. Since $e$ is the $\mathfrak m$-adic order of $p$, $p\notin\mathfrak m^p$ means $e\leq p-1$.

To apply this to a generic abelian variety, first spread it out over a nonempty open of $\Spec R$. There are only finitely many height-one points outside this open: each lies over a height-one component of its closed complement. A model over $R_{\mathfrak p}$ spreads out to a neighbourhood of $\mathfrak p$, since it and its group structure have finite presentation. The models glue by uniqueness of homomorphism extension over normal bases; a precise form is \citet[Proposition 3.3.1.5]{c_lan2021}. They give an abelian scheme on $U$. The purity theorem then applies. The equivalence of good reduction and the indicated unramifiedness is SGA~7~I, IX, Corollary 2.2.9 \citep{c_sga7}; the criterion is pointwise over the height-one discrete valuation rings, with $\ell$ invertible there. This proves the transfer required in the question, rather than silently assuming a model already existed on $U$.
\end{proof}

\begin{theorem}[An explicit abelian-surface counterexample at $p=2$]
\label{c:explicit-bad-reduction}
\status{Proved refutation of the unrestricted criterion}
Let $k=\overline\F_2$ and
\[
 R=W(k)[[u,v]]/(2-uv),\qquad K=\operatorname{Frac}(R),\qquad
 \mathfrak m=(u,v).
\]
Consider the elliptic schemes
\begin{align*}
 E_1:&\quad Y^2Z+(u+v^2)XYZ+YZ^2=X^3,\\
 E_2:&\quad Y^2Z+(v+u^2)XYZ+YZ^2=X^3,
\end{align*}
with origins $[0:1:0]$, and put $A=E_1\times_R E_2$. The sections
\[
 P_v=[-v:1:v^3]\in E_1(R),\qquad
 P_u=[-u:1:u^3]\in E_2(R)
\]
have order two on the generic fibre. Let $C_K\subset A_K$ be the order-two subgroup generated by $(P_v,P_u)$, and define the abelian surface
\[
 B_K=A_K/C_K.
\]
Then $B_K$ has good reduction at every height-one prime of $R$, hence its $T_\ell$ is unramified at every height-one point of $S_\ell=\Spec R[1/\ell]$ for every prime $\ell$, but it does not extend to an abelian scheme over $R$. The schematic closure of $C_K$ in $A$ is explicitly nonflat.
\end{theorem}

\begin{proof}
The ambient ring $W(k)[[u,v]]$ is regular of dimension three, and $2-uv$ has nonzero linear term $2$ in its maximal ideal. The quotient is regular local of dimension two, with maximal ideal $(u,v)$ and residue field $k$. It is consequently a domain of mixed characteristic; $2=uv$ has order two in $\mathfrak m$. For $E_a:Y^2Z+aXYZ+YZ^2=X^3$, the discriminant is $a^3-27$. For both displayed values of $a$ it is $1$ modulo $\mathfrak m$, hence a unit. Thus both equations define smooth projective elliptic schemes over $R$.

For $s=v,r=u$, or $s=u,r=v$, we have $rs=2$ and $a=r+s^2$. Substitution shows that $[-s:1:s^3]$ lies on $E_a$. The inverse formula sends its $Y$-coordinate to $-1+as-s^3=-1+rs=1$, so the section equals its inverse. Its generic fibre is not the origin, hence has exact order two.

Define two finite flat group schemes
\[
 G_u=\Spec R[x]/(x^2+vx),\qquad
 G_v=\Spec R[z]/(z^2+uz),
\]
with laws $x*x'=x+x'+uxx'$ and $z*z'=z+z'+vzz'$. These are integral Hopf laws: for example $1+ux$ is group-like and $(1+ux)^2=1$, by $uv=2$ and the relation. The inverse is $x\mapsto x$; direct substitution verifies the identity. The maps
\[
 G_u\longrightarrow E_1,\quad x\longmapsto[x:1:-v^2x],\qquad
 G_v\longrightarrow E_2,\quad z\longmapsto[z:1:-u^2z]
\]
are closed immersions. On the chart $Y=1$, the coordinate $X/Y$ pulls back to $x$, respectively $z$, proving the closed-immersion assertion there; their finite source makes the images closed in the entire elliptic scheme. Their images satisfy the elliptic equations: in the first case the substitution is
\[
 (-x+v-uv^2)(x^2+vx)+v^2x(uv-2)=0,
\]
and the other case interchanges $u,v$. Generically they are the subgroups $\{0,P_v\}$ and $\{0,P_u\}$. Equality of the group-law morphisms then holds over $R$ because the source and its square are flat over the domain $R$ and the target is separated. Put $H=G_u\times G_v\subset A$.

The constant order-two group may be written
\[
 D=\Spec R[t]/(t^2+t),\qquad t*t'=t+t'+uvtt'.
\]
There is a group homomorphism $\alpha:D\to H$ defined on coordinate rings by $x\mapsto vt$, $z\mapsto ut$. The polynomial relations and the coproducts agree directly. On $D(u)$ or $D(v)$ one of these two coordinate maps recovers $t$, so $\alpha_U$ is a closed immersion. These two opens cover $U$. On the generic fibre it is exactly $C_K$.

Its schematic image has coordinate algebra
\[
 R[ut,vt]\subset R[t]/(t^2+t),\qquad
 R[ut,vt]=R\oplus\mathfrak m t
\]
as an $R$-module. To verify this equality, every nonconstant monomial in $ut,vt$ reduces using $t^2=-t$ to an element of $\mathfrak m t$, while $ut,vt$ generate that summand. This algebra is torsion-free and has generic rank two, hence it is the schematic closure of $C_K$. Its closed fibre has dimension
\[
 \dim_k\bigl((R\oplus\mathfrak m t)/\mathfrak m(R\oplus\mathfrak m t)\bigr)
 =1+\dim_k(\mathfrak m/\mathfrak m^2)=3.
\]
A finite flat module over the local ring $R$ would be free of its generic rank, and would have closed-fibre dimension two. Thus the closure is not flat.

The quotient $B_U=A_U/\alpha_U(D_U)$ is an abelian scheme, since its kernel is finite locally free. Every height-one point is in $U$, which proves good reduction there. Suppose that $B_U$ extended to an abelian scheme $B/R$. The quotient isogeny $A_U\to B_U$ and its inverse up to multiplication by $2$ extend uniquely to homomorphisms $f:A\to B$ and $g:B\to A$ by the normal-base extension theorem cited above. Their composites equal multiplication by $2$, since they do so on the dense generic fibre. Thus $f$ has finite kernel, contained in $A[2]$, and is fibrewise a surjective isogeny. It is finite and flat: properness and finite fibres give finiteness, and fibrewise flatness of isogenies between abelian varieties, together with flatness over $R$, gives flatness by the fibre criterion. Therefore $\ker f$ is finite flat. Its generic fibre is $C_K$, so torsion-freeness identifies it with the schematic closure just computed. This contradicts the closed-fibre dimension three. No such $B$ exists.

All substitutions, discriminants, inverse formulas and Hopf-map polynomial identities have been checked over $\Z$ by \texttt{check\_explicit\_bad\_reduction.py}; no calculation in that certificate inverts $2$. The example is an explicit instance of the obstruction in \citet[Lemma 27 and Theorem 28(ii)]{c_vasiuzink2010}, including the prime $2$.
\end{proof}

\begin{remark}[What fails to extend]
\label{c:subgroup-versus-object}
\status{Proved clarification}
The abstract finite group $D_U$ already extends as $D$. It is its closed immersion into $A_U$, and consequently the exact sequence defining the isogeny, that cannot extend as a finite flat subgroup. Indeed finite locally free group schemes on a punctured regular local surface extend as objects by extending their coordinate bundles to reflexive modules, which are free in dimension two; multiplication and comultiplication extend by Hartogs. Exactness and closed-immersion conditions need not survive. This is the precise content of the finite-subgroup reduction in R6, rather than a failure to extend the abstract group. The source proof of Vasiu--Zink Lemma 27 reverses the Cartier-dual arrows in its printed embedding sentence; our displayed maps avoid that typing error. Lead C5 is therefore valid as an identification of the extension obstruction, while the unrestricted good-reduction claim is false.
\end{remark}

\subsection{C.6: what independence of \texorpdfstring{$\ell$}{ell} actually supplies}

\begin{proposition}[The degree-one and semistable abelian-type cases]
\label{c:ell-independence-cases}
\status{Proved in the stated cases}
Let $K$ be a number field, $v$ a finite place of residue characteristic $p$, and $\ell\ne p$.
For $M=h^1(A)$, unramifiedness of $M_\ell$ at $v$ is independent of $\ell$ and is equivalent to good reduction of $A$. It follows for every tensor construction and motivated summand of abelian varieties already having good reduction at $v$.

More generally suppose that $A/K$ has semistable reduction at $v$, that its Galois representations take values in its Mumford--Tate group $G$ over $K$, and that a motivated abelian-type object $M$ in the category it generates corresponds to a rational representation $r:G\to\operatorname{GL}(V)$. Then unramifiedness of $M_\ell$ at $v$ is independent of $\ell\ne p$. The same conclusion holds for potentially CM motives without a semistability assumption.
\end{proposition}

\begin{proof}
For the first assertion we use SGA~7~I, IX, Proposition 2.2.5 and Corollary 2.2.9. In its notation, the dimension of inertia invariants is $\mu+2\alpha$, whereas the full dimension is $2(\mu+\lambda+\alpha)$. Equality holds if and only if $\mu=\lambda=0$, which is exactly good reduction and is a geometric condition independent of $\ell$. Smooth proper base change then makes all tensor constructions unramified, and restriction to an invariant direct summand preserves the identity action.

For the semistable assertion, \citet[Theorem 1.2]{c_kisinzhou2025} supplies a common $G$-valued Weil--Deligne representation, up to $G(\C)$-conjugacy, for every prime, including $p$. Apply the rational representation $r$. In the semistable case the finite inertial type is trivial and the actual inertia action is $\exp(t_\ell(\sigma)\,dr(N))$. It is trivial exactly when $dr(N)=0$, since the tame character has infinite image and the exponential of a nonzero nilpotent cannot be identically one in characteristic zero. This vanishing is invariant under conjugation and extension of the coefficient field, proving the assertion. No descent from a ramified field extension has been used.

For potentially CM motives, the tannakian identification with the Taniyama group gives their realizations from one representation of the Weil group; the strict-compatibility statement is \citet[Remark 3, p.~262]{c_delignetaniyama1982}. On inertia the algebraic Hecke-character factors have a common finite-order part and the Artin factors are defined over a fixed number field of coefficients. Triviality of those matrices is unchanged by any embedding of that field into $\overline\Q_\ell$. Hence unramifiedness is independent of $\ell\ne p$ in this case as well.
\end{proof}

\begin{lemma}[A generator with potentially good reduction]
\label{c:potentially-good-generator}
\status{Proved}
Let $M$ be a rational motive of abelian type over a number field $K$, in either
the absolute-Hodge or motivated category.  Let $v$ have residue characteristic
$p$.  If $M_{\ell_0}$ is potentially unramified at $v$ for one $\ell_0\ne p$,
then, after a finite Galois extension $E/K$, the object $M_E$ belongs to the
tensor category generated by finitely many abelian varieties over $E$ which
have good reduction at every place above $v$.
\end{lemma}

\begin{proof}
We first work over $\overline\Q$, using the chosen complex embedding.
The weak Mumford--Tate lifting construction of
\citet[Theorem 2.12 and Corollary 2.15, pp.~15--20]{c_noot2006} gives an
abelian variety $B$ with a central isogeny
\[
 \pi^{\mathrm{der}}:\operatorname{MT}(B)^{\mathrm{der}}
       \longrightarrow\operatorname{MT}(M)^{\mathrm{der}}
\]
identifying their adjoint Hodge homomorphisms.  Here descent to
$\overline\Q$ is part of the argument, not an assumption: start with an
abelian generator $A/\overline\Q$ of $M$, take its Mumford--Tate decomposed
weak lift $B_0$, and descend $B_0$ by Noot's Proposition 3.1 and Corollary 3.2,
pp.~20--21.  The proof of Corollary 2.15 selects the factors of $B_0$ whose
$\Q$-simple derived groups survive in $\operatorname{MT}(M)$.  These factors
descend up to isogeny because their rational endomorphism idempotents descend
under an algebraically closed field extension.  Rational motives only require
this isogeny-level descent.

The proof of Noot's Proposition 3.4, pp.~22--23, applies with the Hodge
realization of $M$ in place of $H^1(A)$ and gives
\begin{equation}
\label{c:weak-lift-generates}
 M_{\overline\Q}\ \in\
 \langle h^1(B),\text{CM motives}\rangle^{\otimes}.
\end{equation}
To make this extension of the source statement explicit, use the common
group formed from the central tori and $\operatorname{MT}(B)^{\mathrm{der}}$.
Its representation on $\omega_B(M)$ is generated by the faithful
representation on $H^1(B)$ and representations of its torus quotient.
For an irreducible constituent $W_2$ of $\omega_B(M)$, choose an irreducible
constituent $W_1$ of a tensor construction on $H^1(B)$ and a torus
representation $W_3$ with $W_2\subset W_3\otimes W_1$.
The nonzero contraction $W_2\otimes W_1^\vee\to W_3$ shows that $W_3$ itself
is an actual polarizable Hodge subquotient, of CM type, with integral
Hodge degrees.  Full faithfulness of Betti realization for abelian motives
and the realization of CM Hodge structures by CM abelian varieties prove
\eqref{c:weak-lift-generates}.  In the motivated category the same step uses
Andr\'e's Theorem 0.6.2.  Only finitely many CM varieties and morphisms are
needed for this particular $M$, so all these data descend to a finite
extension of $K$.

Put $H=\operatorname{MT}(M_{\overline\Q}\oplus h^1(B))$.
Its two projections have the same adjoint Hodge datum.  The Hodge
homomorphism defining $H$ lies in the fibre product over that adjoint group,
and projects onto both Mumford--Tate groups.  The kernel of the projection
$H\to\operatorname{MT}(M)$ lies in $\{1\}\times Z(\operatorname{MT}(B))$,
so is central in $H$ and meets its semisimple derived group in a finite
group.  Thus $H^{\mathrm{der}}\to\operatorname{MT}(M)^{\mathrm{der}}$ is a
central isogeny.  After a finite extension, the arithmetic realizations of
$M\oplus h^1(B)$ take values in $H(\Q_\ell)$ for every $\ell$: finitely many
absolute Hodge tensors cut out $H$, and one extension defines all of them.

Fix a place above $v$ and take a further finite local extension over which
$M_{\ell_0}$ is unramified and $B$ is semistable.  An inertia element acts on
$M_{\ell_0}\oplus H^1(B,\Q_{\ell_0})$ as
\[
 u=(1,\rho_{B,\ell_0}(\tau))\in H(\Q_{\ell_0}),
\]
and $u$ is unipotent by the semistable monodromy theorem.  Its image in the
torus $H/H^{\mathrm{der}}$ is trivial, so $u\in H^{\mathrm{der}}$.
Its image in $\operatorname{MT}(M)^{\mathrm{der}}$ is trivial, and hence
$u$ lies in a finite kernel.  A finite-order unipotent in characteristic
zero is the identity.  Thus $H^1(B,\Q_{\ell_0})$ is unramified, and
N\'eron--Ogg--Shafarevich gives good reduction of $B$ over this local
extension.  The same reasoning applies at \emph{every} place above $v$,
since potential unramifiedness of $M_{\ell_0}$ transfers to all those
completions.

The auxiliary CM abelian varieties are potentially good as well: after
fixing their Hodge tensors and taking a semistable extension, their inertia
is both unipotent and contained in a torus, hence trivial.  The same
criterion applies.  There are only finitely many varieties and places.
Choose finite extensions of the completions giving good reduction and
realize them within one finite global extension by weak approximation and
Krasner's lemma.  Pass to its normal closure over $K$; good reduction is
preserved by further extension.  This gives the asserted $E$ and proves
the direction of transfer required in the lemma.
\end{proof}

\begin{theorem}[Independence of unramifiedness for all abelian-type motives]
\label{c:ell-independence-abelian}
\status{Proved, using the October 2026 rationality theorem}
Let $M$ be a rational absolute-Hodge or motivated motive of abelian type
over a number field $K$, and let $v$ have residue characteristic $p$.
If $M_{\ell_0}$ is unramified at $v$ for one $\ell_0\ne p$, then
$M_\ell$ is unramified at $v$ for every $\ell\ne p$.
The statement includes $p=2$, and includes $\ell=2$ when $p\ne2$.
\end{theorem}

\begin{proof}
Apply Lemma~\ref{c:potentially-good-generator}.  Over a finite Galois
extension $E/K$, the motive $M_E$ is a summand of a finite direct sum $P$
of motives $h^{i_a}(A_a)(m_a)$, where the $A_a$ are powers or products of
abelian varieties with good reduction at \emph{every} place above $v$.
Indeed, tensor constructions on $h^1$, their duals and Tate twists occur
as K\"unneth summands in such powers.  A finite direct sum permits different
degrees $i_a$ and twists $m_a$.

Restriction and induction give morphisms
\[
 M\longrightarrow\operatorname{Ind}_E^K(M_E)\longrightarrow M
\]
whose composite is $[E:K]$.  After division by this nonzero rational
integer, $M$ is a summand of $\operatorname{Ind}_E^K(P)$.
These are finite \'etale correspondence maps in either motivic category;
after extension to $E$ they are the diagonal and summation maps over
$\Gal(E/K)$.  Thus the descent assertion is established over the original
$K$, without inferring it from trivial inertia over a ramified extension.

Write $e$ for the resulting rational Hodge, equivalently motivated,
projector.  The induced ambient varieties are finite disjoint unions of
the conjugate $A_a$.  Fix the chosen place of $\overline\Q$ above $v$ and
take a common finite Galois local extension $L/K_v$ splitting the finite
extensions of induction and giving good models for all components.
Smooth proper base change shows, simultaneously for every $\ell\ne p$,
that inertia on the ambient object factors through the same finite
quotient
\[
 \Delta=I_v/(I_v\cap G_L).
\]
It follows that this also holds on $M_\ell$.  Tate twists are trivial on
this inertia group when $\ell\ne p$.

By uniqueness of the abelian scheme models, the descent maps extend to
them.  On the geometric special fibres, each $\tau\in\Delta$ acts by a
permutation of components and origin-preserving algebraic isomorphisms
$u:X_0\to Y_0$ between conjugate abelian varieties.  Inertia fixes the
algebraic residue field, so these maps are actual homomorphisms over that
field.  Their pullbacks describe the cohomological inertia action.

The graph of every such $u$ is a rational Lefschetz class on $X_0\times Y_0$.
If $d=\dim Y_0=\dim X_0$ and $f(x,y)=y-u(x)$, then in every Weil
cohomology,
\begin{equation}
\label{c:inertia-graph-lefschetz}
 [\Gamma_u]=f^*[0_{Y_0}]
   =f^*\left(\frac{\theta_{Y_0}^{d}}
                         {\deg(\theta_{Y_0}^{d})}\right),
\end{equation}
where $\theta_{Y_0}$ is any ample divisor class.  This uses the common
degree normalization of top classes.  The expression on the right is one
rational polynomial in divisor classes, with the same cycle realizations
for all $\ell$.

Use the correspondence convention
$Z_*(x)=(\mathrm{pr}_2)_*([Z]\smile\mathrm{pr}_1^*x)$.
Then $\Gamma_u$ acts by $u_*=(u^{-1})^*$, and the alternating trace of
$e_*\circ u^*$ is the intersection pairing of $e$ with $\Gamma_u$.
For a projector supported in degree $i$, the ordinary trace is therefore
$(-1)^i$ times that pairing.  On a disjoint union, the same identity pairs
the matrix entries of $e$ with the graphs indexed by the component
permutation.  For a direct sum of different degrees and twists, only
the diagonal blocks relative to those Galois-stable summands contribute
to the trace.  Consequently
\[
 \operatorname{Tr}(\tau\mid M_\ell)
\]
is a finite sum, with the indicated signs, of pairings of specialized
rational Hodge classes with complementary rational Lefschetz classes on
products of abelian varieties with good reduction.  The generic matrix
entries of $e$ are Hodge classes on these products; realization and smooth
proper specialization preserve their actions, tensor operations, duality
and K\"unneth decompositions.

The rationality theorem \citep[Theorem 1.1]{c_rationality2026}, also
\citep[Theorem 1.1, pp.~1--2]{c_rationalityexposition2026}, says that every
such pairing is one rational number in every $\ell$-adic realization
$\ell\ne p$, including residue characteristic $2$.  Hence for each
$\tau\in\Delta$ there is a common rational trace $t_\tau$.
The assumed unramifiedness at $\ell_0$ gives
$t_\tau=\dim_\Q\omega_B(M)$.  In every other realization the operator
$\tau$ has finite order.  Its eigenvalues are roots of unity, and under a
complex embedding their sum equals their number.  Equality in the triangle
inequality forces every eigenvalue to be $1$.  Finite-order operators in
characteristic zero are semisimple, so $\tau$ is the identity on $M_\ell$.
This holds for every $\tau$, proving the theorem.
\end{proof}

\begin{remark}[The exact role of the new input]
\label{c:ell-independence-new-input}
\status{Proved scope clarification}
The proof does not apply N\'eron--Ogg--Shafarevich to an arbitrary original
generator: such a generator need not have good reduction when its summand
is unramified.  Lemma~\ref{c:potentially-good-generator} changes the
generator and proves its potential good reduction.  The finite descent
action is then recovered over $K$ by the common trace pairing.  Since
\eqref{c:inertia-graph-lefschetz} is Lefschetz, the proof uses the
rationality theorem alone; it does not use the CM Hodge theorem or assume
algebraicity of the Hodge projector.  It proves independence of
unramifiedness and, for potentially unramified objects, independence of
the finite inertia character.  It does not assert full Weil--Deligne
compatibility for objects with nonzero monodromy.
\end{remark}

% Supplementary reduction for C.6.  This does not settle the general iff.
\begin{proposition}[A purity reduction for homogeneous Chow motives]
\label{c:chow-unramified-purity-reduction}
\status{Proved conditional implication}
Let $K$ be a number field, $v$ a finite place with residue field of
cardinality $q$, and $p$ its residue characteristic.  Let
$M=(X,e,m)$ be a rational Chow motive over $K$, where $X$ is smooth
projective of dimension $d$ and the rational algebraic projector
$e\in\CH^d(X\times_KX)_{\Q}$ acts only in cohomological degree $i$.
Suppose that, for one prime $\ell_0\ne p$, the representation
$M_{\ell_0}$ is unramified at $v$ and the eigenvalues of geometric
Frobenius on it are algebraic numbers all of whose complex conjugates
have one common absolute value.  Then $M_\ell$ is unramified at $v$
for every $\ell\ne p$.

In particular, the asserted independence follows for such a motive if
the weight--monodromy conjecture holds for
$H^i(X_{\overline{K_v}},\Q_{\ell_0})$: its unramified direct summand
$M_{\ell_0}$ is then pure of weight $i-2m$.
\end{proposition}
\begin{proof}
Write $V_\ell=eH^i(X_{\overline{K_v}},\Q_\ell)(m)$.
Saito's Theorem~0.1, in its more general form Theorem~3.1,
gives a rational alternating trace independent of $\ell$ for the
composition of an algebraic correspondence and any Weil element
\citep[Theorems~0.1 and 3.1]{c_saito2003weights}.
The theorem extends from integral to rational correspondences by
linearity.  Taking the correspondence $e$, its action on all other
degrees is zero; since it is idempotent and commutes with the Galois
action, one obtains the exact equality
\[
 \operatorname{Tr}(\sigma\mid V_\ell)
 =(-1)^i q^{-m n(\sigma)}
   \sum_a(-1)^a\operatorname{Tr}
   \bigl(e\circ\sigma\mid H^a(X_{\overline{K_v}},\Q_\ell)\bigr).
\]
Thus $\operatorname{Tr}(\sigma\mid V_\ell)$ is a common rational
number for each $\sigma\in W_{K_v}$.  Apply this also to all positive
powers of $\sigma$.  Newton identities show that the characteristic
polynomial of every such $\sigma$ is in $\Q[T]$ and independent of
$\ell$.  In particular the dimension $r$ is common; if $r=0$ there
is nothing to prove.

Fix a lift $\phi\in W_{K_v}$ of geometric Frobenius.  The local
monodromy theorem supplies, for every $\ell\ne p$, a nilpotent operator
$N_\ell$ and an open inertia subgroup on which the action is
$\exp(t_\ell N_\ell)$; it also gives
\[
 F_\ell N_\ell F_\ell^{-1}=q^{-1}N_\ell,
 \qquad F_\ell=\rho_\ell(\phi).
\]
The common characteristic polynomial of $F_\ell$ has roots of the
single absolute value in the hypothesis.  Over an algebraic closure
of the coefficient field, $N_\ell$ maps a generalized
$\lambda$-eigenspace of $F_\ell$ into its generalized
$q^{-1}\lambda$-eigenspace.  The latter eigenvalue cannot occur:
under any embedding of the common algebraic splitting field into
$\C$, multiplication by $q^{-1}$ changes its absolute value.
Consequently $N_\ell=0$ for every $\ell$.

The image of inertia in each $\rho_\ell$ is therefore finite.  For
each $\tau\in I_{K_v}$ the common trace is
$\operatorname{Tr}(\rho_{\ell_0}(\tau))=r$.  A finite-order matrix
in characteristic zero is semisimple with roots of unity as
eigenvalues.  Their sum can equal the positive integer $r$ only if
all of them are $1$, as follows by embedding their cyclotomic field
into $\C$ and using equality in the triangle inequality.  Hence
$\rho_\ell(\tau)=1$ for every $\tau$ and every $\ell\ne p$.

For the last assertion, the weight--monodromy filtration is functorial
for the algebraic projector.  A pure Weil--Deligne direct summand
whose monodromy is zero is concentrated in monodromy degree zero,
so its Frobenius roots all have weight $i-2m$.
\end{proof}

\begin{remark}[Exact scope of this reduction]
\label{c:chow-purity-scope}
\status{Proved scope clarification; the general claim remains unresolved}
For the full group $H^i(X)$, the purity assumption in the proposition
is Serre's unramified-purity conjecture C7
\citep[\S2.3, printed p.~19-07]{c_serre1970local}.
For a prescribed proper summand it is the corresponding
\emph{motivic extension} of that condition; the literal original C7
must not be silently enlarged.  The named weight--monodromy
conjecture for the ambient cohomology is a sufficient alternative.
The use of Saito's theorem requires an algebraic correspondence.
It cannot be applied to an arbitrary absolute-Hodge or motivated
projector without an additional comparison result.  No converse to
any of these named conjectures is proved here, and this proposition
does not close the general-motive requirement in C.6.
\end{remark}

\begin{corollary}[Surface and low-degree Chow constructions]
\label{c:surface-lowdegree-chow-unramified}
\status{Proved}
Let $M$ be a pure rational Chow motive over a number field.  Suppose
that it is a homogeneous algebraic Chow summand of the cohomology of
a finite product of smooth projective surfaces, or more generally
belongs to the tensor constructions on supplied homogeneous algebraic
Chow summands supported in degrees $0,1,2$ of smooth projective
varieties, allowing finite sums, duals, Tate twists and homogeneous
algebraic projectors.  Then, at each finite place of residue
characteristic $p$, unramifiedness of $M_\ell$ is independent of
$\ell\ne p$.
\end{corollary}
\begin{proof}
We use the complete low-degree result
\citep[Lemma~3.9, author p.~42]{c_saito2003weights}; we spell out its
splitting step.  If $X$ is smooth projective of dimension $d\ge r$,
choose a smooth complete intersection $\iota:Y\hookrightarrow X$ of
dimension $r\le2$, after a finite extension if necessary.  For an ample
class $L$, with the degrees of the chosen hypersurfaces included in
its nonzero scalar normalization, the projection formula gives
\[
 \iota_*\iota^*=L^{d-r}:
 H^r(X)\longrightarrow H^{2d-r}(X)(d-r).
\]
Hard Lefschetz makes this an isomorphism, so
$(L^{d-r})^{-1}\iota_*$ retracts $\iota^*$ as a Galois map.  No
algebraicity of this inverse is needed or asserted here.  For curves
and surfaces, semistable alteration and the degree-divided trace
reduce purity to the Rapoport--Zink surface theorem
\citep[Satz~2.13]{c_rapoportzink1982}.  Purity descends through the
finite extension.  The cases $d<r\le2$ follow from duality or are zero.
Thus $H^r(X)$ is pure as a Weil--Deligne representation for $r\le2$.
Poincar\'e duality gives the corresponding statement in degrees
$2d,2d-1,2d-2$, and therefore in every degree of a surface.

Purity is preserved by tensor products, direct sums of the same
weight, duals, Tate twists and direct summands.  K\"unneth gives the
required purity for the stated constructions.  For finite sums, insert
the algebraic unit and point projectors of projective spaces to put
the summands into one cohomological degree and one common Tate twist
on a disjoint union of smooth projective varieties of the same
dimension.  Proposition~\ref{c:chow-unramified-purity-reduction}
then applies to the resulting homogeneous algebraic projector.  In particular it preserves the
finite inertia information at the original place; there is no
inference of unramifiedness by ramified descent.
\end{proof}

This supplies the expected local C.4 bound
$\operatorname{weights}(M_\ell^{I_v})\le j$ for a pure weight-$j$
coefficient in these tensor constructions: inertia invariants lie
in $\ker N$, which consists of the bottom vectors of pure monodromy
strings.  The first remaining general homogeneous Chow purity case
is $H^3$ of arbitrary smooth projective threefolds.  The assumed
homogeneous projectors on arbitrary higher-dimensional varieties
are part of the statement; none is manufactured from a possibly
nonalgebraic hard-Lefschetz inverse.

\begin{corollary}[Chow tensor summands of hyper-K\"ahler varieties]
\label{c:hyperkahler-chow-unramified}
\status{Proved using the stated 6 October 2026 purity theorem}
Let $M$ be a pure rational Chow motive over a number field, obtained
from products of hyper-K\"ahler varieties by homogeneous algebraic
projectors, finite sums, duals and Tate twists.  At a finite place
$v$ of residue characteristic $p$, if $T_{\ell_0}(M)$ is unramified
for one $\ell_0\ne p$, then $T_\ell(M)$ is unramified for every
$\ell\ne p$.
\end{corollary}

\begin{proof}
\citet[Theorem~4.4, with Lemmas~4.2--4.3]{c_itozou2026}
prove weight--monodromy in every degree for every hyper-K\"ahler
variety over a $p$-adic field, including $\ell=p$.  Their theorem
is over the original local field: the finite extension in its proof
does not remain in its conclusion, because purity descends through
finite extensions.  Pure Weil--Deligne representations are stable
under tensor products, duals, Tate twists and direct summands.
K\"unneth therefore gives the ambient purity hypothesis in Proposition~\ref{c:chow-unramified-purity-reduction}.  At $\ell_0$, zero monodromy gives
a single Frobenius weight on the chosen pure summand.  That reduction,
using Saito's algebraic-correspondence trace theorem, proves the
assertion at the original place $v$.
\end{proof}

This consequence uses the purity theorem, not a descent of
unramifiedness through a ramified extension.  The stronger
group-valued compatibility statement in
\citet[Theorem~5.8 and Remark~5.9]{c_itozou2026} assumes connected
monodromy and a semistable Kuga--Satake variety over the chosen base;
Corollary~5.10 obtains its general local Frobenius comparison after
a finite extension.  The displayed corollary retains algebraic
projectors.  The same purity theorem gives the local C.4 bound
$\operatorname{weights}(\ker N)\le j$ for pure weight-$j$
tensor summands of hyper-K\"ahler cohomology.  It does not assert the missing trace comparison for
arbitrary absolute-Hodge projectors outside the abelian-type case.

\begin{proposition}[A finite-field rank conjecture controls algebraic homological projectors]
\label{c:chow-unramified-rank-reduction}
\status{Proved conditional implication; no converse is asserted}
Fix a prime $p$. Assume independence of $\ell$ of homological
equivalence for rational cycles on every smooth projective variety
over $\overline{\F}_p$. Equivalently, assume independence of the
ranks of all rational algebraic correspondences on their individual
cohomology groups; this is the exact equivalence of
\citet[Theorem~1, Definition~4.1 and Theorem~4.6]{c_vandobben2020}.

Let $K$ be a number field, $v\mid p$, and let $X/K$ be smooth
projective of pure dimension $d$. Let
$e\in\CH^d(X\times_KX)_\Q$ be a rational algebraic correspondence
whose Betti action is idempotent and vanishes outside degree $i$.
For an integer $m$, put
$M_\ell=\operatorname{im}(e:H^i(X_{\bar K},\Q_\ell)(m)
\to H^i(X_{\bar K},\Q_\ell)(m))$.
Betti--\'etale comparison gives the same homogeneous idempotent
condition for every $\ell$. We require no identity $e^2=e$ in
the Chow ring; the assertion therefore covers algebraic
homological projectors as well as rational Chow motives.
Then the ranks of all powers of the local monodromy operator on
$M_\ell$ are independent of $\ell\ne p$.
Moreover,
\[
 M_{\ell_0}\text{ unramified at }v\text{ for one }\ell_0\ne p
 \quad\Longrightarrow\quad
 M_\ell\text{ unramified at }v\text{ for every }\ell\ne p.
\]
This implication does not assume weight--monodromy.
\end{proposition}

\begin{proof}
We work over the local field $K_v$ and first suppress the Tate twist:
for $\ell\ne p$ it has trivial inertia and zero monodromy.
There are three linear-algebra and spectral-sequence points to verify.

\emph{1. Ranks on cohomology are finite block-matrix ranks.}
Let $f:C^\bullet\to D^\bullet$ be a map of finite complexes over a
field. For every $a$,
\begin{equation}
\label{eq:c-induced-cohomology-rank}
 \operatorname{rank}H^a(f)
 =\operatorname{rank}
 \begin{pmatrix}
  d_C^a&0\\ f^a&d_D^{a-1}
 \end{pmatrix}
 -\operatorname{rank}d_C^a-\operatorname{rank}d_D^{a-1}.
\end{equation}
Indeed, projection of the block image onto $C^{a+1}$ has image
$\operatorname{im}d_C^a$ and kernel
$f^a(\ker d_C^a)+\operatorname{im}d_D^{a-1}$.
Dividing that kernel by $\operatorname{im}d_D^{a-1}$ gives exactly
the image on cohomology, since $f$ carries boundaries to boundaries.

The rank conjecture applies to these block matrices whenever their
entries are algebraic correspondences between finite sums of twisted
stratum cohomology groups of one weight. To check this assertion
explicitly, write a summand as $H^{j_s}(Z_s)(t_s)$ with
$j_s-2t_s=w$. Choose $T\geq\max_s t_s$ and put
$r_s=T-t_s\geq0$. The summand identifies with the part of
\[
 H^{w+2T}(Z_s\times\mathbf P^{r_s})(T)
\]
selected by the algebraic top-degree projector of the projective
space. Exterior products of the original correspondences with the
projective-space correspondences implement every block.
Choose a common dimension $D$ at least as large as every
$\dim Z_s+r_s$, and further multiply the $s$th component by
$\mathbf P^{D-\dim Z_s-r_s}$ with its algebraic degree-zero
projector. This leaves the selected cohomological degree unchanged
and makes all components equidimensional of dimension $D$.
Disjoint unions implement finite direct sums, and the top-degree
and degree-zero projectors kill the unwanted projective-space
components. Thus the block map is a single codimension-$D$
rational algebraic correspondence
on a fixed cohomological degree of smooth projective schemes.
Its rank is independent of $\ell$ by the stated conjecture.
No algebraic projector onto a kernel of a differential is required.

\emph{2. A semistable projector has common algebraic $E_1$ maps.}
Choose a finite Galois extension $L/K_v$, a projective strictly
semistable scheme $\mathcal Y/\mathcal O_L$, and a proper
generically finite map $f:Y\to X_L$ of positive degree $\delta$.
Such an alteration is supplied by
\citet[Corollary~1.2.6, author p.~16]{c_saito2003weights}.
For geometric components, take disjoint unions of alterations and,
if necessary, repeat components to make the generic degree common.
The correspondence
\[
                 e_Y=\delta^{-1}f^* e f_*
\]
is algebraic and acts idempotently on cohomology. The identities $f_*f^*=\delta$ and
the projection formula identify its realized image with that of
$e$, compatibly with Galois action and monodromy.
Passing from $K_v$ to $L$ multiplies the normalized monodromy
operator by the nonzero ramification index, so none of its ranks
changes.

Use Saito's convention that $Y^{(s)}$ is the disjoint union of the
$(s+1)$-fold intersections of geometric components. The weight
spectral sequence is
\[
 E_{1,\ell}^{a,b}
 =\bigoplus_{j\geq\max(0,-a)}
 H^{b-2j}(Y^{(a+2j)}_{\overline k},\Q_\ell)(-j)
 \ \Longrightarrow\ H^{a+b}(Y_{\overline L},\Q_\ell).
\]
Its differentials are fixed signed restriction and Gysin
correspondences. The action of $e_Y$ on this spectral sequence
has $E_1$ components induced by fixed rational algebraic
correspondences on the strata, independent of $\ell$.
This is precisely
\citet[Lemma~2.4.2, Corollary~2.4.3 and Proposition~2.4.4,
complete proofs, author pp.~34--36]{c_saito2003weights}:
one resolves the closure of a cycle on the regular semistable
product blowup by vector bundles, restricts their Chern characters
to the strata, and divides by the fixed factorial.
The resulting rational cycles are chosen before $\ell$.
The $E_1$ endomorphism need not be idempotent; its induced action
on $E_2$ and the abutment is the required projector.

On $E_1$, the monodromy map
\[
 N:E_{1,\ell}^{a,b}\longrightarrow
       E_{1,\ell}^{a+2,b-2}(-1)
\]
is the identity on matching summands and zero on the others.
It sends $(a,j)$ to $(a+2,j-1)$ and hence preserves the stratum
index $a+2j$. Consequently the stratum correspondence blocks
commute with this map. The assertion for $N$ follows from
\citet[Corollary~2.2.4(2), author pp.~21--22]{c_saito2003weights}:
the displayed map for $T-1$ differs on the graded pieces from
$N$ only by its nonzero tame scalar. Higher logarithm terms drop
the filtration further and do not change this leading map.

Fix $b$ and an integer $r\geq1$. Regard
$C^a=E_{1,\ell}^{a,b}$ and
$D^a=E_{1,\ell}^{a+2r,b-2r}(-r)$ as complexes in $a$.
All their summands have weight $b$, and
$N^r e_Y:C^\bullet\to D^\bullet$ has rational algebraic
stratum blocks. Formula~\eqref{eq:c-induced-cohomology-rank}
therefore shows that the rank of its induced $E_2$ map is
independent of $\ell$.
The spectral sequence degenerates at $E_2$, since its $E_1$
rows are pure of their distinct weights. Thus these are the ranks
of all graded maps
\[
 N^r e_Y:\operatorname{Gr}^W_b H^i(Y)(r)
       \longrightarrow\operatorname{Gr}^W_{b-2r}H^i(Y),
\]
or, equivalently, of their simultaneous untwisted form with
$(-r)$ on the target.

\emph{3. Graded monodromy detects actual monodromy here.}
Let $F$ be a geometric Frobenius lift over $L$ and let $q_L$
be its residue cardinality. Each $\operatorname{Gr}^W_b H^i(Y)$
is pure of weight $b$, by
\citet[Lemma~3.7, author p.~42]{c_saito2003weights}.
Distinct weights have disjoint Frobenius spectra. Consequently
the filtration has a unique $F$-equivariant splitting by the
generalized Frobenius weight spaces
\[
             H^i(Y)=\bigoplus_b V_b,\qquad
             W_b H^i(Y)=\bigoplus_{c\leq b}V_c.
\]
This follows from primary decomposition for $F$ and does not
assume that $F$ is semisimple. The relation
$FNF^{-1}=q_L^{-1}N$ carries $V_b$ into $V_{b-2}$,
while $e_Y$ preserves every $V_b$. Therefore
\[
 \operatorname{rank}(N^r e_Y\mid H^i(Y))
 =\sum_b\operatorname{rank}
       (\operatorname{Gr}^W_b(N^r e_Y)).
\]
In particular zero on all graded pieces is equivalent to zero
on the entire space. Because $e_Y$ is a projector commuting with
$N$, the left side is the rank of $N^r$ on its image.
This proves the first assertion, including the independence
of vanishing of $N$.

Finally suppose $M_{\ell_0}$ is unramified. Its monodromy
vanishes, hence the preceding argument makes monodromy vanish
on $M_\ell$ for every $\ell\ne p$. Strict semistability of
$\mathcal Y$ then makes $I_L$ act trivially on every such image.
Thus the original inertia actions factor through one finite
quotient of $I_{K_v}$.
By Saito's Theorem~0.1, the alternating trace of $e\sigma$ is
rational and independent of $\ell$ for every
$\sigma\in I_{K_v}$. Homogeneity of $e$ identifies this with
$(-1)^i\operatorname{Tr}(\sigma\mid M_\ell)$; the Tate twist
is trivial on inertia. At $\ell_0$ the trace is the common
dimension of $M$, so it is that dimension for every $\ell$.
Averaging these traces over the finite inertia quotient gives
the full dimension for its invariant subspace. Finite-group
representations in characteristic zero are semisimple, hence
the inertia action is trivial for every $\ell$.
\end{proof}

\begin{remark}[The precise conditional and categorical boundary]
\label{c:chow-rank-reduction-scope}
\status{Proved sufficient-conjecture reduction; the general equivalence is open}
The hypothesis is a published, independently formulated
independence-of-homological-equivalence conjecture over
$\overline{\F}_p$, not a renamed version of C.6.
The result proves only its implication to the stated C.6 assertion.
It does not prove the converse: realizing an arbitrary finite-field
correspondence as local monodromy of a number-field variety would
require an additional geometric construction which is not supplied.
The representative must be algebraic and homogeneous, but only
its cohomological action must be idempotent. Inspecting the proof,
idempotence is used on the realization and its $E_2$ graded pieces,
never in the Chow group. No extension to an arbitrary
absolute-Hodge, motivated, or numerical projector is inferred
from Saito's algebraic-correspondence construction alone.
The following theorem treats motivated projectors by an additional
finite algebraic-and-Lefschetz construction.
Thus this proposition gives a second exact sufficient hypothesis,
alongside the purity reduction, but does not satisfy the requested
general-motive proof or a two-way named-conjecture equivalence.
\end{remark}


\begin{theorem}[The finite-field rank hypothesis also controls motivated motives]
\label{c:motivated-unramified-rank-reduction}
\status{Proved conditional implication; the general named-conjecture equivalence remains open}
Fix a prime $p$, and assume independence of $\ell$ of homological
equivalence for rational cycles on every smooth projective variety
over $\overline{\F}_p$, equivalently the correspondence-rank
conjecture of
\citet[Theorems~4.6 and~5.1 and Corollary~5.2]{c_vandobben2020}.
Let $M$ be an Andr\'e motivated motive with rational coefficients
over a number field $K$, and let $v\mid p$.
Then the ranks of all powers of its local monodromy operator are
independent of $\ell\ne p$. Moreover,
\[
 M_{\ell_0}\text{ is unramified at }v\text{ for some }\ell_0\ne p
 \quad\Longleftrightarrow\quad
 M_\ell\text{ is unramified at }v\text{ for every }\ell\ne p.
\]
Here the equivalence concerns the coefficient primes under the
stated conjectural hypothesis; it is not a converse from C.6 to
the finite-field conjecture. No weight--monodromy hypothesis or
algebraicity of the motivated projector is assumed.
\end{theorem}

\begin{proof}
We give the operator and geometric constructions, including the
points needed to avoid a tensor-product assertion for nearby cycles.
Finite direct sums and Tate twists cause no difficulty, so realize
$M$ as the image of a motivated idempotent $e$ on the cohomology
of a smooth projective variety $U/K$. We work degree by degree
and suppress an overall Tate twist, which has trivial inertia
at $v$ for $\ell\ne p$.

\emph{1. Motivated operators are finite inverse expressions on whole products.}
For a smooth projective variety $Z$ of dimension $d$, write $L$
for an ample Lefschetz operator. On the string generated by a
primitive vector $z$ of degree $j\leq d$, put $n=d-j$.
Andr\'e's Hodge involution $T$ satisfies
\[
 T(L^a z)=(-1)^{j(j+1)/2}\frac{a!}{(n-a)!}L^{n-a}z,
 \qquad
 \Lambda(L^a z)=a(n-a+1)L^{a-1}z.
\]
Direct substitution gives $T^2=1$, $TLT=\Lambda$, and
$[\Lambda,L]|_{H^m(Z)}=(d-m)\operatorname{id}$.
These are the definitions in
\citet[\S1.1, printed p.~10]{c_andre1996}.
We use this direct computation; the different normalization of
the Lefschetz involution is not substituted for $T$.

On $Z=X\times Y$ with product ample class,
\citet[Lemma~1.3.2, printed p.~13]{c_andre1996} gives
\[
 T_Z|_{H^m(X)\otimes H^n(Y)}
       =(-1)^{mn}(T_X\otimes T_Y).
\]
The cross-parity sign commutes with partial Lefschetz operators
and both factor involutions. Consequently
$T_ZL_XT_Z=\Lambda_X\otimes1$.
Put $h_X=[T_ZL_XT_Z,L_X]$. The polynomial
\[
 \pi_m^X=
 \prod_{\substack{0\leq r\leq2\dim X\\r\ne m}}
       \frac{h_X-(\dim X-r)}{r-m}
\]
projects onto factor degree $m$. This is an operator expression,
not an assumption that the generic K\"unneth projector is algebraic.
The finite Weyl expression
\[
 w_X=\exp(L_X)\exp(-T_ZL_XT_Z)\exp(L_X)
\]
satisfies
\[
 T_X\otimes1
     =w_X\sum_m(-1)^{m(m+1)/2}\pi_m^X.
\]
The sign acts on the input and is therefore on the right.
Both identities follow on each primitive Lefschetz string;
all exponentials are finite rational polynomials.
Thus every factor involution is a finite rational expression
in partial algebraic cup products and a global involution.

A global involution itself is a finite expression in algebraic
Lefschetz maps and inverses of Hard-Lefschetz isomorphisms.
For $2\leq j\leq d$, the primitive projector in degree $j$ is
\[
 1-L(L^{d-j+2})^{-1}L^{d-j+1}.
\]
The composition defining its complementary projector is typed as
\[
\begin{aligned}
 H^j(Z)&\xrightarrow{L^{d-j+1}}
 H^{2d-j+2}(Z)(d-j+1)\\
 &\xrightarrow{(L^{d-j+2})^{-1}(-1)}
 H^{j-2}(Z)(-1)\xrightarrow{L}H^j(Z).
\end{aligned}
\]
Here the inverse Hard-Lefschetz isomorphism is tensored by $\Q(-1)$.
The complementary expression recovers the lower Lefschetz component;
iteration gives the entire decomposition.
For $j>d$, first invert
$L^{j-d}:H^{2d-j}(Z)\to H^j(Z)(j-d)$.
The coefficient formula for $T$ then supplies the involution.
Only finitely many degrees and inverses occur.

To see explicitly that this handles a motivated correspondence,
use Andr\'e's definition with
\[
 e=\operatorname{pr}_{12*}(\alpha\cup T_P\beta),
 \qquad P=U_1\times U_2\times V,
\]
term by term, where $\alpha,\beta$ are rational algebraic cycles
and all data are defined over $K$
\citep[Definition~1, \S2.1, printed p.~14]{c_andre1996}.
Here the setup of \S2.1 takes its auxiliary objects to be smooth
projective $K$-schemes; we use that definition over $K$, rather
than descending a presentation chosen only over $\bar K$.
Let $d=\dim U$, $D=\dim P$, and let their codimensions be
$a,b$. The degree-zero endomorphism condition is $b=a+d$.
On $W=U_0\times P$, exterior multiplication by $\beta$ sends
$H^i(U_0)$ to $H^{i+2b}(W)(b)$. Apply $1\otimes T_P$, then
pull back along $P\to W$ identifying $U_0$ and $U_2$, cup with
$\alpha$, and push forward to $U_1$. The final degree and twist are
\[
 i+2(D-b)+2a-2(D-d)=i,
 \qquad (D-b)+a-(D-d)=0.
\]
This is exactly the action of $e$.
The factor involution is replaced by the preceding whole-product
expression. All remaining elementary arrows are algebraic
correspondences between whole smooth projective varieties, or
inverses of specified Hard-Lefschetz isomorphisms on those varieties.
The construction never requires a tensor comparison of $E_1$ pages.

Every intermediate vertex may be written
$H^{i+2r}(Z)(r)$. Keeping these finitely many vertices separately
also handles the degree changes inside the polynomial expressions.
In particular no inverse of an unspecified inhomogeneous operator
is taken.

\emph{2. Simultaneous semistable alterations with the original inertia action.}
The expression uses only finitely many varieties $Z$.
By \citet[Corollary~1.2.6, author p.~16]{c_saito2003weights},
choose one finite Galois extension $L/K_v$ such that each $Z_L$
has a projective strictly semistable alteration
$f_{Z,0}:Y_{Z,0}\to Z_L$.
The corollary allows every finite extension containing a suitable
one, so a common Galois extension is permitted.
For disconnected varieties, work componentwise and repeat covers
to arrange a common degree on all geometric components.

Put $G=\operatorname{Gal}(L/K_v)$ and replace each alteration
and its integral model by the disjoint union of all its $G$-conjugates.
Denote the resulting map by $f_Z:Y_Z\to Z_L$ and its degree by
$\delta_Z$. The union has a canonical semilinear $G$-action and
$f_Z$ is equivariant, since $Z$ descends to $K_v$.
Each conjugate integral model is strictly semistable over
$\mathcal O_L$. No descent of the integral model to
$\mathcal O_{K_v}$ is asserted or needed.
At the generic fibre this is the ordinary conjugate-coproduct
descent construction: a variety over $L$, viewed over $K_v$ and
then base changed to $L$, is the coproduct of its conjugates.

Let $i_Z=f_Z^*$, $p_Z=\delta_Z^{-1}f_{Z*}$ and
$q_Z=i_Zp_Z$. The identity $p_Zi_Z=1$ identifies every original
cohomology vertex with $\operatorname{im}q_Z$.
An algebraic elementary arrow $a:Z\to Z'$ in the cohomological
diagram transfers to
\[
                    \widetilde a=i_{Z'}a p_Z.
\]
These are rational algebraic correspondences between the
semistable generic fibres. Saito's construction gives fixed
rational algebraic $E_1$ blocks for them and for the $q_Z$
\citep[Lemmas~2.4.1--2.4.2, Corollary~2.4.3 and
Proposition~2.4.4, author pp.~33--36]{c_saito2003weights}.
Only their induced relations on $E_2$ will be used;
their chosen $E_1$ representatives need not be projectors.

If $\sigma\in I_{K_v}$ has image $g\in G$, then $g$ acts
trivially on the residue field of $L$. Its semilinear action
on the union of integral models therefore induces algebraic
isomorphisms and permutations of their closed strata, over
that residue field. Ordering intersection components introduces
only fixed signs. These give a common algebraic action on $E_1$.
This is the transport-of-structure spectral-sequence map in
\citet[Lemma~3.2, complete proof on author pp.~36--38;
Lemma~3.3 and proof, pp.~38--39]{c_saito2003weights},
with $n(\sigma)=0$. Inertia of $L$ acts trivially on $E_1$,
so this action depends only on $g$ there.
On the abutment it is the original $\sigma$ action on the
conjugate coproduct, and the embeddings $i_Z$ respect it.

\emph{3. Inverses require only graph constraints on $E_2$.}
Fix a normalized weight $b$.
For the vertex $H^{i+2r}(Z)(r)$ use the row complex
\[
 C_{Z,r,b}^a=E_{1,Z}^{a,b+2r}(r).
\]
Its cohomology in the common index $a=i-b$ is the weight-$b$
piece of that vertex. All its terms are pure of weight $b$,
and all transferred elementary arrows are chain maps between
such complexes with common algebraic stratum blocks.

Impose $(q_Z-1)x=0$ on every variable representing an original
vertex. For an elementary arrow impose $y-\widetilde a x=0$;
for addition use the corresponding linear equation. For an inverse
of a Hard-Lefschetz map $a$, impose $\widetilde a y=x$, with
the image constraints on both variables. This inverse is required
only on the $q$-images, not on the entire alteration cohomology.
It is indeed invertible on each weight piece: the original
Hard-Lefschetz map is a Galois-equivariant isomorphism with the
displayed twists, and the generalized Frobenius weight spaces
give a canonical splitting of the weight filtration.
The map and its inverse therefore preserve the weight pieces.

These equations describe the graph of the intended motivated
operator on the specified images. More generally, if $R:E\to F$
collects such constraints and $\pi:E\to V$ is the output
projection, elementary rank-nullity gives
\[
       \dim\pi(\ker R)
       =\operatorname{rank}\begin{pmatrix}R\\\pi\end{pmatrix}
          -\operatorname{rank}R.
\]
Internal solutions need not be unique, and the input may itself
be restricted by $q x=x$. This formula consequently applies
without a chain-level inverse or a chain-level image projector.

Both block maps on the right are induced on cohomology by
finite block chain maps. Their ranks are computed by
\eqref{eq:c-induced-cohomology-rank}.
All resulting $E_1$ matrices have rational algebraic entries and
one weight. The explicit projective-space padding in the proof of
Proposition~\ref{c:chow-unramified-rank-reduction} realizes each
as a single correspondence in one degree between equidimensional
smooth projective schemes. The finite-field rank hypothesis
therefore makes every such rank independent of $\ell$.

Append $N^s$ as the final arrow to the motivated expression.
On $E_1$ it is identity or zero on matching summands. Its output
row is shifted by $2s$ and twisted by $-s$:
\[
           E_{1,U}^{a+2s,b-2s+2r}(r-s).
\]
Here $E_{1,U}$ denotes the spectral sequence of the chosen
alteration for $U$, with the image constraints already imposed.
This row still has weight $b$ and is a complex in the same
index $a$. Thus the preceding graph argument gives independence
of the ranks of every $\operatorname{Gr}^W(N^s e)$.
In this step $N$ is the spectral-sequence operator; it has not
been treated as a generic degree-zero motivic correspondence.

The spectral sequence degenerates at $E_2$ by its distinct
row weights. By primary decomposition for Frobenius, those
graded weights split the abutment into generalized weight spaces.
The relation $FNF^{-1}=|k_L|^{-1}N$ sends weight $b$ exactly to
weight $b-2$. Since $e$ commutes with Frobenius, the rank of
$N^s e$ on the abutment is the sum of the graded ranks.
This proves independence of all monodromy ranks on $M$.
Changing the local field multiplies $N$ by a nonzero ramification
index, which does not affect ranks.

\emph{4. Descend triviality of inertia to the original place.}
Suppose $M_{\ell_0}$ is unramified. The first assertion gives
$N=0$ on every $M_\ell$. Strict semistability then makes
$I_L$ act trivially on all of them, so their inertia actions
factor through the same finite quotient of $I_{K_v}$.
For each representative $\sigma$ of that quotient, append the
common algebraic $E_1$ operator $\sigma-1$ to the graph
expression for $e$. The same rank argument shows independence
of the rank of every $\operatorname{Gr}^W((\sigma-1)e)$.
These ranks are zero at $\ell_0$, hence zero for every $\ell$.
On $M_\ell$, therefore, $\sigma$ is the identity on the
associated graded of a finite filtration and is unipotent.
It also has finite order. In characteristic zero a finite-order
unipotent operator is the identity. This proves the assertion
at the original place, without a trace formula for a motivated
projector.
\end{proof}

\begin{corollary}[Abelian strata suffice without the rank conjecture]
\label{c:motivated-abelian-strata-unramified}
\status{Proved special geometric case using the current abelian Tate theorem}
In the preceding construction, suppose every geometric stratum
of the chosen simultaneous semistable alterations has rational
Chow motive of abelian type, with actual algebraic presentations
in that category. Then the monodromy-rank and unramifiedness
conclusions hold without the finite-field rank conjecture.
The condition applies to all auxiliary varieties required by
the finite motivated presentation, not merely to its original
ambient variety.
\end{corollary}
\begin{proof}
Every block is a rational algebraic correspondence among the
strata and their projective-space factors.
The current CM Hodge theorem and its abelian Tate consequence
give homological equals numerical equivalence on all products
in their abelian-motive category
\citep[Theorem~I and \S1.1.3]{AvanHoftenSempliner2026}
\citep[Theorem~3.3(a)]{DMilne2002Polarizations}.
The common numerical category is semisimple
\citep[Theorem~1, pp.~448--449]{c_jannsen1992semisimple}; every block has
an image projector there, represented by one rational
algebraic correspondence. Its realizations are projectors
with that same image. Their ranks agree, for example by the
common algebraic trace after selecting the one padded degree.
Thus every finite block-rank input in the proof is available.
This argument uses ranks and semisimplicity, and does not
transport a geometric primitive form to a positive polarization.
\end{proof}

\begin{remark}[Boundary of the motivated reduction]
\label{c:motivated-rank-reduction-scope}
\status{Proved scope clarification; no general-motive closure}
The theorem removes the algebraic-projector limitation of the
earlier sufficient-conjecture implication by using Andr\'e's
finite algebraic-and-Lefschetz presentation. An arbitrary
absolute-Hodge projector is not thereby proved motivated,
and a numerical projector is not thereby supplied with a
realization-compatible presentation. No converse from
unramifiedness independence to the finite-field rank conjecture
has been proved. The geometric construction needed to encode
an arbitrary finite-field correspondence in the monodromy of
a number-field variety remains separate.
\end{remark}


\begin{proposition}[Finite-prime approximation by a number-field motive]
\label{c:local-global-motive-approximation}
\status{Proved; a globalization reduction, not a proof of general independence}
Let $K/\Q_p$ be finite, and let $M$ be a homogeneous rational
motive over $K$ in either the algebraic homological category or
Andr\'e's motivated category.
For every nonempty finite set $\Sigma$ of primes different from $p$
there exist a number field $F\subset K$, a place $w$ of $F$ with
$F_w=K$, and a homogeneous motive $M'$ of the same kind and weight
over $F$, such that
\[
 H_\ell(M')|_{G_{F_w}}\simeq H_\ell(M)
 \qquad(\ell\in\Sigma)
\]
as representations of $G_K$, under the specified identification
$F_w=K$.
Consequently, universal independence of unramifiedness for these
motives over number fields is equivalent to the same assertion over
all finite extensions of $\Q_p$, for all $p$.
The number-field approximation may depend on $\Sigma$.
\end{proposition}

\begin{proof}
First choose a number field $F_0\subset K$ whose completion at the
induced place is $K$.
For completeness, choose a primitive element of $K/\Q_p$ and
approximate its separable minimal polynomial by a monic polynomial
over $\Q$ sufficiently closely. Root continuity and Krasner's lemma
give a root $b\in K$ of the latter polynomial with
$\Q_p(b)=K$. Thus $F_0=\Q(b)$ is dense in $K$.
This is also the separable case of
\citet[Theorem~2.1 and its proof, pp.~2--3]
{c_conrad_localglobal}.

Present $M$ as the homogeneous image of a rational projector on the
cohomology of smooth projective varieties, allowing Tate twists and
finite direct sums. In the algebraic homological case the projector
is represented by rational algebraic correspondences.
In the motivated case choose the finite algebraic-cycle and
Lefschetz-involution presentation in Andr\'e's definition.
The finite operator presentation established in the proof of
Theorem~\ref{c:motivated-unramified-rank-reduction} expresses these
involutions using algebraic maps and inverses of Hard Lefschetz
isomorphisms on whole products.
This use of that operator presentation does not use its
finite-field rank hypothesis.

All the varieties, cycles, maps, line bundles, and finitely many
rational coefficients in this presentation are defined over a
finitely generated subfield $E\subset K$ containing $F_0$.
Enlarge $F_0$ to its algebraic closure in $E$, which is still a
number field contained in $K$ and still dense there.
Then $E/F_0$ is regular. Spread the data over a smooth geometrically
integral $F_0$-scheme $S$ with function field $E$.
After shrinking $S$, all the chosen varieties form smooth projective
families, all chosen line bundles are relatively ample, and the
cycles and maps defining the finite presentation extend.
The embedding $E\hookrightarrow K$ gives a point $x\in S(K)$.

For each $\ell\in\Sigma$, relative smooth proper cohomology and
the extended algebraic correspondences give morphisms of lisse
$\Q_\ell$-sheaves on $S$. Relative Hard Lefschetz gives the
required isomorphisms of lisse sheaves, so their inverses are also
morphisms of such sheaves. Hence the chosen expression defines
an endomorphism $e_\ell$ of a finite direct sum of relative
cohomology sheaves.
The identity $e_\ell^2=e_\ell$, and the identities saying that
its image has the prescribed homogeneous degree and twist, hold
at the geometric fibre over $x$.
A morphism of lisse sheaves on a connected scheme is zero if its
map at one geometric point is zero. These identities therefore
hold on $S$. In particular
\[
 \mathcal M_\ell=\operatorname{im}(e_\ell)
\]
is a lisse $\Q_\ell$-sheaf whose fibre at $x$ is $H_\ell(M)$.
Its monodromy is compact and therefore preserves a full
$\Z_\ell$-lattice, giving a lisse finite-rank free
$\Z_\ell$-sheaf after a choice of lattice.

Apply Kisin's local constancy theorem to this projected local
system on the rigid analytification of $S_K$:
\citet[Theorem~5.1, author p.~16; Proposition~5.2(2) and proof,
author pp.~17--18]{c_kisin1999localconstancy}.
For every $\ell\in\Sigma$ it gives a $K$-adic neighbourhood
$U_\ell$ of $x$ such that the fibres of $\mathcal M_\ell$ at
$x$ and at any $y\in U_\ell(K)$ are isomorphic as $G_K$-modules.
We apply the theorem to the image local system itself; no
arbitrary isomorphism of ambient cohomology is presumed to
preserve the projector.
The finite intersection $U=\bigcap_{\ell\in\Sigma}U_\ell$
is a neighbourhood of $x$.

There is a closed algebraic point of $S$ whose associated
$K$-point lies in $U$ and whose residue field has completion
exactly $K$. Here are the details of this approximation.
Choose a smooth coordinate chart around the image of $x$ and
an \'etale map $S^\circ\to\mathbb A^n_{F_0}$.
The nonarchimedean inverse-function theorem gives a local
analytic inverse at $x$. For a point
$a\in F_0^n$ sufficiently close to the image of $x$, its lift
$y\in S^\circ(K)$ lies in $U$.
Such $a$ exist because $F_0$ is dense in $K$.
Since the chosen chart is \'etale, the fibre over $a$ is
zero-dimensional of finite type over $F_0$; the residue field
$F=\kappa(y)$ is consequently finite over $F_0$ and is embedded
in $K$ by this lift. Its closure in $K$ contains the dense
subfield $F_0$, so $F_w=K$ for the induced place $w$.
If $n=0$, the point is already algebraic and the same closure
argument applies. Alternatively one may replace $S$ by
$S\times\mathbb A^1$ and pull back the entire presentation.

Specialize the finite cycle and polarization presentation at
$y$ to obtain $M'$ over $F$.
It is a presentation of the required kind: the cycles remain
algebraic, the chosen polarizations remain ample, and the
resulting endomorphism is a homogeneous projector. In the
motivated case this last assertion follows in the defining
homological category because its square-minus-itself has zero
$\ell$-adic realization for any chosen $\ell\in\Sigma$;
Betti--\'etale comparison over the number field identifies this
with the defining Betti equality. Thus no Chow-idempotent
lifting or spread of an arbitrary absolute-Hodge class is used.
Kisin's theorem now gives precisely the isomorphisms in the
statement. The homogeneous degree and twist also show that
the weight is unchanged.

The local universal assertion clearly implies its restriction
to completions of number fields.
Conversely, assume the number-field assertion and fix a local
motive $M/K$ and two primes $\ell_1,\ell_2\ne p$.
Apply the proved approximation with
$\Sigma=\{\ell_1,\ell_2\}$.
The number-field assertion for $M'$ at $w$ and the two local
isomorphisms imply that $H_{\ell_1}(M)$ is unramified if and
only if $H_{\ell_2}(M)$ is unramified.
Since the pair was arbitrary, this is the local assertion.
Finite direct sums of different homogeneous weights cause no
additional issue: inertia is trivial on a direct sum if and
only if it is trivial on every summand.
\end{proof}

\begin{remark}[The remaining purity step in the analytic encoding]
\label{c:bgv-pure-local-boundary}
\status{Exact boundary; no converse to the finite-field conjecture is claimed}
The preceding proposition removes number-field globalization as
an additional obstruction once a pure algebraic or motivated
local realization has been constructed.
It does not turn an arbitrary compact rigid analytic motive
into such a pure motive. In particular, the square-zero
monodromy construction from
\citet[Corollary~4.14]{c_bgv2024monodromy}
still requires a pure local realization, or a proved reduction
to finitely many pure local motives preserving the relevant
unramifiedness flag. The finite-prime approximation is sufficient
for the universal two-prime assertion; an approximation
simultaneously preserving every coefficient prime has not been
asserted.
\end{remark}


\begin{remark}[The remaining general-motive claim]
\label{c:ell-independence-gate}
\status{Unresolved in the general motivic category}
Theorem~\ref{c:ell-independence-abelian} settles the whole abelian-type assertion, including its finite-inertia descent. The direct argument from an arbitrary original generator would still be invalid: $\bigwedge^2H^1(E)=\Q_\ell(-1)$ is unramified at a place of residue characteristic different from $\ell$, even when $E$ has bad reduction. The weak-lift lemma and the trace argument are what repair that inference.

For Andr\'e motivated motives, Theorem~\ref{c:motivated-unramified-rank-reduction} supplies a second named sufficient hypothesis: independence of homological equivalence for rational cycles on all smooth projective varieties over $\overline{\F}_p$. It controls the monodromy ranks and the finite inertia action, but remains conditional on that finite-field conjecture. Corollary~\ref{c:motivated-abelian-strata-unramified} gives the stated unconditional geometric case only when all auxiliary strata have the required algebraic abelian-type presentations. Neither result proves a converse to the finite-field conjecture.

For general motives, another precise named sufficient conjecture is Fontaine's motivic Weil--Deligne compatibility conjecture: \citet[\S2.4.3, Conjecture $C_{\mathrm{WD}}(X,m)$, p.~342, and its motivic formulation in \S2.4.7, pp.~344--345]{c_fontaine1994pst}. The motivic formulation is required to retain the specified projector; compatibility of an unlabelled ambient cohomology representation alone does not identify every prescribed summand. Write the local Weil--Deligne data as $(r_\ell,N_\ell)$, with $r_\ell$ of finite inertial image. On an open subgroup of inertia the actual action is $\exp(t_\ell(\tau)N_\ell)$, where the tame character has infinite image. Consequently
\[
 M_\ell\text{ is unramified at }v
 \quad\Longleftrightarrow\quad
 N_\ell=0\ \text{ and }\ r_\ell|_{I_v}=1.
\]
Both conditions are invariant under extension of coefficient fields and isomorphism of Weil--Deligne representations. Thus Fontaine's compatibility, which only requires isomorphism after common coefficient extension and not a chosen rational form, implies the requested independence. His weaker monodromy-graded character formulation in \S2.4.5 also suffices: the characters at the identity determine whether the monodromy filtration is concentrated in degree zero, equivalently $N=0$; after that, the finite inertia character detects triviality.

This proves an implication, not the exact equivalence requested in the question. Algebraically, the representations trivial on a subgroup $J\subset G$ detect its closed normal closure, not the whole inertia representation. Even all tensor tests do not formally recover the nilpotent's Jordan type: in $\mathfrak{sl}_3$, both $E_{12}$ and $E_{12}+E_{23}$ act nontrivially in every nontrivial algebraic representation, since the differential's kernel is an ideal in the simple Lie algebra. Their ranks differ. The accompanying notes give compatible abstract Weil--Deligne pairs and a matrix certificate; no such pair is asserted to come from a global motive, and it is not a motivic counterexample.

The theorem of \citet[Theorem 5.1]{c_commelin2019} concerns abelian motives and quasi-compatibility; by itself it does not supply finite-inertia traces for an arbitrary motive at a bad place. Similarly the $G$-valued expectation of \citet[Conjecture 1.1]{c_kisinzhou2025} is stronger than the asserted triviality test. A proof in the general motivic category, or an exact equivalence with a published named conjecture of precisely the required strength, remains an outstanding assertion, recorded separately from the now proved abelian-type theorem.
\end{remark}
}{}
\IfFileExists{sections/part_d.tex}{\section{The Serre torus, CM types, and the limits of the class-field slogan}
\label{sec:part-d}

\providecommand{\Res}{\operatorname{Res}}
\providecommand{\MT}{\operatorname{MT}}
\providecommand{\Cl}{\operatorname{Cl}}
\providecommand{\CM}{\operatorname{CM}}
\providecommand{\Tors}{\operatorname{Tors}}
\providecommand{\Isom}{\operatorname{Isom}}
\providecommand{\cor}{\operatorname{cor}}
\providecommand{\res}{\operatorname{res}}
\providecommand{\ord}{\operatorname{ord}}

The distinctions between a finite-modulus Serre group, the connected Serre
protorus, and the Taniyama group are essential here.  We use $S^K$ for the
standard connected Serre torus attached to a Galois CM field $K$, $S$ for its
projective limit, and $S_{\mathfrak m}$ for Serre's possibly disconnected
finite-modulus group.  The latter is a commutative group; the Taniyama group
is generally not.  Statements about the wording of 15:157 and 15:159 below
refer to the supplied AI transcriptions.  The claimed earlier visual check
of ``m\'etab\'eliens'' cannot be independently repeated without the scans.

\subsection{D.1: the integral model and the actual direction of reduction}
\label{subsec:d1}

\subsubsection{The schematic quotient and the narrow class group}

Let $k$ be a number field, $A=\mathcal O_k$, $E^+=A^{\times,+}$, and
$\mathcal T=\Res_{A/\Z}\Gm$.  Let $\mathcal R$ be the schematic closure
of the sections $E^+\subset\mathcal T(\Z)$.  This is the meaning of
``adh\'erence de l'union des sections'' in 15:159.  It is not the profinite
closure occurring in the adjacent class-field-theoretic formula.

\begin{proposition}[Proved: integral finite-modulus Serre group]
\label{prop:d-integral-serre}
The fppf quotient $\mathcal G^0=\mathcal T/\mathcal R$ is a separated,
flat group scheme of finite type over $\Z$.  Its geometric fibres are
connected.  There is an extension of group schemes
\[
  1\longrightarrow\mathcal G^0\longrightarrow\mathcal G
   \longrightarrow\underline{\Cl^+(k)}\longrightarrow1
\]
whose generic fibre is Serre's $S_{\mathfrak m}$ for
$\mathfrak m=\prod_{v\text{ real}}v$.  In particular
$(\mathcal G^0)_\Q=T_{\mathfrak m}
=\Res_{k/\Q}\Gm/\overline{E^+}^{\,\mathrm{Zar}}$.
The assertion does not say that these integral models are tori at ramified
primes or that their representations comprise all CM motives over $k$.
\end{proposition}

\begin{proof}
The coordinate ring of the closure embeds in the product of copies of
$\Z$ obtained by evaluation at the sections.  It is therefore
$\Z$-torsion-free and flat.  Equivalently, $\mathcal R$ is the schematic
closure of the closed generic-fibre subgroup
$\overline{E^+}^{\,\mathrm{Zar}}$.  Multiplication and inversion preserve
this closure: equality on the generic fibre extends over $\Z$, since the
source schemes are flat.  Thus $\mathcal R$ is a flat closed subgroup of
$\mathcal T$.  Representability and finite type of the fppf quotient follow
from \citet[Theorem 4.C, p.~53]{DAnantharaman1973}; separatedness follows
from closedness of $\mathcal R$.  The map $\mathcal T\to\mathcal G^0$ is
faithfully flat, and flatness descends along it.  For an algebraically closed
residue field $\Omega$, the unit group of $A\otimes\Omega$ is a product
of groups of the form $\Gm\times\mathbb A^n$ as varieties, hence is
connected.  Its quotient $\mathcal G^0_\Omega$ is connected as well.

Here is an integral construction of the class-group extension, including
its multiplication.  For every $c\in\Cl^+(k)$ choose an invertible
fractional ideal $\mathfrak a_c$, with $\mathfrak a_0=A$.  Put
\[
 \mathcal P_c=\Isom_{A}(\mathfrak a_c,A),\qquad
 \mathcal Q_c=\mathcal P_c/\mathcal R.
\]
The first is a $\mathcal T$-torsor and the second a
$\mathcal G^0$-torsor; representability follows either from the same
quotient theorem or from fppf descent.  Choose totally positive
$t_{c,d}\in k^\times$ with
$\mathfrak a_c\mathfrak a_d=t_{c,d}\mathfrak a_{c+d}$.  The map
\[
 (u,v)\longmapsto (u\otimes v)\circ[t_{c,d}]
\]
on ideal isomorphisms defines
$\mathcal Q_c\times\mathcal Q_d\to\mathcal Q_{c+d}$.  Any two choices
of $t_{c,d}$ differ by $E^+$, so the map is independent of the choice.
The ratios occurring in the associativity identity are also positive
units; the identity and inverse maps are constructed in the same manner.
Consequently $\mathcal G=\coprod_c\mathcal Q_c$ is the required group
scheme.  Changing the representatives gives the canonically isomorphic
construction by the same ideal-isomorphism torsors.

Over $\Q$ all these torsors are trivial.  Choosing their trivializations
gives precisely the pushout cocycle of
$k^\times/E^+\to I_{\mathfrak m}\to\Cl^+(k)$ under
$k^\times/E^+\to T_{\mathfrak m}(\Q)$.  This is Serre's definition:
\citet[Chapter II, \S1.3 and \S2.2, pp.~II-3--II-9]{DSerre1968}.
\end{proof}

\paragraph{Verification of R4 and the finite-modulus qualification.}
The finite-idelic class-field exact sequence is
\[
 1\longrightarrow\widehat A^\times/\overline{E^+}
 \longrightarrow\Gal(k^{\mathrm{ab}}/k)
 \longrightarrow\Cl^+(k)\longrightarrow1.
\]
Indeed the valuation map on finite ideles has kernel
$\widehat A^\times$, and quotienting by principal ideles with positive
archimedean components gives exactly the narrow ideal class group.
Using all principal ideles instead leaves the archimedean sign group in
the kernel and produces $\Cl(k)$.  This proves the correction asserted in
R4.  Here the bar denotes topological closure; in
Proposition~\ref{prop:d-integral-serre} it denotes schematic closure.

The finite-modulus torus need only be isogenous to the standard connected
Serre torus.  For example, when $k=\Q(i)$ and $\mathfrak m$ has empty
finite support,
\[
 T_{\mathfrak m}=\Res_{k/\Q}\Gm/\mu_4,
 \qquad S^k=\Res_{k/\Q}\Gm.
\]
Their character lattices, with complex conjugation swapping the two
coordinates, are respectively
\[
 \{(a,b)\in\Z^2:a-b\equiv0\pmod4\},\qquad\Z^2.
\]
The first has basis $(1,1),(2,-2)$ and is the direct sum of the trivial
and sign lattices.  In the second, the sum of the invariant and
anti-invariant sublattices has index two.  Hence the lattices, and the
tori, are not isomorphic.  Allowing a sufficiently small congruence
subgroup of the units removes this finite discrepancy.  A fixed modulus
also imposes a fixed conductor condition; taking only the real places
does not include CM motives ramified at finite primes.

\subsubsection{Frobenius, reflex norms, and two false literal assertions}

\begin{proposition}[Refuted literal Frobenius formula; proved correction]
\label{prop:d-frobenius-correction}
The CM type and a uniformizer do not, by the reflex norm alone,
determine the Frobenius element of a CM motive.  They determine its
algebraic Hecke-character part and its principal ideal; the finite-order
character data are also required to determine the element.
\end{proposition}

\begin{proof}
The precise Shimura--Taniyama assertion is an ideal or an idelic
reciprocity assertion.  If $B$ has CM by $E$ of type $\Phi$, all this
endomorphism structure is defined over $k$, and $B$ has good reduction at
$v$, then the Frobenius $\pi_v$ belongs to $E$, its principal ideal is
the appropriate reflex norm of $N_{k/E^*}(v)$, and its slopes satisfy
\begin{equation}
 \frac{\ord_w(\pi_v)}{\ord_w(q_v)}
 =\frac{\#\{\phi\in\Phi_p:\phi\text{ induces }w\}}
 {[E_w:\Q_p]}.
 \label{eq:d-st-slopes}
\end{equation}
Equivalently, the actual Frobenius is the value of the algebraic Hecke
character of $B$ on a uniformizer idele.  The finite-order part of this
character is necessary.  The reflex norm specifies its algebraic part;
it does not by itself remove the unit or root-of-unity ambiguity.
See \citet[Theorem 8.1, Remark 8.6(b), and Theorems 9.3 and 9.10]{DMilneCM2020}
and the particularly explicit
\citet[Theorem 5.1(ii)--(v)]{DCCO2008}.
Here $\pi_v$ is the $q_v$-power Frobenius endomorphism, whose
eigenvalues are those of geometric Frobenius on $H^1$.  The cited
Hecke-character formula uses arithmetic reciprocity on the Tate
module, as specified in \S1.7 of that source.  Dualizing to
cohomology accounts for this convention without changing the
displayed eigenvalues or slopes.

An explicit check is provided by the two curves over $\Q(i)$
\[
 B:y^2=x^3-x,
 \qquad B^{(2)}:y^2=x^3-4x.
\]
They have the same CM type and endomorphism ring $\Z[i]$, and both have
good reduction at the primes above $5$.  Their discriminants are $64$
and $4096$, so good reduction here follows directly from the equations.
Counting points gives
$\#B(\F_5)=8$ and $\#B^{(2)}(\F_5)=4$, and hence Frobenius traces
$-2$ and $2$.  Their Frobenius eigenvalues are
$-1\pm2i$ and $1\pm2i$.  Under the identification of endomorphisms
induced by the quadratic twisting isomorphism, Frobenius changes sign.
Thus the CM type and the chosen prime give the same reflex-norm ideal
but not the same Frobenius element.  The exact enumeration is attached
in \texttt{part\_d/check\_part\_d.py} and its recorded output.
\end{proof}

\begin{proposition}[Refuted: special-fibre gerbe bridge]
\label{prop:d-special-fibre-false}
The proposed morphism from the finite-field motivic gerbe to the
special fibre of the integral group in Lead D1 is not a morphism of
gerbes over the coefficient field, and good reduction does not make
the Frobenius an integral point of that group.
\end{proposition}

\begin{proof}
There is an even simpler obstruction to the proposed special-fibre
bridge in Lead D1.  For $k=\Q$ the integral group above is $\Gm$.
The Tate motive $\Q(-1)$ has good reduction everywhere and Frobenius
eigenvalue $p$ at $p$.  This is not a $\Z_p$-point of $\Gm$, and its
reduction is not a point of $\Gm_{\F_p}$.  Reversing the Frobenius
convention replaces $p$ by $p^{-1}$, which is not integral.  Moreover,
the motivic gerbe of a finite field is a gerbe over the coefficient
field $\Q$, whereas $\mathcal G_{\F_p}$ is a group scheme over
$\F_p$.  These are not two bands over a common base to which the
asserted morphism applies.  Lead D1 in its special-fibre formulation
is therefore false.
\end{proof}

\begin{proposition}[Proved with the stated CM Hodge input: reduction of gerbes]
\label{prop:d-cm-gerbe-reduction}
Reduction of CM motives induces a morphism of gerbes over $\Q$
from the abelian motivic gerbe over $\overline\F_p$ to the CM
motivic gerbe over $\overline\Q$, bound by $P\to S$.  It satisfies
the Langlands--Rapoport local conditions.  These local conditions
alone do not canonically specify its global comparison with a
separately chosen pseudomotivic morphism.
\end{proposition}

\begin{proof}
Fix an extension of $v$ to $\overline\Q$.  Reduction of CM abelian
motives gives an exact tensor functor over $\Q$
\[
 R:\CM(\overline\Q)\longrightarrow
       \operatorname{Mot}_{\mathrm{ab}}(\overline\F_p),
\]
and therefore, by precomposition on fibre functors, a morphism
\begin{equation}
 R^\vee:\operatorname{Mot}_{\mathrm{ab}}(\overline\F_p)^\vee
    \longrightarrow\CM(\overline\Q)^\vee,
 \qquad\text{with band map }\rho:P\longrightarrow S.
 \label{eq:d-reduction-gerbe}
\end{equation}
This is the direction that exists.  Both sides remain over $\Q$.
The construction and its local properties are proved in
\citet[\S7, pp.~39--40]{AMilne2003}.  The Hodge-conjecture hypothesis
on CM abelian varieties used there is the conclusion of
\citet[Theorem I]{AvanHoftenSempliner2026}; this is an explicitly
identified 2026 preprint input, rather than an assumption about general
Hodge classes.  It permits algebraic specialization of the
correspondences and supplies the required abelian Tate and positivity
statements through the implications cited in that source.

The map \eqref{eq:d-reduction-gerbe} satisfies the Langlands--Rapoport
local conditions: the prime-to-$p$ comparison torsors are trivial,
the $p$-adic map is the isocrystal map, and the real map is the
polarized real realization.  These are the statements immediately
preceding \citet[Theorem 7.2]{AMilne2003}.  They do not uniquely
identify a global comparison morphism.  That theorem says that the
classes of compatible $(P,S)$-gerbes form a principal homogeneous
space under $\varprojlim_K\Cl(K)$; Remark 7.3 gives equality at each
finite CM level, and Remark 7.4 warns that the chosen local objects add
further comparison data.  Thus the motivic reduction is a realization
of the Langlands--Rapoport type of morphism.  Identifying it with a
particular separately chosen pseudomotivic morphism requires the
corresponding global comparison choice.  The local statements alone
cannot prove that stronger assertion.
\end{proof}

\paragraph{Weil numbers: the field of reduction matters.}
At the level of Weil-number germs, Honda--Tate and CM lifting after
a finite residue-field extension and isogeny already suffice.  This
is precisely the formulation of \citet[Theorem (1.7)]{DOort2009};
its Theorem (1.9) recalls the Honda--Tate classification.  For exact
eigenvalues over a fixed finite field one must distinguish lifting
the simple isogeny class itself from obtaining its eigenvalue with
larger multiplicity.  The first can fail, whereas the second always
works when the number field and CM type may vary.  We prove this
stronger and precise statement after explaining the obstruction.

For an isotypic $B_0/\F_q$ with a specified CM field
$E\hookrightarrow\End^0(B_0)$ of degree $2\dim B_0$, and a
$p$-adic CM type $\Phi_p$, the \emph{residual reflex condition} is
\eqref{eq:d-st-slopes} together with the requirement that the residue
field of the reflex field at the specified $p$-adic place embed in
$\F_q$.  The existence, up to an $E$-linear $\F_q$-isogeny, of a CM
lifting over a normal characteristic-zero local domain with exactly
that residue field is equivalent to this condition:
\citet[Theorem (4.4)]{DOort2009}.  More precisely,
\citet[Theorem 1.6 and Proposition 5.4]{DCCO2008} produce an
abelian variety over a number field, with good reduction and exactly
the prescribed residue field.  This is a finite criterion once the
CM structure is specified: enumerate the CM types, compute their
stabilizers and local reflex fields, and test the displayed slope
identities and residue degrees.  The squarefree-Frobenius version over
a number field is \citet[Definition 2.14 and Theorem 2.16]{DBergstromEtAl2023},
which cites Chai--Conrad--Oort, Theorems 2.1.6 and 2.5.3.  The latter
monograph's full text was not obtained; the corresponding general
theorem and proof were opened in the authors' 26 July 2008 version,
Theorem 1.6 and Proposition 5.4, pp.~6 and 31--33.  The author-uploaded
copy has erroneous hosting-page metadata; the date and authors cited
here are those printed on its first page.

For example, put $q=4$ and $\pi=2\zeta_5$.  Honda--Tate gives a simple
abelian surface over $\F_4$ with this Frobenius and
$\End^0(B_0)=\Q(\zeta_5)$.  Indeed the sole place above $2$ has
degree four, its slope is $1/2$, so its division-algebra invariant is
$4/2=0$ modulo $\Z$; all other invariants vanish.  The division algebra
is the field itself and the dimension is $4/2=2$.  The four CM types
on $\Q(\zeta_5)$ have trivial stabilizer, so all their reflex fields
are $\Q(\zeta_5)$.  Since $2$ has order four modulo five, the residue
field is $\F_{16}$, which cannot embed in $\F_4$.  Thus this isogeny
class, with its necessarily full CM endomorphism field, has no such
lifting.  This is the fully specified example in
\citet[Example (3.3)]{DOort2009}; the CM-type enumeration and the
polynomial $x^4+2x^3+4x^2+8x+16$ are independently checked in the
attached program.  Allowing a nonnormal domain changes the question:
\citet[Theorem (5.2)]{DOort2009} gives lifting up to isogeny without
extending the finite field in that broader setting.

\begin{lemma}[Proved: a power admits an exact-residue-field CM lifting]
\label{lem:d-power-cm-lift}
For every simple abelian variety $A_0/\F_q$ there exist $m\geq1$,
a number field $k$, a place $v$ of $k$ with $\kappa(v)=\F_q$, and a
CM abelian variety $B/k$, with its full CM field acting over $k$,
such that $B$ has good reduction at $v$ and
$B_v\sim_{\F_q}A_0^m$.  The integer $m$ is not required to equal one.
\end{lemma}

\begin{proof}
Write $D_0=\End^0(A_0)$, $F_0=\Q(\pi)$, and
$d^2=[D_0:F_0]$.  Honda--Tate gives
$2\dim A_0=d[F_0:\Q]$ and a maximal commutative CM subfield
$L_0\subset D_0$ with $[L_0:F_0]=d$.  These precise facts, including
the CM choice, are recalled at the beginning of
\citet[\S2, p.~7]{DCCO2008}.  Fix a Galois CM field $E_0$ containing
$L_0$ and also the imaginary quadratic field
$J=\Q(\sqrt{1-4p})$.  The polynomial $x^2-x+p$ has two distinct
roots modulo $p$, so $p$ splits in $J$.  Consequently complex
conjugation belongs to no decomposition group at $p$ in $E_0$.

Let $N$ be a common denominator of all slopes
$s_w=\ord_w(\pi)/\ord_w(q)$.  These belong to $[0,1]$ and satisfy
$s_{cw}=1-s_w$.  Start with the completion of $E_0$ at the fixed
$p$-adic place and enlarge it to a finite Galois extension $L_1/\Q_p$
containing $p^{1/N}$ and the $N$-th roots of unity.  Its ramification
index is divisible by $N$.  Let $f$ be its residue degree and choose
$\sigma\in\Gal(L_1/\Q_p)$ inducing arithmetic Frobenius on its
residue field.  Put $n=\operatorname{ord}(\sigma)$; then $f\mid n$.
If $U/\Q_p$ is unramified of degree $n$, then
$L=L_1U$ has the same ramification index and
\[
 \Gal(L/\Q_p)=I\rtimes\langle\tau\rangle,
 \qquad\operatorname{ord}(\tau)=n,
\]
where $\tau=(\sigma,\operatorname{Frob}_U)$ in the fibre product
over the cyclic residue group of $L_1$.  This proves the splitting
directly; it is also the splitting construction of
\citet[Remark 4.12 and Lemma 6.1]{DCCO2008}.

We need an enlargement which is still globally CM.  There exists a
totally real Galois number field $R$ whose completion at a specified
place is exactly $L$.  To prove this, take a primitive element of
the Galois extension $L/\Q_p$.  Approximate its minimal polynomial
$p$-adically and, simultaneously, a monic real polynomial with
distinct real roots, coefficient by coefficient by rational numbers.
For sufficiently close approximation Hensel's lemma supplies all
the roots inside $L$ and Krasner's lemma says that each root
generates $L$ over $\Q_p$.  At the real place all the roots remain
real.  The splitting field over $\Q$ is the required $R$.
Set $E=E_0R$.  It is Galois CM; its completion is $L$ and its
decomposition group is $\Delta=I\rtimes\langle\tau\rangle$.
Its inertia order $e=|I|$ is divisible by $N$, and $c\notin\Delta$.

Identify the embeddings of $E$ into $\overline\Q_p$ with
$G=\Gal(E/\Q)$.  The embeddings inducing one $p$-adic place form
a left coset $\Delta g$, of size $en$.  Such cosets occur in
disjoint conjugate pairs $\Delta g,c\Delta g$.  In $\Delta g$
choose exactly $e s_g$ orbits for the left action of
$\langle\tau\rangle$, each of size $n$.  This is possible since
$e s_g$ is an integer between zero and $e$.  On $c\Delta g$ choose
the complement of the conjugate choice.  The resulting subset
$\Phi\subset G$ is a CM type, is fixed by $\tau$, and satisfies
\[
 \frac{|\Phi\cap\Delta g|}{|\Delta|}=s_g.
\]
Thus it has exactly the required slopes.  If
$H=\operatorname{Stab}_G(\Phi)$, its reflex field is $E^H$.
Since $\tau\in H$, one has $I(\Delta\cap H)=\Delta$; the residue
degree of the induced place of $E^H$ is
$[\Delta:I(\Delta\cap H)]=1$.  Its residue field is therefore
$\F_p$, which embeds in the prescribed $\F_q$.

Finally $E$ splits $D_0$, since it contains $L_0$.  Put
$m=[E:F_0]/d$.  The splitting
$D_0\otimes_{F_0}E\simeq M_d(E)$ makes $E^d$ a right
$D_0$-module of rank $m$, with a commuting faithful left $E$-action.
It therefore embeds $E$ as a maximal commutative subfield of
$M_m(D_0)=\End^0(A_0^m)$.  The CM structure $(A_0^m,E,\Phi)$
satisfies both parts of the residual reflex condition.  Applying
\citet[Theorem 1.6, equivalently Proposition 5.4]{DCCO2008}
gives the asserted number-field CM lift with residue field
\emph{exactly} $\F_q$.
\end{proof}

\begin{corollary}[Proved: all Weil numbers occur, with variable CM data]
\label{cor:d-all-weil}
Let $\Gamma_q$ be the group of algebraic $p$-units whose complex
conjugates have absolute value $q^{w/2}$ for some integer $w$.
Every element of $\Gamma_q$ occurs as a Frobenius eigenvalue of the
good reduction of a tensor CM motive over some number field at a
place with residue field $\F_q$.  Every integral weight-one member
already occurs in $H^1$ of a CM abelian variety with such reduction.
No single number field, finite modulus, or primitive CM type is
asserted to work for all members.
\end{corollary}

\begin{proof}
For an integral weight-one Weil number, apply Honda--Tate and
Lemma~\ref{lem:d-power-cm-lift}.  For arbitrary $\alpha\in\Gamma_q$
of weight $w$, choose $a\geq0$ sufficiently large that
$\beta=\alpha q^a$ is integral and $b=w+2a>0$.  Any root
$\gamma=\beta^{1/b}$ is an algebraic integer, because it is
integral over the ring of algebraic integers, and every conjugate
of $\gamma$ has absolute value $q^{1/2}$.  It is therefore an
integral weight-one Weil number.  Choose the CM abelian variety
just constructed with $\gamma$ as a reduction eigenvalue.  Then
$(H^1(B))^{\otimes b}(a)$ has eigenvalue
$\gamma^bq^{-a}=\alpha$.  Conversely, tensor CM motives have Weil
eigenvalues with precisely these archimedean and non-$p$ integrality
properties, by good reduction and the Weil theorem for abelian
varieties, together with tensor products, duals and Tate twists.
\end{proof}

The earlier surface example also illustrates the power distinction
concretely.  For $q=4$, enlarge its CM field to
$E=\Q(\zeta_5,\sqrt{-3})$.  It embeds in
$M_2(\Q(\zeta_5))=\End^0(A_0^2)$.  Take the CM type induced from
$\Q(\sqrt{-3})$.  At $2$ its slope is $1/2$ at each place and its
reflex residue field is $\F_4$, so the residual reflex condition
holds.  Thus $A_0^2$ lifts up to $\F_4$-isogeny, although $A_0$
does not.  The corresponding finite-group calculation is included
in the attached verification program.

\subsection{D.2: toric identity components and the Taniyama group}
\label{subsec:d2}

\subsubsection{The exact general Hodge-class gate}

For either absolute-Hodge or motivated motives over a number field $k$,
write $\CM(k)$ for the Tannakian subcategory generated by Artin motives
and potentially CM abelian varieties.  Fix a complex embedding and the
Betti fibre functor.  The two categories are semisimple and the
motivic group of an object is reductive, possibly disconnected.

\begin{theorem}[Proved: D.2(a) for abelian type and its general equivalent]
\label{thm:d-toric-classification}
For a motive $M$ of abelian type, in either category,
\[
 G_{\mathrm{mot}}(M)^\circ\text{ is commutative}
 \quad\Longleftrightarrow\quad M\in\CM(k).
\]
For arbitrary motives the same classification is equivalent to the
following restricted Hodge-class assertion:

\emph{CM comparison case of the absolute-Hodge (respectively motivated
Hodge) conjecture.}  If $G_{\mathrm{mot}}(M)^\circ$ is commutative and
$N$ is a CM abelian motive over $\overline k$, every rational Hodge
morphism $H_B(M_{\overline k})\to H_B(N)$ is induced by an absolute
Hodge (respectively motivated) correspondence.

It is enough in this assertion to require one Hodge isomorphism to
some such $N$ to lift.  In the absolute-Hodge case this is a precise
restriction of \citet[Conjecture 8.24]{c_delignehodge1982}, rather
than a claim of equivalence with the unrestricted conjecture.
\end{theorem}

\begin{proof}
The Mumford--Tate group $\MT(H_B(M))$ is a connected subgroup of
$G_{\mathrm{mot}}(M)$, because every defining motivic tensor is a
Hodge tensor.  It is therefore contained in the identity component.
If that component is commutative, it is a torus and $H_B(M)$ is a
polarizable Hodge structure of CM type.  Every such Hodge structure
belongs to the tensor category generated by $H^1$ of CM abelian
varieties.  One way to see this explicitly is to use the character
lattice of $S$: the CM-type characters generate it integrally, as
proved below in the verification of R2.  Thus there exists a CM
abelian motive $N$ with a rational Hodge isomorphism
$H_B(M_{\overline k})\simeq H_B(N)$.

For abelian motives, Hodge realization is fully faithful by
\citet[Theorem 8.25, restating Main Theorem 2.11]{c_delignehodge1982}
in the absolute-Hodge category, and by
\citet[Theorem 0.6.2 and \S6.3]{c_andre1996} in the motivated
category.  The isomorphism therefore lifts to an isomorphism of
motives.  All the varieties and correspondences involved are defined
over some finite extension $k'/k$.  The adjunction between base
change and restriction of scalars has unit followed by trace equal
to $[k':k]$; dividing by this degree shows that $M$ is a direct summand
of $\operatorname{Ind}_{k'/k}(M_{k'})$.  Restriction of scalars of a
CM abelian variety is potentially CM, and induction of a Tate object
is Artin--Tate.  Hence this induced motive, and its summand $M$, lie
in $\CM(k)$.

Conversely, over $\overline k$ the category of CM abelian motives is
the category of representations of the Serre protorus.  Over $k$ it
is extended by the profinite Artin group.  Thus the identity component
of every object group in $\CM(k)$ is a quotient of a torus and is
commutative.

For general $M$ the same proof applies as soon as the comparison
isomorphism to $N$ lifts, and so the restricted conjecture implies
the classification.  If the classification holds, both $M$ and $N$
are of abelian type and the already cited full-faithfulness theorem
lifts every Hodge morphism between them.  This proves both directions
of the claimed equivalence, without assuming that a general geometric
motivic group is connected.  Such connectedness is itself conditional
in \citet[Propositions 8.22(a) and 8.23(b)]{c_delignehodge1982}.
\end{proof}

The 2026 theorem on Hodge classes of CM abelian varieties does not
settle this last comparison assertion: the classes here can lie on
$X\times B$ with $X$ an arbitrary smooth projective variety and $B$
CM abelian.  This distinction is reflected explicitly in
\citet[Remark 2.26(i), following Theorem 2.24]{DKahnOtsubo2026},
which treats the corresponding comparison condition for invertible
motives.  Thus D.2(a) is answered by the requested exact equivalence;
no proof of the unrestricted Hodge conjecture is being claimed.

\subsubsection{D.2(b): a non-metabelian Artin counterexample}

\begin{proposition}[Refuted: metabelian motives are all Taniyama quotients]
\label{prop:d-metabelian-false}
The category of Taniyama quotients is $\CM(\Q)$ and contains
motives with non-metabelian Galois group.  Its metabelian objects
are precisely the representations of the maximal metabelian quotient
of the Taniyama group.  Among motives of abelian type, this is
exactly the category with metabelian motivic Galois group.
\end{proposition}

\begin{proof}
The Taniyama group sits in
\begin{equation}
 1\longrightarrow S\longrightarrow\mathcal T
 \longrightarrow\Gal(\overline\Q/\Q)\longrightarrow1.
 \label{eq:d-taniyama}
\end{equation}
The quotient is the full absolute Galois group, not its abelianization.
Deligne's unnumbered main theorem on p.~262 identifies this group with
the motivic group of $\CM(\Q)$, preserving the Artin quotient, Hodge
cocharacter and finite-adelic splitting
\citep[pp.~261--264, main theorem and Lemmas 1--2]{c_delignetaniyama1982}.

Let $f(x)=x^4-x-1$, $\alpha$ a root and
$M=h^0(\Spec\Q(\alpha))$.  Reduction modulo two gives the irreducible
$x^4+x+1$: it has no linear factor and is not divisible by the only
irreducible quadratic $x^2+x+1$.  Thus $f$ is irreducible over $\Q$.
Its discriminant is $-283$.  Modulo seven it factors as
\[
 (x-3)(x^3+3x^2+2x-2),
\]
with the cubic irreducible.  The splitting-field group contains a
four-cycle and a three-cycle.  Its order is divisible by twelve,
and the four-cycle is odd, so the subgroup is $S_4$, not $A_4$.
The permutation representation on the four embeddings is faithful.
Consequently $G_{\mathrm{mot}}(M)=S_4$ in both categories, by
\citet[Proposition 8.17]{c_delignehodge1982} and the same Artin
realization in motivated motives.  But
\[
 [S_4,S_4]=A_4,\qquad[A_4,A_4]=V_4\ne1.
\]
All these finite calculations are checked by the attached program.
This is an Artin motive and hence a quotient representation of
$\mathcal T$, but its motivic group is not metabelian.  Its identity
component is trivial, so it also refutes the inference ``commutative
identity component implies metabelian'' in the supplied reading of R4.

The exact replacement is the category $\CM(\Q)=\Rep_\Q(\mathcal T)$.
Inside it, the objects with metabelian group are exactly
\[
 \Rep_\Q\!\left(\mathcal T/
       \overline{[[\mathcal T,\mathcal T],[\mathcal T,\mathcal T]]}\right),
\]
where the bar means the closed normal subgroup generated by the
indicated commutators, equivalently the intersection of the kernels
of all metabelian algebraic quotients.  This equality follows directly
by factoring a representation through its scheme-theoretic image.
For motives of abelian type, a metabelian reductive connected group
is a torus, so Theorem~\ref{thm:d-toric-classification} places them
inside this category.  For arbitrary motives, the additional step is
exactly the comparison gate in that theorem.
\end{proof}

\subsubsection{D.2(c): comparison of the Hodge structures at two places}

\begin{proposition}[Known and rechecked: the Taniyama comparison]
\label{prop:d-hodge-comparison}
The two Betti fibre functors on geometric CM motives at
$\widetilde\xi$ and $\widetilde\xi'=\sigma\widetilde\xi$
are related by the fibre of the Taniyama group above
$\sigma|_{\overline\Q}$.  Their Hodge types change by the
$\Aut(\C)$ action and their finite-adelic comparison is supplied
by its canonical adelic section.  Formula~\eqref{eq:d-conjugate-cm-type}
gives the explicit CM-type and polarization change.
\end{proposition}

\begin{proof}
Let $\widetilde\xi:k\hookrightarrow\C$ and
$\widetilde\xi'=\sigma\widetilde\xi$, where
$\sigma\in\Aut(\C)$, and let $\tau$ be its restriction to
$\overline\Q$.  The fibre $\mathcal T_\tau$ of
\eqref{eq:d-taniyama} is an $S$-torsor.  It is the tensor-isomorphism
torsor between the Betti fibre functors at the two embeddings on
geometric CM motives.  This identification is
\citet[Proposition 8.23(c)--(d)]{c_delignehodge1982} specialized by
Deligne's Taniyama theorem.  A tensor trivialization transports the
underlying rational vector spaces, and the Hodge cocharacter is
transported to its $\sigma$-conjugate.  Such a rational trivialization
need not be canonical, or even exist; the adelic splitting supplies
the canonical finite-adelic comparison.  In particular, the answer
does not identify the two rational Hodge structures without changing
the CM type.

More explicitly, use homology and write a polarized CM abelian
variety at $\widetilde\xi$ in the form
$(E,\Phi;\mathfrak a,t)$.  For
$f\in f_\Phi(\sigma)\subset\mathbb A_{f,E}^{\times}/E^\times$,
its conjugate has type
\begin{equation}
 (E,\sigma\Phi;
       f\mathfrak a,\ t\chi_{\mathrm{cyc}}(\sigma)/(f\overline f)).
 \label{eq:d-conjugate-cm-type}
\end{equation}
There is an $E$-linear rational homology isomorphism $\alpha$ for
which $(\alpha\otimes1)f$ is the finite-adelic comparison, with
the displayed polarization factor.  These are
\citet[Theorems 1.1 and 1.3]{DMilnePedestrians}.
Applying tensor products, duals and projectors gives the comparison
for every CM motive.  On cohomology one dualizes the formula; this
accounts for the sign of the norm weight in 15:165.

At a real embedding the compatible de Rham structure is the fixed
real form of the Betti comparison under
$F_\infty\otimes$ complex conjugation.  This verifies the real-place
condition used from R4: the involution acts on the Betti factor and
coefficient conjugation acts on $\C$, so their fixed subspace is
the real de Rham realization.  It does not remove the torsor ambiguity
of arbitrary rational fibre functors examined next.
\end{proof}

\subsection{D.3: the finite-place kernel for the actual Serre torus}
\label{subsec:d3}

Let $K/\Q$ be a finite Galois CM extension, $F=K^+$, and
$r=[F:\Q]$.  Set
\[
 \Sha^1_f(\Q,S^K)=
 \ker\!\left(H^1(\Q,S^K)\longrightarrow
                      \prod_{\ell<\infty}H^1(\Q_\ell,S^K)\right).
\]
The subscript $f$ is essential: this is not the ordinary
Tate--Shafarevich group, which also tests the real place.

\begin{theorem}[Proved: finite-place kernel and fibre functors]
\label{thm:d-serre-finite-kernel}
There is a canonical identification
\[
 \Sha^1_f(\Q,S^K)
 \simeq
 \left\{(\epsilon_v)_{v\mid\infty}\in\F_2^r:
                            \sum_v\epsilon_v=0\right\}.
\]
For the geometric neutral category of CM motives whose CM fields
are subfields of $K$, with Betti functor as origin, this group
classifies the isomorphism classes of rational fibre functors that
become isomorphic to Betti at every finite completion.  It has
$2^{r-1}$ elements.  Thus finite comparisons single out Betti only
when $r=1$.  Adding a real comparison singles it out for every $r$.
\end{theorem}

\begin{proof}
There is an exact sequence of tori
\begin{equation}
 1\longrightarrow\Res_{F/\Q}\Gm
 \xrightarrow{a\mapsto(a,N_{F/\Q}(a)^{-1})}
 \Res_{K/\Q}\Gm\times\Gm
 \longrightarrow S^K\longrightarrow1.
 \label{eq:d-serre-resolution}
\end{equation}
Its character sequence says exactly that the coefficients $f_\sigma$
have constant pair sums $f_\sigma+f_{c\sigma}$.  Hilbert 90 gives
\begin{equation}
 H^1(\Q,S^K)=
 \ker\left(\Br(F)\xrightarrow{(\res,-\cor)}
                         \Br(K)\oplus\Br(\Q)\right).
 \label{eq:d-brauer-kernel}
\end{equation}
The same calculation holds over every completion.  This is
\citet[Equation (23) and Lemma 3.1]{AMilne2003}, with the sign
of the second map arising from the inverse norm.

An element in the finite-place kernel corresponds under
\eqref{eq:d-brauer-kernel} to a Brauer class $b$ whose finite local
invariants vanish.  Its real invariants are
$\operatorname{inv}_v(b)=\epsilon_v/2$, with
$\epsilon_v\in\{0,1\}$.  Global Brauer reciprocity imposes
$\sum_v\epsilon_v=0$ in $\F_2$.  Conversely, every such even
pattern is realized by a unique global Brauer class.  Its restriction
to $K$ vanishes because the real places become complex and the
finite invariants are zero.  Its corestriction to $\Q$ also vanishes:
its only possible nonzero invariant is at infinity and is the sum
of the displayed real invariants.  This proves both existence and
uniqueness in the formula.  It also proves that restriction to
$H^1(\R,S^K)$ is an isomorphism on this kernel and that the kernel
over \emph{all} places is zero, in agreement with
\citet[Proposition 3.2]{AMilne2003}.

The geometric CM category under discussion has Tannaka group $S^K$:
all its Hodge characters satisfy the pair-sum condition, and CM
types generate that lattice integrally.  The absolute-Hodge and
motivated versions agree here by the full-faithfulness results cited
in Theorem~\ref{thm:d-toric-classification}.  The usual torsor
classification sends a fibre functor $\omega$ to
$\Isom^\otimes(\omega_B,\omega)$ in $H^1(\Q,S^K)$; conversely an
$S^K$-torsor twists $\omega_B$.  The local comparison requirement
is exactly the displayed kernel.  Since $S^K$ is commutative these
isomorphism classes form a group after choosing Betti as origin,
or intrinsically a torsor under that group.
\end{proof}

\subsubsection*{Verification of R3 and the missing comparison in Lead D3}
Put $U=\Res_{F/\Q}(\Res^1_{K/F}\Gm)$.  There is an exact sequence
\[
 1\longrightarrow\Gm\longrightarrow S^K
       \xrightarrow{(x,z)\mapsto x/\overline x}U\longrightarrow1.
\]
Thus $H^1(\Q,S^K)$ injects into
$H^1(\Q,U)=F^\times/N_{K/F}K^\times$.  The connecting map to
$\Br(\Q)$ is the corestriction of the cyclic algebra
$(K/F,a)$.  A norm class locally trivial at every finite place has
only real invariants; reciprocity imposes the same even-sign
condition, and then its corestriction vanishes.  Therefore this
injection induces an isomorphism on the finite-place kernels.
This derives R3 by the cyclic norm/Brauer exact sequence and proves
the additional step required to use it for $S^K$.  Identifying the
two tori themselves would be incorrect.

\begin{example}[Verified: a non-Betti form with all finite comparisons]
\label{ex:d-non-betti}
Take $F=\Q(\sqrt5)$ and $K=F(i)$.  The quaternion algebra
$(-1,-1)_F$ is nonsplit at both real embeddings.  It is split at
every finite place: over $\Q$ Hamilton's algebra is ramified only
at $2$ and infinity, and $2$ is inert in $F$ because
$x^2-x-1$ is irreducible modulo two.  The degree-two local extension
at $2$ multiplies the invariant $1/2$ by two.  The algebra is split
by $K$, and its corestriction to $\Q$ is trivial since its two real
invariants sum to zero in $\Q/\Z$.  It is the nonzero element of
$\Sha^1_f(\Q,S^K)\simeq\Z/2$.  Equivalently, the corresponding
$U$-torsor is $x^2+y^2=-1$ over $F$.  Twisting Betti by its lift to
an $S^K$-torsor supplies the asserted different rational fibre
functor.  The irreducibility calculation is recorded in the program.
\end{example}

The word ``geometric'' in the theorem cannot be omitted.  Over a
number field the arithmetic CM category also contains the Artin
quotient in \eqref{eq:d-taniyama}.  Its fibre functors are torsors
under that full arithmetic group, or under the particular quotient
generated by the chosen motives; replacing it by $S^K$ without
fixing the Artin data changes the question.  Likewise a smaller CM
subcategory has its actual quotient torus as band.  The computation
above addresses the full geometric category and finite $K$, not
an unexamined inverse limit over all CM fields.

Finally, this is an ambiguity in $H^1$ between fibre functors of an
already neutral category.  The $\mu_2$ at infinity in 15:11 is an
$H^2$ obstruction to neutralizing a different finite-field gerbe.
The former category has the Betti real fibre functor, while the
latter may have no ordinary real fibre functor at all.  Their common
archimedean origin does not identify their cohomology classes.  The
same distinction must be retained when comparing the real functors
of SA06-10: once a real functor exists its forms are $H^1$-classes;
existence itself is the gerbe's $H^2$ question.

\subsection{D.4: the image of the type norm, including its integral lattice}
\label{subsec:d4}

Let $I_E=\Hom(E,\overline\Q)$ and
$I_K=\Hom(K,\overline\Q)$.  Let $\Sigma\subset I_E\times I_K$
be the Galois-invariant CM structure in 15:147--151, and use the
convention in which its type norm has exponent matrix
\[
 A=(a_{\sigma,\tau})_{\sigma\in I_E,\tau\in I_K},
 \qquad a_{\sigma,\tau}=1_\Sigma(\sigma,\tau).
\]
Inverting all exponents to switch between homology and cohomology
does not change the image.  The source map $T_K\to T_E$ factors
through $S_K^0$ as in R2.

\begin{theorem}[Proved: the complete lattice and rank answer to 15:151]
\label{thm:d-reflex-image}
The image $I_\Sigma$ of $\psi_\Sigma$ is the connected $\Q$-torus
with character and cocharacter lattices
\begin{align*}
 X^*(I_\Sigma)&=\Z^{I_E}/\ker(A^{\mathsf t})
                  \simeq\im(A^{\mathsf t})\subset\Z^{I_K},\\
 X_*(I_\Sigma)&=(\im A\otimes\Q)\cap\Z^{I_E}.
\end{align*}
In particular $\dim I_\Sigma=\operatorname{rank}_\Q(A)$.
If $E$ is CM of degree $2g$, put
\[
 C_E=\{x\in\Res_{E/\Q}\Gm:
              x\overline x\in\Gm\subset\Res_{E^+/\Q}\Gm\}.
\]
Then $I_\Sigma\subset C_E$, and the following conditions are equivalent:
\begin{enumerate}[label=(\roman*)]
 \item $I_\Sigma=C_E$;
 \item $\operatorname{rank}_\Q(A)=g+1$;
 \item every rational row relation of $A$ is generated by
 $e_\sigma+e_{c\sigma}-e_\tau-e_{c\tau}$.
\end{enumerate}
These are effective necessary and sufficient conditions for every CM
type, including non-Galois and imprimitive types.  For the CM type
$\Phi$ represented by $\Sigma$, $I_\Sigma=\MT(E,\Phi)$.
\end{theorem}

\begin{proof}
Over $\overline\Q$, the map is
$y_\sigma=\prod_\tau x_\tau^{a_{\sigma,\tau}}$, so its character map
is $A^{\mathsf t}$.  The characters vanishing on its image are
exactly $\ker A^{\mathsf t}$.  That kernel is saturated, because the
target $\Z^{I_K}$ is torsion-free.  Hence the displayed quotient
is a free lattice and describes a connected torus.  Dually, the
cocharacters of a subtorus are the saturation of the image on
cocharacters, which gives the second formula.  Omitting this
saturation can give an incorrect integral lattice even though it
preserves the rank.

For a CM structure each conjugate row pair has sum the all-ones
row.  Therefore every vector
$e_\sigma+e_{c\sigma}-e_\tau-e_{c\tau}$ lies in the character
kernel.  These are exactly the defining relations of $C_E$, whose
dimension is $g+1$.  Inclusion follows.  A connected closed subtorus
of the same dimension is the whole torus, proving (i)$\Leftrightarrow$(ii).
The relation space just described has dimension $g-1$; equality of
its dimension with the full row-kernel dimension proves (iii).
Over $\Z$ the exact answer remains the displayed kernel, computable
by Smith normal form, rather than just a choice of rational basis.

The Hodge cocharacter $\mu_\Phi$ of the rank-one $E$-module has
coordinates $1_\Phi$.  The columns of the reflex-norm matrix are its
Galois conjugates, with repetitions if the chosen source contains
the reflex field properly.  Hence the image is the smallest rational
subtorus containing $\mu_\Phi$.  A rational algebraic subgroup
containing the Hodge homomorphism contains $\mu_\Phi$ and every
Galois conjugate; conversely the torus they generate contains both
$\mu_\Phi$ and its complex conjugate and therefore contains the
Hodge homomorphism.  This is the defining minimality property of
the Mumford--Tate group.  It proves the equality without assuming
nondegeneracy.  The reflex norm identity is stated in
\citet[Proposition 1.23 and \S4, p.~35]{DMilneCM2020}; its
universal cocharacter property is Proposition 4.19, and the CM
criterion for the Mumford--Tate group is Corollary 4.7 there.
\end{proof}

\begin{corollary}[Proved: the compact part in the manuscript's reduction]
\label{cor:d-compact-image}
Let $V'_K=\ker(N_{K/\Q}:S^K\to\Gm)$, and let
$\nu:C_E\to\Gm$ send $x$ to its common product
$x_\sigma x_{c\sigma}$.  Then
\[
 \psi_\Sigma(V'_K)=I_\Sigma\cap\ker\nu,\qquad
 \dim\psi_\Sigma(V'_K)=\operatorname{rank}_\Q A-1,
\]
and the character lattice of this image is
\[
 \im(A^{\mathsf t})/\Z\mathbf1_{I_K}.
\]
The image is all of the norm-one torus of $E/E^+$ precisely when
$\operatorname{rank}_\Q A=g+1$.
\end{corollary}

\begin{proof}
The pair-sum identity gives $\nu\circ\psi_\Sigma=N_{K/\Q}$.
If an element of $I_\Sigma$ has $\nu=1$, any fppf-local lift to
$S^K$ consequently has norm one.  This proves equality of the
fppf images.  The pullback of $\nu$ is the all-ones vector,
which is primitive even in $\Z^{I_K}$, hence also in
$\im(A^{\mathsf t})$.  Its kernel in $I_\Sigma$ is therefore
connected, and the exact character sequence gives the formula
and rank.  The norm-one torus of $E/E^+$ has dimension $g$;
equality follows from the same dimension criterion as above.
Over $\R$ it is compact, which explains the manuscript's reduction
to this part.  The weight cocharacter supplies the omitted scalar
direction; its intersection with the compact part is $\mu_2$,
so this decomposition is an isogeny, not an asserted integral
direct-product decomposition.
\end{proof}

\paragraph{A complete Fourier test in the abelian Galois case.}
If $E/\Q$ is abelian Galois with group $G$ and central conjugation
$c$, use all translates of $1_\Phi$ as columns.  Diagonalizing the
convolution matrix by the characters of $G$ gives
\[
 \dim\MT(E,\Phi)=1+
 \#\left\{\chi\in\widehat G:\chi(c)=-1,
                         \ \sum_{\phi\in\Phi}\chi(\phi)\ne0\right\}.
\]
Indeed the trivial Fourier coefficient is $g$; every nontrivial
even character has sum zero on $\Phi$, since doubling that sum
gives its sum on $G$; and the remaining eigenvalues are precisely
the displayed odd sums.  Thus the full norm-compatible torus
occurs exactly when every odd character has nonzero type sum.
This criterion, together with Theorem~\ref{thm:d-reflex-image}, is
a classification test for all cases, rather than a list of only
the nondegenerate examples.

For example the imaginary quadratic type has rank two, the cyclic
quartic type $\{0,1\}\subset\Z/4$ has rank three, the cyclic sextic
type $\{0,1,2\}\subset\Z/6$ has rank four, and the induced sextic
type $\{0,2,4\}$ has rank two.  For a biquadratic quartic field every
type is induced from an imaginary quadratic subfield and has rank
two, less than $\dim C_E=3$.  The attached program prints all the
matrices, exact ranks, kernels and Smith invariants in these cases.
If $E$ itself is not CM but the CM structure is pulled back from
a minimal CM subfield $E_0$, the same matrix formula applies and
the image is the diagonal image of $\MT(E_0,\Phi_0)$ in $T_E$.
The notation $C_E$ above is used only when $E$ is CM.

\paragraph{Verification of the portions of R2 used here.}
For a CM field with $r$ conjugate pairs, the character lattice of
the connected Serre quotient is
\citep[\S1, equations (1.1)--(1.3)]{DMilneShih1982}
\[
 L_S=\{(n_\sigma): n_\sigma+n_{c\sigma}
                            \text{ is independent of }\sigma\}.
\]
Choose one CM type $\Phi$.  A basis of this lattice is $1_\Phi$
together with the $r$ vectors $e_\sigma-e_{c\sigma}$, one for each
pair with $\sigma\in\Phi$.  Each latter vector is the difference
between two CM-type indicators obtained by flipping one pair.
Consequently all CM types generate $L_S$ over $\Z$, not merely
over $\Q$.  This proves the simultaneous faithfulness invoked in
R2 and in Theorem~\ref{thm:d-serre-finite-kernel}.  The same pair-sum
identity proves that $\psi_\Sigma$ annihilates the norm-one subtorus
from the maximal real subfield and factors through $S_K^0$.
The Galois orbit of a CM type is a subset of the $2^r$ possible
choices, so its reflex field has degree at most $2^r$.  Finally the
character argument for $\MT$ above proves the part of R2 concerning
the universal Hodge cocharacter; it does not identify a fixed
$T_{\mathfrak m}$ with $S^K$ in spite of the finite unit quotient.
}{}
\section{Statement of results and verification record}
The machine-readable results table and the dated progress ledger are part of this research package. They distinguish proved assertions from proof gates and refuted readings. The bibliography records whether each source was opened. The attached archive is the witness for manuscript references; all wording and mathematical content remain subject to that transcription limitation.
\bibliographystyle{apalike}
\bibliography{RESEARCH_SA05_references}
\end{document}
